\documentclass[12pt]{book}

% Idioma y codificación
\usepackage[spanish, es-tabla, es-notilde]{babel}       %es-tabla para que se titule "Tabla"
\usepackage[utf8]{inputenc}

% Márgenes
\usepackage[a4paper,top=3cm,bottom=2.5cm,left=3cm,right=3cm]{geometry}

% Comentarios de bloque
\usepackage{verbatim}

% Paquetes de links
\usepackage[hidelinks]{hyperref}    % Permite enlaces
\usepackage{url}                    % redirecciona a la web

% Más opciones para enumeraciones
\usepackage{enumitem}

% Personalizar la portada
\usepackage{titling}

% Paquetes de tablas
\usepackage{multirow}

% Para añadir el símbolo de euro
\usepackage{eurosym}

% Paquete para subrayar con diferentes estilos
\usepackage[normalem]{ulem}


%------------------------------------------------------------------------

%Paquetes de figuras
\usepackage{caption}
\usepackage{subcaption} % Figuras al lado de otras
\usepackage{float}      % Poner figuras en el sitio indicado H.


% Paquetes de imágenes
\usepackage{graphicx}       % Paquete para añadir imágenes
\usepackage{transparent}    % Para manejar la opacidad de las figuras

% Paquete para usar colores
\usepackage[dvipsnames, table, xcdraw]{xcolor}
\usepackage{pagecolor}      % Para cambiar el color de la página

% Habilita tamaños de fuente mayores
\usepackage{fix-cm}

% Para los gráficos
\usepackage{tikz}
\usepackage{forest}

% Para poder situar los nodos en los grafos
\usetikzlibrary{positioning}

% Para evitar problemas entre babel y tikz
\usetikzlibrary{babel}


%------------------------------------------------------------------------

% Paquetes de matemáticas
\usepackage{mathtools, amsfonts, amssymb, mathrsfs}
\usepackage[makeroom]{cancel}     % Simplificar tachando
\usepackage{polynom}    % Divisiones y Ruffini
\usepackage{units} % Para poner fracciones diagonales con \nicefrac

\usepackage{pgfplots}   %Representar funciones
\pgfplotsset{compat=1.18}  % Versión 1.18

\usepackage{tikz-cd}    % Para usar diagramas de composiciones
\usetikzlibrary{calc}   % Para usar cálculo de coordenadas en tikz

%Definición de teoremas, etc.
\usepackage{amsthm}
%\swapnumbers   % Intercambia la posición del texto y de la numeración

\theoremstyle{plain}

\makeatletter
\@ifclassloaded{article}{
  % --- Configuración para ARTICLE ---
  \newtheorem{teo}{Teorema}[section]
  \theoremstyle{definition}
  \newtheorem{definicion}{Definición}[section]
  \newtheorem{ejercicio}{Ejercicio}[section]

  \renewcommand{\theteo}{%
    \ifnum\value{section}=0 \arabic{teo}\else \thesection.\arabic{teo}\fi
  }
  \renewcommand{\thedefinicion}{%
    \ifnum\value{section}=0 \arabic{definicion}\else \thesection.\arabic{definicion}\fi
  }
  \renewcommand{\theejercicio}{%
    \ifnum\value{section}=0 \arabic{ejercicio}\else \thesection.\arabic{ejercicio}\fi
  }
}{
  % --- Configuración para BOOK / REPORT ---
  \newtheorem{teo}{Teorema}[chapter]
  \theoremstyle{definition}
  \newtheorem{definicion}{Definición}[chapter]
  \newtheorem{ejercicio}{Ejercicio}[section]

  % Si no hay sección: Capítulo.Teorema (ej. 2.1)
  % Si hay sección:    Capítulo.Sección.Teorema (ej. 2.3.1)
  \renewcommand{\theteo}{%
    \ifnum\value{section}=0
      \thechapter.\arabic{teo}%
    \else
      \thesection.\arabic{teo}% (\thesection ya contiene el número de capítulo)
    \fi
  }
  \renewcommand{\thedefinicion}{%
    \ifnum\value{section}=0
      \thechapter.\arabic{definicion}%
    \else
      \thesection.\arabic{definicion}%
    \fi
  }
  \renewcommand{\theejercicio}{%
    \ifnum\value{section}=0
      \arabic{ejercicio}%
    \else
      \thesection.\arabic{ejercicio}% Formato por defecto: ej. 2.2.1
    \fi
  }
}
\makeatother

% Teoremas dependientes o con contador compartido
\newtheorem{coro}{Corolario}[teo]
\newtheorem{prop}[teo]{Proposición}
\newtheorem{lema}[teo]{Lema}
\newtheorem*{lema*}{Lema}

\theoremstyle{remark}
\newtheorem*{observacion}{Observación}

\theoremstyle{definition}
\newtheorem*{notacion}{Notación}
\newtheorem*{ejemplo}{Ejemplo}
\newtheorem*{ejercicio*}{Ejercicio}


% \renewcommand\qedsymbol{$\blacksquare$}         % Cambiar símbolo QED
%------------------------------------------------------------------------

% Paquetes para encabezados
\usepackage{fancyhdr}
\pagestyle{fancy}
\fancyhf{}

\newcommand{\helv}{ % Modificación tamaño de letra
\fontfamily{}\fontsize{12}{12}\selectfont}
\setlength{\headheight}{15pt} % Amplía el tamaño del índice


%\usepackage{lastpage}   % Referenciar última pag   \pageref{LastPage}
%\fancyfoot[C]{%
%  \begin{minipage}{\textwidth}
%    \centering
%    ~\\
%    \thepage\\
%    \href{https://losdeldgiim.github.io/}{\texttt{\footnotesize losdeldgiim.github.io}}
%  \end{minipage}
%}
\fancyfoot[C]{\thepage}
\fancyfoot[R]{\href{https://losdeldgiim.github.io/}{\texttt{\footnotesize losdeldgiim.github.io}}}

%------------------------------------------------------------------------

% Conseguir que no ponga "Capítulo 1". Sino solo "1."
\makeatletter
\@ifclassloaded{book}{
  \renewcommand{\chaptermark}[1]{\markboth{\thechapter.\ #1}{}} % En el encabezado
    
  \renewcommand{\@makechapterhead}[1]{%
  \vspace*{50\p@}%
  {\parindent \z@ \raggedright \normalfont
    \ifnum \c@secnumdepth >\m@ne
      \huge\bfseries \thechapter.\hspace{1em}\ignorespaces
    \fi
    \interlinepenalty\@M
    \Huge \bfseries #1\par\nobreak
    \vskip 40\p@
  }}
}
\makeatother

%------------------------------------------------------------------------
% Paquetes de cógido
\usepackage{minted}
\renewcommand\listingscaption{Código fuente}

\usepackage{fancyvrb}
% Personaliza el tamaño de los números de línea
\renewcommand{\theFancyVerbLine}{\small\arabic{FancyVerbLine}}

% Estilo para C++
\newminted{cpp}{
    frame=lines,
    framesep=2mm,
    baselinestretch=1.2,
    linenos,
    escapeinside=||
}

% para minted
\definecolor{LightGray}{rgb}{0.95,0.95,0.92}
\setminted{
    linenos=true,
    stepnumber=5,
    numberfirstline=true,
    autogobble,
    breaklines=true,
    breakautoindent=true,
    breaksymbolleft=,
    breaksymbolright=,
    breaksymbolindentleft=0pt,
    breaksymbolindentright=0pt,
    breaksymbolsepleft=0pt,
    breaksymbolsepright=0pt,
    fontsize=\footnotesize,
    bgcolor=LightGray,
    numbersep=10pt
}


\usepackage{listings} % Para incluir código desde un archivo

\renewcommand\lstlistingname{Código Fuente}
\renewcommand\lstlistlistingname{Índice de Códigos Fuente}

% Definir colores
\definecolor{vscodepurple}{rgb}{0.5,0,0.5}
\definecolor{vscodeblue}{rgb}{0,0,0.8}
\definecolor{vscodegreen}{rgb}{0,0.5,0}
\definecolor{vscodegray}{rgb}{0.5,0.5,0.5}
\definecolor{vscodebackground}{rgb}{0.97,0.97,0.97}
\definecolor{vscodelightgray}{rgb}{0.9,0.9,0.9}

% Configuración para el estilo de C similar a VSCode
\lstdefinestyle{vscode_C}{
  backgroundcolor=\color{vscodebackground},
  commentstyle=\color{vscodegreen},
  keywordstyle=\color{vscodeblue},
  numberstyle=\tiny\color{vscodegray},
  stringstyle=\color{vscodepurple},
  basicstyle=\scriptsize\ttfamily,
  breakatwhitespace=false,
  breaklines=true,
  captionpos=b,
  keepspaces=true,
  numbers=left,
  numbersep=5pt,
  showspaces=false,
  showstringspaces=false,
  showtabs=false,
  tabsize=2,
  frame=tb,
  framerule=0pt,
  aboveskip=10pt,
  belowskip=10pt,
  xleftmargin=10pt,
  xrightmargin=10pt,
  framexleftmargin=10pt,
  framexrightmargin=10pt,
  framesep=0pt,
  rulecolor=\color{vscodelightgray},
  backgroundcolor=\color{vscodebackground},
}

%------------------------------------------------------------------------
% La etiqueda IA sirve para indicar a partir de donde le pasamos el preambulo a la IA para el escaneo
% IA

% Comandos definidos
\newcommand{\bb}[1]{\mathbb{#1}}
\newcommand{\cc}[1]{\mathcal{#1}}

% I prefer the slanted \leq
\let\oldleq\leq % save them in case they're every wanted
\let\oldgeq\geq
\renewcommand{\leq}{\leqslant}
\renewcommand{\geq}{\geqslant}

% Si y solo si
\newcommand{\sii}{\iff}

% MCD y MCM
\DeclareMathOperator{\mcd}{mcd}
\DeclareMathOperator{\mcm}{mcm}

% Signo
\DeclareMathOperator{\sgn}{sgn}

% Imagen
\DeclareMathOperator{\Ima}{Im}

% Aplicación identidad
\newcommand{\id}{\mathrm{Id}}

% Letras griegas
\newcommand{\eps}{\epsilon}
\newcommand{\veps}{\varepsilon}
\newcommand{\lm}{\lambda}

\newcommand{\ol}{\overline}
\newcommand{\ul}{\underline}
\newcommand{\wt}{\widetilde}
\newcommand{\wh}{\widehat}

\let\oldvec\vec
\renewcommand{\vec}{\overrightarrow}

% Derivadas parciales
\newcommand{\del}[2]{\frac{\partial #1}{\partial #2}}
\newcommand{\Del}[3]{\frac{\partial^{#1} #2}{\partial #3^{#1}}}
\newcommand{\deld}[2]{\dfrac{\partial #1}{\partial #2}}
\newcommand{\Deld}[3]{\dfrac{\partial^{#1} #2}{\partial #3^{#1}}}


\newcommand{\AstIg}{\stackrel{(\ast)}{=}}
\newcommand{\Hop}{\stackrel{L'H\hat{o}pital}{=}}

\newcommand{\red}[1]{{\color{red}#1}} % Para integrales, destacar los cambios.

% Método de integración
\newcommand{\MetInt}[2]{
    \left[\begin{array}{c}
        #1 \\ #2
    \end{array}\right]
}

% Declarar aplicaciones
% 1. Nombre aplicación
% 2. Dominio
% 3. Codominio
% 4. Variable
% 5. Imagen de la variable
\newcommand{\Func}[5]{
    \begin{equation*}
        \begin{array}{rrll}
            \displaystyle #1:& \displaystyle  #2 & \longrightarrow & \displaystyle  #3\\
               & \displaystyle  #4 & \longmapsto & \displaystyle  #5
        \end{array}
    \end{equation*}
}

%------------------------------------------------------------------------


\newcommand{\myhref}[2]{\href{#1}{\color{blue}\uline{#2}}}

\let\oldRe\Re % save them in case they're every wanted
\let\oldIm\Im
\renewcommand{\Re}{\operatorname{Re}} % redefine them
\renewcommand{\Im}{\operatorname{Im}}
\DeclareMathOperator{\Log}{Log}
\DeclareMathOperator{\Arg}{Arg}
\DeclareMathOperator{\ord}{ord}
\DeclareMathOperator{\Ind}{Ind}
\DeclareMathOperator{\Fr}{Fr}
\DeclareMathOperator{\Res}{Res}
\DeclareMathOperator{\ind}{I}

\usetikzlibrary{arrows.meta, decorations.markings} % Cargar las bibliotecas necesarias

% Configuración para las flechas
\tikzset{
    arrow at 1/3/.style={postaction={decorate},
        decoration={markings, mark=at position 0.33 with {\arrow{Stealth}}}},
    arrow at 2/3/.style={postaction={decorate},
        decoration={markings, mark=at position 0.66 with {\arrow{Stealth}}}}
}

\begin{document}

    % 1. Foto de fondo
    % 2. Título
    % 3. Encabezado Izquierdo
    % 4. Color de fondo
    % 5. Coord x del titulo
    % 6. Coord y del titulo
    % 7. Fecha
    % 8. Autor

    % 1. Foto de fondo
% 2. Título
% 3. Encabezado Izquierdo
% 4. Color de fondo
% 5. Coord x del titulo
% 6. Coord y del titulo
% 7. Fecha
% 8. Autor

\newcommand{\portada}[8]{

    \portadaBase{#1}{#2}{#3}{#4}{#5}{#6}{#7}{#8}
    \portadaBook{#1}{#2}{#3}{#4}{#5}{#6}{#7}{#8}
}

\newcommand{\portadaExamen}[8]{

    \portadaBase{#1}{#2}{#3}{#4}{#5}{#6}{#7}{#8}
    \portadaArticle{#1}{#2}{#3}{#4}{#5}{#6}{#7}{#8}
}




\newcommand{\portadaBase}[8]{

    % Tiene la portada principal y la licencia Creative Commons
    
    % 1. Foto de fondo
    % 2. Título
    % 3. Encabezado Izquierdo
    % 4. Color de fondo
    % 5. Coord x del titulo
    % 6. Coord y del titulo
    % 7. Fecha
    % 8. Autor
    
    
    \thispagestyle{empty}               % Sin encabezado ni pie de página
    \newgeometry{margin=0cm}        % Márgenes nulos para la primera página
    
    
    % Encabezado
    \fancyhead[L]{\helv #3}
    \fancyhead[R]{\helv \nouppercase{\leftmark}}
    
    
    \pagecolor{#4}        % Color de fondo para la portada
    
    \begin{figure}[p]
        \centering
        \transparent{0.3}           % Opacidad del 30% para la imagen
        
        \includegraphics[width=\paperwidth, keepaspectratio]{../../_assets/#1}
    
        \begin{tikzpicture}[remember picture, overlay]
            \node[anchor=north west, text=white, opacity=1, font=\fontsize{60}{90}\selectfont\bfseries\sffamily, align=left] at (#5, #6) {#2};
            
            \node[anchor=south east, text=white, opacity=1, font=\fontsize{12}{18}\selectfont\sffamily, align=right] at (9.7, 3) {\href{https://losdeldgiim.github.io/}{\textbf{Los Del DGIIM}, \texttt{\footnotesize losdeldgiim.github.io}}};
            
            \node[anchor=south east, text=white, opacity=1, font=\fontsize{12}{15}\selectfont\sffamily, align=right] at (9.7, 1.8) {Doble Grado en Ingeniería Informática y Matemáticas\\Universidad de Granada};
        \end{tikzpicture}
    \end{figure}
    
    
    \restoregeometry        % Restaurar márgenes normales para las páginas subsiguientes
    \pagecolor{white}       % Restaurar el color de página
    
    
    \newpage
    \thispagestyle{empty}               % Sin encabezado ni pie de página
    \begin{tikzpicture}[remember picture, overlay]
        \node[anchor=south west, inner sep=3cm] at (current page.south west) {
            \begin{minipage}{0.5\paperwidth}
                \href{https://creativecommons.org/licenses/by-nc-nd/4.0/}{
                    \includegraphics[height=2cm]{../../_assets/Licencia.png}
                }\vspace{1cm}\\
                Esta obra está bajo una
                \href{https://creativecommons.org/licenses/by-nc-nd/4.0/}{
                    Licencia Creative Commons Atribución-NoComercial-SinDerivadas 4.0 Internacional (CC BY-NC-ND 4.0).
                }\\
    
                Eres libre de compartir y redistribuir el contenido de esta obra en cualquier medio o formato, siempre y cuando des el crédito adecuado a los autores originales y no persigas fines comerciales. 
            \end{minipage}
        };
    \end{tikzpicture}
    
    
    
    % 1. Foto de fondo
    % 2. Título
    % 3. Encabezado Izquierdo
    % 4. Color de fondo
    % 5. Coord x del titulo
    % 6. Coord y del titulo
    % 7. Fecha
    % 8. Autor


}


\newcommand{\portadaBook}[8]{

    % 1. Foto de fondo
    % 2. Título
    % 3. Encabezado Izquierdo
    % 4. Color de fondo
    % 5. Coord x del titulo
    % 6. Coord y del titulo
    % 7. Fecha
    % 8. Autor

    % Personaliza el formato del título
    \pretitle{\begin{center}\bfseries\fontsize{42}{56}\selectfont}
    \posttitle{\par\end{center}\vspace{2em}}
    
    % Personaliza el formato del autor
    \preauthor{\begin{center}\Large}
    \postauthor{\par\end{center}\vfill}
    
    % Personaliza el formato de la fecha
    \predate{\begin{center}\huge}
    \postdate{\par\end{center}\vspace{2em}}
    
    \title{#2}
    \author{\href{https://losdeldgiim.github.io/}{Los Del DGIIM, \texttt{\large losdeldgiim.github.io}}
    \\ \vspace{0.5cm}#8}
    \date{Granada, #7}
    \maketitle
    
    \tableofcontents
}




\newcommand{\portadaArticle}[8]{

    % 1. Foto de fondo
    % 2. Título
    % 3. Encabezado Izquierdo
    % 4. Color de fondo
    % 5. Coord x del titulo
    % 6. Coord y del titulo
    % 7. Fecha
    % 8. Autor

    % Personaliza el formato del título
    \pretitle{\begin{center}\bfseries\fontsize{42}{56}\selectfont}
    \posttitle{\par\end{center}\vspace{2em}}
    
    % Personaliza el formato del autor
    \preauthor{\begin{center}\Large}
    \postauthor{\par\end{center}\vspace{3em}}
    
    % Personaliza el formato de la fecha
    \predate{\begin{center}\huge}
    \postdate{\par\end{center}\vspace{5em}}
    
    \title{#2}
    \author{\href{https://losdeldgiim.github.io/}{Los Del DGIIM, \texttt{\large losdeldgiim.github.io}}
    \\ \vspace{0.5cm}#8}
    \date{Granada, #7}
    \thispagestyle{empty}               % Sin encabezado ni pie de página
    \maketitle
    \vfill
}
    \portada{ffccA4.jpg}{Inferencia\\Estadística}{Inferencia Estadística}{MidnightBlue}{-8}{28}{2025}{José Juan Urrutia Milán}

    \fancyhead[R]{\helv \nouppercase{\rightmark}}
    \chapter{Ejercicios de clase}
    \noindent
    Este capítulo tiene el propósito de recoger todos los ejercicios propuestos en clase por parte de la profesora que acompañan a los contenidos teóricos. No se pretende que sean ejercicios de examen, sino que son ejercicios de práctica para afianzar los conceptos teóricos.\\
    \section{Estadísticos muestrales}
\begin{ejemplo}
    Obtener la función masa de probabilidad conjunta de una m.a.s. de $X\rightsquigarrow \cc{B}(k_0,p)$ y la función de densidad de una m.a.s. de $X\rightsquigarrow U(a,b)$.\\

    \noindent
    Recordamos que si $X\rightsquigarrow \cc{B}(k_0, p)$, entonces:
    \begin{equation*}
        P[X = x] = \binom{k_0}{x} p^x{(1-p)}^{n-x} \qquad \forall x\in \{0,\ldots,k_0\}
    \end{equation*}
    Por lo que si tenemos una m.a.s. de $n$ variables independientes e idénticamente distribuidas a $X$, $(X_1, \ldots, X_n)$, su función de densidad vendrá dada por:
    \begin{align*}
        P[X_1 &= x_1, \ldots, X_n = x_n] \stackrel{\text{indep.}}{=} \prod_{i=1}^{n}P[X_i = x_i]\stackrel{\text{id. d.}}{=} \prod_{i=1}^{n} P[X=x_i] \\
              &= \prod_{i=1}^{n} \binom{k_0}{x_i} p^{x_i} {(1-p)}^{k_0-x_i} = p^{\sum\limits_{i=1}^{n}x_i} {(1-p)}^{nk_0 - \sum\limits_{i=1}^{n}x_i} \prod_{i=1}^{n}\binom{k_0}{x_i} \\
              & \qquad \forall x_i \in \{0,\ldots,k_0\}
    \end{align*}

    \noindent
    Si ahora $X\rightsquigarrow U(a,b)$ para ciertos $a,b\in \mathbb{R}$ con $a<b$, entonces:
    \begin{equation*}
        f_X(x) = \dfrac{1}{b-a} \qquad \forall x\in [a,b]
    \end{equation*}

    de donde:
    \begin{equation*}
        f_{(X_1, \ldots, X_n)}(x_1, \ldots, x_n) \stackrel{\text{indep.}}{=} \prod_{i=1}^{n} f_{X_i}(x_i)\stackrel{\text{id. d.}}{=} \prod_{i=1}^{n} f_X(x_i) = \prod_{i=1}^{n} \dfrac{1}{b-a} = \dfrac{1}{{(b-a)}^{n}} \qquad \forall x\in [a,b]
    \end{equation*}
\end{ejemplo}

\begin{ejemplo}
    Para cada realización muestral, $(x_1, \ldots, x_n)\in \cc{X}^n$, $F_{x_1,\ldots,x_n}^\ast$ es una función de distribución en $\mathbb{R}$. En particular es una función a saltos, con saltos de amplitud $\nicefrac{1}{n}$ en los sucesivos valores muestrales ordenados de menor a mayor, supuestos que sean distintos, y de saltos múltiplos en el caso de que varios valores muestrales coincidieran.\\

    \noindent
    En las condiciones del enunciado, es decir, suponiendo que $x_1, \ldots, x_n$ están ordenados de menor a mayor y son distintos, entonces es fácil ver que:
    \begin{equation*}
        F_{x_1, \ldots, x_n}^\ast (x) = \left\{\begin{array}{ll}
                0 & \text{si\ } x < x_1 \\
                \nicefrac{1}{n} & \text{si\ } x_1 \leq x < x_2 \\
                                &\vdots \\
                            1 & \text{si\ } x > x_n
        \end{array}\right. \qquad \forall x\in \mathbb{R}
    \end{equation*}
    Por lo que es claro que $F_{x_1, \ldots, x_n}^\ast $ es no decreciente, continua por la derecha, con límite 0 en $-\infty$ y con límite 1 en $+\infty$.
\end{ejemplo}

\begin{ejemplo}
    Para todo $x\in \mathbb{R}$, $F_{X_1, \ldots, X_n}^\ast(x)$ es una variable aleatoria tal que \newline $nF_{X_1, \ldots, X_n}^\ast(x) \rightsquigarrow \cc{B}(n,F(x))$ y:
    \begin{equation*}
        E[F_{X_1, \ldots, X_n}^\ast (x)] = F(x), \qquad Var[F_{X_1, \ldots, X_n}^\ast (x)] = \dfrac{F(x)(1-F(x))}{n}
    \end{equation*}
    donde $F(x)$ es la función de distribución de $X$.\\

    \noindent
    Recordamos que:
    \begin{equation*}
    F_{X_1, \ldots, X_n}^\ast (x) = \dfrac{1}{n}\sum\limits_{i=1}^{n}I_{\left]-\infty,x\right]}(X_i) \qquad \forall x\in \mathbb{R}
    \end{equation*}
Fijado $x\in \mathbb{R}$, tenemos que $I_{\left]-\infty,x\right]}(X)\rightsquigarrow \cc{B}(1,P[X\leq x]) \equiv \cc{B}(1,F(x))$, por lo que por la propiedad reproductiva de la binomial tenemos que:
    \begin{equation*}
        nF_{X_1, \ldots, X_n}^\ast (x) \rightsquigarrow \cc{B}(n,F(x))
    \end{equation*}
    Por lo que:
    \begin{equation*}
        nE[F_{X_1, \ldots, X_n}^\ast (x)] = E[nF_{X_1, \ldots, X_n}^\ast (x)] = nF(x)
        \Longrightarrow
        E[F_{X_1, \ldots, X_n}^\ast (x)] = F(x)
    \end{equation*}
    Para la varianza:
    \begin{multline*}
        n^2Var[F_{X_1, \ldots, X_n}^\ast (x)] = Var[nF_{X_1, \ldots, X_n}^\ast (x)] = nF(x)(1-F(x))
        \Longrightarrow \\ \Longrightarrow
        Var[F_{X_1, \ldots, X_n}^\ast (x)] = \dfrac{F(x)(1-F(x))}{n}
    \end{multline*}
\end{ejemplo}

\begin{ejemplo}
    Para valores grandes de $n$, en virtual del Teorema Central del Límite:
    \begin{equation*}
        F_{X_1, \ldots, X_n}^\ast (x) \rightsquigarrow \cc{N}\left(F(x), \dfrac{F(x)(1-F(x))}{n}\right)
    \end{equation*}

    \noindent
    Sea $(X_1, \ldots, X_n)$ una m.a.s. de $n$ muestras, sea:
    \begin{equation*}
        S_n = \sum\limits_{i=1}^{n} I_{\left]-\infty,x\right]}(X_i) \qquad \forall n\in \mathbb{N}
    \end{equation*}
    Por el Teorema Central del Límite tenemos que:
    \begin{equation*}
        \dfrac{S_n - E[S_n]}{\sqrt{Var[S_n]}} \stackrel{n\to \infty}{\rightsquigarrow} \cc{N}(0,1) \Longrightarrow S_n \stackrel{n\to \infty}{\rightsquigarrow} \cc{N}\left(E[S_n], Var[S_n]\right)
    \end{equation*}
    Como $S_n \rightsquigarrow \cc{B}(n,F(x))$, entonces tenemos que:
    \begin{align*}
        E[S_n] &= nF(x) \\
        Var[S_n] &= nF(x)(1-F(x))
    \end{align*}
    Por lo que:
    \begin{equation*}
        F_{X_1, \ldots, X_n}^\ast (x) = \dfrac{1}{n}S_n \stackrel{n\to \infty}{\rightsquigarrow} \cc{N}\left(F(x), \dfrac{F(x)(1-F(x))}{n}\right)
    \end{equation*}
\end{ejemplo}

\begin{ejemplo}
    Dada una muestra aleatoria simple formada por las observaciones $(3, 8, 5, 4, 5)$, obtener su función de distribución muestral y realizar la representación gráfica.\\

    \noindent
    Aplicando la definición de la función de distribución muestral obtenemos que:
    \begin{equation*}
        F_{(3, 8, 5, 4, 5)}^\ast(x) = \left\{\begin{array}{ll}
                0 & \text{si\ } x < 3 \\
                \nicefrac{1}{5} & \text{si\ } 3 \leq x < 4 \\
                \nicefrac{2}{5} & \text{si\ } 4 \leq x < 5 \\
                \nicefrac{4}{5} & \text{si\ } 5 \leq x < 8 \\
                1 & \text{si\ } x \geq 8 
        \end{array}\right.
    \end{equation*}

    \begin{figure}[H]
        \centering
        \begin{tikzpicture}
        \begin{axis}[
            axis lines=middle,
            xlabel={$x$},
            ylabel={$F_{(3,8,5,4,5)}^\ast (x)$},
            ymin=-0.5, ymax=6,
            xmin=0, xmax=10,
            xtick={3,4,5,8},
            ytick={0,1,2,4,5},
            grid=both,
            width=10cm,
            height=6cm
        ]

        % Tramos de la función
        \addplot[domain=0:3, thick] {0};

        \addplot[domain=3:4, thick] {1};
        \addplot[domain=4:5, thick] {2};
        \addplot[domain=5:8, thick] {4};
        \addplot[domain=8:10, thick] {5};

        % Puntos abiertos y cerrados
        \addplot[only marks, mark=o] coordinates {(3,0) (4,1) (5,2) (8,4)};
        \addplot[only marks, mark=*] coordinates {(3,1) (4,2) (5,4) (8,5)};

        \end{axis}
        \end{tikzpicture}
        \caption{Representación gráfica de $F_{(3,8,5,4,5)}^\ast (x)$.}
    \end{figure}
\end{ejemplo}

\begin{ejemplo}
    Sea $X$ una variable aleatoria con distribución $\cc{B}(1,p)$ con $p\in (0,1)$. Se toma una muestra de tamaño 5, $(X_1, X_2, X_3, X_4, X_5)$, y se obtiene la siguiente observación $(0,1,1,0,0)$. Determinar el valor de los estadísticos estudiados en la observación.\\

    \noindent
    Aplicando las fórmulas vistas en clase obtenemos:
    \begin{itemize}
        \item Media:
        \begin{equation*}
            \overline{X} = \dfrac{0+1+1+0+0}{5} = 0.4
        \end{equation*}
        \item Varianza muestral:
        \begin{equation*}
            \sigma^2 = \dfrac{(0-0.4)^2 + (1-0.4)^2 + (1-0.4)^2 + (0-0.4)^2 + (0-0.4)^2}{5} = \dfrac{1.2}{5} = 0.24
        \end{equation*}
        \item Cuasivarianza:
        \begin{equation*}
            s^2 = \dfrac{(0-0.4)^2 + (1-0.4)^2 + (1-0.4)^2 + (0-0.4)^2 + (0-0.4)^2}{5-1} =
            \dfrac{1.2}{4} = 0.3
        \end{equation*}
        \item Estadísticos ordenados:
        \begin{equation*}
            x_{(1)} = 0, \quad x_{(2)} = 0, \quad x_{(3)} = 0, \quad x_{(4)} = 1, \quad x_{(5)} = 1
        \end{equation*}
    \end{itemize}
\end{ejemplo}

\begin{ejemplo}
    Sea $(X_1, \ldots, X_n)$ una m.a.s. y $\overline{X} = \dfrac{1}{n}\sum\limits_{i=1}^{n}X_i$, entonces:
    \begin{equation*}
        M_{\overline{X}}(t) = {(M_X(\nicefrac{t}{n}))}^{n}
    \end{equation*}
    \begin{equation*}
        M_{\overline{X}}(t) = E\left[e^{t\overline{X}}\right] = E\left[e^{\frac{t}{n}\sum\limits_{i=1}^{n}X_i}\right] = M_{\sum\limits_{i=1}^{n}X_i}\left(\frac{t}{n}\right) \stackrel{\text{indep.}}{=} \prod_{i=1}^{n} M_{X_i}\left(\frac{t}{n}\right) \stackrel{\text{id. d.}}{=} {\left(M_X\left(\frac{t}{n}\right)\right)}^{n}
    \end{equation*}
\end{ejemplo}

\begin{ejemplo}\label{ej:dist_muestral_media}
    Obtener la distribución muestral de $\overline{X}$ para $(X_1, \ldots, X_n)$ una m.a.s. de $X\rightsquigarrow \cc{N}(\mu, \sigma^2)$.

    \begin{equation*}
        M_{\overline{X}}(t) = {\left(M_X\left(\frac{t}{n}\right)\right)}^{n} = {\left(e^{\frac{\mu t}{n} + \frac{\sigma^2 t^2}{2n^2}}\right)}^{n} = e^{\mu t + \frac{\sigma^2 t^2}{2n}}
    \end{equation*}
    Luego $\overline{X}\rightsquigarrow \cc{N}\left(\mu, \frac{\sigma^2}{n}\right)$, ya que la función generatriz de momentos caracteriza la distribución.
\end{ejemplo}

\begin{prop}
    Si tenemos una m.a.s. $(X_1, \ldots, X_n)$, entonces:
    \begin{align*}
        F_{X_{(n)}}(x) &= {(F_X(x))}^{n} \qquad \forall x\in \mathbb{R} \\
        F_{X_{(1)}}(x) &= 1-{(1-F_X(x))}^{n}
    \end{align*}
    \begin{proof}
        Para la distribución del máximo:
        \begin{align*}
            F_{X_{(n)}}(x) &= P[X_{(n)} \leq x] = P[X_1 \leq x, \ldots, X_n \leq x] \stackrel{\text{indep.}}{=} \prod_{i=1}^{n} P[X_i \leq x] \\ &\stackrel{\text{id. d.}}{=} \prod_{i=1}^{n} P[X\leq x] = {(F_X(x))}^{n}
        \end{align*}
        Para la del mínimo:
        \begin{align*}
            F_{X_{(1)}}(x) &= P[X_{(1)}\leq x] = 1-P[X_{(1)} > x] = 1-P[X_1 > x, \ldots, X_n > x] \\ &\stackrel{\text{indep.}}{=} 1 - \prod_{i=1}^{n}P[X_i > x] \stackrel{\text{id. d.}}{=} 1-{(P[X>x])}^{n} = 1-{(1-F_X(x))}^{n}
        \end{align*}
    \end{proof}
\end{prop}

\begin{ejemplo}
    Obtener las distribuciones muestrales de $X_{(1)}$ y  $X_{(n)}$ para \newline $X\rightsquigarrow U(a,b)$.\\

    \noindent
    Si $X\rightsquigarrow U(a,b)$, entonces:
    \begin{equation*}
        F_X(x) = \dfrac{x-a}{b-a} \qquad \forall x\in [a,b]
    \end{equation*}
    Por lo que aplicando la Proposición superior:
    \begin{align*}
        F_{X_{(n)}}(x) &= {(F_X(x))}^{n} = {\left(\dfrac{x-a}{b-a}\right)}^{n} \qquad \forall x\in [a,b] \\
        F_{X_{(1)}}(x) &= 1 - {(1-F_X(x))}^{n} = 1-{\left(1-\dfrac{x-a}{b-a}\right)}^{n} 
        = 1-{\left(\dfrac{b-x}{b-a}\right)}^{n} \qquad \forall x\in [a,b]
    \end{align*}
\end{ejemplo}

\begin{ejemplo}
    Sea $(X_1, \ldots, X_n)$ una m.a.s. de $X$. Entonces, para los momentos no centrados de orden $k$ de la m.a.s. se tiene que:
    \begin{align*}
        E\left[A_k\right] &= E[X^k] \\
        Var\left[A_k\right] &= \dfrac{Var[X^k]}{n}
    \end{align*}

    Demostremos la primera igualdad:
    \begin{align*}
        E\left[A_k\right] &= E\left[\dfrac{1}{n}\sum\limits_{i=1}^{n}X_i^k\right] = \dfrac{1}{n}\sum\limits_{i=1}^{n}E[X_i^k]
        \stackrel{\text{id. d.}}{=} \dfrac{1}{n}\sum\limits_{i=1}^{n}E[X^k] = E[X^k]
    \end{align*}

    Demostremos la segunda igualdad:
    \begin{align*}
        Var\left[A_k\right] &= Var\left[\dfrac{1}{n}\sum\limits_{i=1}^{n}X_i^k\right] = \dfrac{1}{n^2}Var\left[\sum\limits_{i=1}^{n}X_i^k\right]
        \stackrel{\text{indep.}}{=} \dfrac{1}{n^2}\sum\limits_{i=1}^{n}Var[X_i^k]
        \stackrel{\text{id. d.}}{=}\\&\stackrel{\text{id. d.}}{=}
        \dfrac{1}{n^2}\sum\limits_{i=1}^{n}Var[X^k] = \dfrac{Var[X^k]}{n}
    \end{align*}
\end{ejemplo}

    
\section{Distribuciones en el muestreo de poblaciones normales}
\begin{prop}
    Sea $X\rightsquigarrow\cc{N}(0,1)$, entonces $X^2\rightsquigarrow \chi^2(1)$.
    \begin{proof}
        Sea $Y = X^2 = h(X)$, entonces $X = \pm \sqrt{Y} = h^{-1}(y)$, por lo que:
        \begin{equation*}
            f_Y(y) = f_X(h_1^{-1}(y)) \left|\dfrac{dh_1^{-1}(y)}{dy}\right| + f_X(h_2^{-1}(y)) + \left|\dfrac{dh_2^{-1}(y)}{dy}\right| 
        \end{equation*}
        Como $X\rightsquigarrow \cc{N}(0,1)$, entonces:
        \begin{equation*}
            f_X(x) = \dfrac{1}{\sqrt{2\pi}} e^{\frac{-x^2}{2}} \qquad \forall x\in \mathbb{R}
        \end{equation*}
        De donde:
        \begin{equation*}
            f_Y(y) = \dfrac{1}{\sqrt{2\pi}} e^{\frac{-{(\sqrt{y})}^{2}}{2}} \left|\dfrac{1}{2\sqrt{y}}\right| + \dfrac{1}{\sqrt{2\pi}} e^{\frac{-{(-\sqrt{y})}^{2}}{2}} \left|\dfrac{-1}{2\sqrt{y}}\right|  = \dfrac{1}{\sqrt{2\pi y}} e^{\frac{-y}{2}} \qquad \forall y>0
        \end{equation*}
        Por lo que $Y\rightsquigarrow\chi^2(1)$.
    \end{proof}
\end{prop}

\begin{ejemplo}
    Calcula el valor de $k$ o la probabilidad inducida:
    \begin{enumerate}[label=\alph*)]
        \item $P[\chi^2(10)\geq k] = 0.005$.

            Consultando la tabla de la distribución $\chi^2$, obtenemos que $k = 25.1881$.
        \item $P[\chi^2(45) \leq k] = 0.005$.
        
            Tenemos que:
            \begin{equation*}
                P[\chi^2(45) \leq k] = 1-P[\chi^2(45) \geq k] = 0.005 \Longrightarrow
                P[\chi^2(45) \geq k] = 0.995
            \end{equation*}
            Consultando la tabla de la distribución $\chi^2$, obtenemos que $k = 24.3110$.
        
        \item $P[\chi^2(14) \geq 21.06]$
        
            En este caso, no podemos consultar de forma directa los encabezados de la tabla de la distribución $\chi^2$ y obtener el valor de la probabilidad, sino que tenemos que buscar qué columna de la tabla contiene el valor buscado. Podemos ver que:
            \begin{equation*}
                P[\chi^2(14) \geq 21.06] \approx 0.1
            \end{equation*}
            
        \item $P[\chi^2(20) \leq 12.44]$
        
            Tenemos que:
            \begin{equation*}
                P[\chi^2(20) \leq 12.44] = 1-P[\chi^2(20) \geq 12.44] \AstIg 1-0.9 = 0.1
            \end{equation*}
            donde en $(\ast)$ hemos consultado la tabla de la distribución $\chi^2$.
    \end{enumerate}
\end{ejemplo}

\begin{ejemplo}
    Calcula el valor de $k$ o la probabilidad inducida:
    \begin{enumerate}[label=\alph*)]
        \item $P[t(26)\geq k] = 0.05$

            Consultando la tabla de la distribución $t$ de Student, obtenemos que $k = 1.706$.
        \item $P[t(20)\leq k] = 0.25$
        
            Por la simetría de la distribución $t$ de Student, tenemos que:
            \begin{equation*}
                0.25 = P[t(20)\leq k] = P[t(20)\geq -k]
            \end{equation*}

            Consultando la tabla de la distribución $t$ de Student, obtenemos que $-k = 0.6870$, luego $k = -0.6870$.
        \item $P[t(26) \geq k] = 0.9$

            Tenemos que:
            \begin{equation*}
                0.9 = P[t(26) \geq k] = 1-P[t(26) \leq k] = 1-P[t(26) \geq -k] \Longrightarrow P[t(26) \geq -k] = 0.1
            \end{equation*}

            Consultando la tabla de la distribución $t$ de Student, obtenemos que $-k = 1.3150$, luego $k = -1.3150$.
        \item $P[t(21) \geq 1.721]$

            Consultando la tabla de la distribución $t$ de Student, obtenemos que:
            \begin{equation*}
                P[t(21) \geq 1.721] \approx 0.05
            \end{equation*}
        \item $P[t(11) \leq 0.697]$
        
            Consultando la tabla de la distribución $t$ de Student, obtenemos que:
            \begin{equation*}
                P[t(11) \leq 0.697] = 1-P[t(11) \geq 0.697] \approx 1-0.25 = 0.75
            \end{equation*}
        \item $P[t(8) \leq -2.306]$
        
            Tenemos que:
            \begin{equation*}
                P[t(8) \leq -2.306] = P[t(8) \geq 2.306] = 0.025
            \end{equation*}
    \end{enumerate}
\end{ejemplo}

\begin{ejemplo}
    Calcula el valor de $k$ o la probabilidad inducida:
    \begin{enumerate}[label=\alph*)]
        \item $P[F(7,3) \leq k] = 0.95$
        
            Consultando de manera directa la tabla de la distribución $F$ de Snedecor, obtenemos que $k = 8.89$.
        \item $P[F(8,4) \geq k] = 0.01$
            \begin{equation*}
                0.01 = P[F(8,4) \geq k] = 1 - P[F(8,4) \leq k] \Longrightarrow P[F(8,4) \leq k] = 0.99
            \end{equation*}
            Consultando la tabla de la distribución $F$ de Snedecor, obtenemos que $k = 14.8$.
        \item $P[F(2,2) \leq 19]$
        
            Consultando la tabla de la distribución $F$ de Snedecor, obtenemos que:
            \begin{equation*}
                P[F(2,2) \leq 19] = 0.95
            \end{equation*}
        \item $P[F(3,5) \geq 12.1]$
            \begin{equation*}
                P[F(3,5) \geq 12.1] = 1-P[F(3,5) \leq 12.1] \AstIg 1-0.99 = 0.01
            \end{equation*}
            donde en $(\ast)$ hemos consultado la tabla de la distribución $F$ de Snedecor.
        \item $P[F(60,40) \leq k] = 0.05$
            
            Consultando la tabla de la distribución $F$ de Snedecor, obtenemos que $k = 0.627$.
    \end{enumerate}
\end{ejemplo}

\subsection{Muestreo en una población normal unidimensional}
En lo que sigue, sea una $(X_1, \ldots, X_n)$ m.a.s. con $X\rightsquigarrow \cc{N}(\mu, \sigma^2)$.

\begin{prop}
    En dichas condiciones, veamos que:
    \begin{equation*}
        \overline{X}, (X_1 - \overline{X}, \ldots, X_n - \overline{X}) \text{\ son independientes}
    \end{equation*}
    \begin{proof}
        Para ellos, usaremos la caracterización de independencia de vectores aleatorios por la función generatriz de momentos conjunta. Es decir, buscamos probar
        \begin{equation*}
            M_{\overline{X},X_1 - \overline{X}, \ldots, X_n - \overline{X}}(t,t_1, \ldots, t_n) \stackrel{\text{?}}{=} M_{\overline{X}}(t) M_{X_1 - \overline{X}, \ldots, X_n - \overline{X}}(t_1, \ldots, t_n)
        \end{equation*}

        Por la reproductividad de la normal, todas las variables aleatorias que aparecen son normales, luego su función característica de momentos está definida en todo $\bb{R}$. Por tanto, para todo $(t, t_1, \ldots, t_n)\in \bb{R}^{n+1}$, tenemos que:
        \begin{align*}
            M_{\overline{X},X_1 - \overline{X}, \ldots, X_n - \overline{X}}&(t, t_1, \ldots, t_n) 
            = E\left[e^{(t, t_1, \ldots, t_n) \cdot (\overline{X},X_1 - \overline{X}, \ldots, X_n - \overline{X})}\right]
            = E\left[e^{t\overline{X} + \sum\limits_{i=1}^{n}(X_i - \overline{X})t_i}\right] 
            =\\&= E\left[e^{\frac{t}{n}\sum\limits_{i=1}^n X_i + \sum\limits_{i=1}^n X_it_i - \sum\limits_{i=1}^n \overline{X}t_i}\right] 
            = E\left[e^{\frac{t}{n}\sum\limits_{i=1}^n X_i + \sum\limits_{i=1}^n X_it_i - \overline{X}\sum\limits_{i=1}^n t_i}\right] 
            =\\&= E\left[e^{\frac{t}{n}\sum\limits_{i=1}^n X_i + \sum\limits_{i=1}^n X_it_i - \frac{1}{n}\sum\limits_{i=1}^n X_i\sum\limits_{i=1}^n t_i}\right] 
            = E\left[e^{\frac{t}{n}\sum\limits_{i=1}^n X_i + \sum\limits_{i=1}^n X_it_i - \sum\limits_{i=1}^n X_i\frac{1}{n}\sum\limits_{i=1}^n t_i}\right] 
            =\\&= E\left[e^{\frac{t}{n}\sum\limits_{i=1}^n X_i + \sum\limits_{i=1}^n X_it_i - \sum\limits_{i=1}^n X_i\overline{t}}\right] 
            = E\left[e^{\sum\limits_{i=1}^n X_i\frac{t}{n} + \sum\limits_{i=1}^n X_it_i - \sum\limits_{i=1}^n X_i\overline{t}}\right] 
            =\\&= E\left[e^{\sum\limits_{i=1}^n X_i \left(\frac{t}{n} + t_i - \overline{t}\right)}\right] 
            = E\left[\prod_{i=1}^{n}e^{X_i\left(\frac{t}{n}+t_i - \overline{t}\right)}\right] 
            \stackrel{\text{indep.}}{=} \prod_{i=1}^{n} E\left[e^{X_i\left(\frac{t}{n}+t_i - \overline{t}\right)}\right]
            =\\&= \prod_{i=1}^{n} M_{X_i}\left(\frac{t}{n}+t_i - \overline{t}\right) 
            \AstIg \prod_{i=1}^{n} \exp\left[\left(\frac{t}{n}+t_i-\overline{t}\right)\mu + {\left(\frac{t}{n}+t_i - \ol{t}\right)}^{2}\frac{\sigma^2}{2}\right]
            =\\&=\exp\left[\sum\limits_{i=1}^n \left[\left(\frac{t}{n}+t_i - \overline{t}\right)\mu + \frac{\sigma^2}{2} \left(\frac{t^2}{n^2} + {\left(t_i - \overline{t}\right)}^{2} + 2\frac{t}{n}(t_i - \overline{t})\right)\right]\right] 
            =\\&= \exp\left[\sum\limits_{i=1}^n \frac{\mu t}{n} + \sum\limits_{i=1}^n \mu (t_i-\overline{t}) + \frac{\sigma^2}{2}\left(\sum\limits_{i=1}^n \frac{t^2}{n^2} + \sum\limits_{i=1}^n {(t_i - \overline{t})}^{2} + 2\sum\limits_{i=1}^n \frac{t}{n}(t_i - \overline{t}) \right)\right]
            \stackrel{(\ast\ast)}{=}\\&\stackrel{(\ast\ast)}{=} \exp\left[\frac{\cancel{n}\mu t}{\cancel{n}} + \mu \cancel{\sum\limits_{i=1}^n (t_i - \overline{t})} + \frac{\sigma^2}{2} \left(\frac{\cancel{n}t^2}{n^{\cancel{2}}} + \sum\limits_{i=1}^n {(t_i - \overline{t})}^{2} + 2\frac{t}{n}\cancel{\sum\limits_{i=1}^n (t_i - \overline{t})}\right)\right]
            =\\&= \exp\left[\mu t + \frac{\sigma^2t^2}{2n} + \frac{\sigma^2}{2}\sum\limits_{i=1}^n {(t_i - \overline{t})}^2{}\right]
        \end{align*}
        donde en $(\ast)$ hemos usado que todas las variables son idénticamente distribuidas a una $\cc{N}(\mu, \sigma^2)$; y en $(\ast\ast)$ hemos usado que el momento centrado de primer orden es cero.
        Sabemos que:
        \begin{align*}
            M_{\overline{X}}(t) &= M_{\left(\overline{X},X_1 - \overline{X}, \ldots, X_n - \overline{X}\right)}(t,0, ..., 0) = e^{\mu t + \frac{\sigma^2t^2}{2n}} \\
            M_{(X_1 - \overline{X}, \ldots, X_n - \overline{X})}(t_1, ..., t_n) &= M_{(\overline{X},X_1 - \overline{X}, \ldots, X_n - \overline{X})}(0,t_1, \ldots, t_n) \\
                                &= e^{\frac{\sigma^2}{2}\sum\limits_{i=1}^n {(t_i - \overline{t})}^{2}}
        \end{align*}

        Por tanto, tenemos que el producto de las funciones generatrices de momentos es la generatriz de mmoentos conjunta, luego las variables son independientes.
    \end{proof}
\end{prop}

\begin{coro}\label{cor:muestreo_1normal}
    En dichas condiciones, tenemos que:
    \begin{enumerate}
        \item $\ol{X}\rightsquigarrow \cc{N}(\mu, \frac{\sigma^2}{n})$
        \item (Lema de Fisher): $\ol{X}$ y $S^2$ son independientes.
        \item $\dfrac{(n-1)S^2}{\sigma^2}\rightsquigarrow \chi^2(n-1)$
        \item $\dfrac{\ol{X}-\mu}{\nicefrac{S}{\sqrt{n}}}\rightsquigarrow t(n-1)$
    \end{enumerate}
\begin{proof}
    Demostramos cada una de las afirmaciones:
    \begin{itemize}
        \item Se vio ya en el ejemplo~\ref{ej:dist_muestral_media}.
        \item Como $S^2$ es función del vector de la Proposición anterior, tenemos que es independiente a $\overline{X}$, ya que las funciones de variables independientes son independientes.
        \item Queremos demostrar que:
            \begin{equation*}
                \dfrac{(n-1)S^2}{\sigma^2} = \dfrac{\sum\limits_{i=1}^n {(X_i - \overline{X})}^{2}}{\sigma^2} \rightsquigarrow\chi^2(n-1)
            \end{equation*}
            
            Tenemos que:
            \begin{align*}
                \dfrac{\sum\limits_{i=1}^n {(X_i - \mu)}^{2}}{\sigma^2} &= \dfrac{\sum\limits_{i=1}^n{(X_i - \overline{X} + \overline{X} - \mu)}^{2}}{\sigma^2} \\ 
                &= \dfrac{\sum\limits_{i=1}^n {(X_i - \overline{X})}^{2} + \sum\limits_{i=1}^n {(\overline{X}-\mu)}^{2} + 2\cancel{\sum\limits_{i=1}^n (X_i - \overline{X})}(\overline{X} - \mu)}{\sigma^2} \\
                &= \dfrac{\sum\limits_{i=1}^n {(X_i - \overline{X})}^{2} }{\sigma^2} + n \dfrac{{(\overline{X}-\mu)}^{2}}{\sigma^2}
            \end{align*}

            Notemos que estamos interesados en conocer la distribución de la primera parte de la suma. Notemos por $A$ a la suma completa, por $B$ a la primera parte y por $C$ a la segunda parte, de forma que $A=B+C$.            
            Veamos la distribución de cada término:
            \begin{itemize}
                \item Por un lado para $A$, tenemos que:
                \begin{equation*}
                    A = \dfrac{\sum\limits_{i=1}^n {(X_i - \mu)}^{2}}{\sigma^2} =
                    \sum\limits_{i=1}^n {\left(\dfrac{X_i - \mu}{\sigma}\right)}^{2} \AstIg \sum\limits_{i=1}^n Z_i^2 \rightsquigarrow \chi^2(n)
                \end{equation*}
                donde en $(\ast)$ hemos usado que $X_i\rightsquigarrow \cc{N}(\mu, \sigma^2)$, por lo que su variable tipificada $Z_i\rightsquigarrow \cc{N}(0,1)$. Además, en la última igualdad hemos empleado la construcción de la distribución $\chi^2$ a partir de variables normales.
                \item Por otro lado, para $C$, del ejemplo~\ref{ej:dist_muestral_media}, tenemos que $\ol{X}\rightsquigarrow \cc{N}(\mu, \frac{\sigma^2}{n})$, por lo que, tipificando 
                \begin{equation*}
                    \dfrac{\ol{X}-\mu}{\sqrt{\nicefrac{\sigma^2}{n}}} = \dfrac{\ol{X}-\mu}{\nicefrac{\sigma}{\sqrt{n}}}=\sqrt{n}\dfrac{\ol{X}-\mu}{\sigma} \rightsquigarrow \cc{N}(0,1)
                \end{equation*}

                Elevando al cuadrado y empleando la construcción de la distribución $\chi^2$, obtenemos que:
                \begin{equation*}
                    C = n\dfrac{{(\ol{X}-\mu)}^{2}}{\sigma^2} \rightsquigarrow \chi^2(1)
                \end{equation*}
            \end{itemize}

            Buscamos la distribución de $B$, y para ello, buscamos su caracterización por la función generatriz de momentos. Tenemos que $B=f(S^2)$, mientras que $C=g(\ol{X})$. Como, por el Lema de Fisher, $\ol{X}$ y $S^2$ son independientes, tenemos que $B$ y $C$ son independientes. Por lo que:
            \begin{equation*}
                M_{A=B+C}(t) \stackrel{\text{indep.}}{=} M_B(t) M_C(t)
            \end{equation*}

            Veamos la función generatriz de momentos de $A$ y $C$:
            \begin{itemize}
                \item Para $A$, tenemos que $A\rightsquigarrow \chi^2(n)$, por lo que:
                \begin{equation*}
                    M_A(t) = \dfrac{1}{{(1-2t)}^{\nicefrac{n}{2}}} \qquad t < \frac{1}{2}
                \end{equation*}
                \item Para $C$, tenemos que $C\rightsquigarrow \chi^2(1)$, por lo que:
                \begin{equation*}
                    M_C(t) = \dfrac{1}{{(1-2t)}^{\frac{1}{2}}} \qquad t < \frac{1}{2}
                \end{equation*}
            \end{itemize}

            Por tanto, tenemos que:
            \begin{equation*}
                M_B(t) = \dfrac{M_A(t)}{M_C(t)} = \dfrac{\dfrac{1}{{(1-2t)}^{\frac{n}{2}}}}{\dfrac{1}{{(1-2t)}^{\frac{1}{2}}}} = \dfrac{1}{{(1-2t)}^{\frac{n-1}{2}}} \qquad  t<\frac{1}{2}
            \end{equation*}
            De aquí, deducimos que $B\rightsquigarrow \chi^2(n-1)$.
        \item Queremos ver que $\dfrac{\overline{X}-\mu}{\nicefrac{S}{\sqrt{n}}}\rightsquigarrow t(n-1)$.

            Para ello, al igual que la $\chi^2$, lo más sencillo es ir a la construcción de $t$:
            \begin{equation*}
                \left.\begin{array}{l}
                        X\rightsquigarrow \cc{N}(0,1) \\
                        Y \rightsquigarrow\chi^2(n) \\
                        \text{$X$ e $Y$ independientes}
                \end{array}\right\}\Longrightarrow \dfrac{X}{\sqrt{\nicefrac{Y}{n}}} \rightsquigarrow t(n)
            \end{equation*}
            
            Por un lado, como $\ol{X}\rightsquigarrow \cc{N}(\mu, \frac{\sigma^2}{n})$, tenemos que:
            \begin{equation*}
                \dfrac{\ol{X}-\mu}{\nicefrac{\sigma}{\sqrt{n}}} \rightsquigarrow \cc{N}(0,1)
            \end{equation*}

            Por otro lado, del punto anterior, tenemos que:
            \begin{equation*}
                \dfrac{(n-1)S^2}{\sigma^2} \rightsquigarrow \chi^2(n-1)
            \end{equation*}

            Como $\ol{X}$ y $S^2$ son independientes por el Lema de Fisher, tenemos que ambas variables aleatorias consideradas son independientes, por lo que podemos aplicar la construcción de la distribución $t$:
            \begin{equation*}
                \dfrac{\dfrac{\overline{X}-\mu}{\nicefrac{\sigma}{\sqrt{n}}}}{\sqrt{\dfrac{\cancel{(n-1)}S^2}{\sigma^2\cancel{(n-1)}}}} = \dfrac{\dfrac{\overline{X}-\mu}{\nicefrac{\cancel{\sigma}}{\sqrt{n}}}}{\dfrac{S}{\cancel{\sigma}}} = \dfrac{\overline{X}-\mu}{\nicefrac{S}{\sqrt{n}}} \rightsquigarrow t(n-1)
            \end{equation*}
    \end{itemize}
\end{proof}
\end{coro}

Gracias a este corolario, podemos inferir sobre $\mu$ y $\sigma^2$ de la población normal a partir de una m.a.s. de la misma. En particular, tenemos que los siguientes resultados, con demostración evidente a partir del corolario anterior:
\begin{enumerate}
    \item Si queremos inferir sobre $\mu$:
    \begin{enumerate}
        \item Si conocemos $\sigma^2_0$ (notaremos con el subíndice $0$ que es un valor conocido), entonces:
        \begin{equation*}
            \dfrac{\ol{X}-\mu}{\nicefrac{\sigma_0}{\sqrt{n}}} \rightsquigarrow \cc{N}(0,1)
        \end{equation*}

        \item Si no conocemos $\sigma^2$, entonces:
        \begin{equation*}
            \dfrac{\ol{X}-\mu}{\nicefrac{S}{\sqrt{n}}} \rightsquigarrow t(n-1)
        \end{equation*}
    \end{enumerate}

    \item Si queremos inferir sobre $\sigma^2$:
    \begin{enumerate}
        \item Si conocemos $\mu_0$, entonces:
        \begin{equation*}
            \dfrac{\sum\limits_{i=1}^n {(X_i - \mu_0)}^{2}}{\sigma^2} \rightsquigarrow \chi^2(n)
        \end{equation*}
        \item Si no conocemos $\mu$, entonces:
        \begin{equation*}
            \dfrac{\sum\limits_{i=1}^n {(X_i - \ol{X})}^{2}}{\sigma^2} \rightsquigarrow \chi^2(n-1)
        \end{equation*}
    \end{enumerate}
\end{enumerate}

\subsection{Muestreo en dos poblaciones normales unidimensionales}

En lo que sigue, sean $(X_1, \ldots, X_{n_1})$ y $(Y_1, \ldots, Y_{n_2})$ m.a.s. independientes con $X\rightsquigarrow \cc{N}(\mu_1, \sigma_1^2)$ y $Y\rightsquigarrow \cc{N}(\mu_2, \sigma_2^2)$.
\begin{teo}[Extensión del Lema de Fisher]
    Los vectores $(\overline{X},\overline{Y})$, $(S_1^2, S_2^2)$ son independientes.
    \begin{proof}
        Usemos la función generatriz de momentos, ya que son independientes si y solo si:
        \begin{equation*}
            M_{\left(\overline{X},\overline{Y},S_1^2, S_2^2\right)}(t_1,t_2,s_1,s_2) \stackrel{\text{?}}{=} M_{\left(\overline{X},\overline{Y}\right)}(t_1,t_2) M_{\left(S_1^2,S_2^2\right)}(s_1,s_2)
        \end{equation*}
        Como $\left(\overline{X},S_1^2\right)= f(X_1, \ldots, X_{n_1})$ y $\left(\overline{Y},S_2^2\right)=g(Y_1, \ldots, Y_{n_2})$, tenemos que ambos vectores son independientes, por lo que:
        \begin{equation*}
            M_{\left(\overline{X},\overline{Y},S_1^2, S_2^2\right)}(t_1,t_2,s_1,s_2) = M_{\left(\overline{X},S_1^2\right)}(t_1,s_1) M_{(\overline{Y},S_2^2)}(t_2,s_2)
        \end{equation*}

        Aplicando ahora el Lema de Fisher a cada uno de los vectores, tenemos que:
        \begin{align*}
            M_{\left(\overline{X},\overline{Y},S_1^2, S_2^2\right)}(t_1,t_2,s_1,s_2) = M_{\overline{X}}(t_1) M_{S_1^2}(s_1) M_{\overline{Y}}(t_2) M_{S_2^2}(s_2)
        \end{align*}

        Nuevamente, como $\ol{X}= f(X_1, \ldots, X_{n_1})$ y $\ol{Y}=g(Y_1, \ldots, Y_{n_2})$, y $S_1^2 = h(X_1, \ldots, X_{n_1})$ y $S_2^2 = k(Y_1, \ldots, Y_{n_2})$, tenemos que $\ol{X}$ y $\ol{Y}$ son independientes, y $S_1^2$ y $S_2^2$ son independientes. Por tanto:
        \begin{align*}
            M_{\left(\overline{X},\overline{Y},S_1^2, S_2^2\right)}(t_1,t_2,s_1,s_2) = M_{\left(\overline{X},\overline{Y}\right)}(t_1,t_2) M_{\left(S_1^2,S_2^2\right)}(s_1,s_2)
        \end{align*}
        
        De aquí deducimos que los vectores $(\overline{X},\overline{Y})$ y $(S_1^2, S_2^2)$ son independientes.
    \end{proof}
\end{teo}

\begin{coro}
    En dichas condiciones, tenemos que:
    \begin{enumerate}
        \item $\displaystyle \dfrac{n_2}{n_1}\dfrac{\sigma_2^2}{\sigma_1^2} \dfrac{\sum\limits_{i=1}^{n_1}{(X_i-\mu_1)}^{2}}{\sum\limits_{i=1}^{n_2}{(Y_i - \mu_2)}^{2}} \rightsquigarrow F(n_1, n_2)$
        \item $\displaystyle \dfrac{\nicefrac{S_1^2}{\sigma_1^2}}{\nicefrac{S_2^2}{\sigma_2^2}} \rightsquigarrow F(n_1-1,n_2-1)$.\\
        En particular, si $\sigma_1 = \sigma_2$, se tiene:
        \begin{equation*}
            \dfrac{S_1^2}{S_2^2}\rightsquigarrow F(n_1-1,n_2-1)
        \end{equation*}
        \item $\displaystyle \dfrac{\ol{X}-\ol{Y} -(\mu_1-\mu_2)}{\sqrt{\dfrac{(n_1-1)S_1^2}{\sigma_1^2} + \dfrac{(n_2-1)S_2^2}{\sigma_2^2}}}\cdot \sqrt{\dfrac{n_1+n_2-2}{\dfrac{\sigma_1^2}{n_1} + \dfrac{\sigma_2^2}{n_2}}} \rightsquigarrow t(n_1+n_2-2)$.\\
        En particular, si $\sigma_1 = \sigma_2$, se tiene:
        \begin{equation*}
            \dfrac{\ol{X}-\ol{Y} -(\mu_1-\mu_2)}{\sqrt{(n_1-1)S_1^2 + (n_2-1)S_2^2}}\cdot \sqrt{\dfrac{n_1+n_2-2}{\dfrac{1}{n_1} + \dfrac{1}{n_2}}} \rightsquigarrow t(n_1+n_2-2)
        \end{equation*}        
    \end{enumerate}
    \begin{proof} Demostramos cada una de las afirmaciones:
    \begin{enumerate}
        \item Partiremos de la construcción de la distribución $F$ de Snedecor:
            \begin{equation*}
                \left.\begin{array}{l}
                    X\rightsquigarrow \chi^2(m) \\
                    Y\rightsquigarrow\chi^2(n)\\
                    \text{$X$ e $Y$ independientes}
                \end{array}\right\} \Longrightarrow \dfrac{\nicefrac{X}{m}}{\nicefrac{Y}{n}} \rightsquigarrow F(m,n)
            \end{equation*}

            Por otro lado, como $X\rightsquigarrow \cc{N}(\mu_1, \sigma_1^2)$ y $Y\rightsquigarrow \cc{N}(\mu_2, \sigma_2^2)$, y haciendo uso de la construcción de la distribución $\chi^2$ a partir de variables normales, tenemos que:
            \begin{align*}
                \dfrac{\sum\limits_{i=1}^{n_1}{(X_i - \mu_1)}^{2}}{\sigma_1^2} &\rightsquigarrow \chi^2(n_1) \qquad \qquad
                \dfrac{\sum\limits_{i=1}^{n_2}{(Y_i - \mu_2)}^{2}}{\sigma_2^2} &\rightsquigarrow \chi^2(n_2) 
            \end{align*}

            Como las m.a.s. son independientes y ambas distribuciones $\chi^2$ son funciones de las m.a.s., tenemos que ambas variables son independientes. Por tanto, podemos aplicar la construcción de la distribución $F$ de Snedecor:
            \begin{equation*}
                \dfrac{\frac{\sum\limits_{i=1}^{n_1}{(X_i-\mu_1)}^{2}}{n_1\sigma_1^2}}{\frac{\sum\limits_{i=1}^{n_2}{(Y_i - \mu_2)}^{2}}{n_2\sigma_2^2}} = \dfrac{n_2}{n_1}\dfrac{\sigma_2^2}{\sigma_1^2} \dfrac{\sum\limits_{i=1}^{n_1}{(X_i-\mu_1)}^{2}}{\sum\limits_{i=1}^{n_2}{(Y_i - \mu_2)}^{2}} \rightsquigarrow F(n_1,n_2)
            \end{equation*}

        \item Seguimos buscando la construcción de $F$ de Snedecor. Por el Colorario~\ref{cor:muestreo_1normal} sabemos que:
            \begin{align*}
                \dfrac{(n_1-1)S_1^2}{\sigma_1^2} \rightsquigarrow\chi^2(n_1-1) \qquad \qquad
                \dfrac{(n_2-1)S_2^2}{\sigma_2^2} \rightsquigarrow\chi^2(n_2-1)
            \end{align*}
            Estas son independientes por ser funciones de ambas m.a.s., que son independientes. Dividimos:
            \begin{equation*}
                \dfrac{\dfrac{\cancel{(n_1-1)}S_1^2}{\sigma_1^2\cancel{(n_1-1)}}}{\dfrac{\bcancel{(n_2-1)}S_2^2}{\sigma_2^2\bcancel{(n_2-1)}}} = \dfrac{\nicefrac{S_1^2}{\sigma_1^2}}{\nicefrac{S_2^2}{\sigma_2^2}} \rightsquigarrow F(n_1-1,n_2-1)
            \end{equation*}
        \item En este caso, buscamos aplicar la construcción de la distribución $t$ de Student:
            \begin{equation*}
                \left.\begin{array}{l}
                        X \rightsquigarrow \cc{N}(0,1) \\
                        Y\rightsquigarrow\chi^2(n)\\
                        \text{$X$ e $Y$ independientes}
                \end{array}\right\} \Longrightarrow \dfrac{X}{\sqrt{\nicefrac{Y}{n}}} \rightsquigarrow t(n)
            \end{equation*}
            Queremos probar:
            \begin{equation*}
                \dfrac{\overline{X}-\overline{Y}-(\mu_1-\mu_2)}{\sqrt{\dfrac{(n_1-1)S_1^2}{\sigma_1^2} + \dfrac{(n_2-1)S_2^2}{\sigma_2^2}}} \sqrt{\dfrac{n_1+n_2-2}{\dfrac{\sigma_1^2}{n_1}+\dfrac{\sigma_2}{n_2}}} \rightsquigarrow t(n_1+n_2-2)
            \end{equation*}
            
            Tenemos que $\ol{X}\rightsquigarrow \cc{N}(\mu_1, \nicefrac{\sigma_1^2}{n_1})$ y $\ol{Y}\rightsquigarrow \cc{N}(\mu_2, \nicefrac{\sigma_2^2}{n_2})$. Además, tenemos que:
            \begin{equation*}
                -\ol{Y}\rightsquigarrow \cc{N}(-\mu_2, \nicefrac{\sigma_2^2}{n_2})
            \end{equation*}

            Como $\ol{X}$ e $-\ol{Y}$ son funciones de m.a.s. independientes, tenemos que son independientes. Por tanto, por la reproductividad de la distribución normal, tenemos que:
            \begin{equation*}
                \overline{X} - \overline{Y} \rightsquigarrow \cc{N}\left(\mu_1 - \mu_2, \frac{\sigma_1^2}{n_1} + \frac{\sigma_2^2}{n_2}\right)
            \end{equation*}

            Tipificando, tenemos que:
            \begin{equation*}
                \dfrac{\overline{X}-\overline{Y}-(\mu_1-\mu_2)}{\sqrt{\frac{\sigma_1^2}{n_1} + \frac{\sigma_2^2}{n_2}}} \rightsquigarrow \cc{N}(0,1)
            \end{equation*}


            Por otro lado, del Corolario~\ref{cor:muestreo_1normal}, tenemos que:
            \begin{align*}
                \dfrac{(n_1-1)S_1^2}{\sigma_1^2} &\rightsquigarrow \chi^2(n_1-1) \qquad \qquad
                \dfrac{(n_2-1)S_2^2}{\sigma_2^2} &\rightsquigarrow \chi^2(n_2-1) 
            \end{align*}

            Como las m.a.s. son independientes, tenemos que estas dos variables aleatorias son independientes. Por la reproductividad de la distribución $\chi^2$, tenemos que:
            \begin{equation*}
                \dfrac{(n_1-1)S_1^2}{\sigma_1^2}  + \dfrac{(n_2-1)S_2^2}{\sigma_2^2}  \rightsquigarrow \chi^2(n_1+n_2-2)
            \end{equation*}
            Por la extensión del Lema de Fisher tenemos que las dos variables aleatorias que hemos calculado son independientes. Procedemos ahora a aplicar la construcción de la $t$ de Student:
            \begin{equation*}
                \dfrac{\dfrac{\overline{X}-\overline{Y}-(\mu_1-\mu_2)}{\sqrt{\frac{\sigma_1^2}{n_1}+\frac{\sigma_2^2}{n_2}}}}{\sqrt{\dfrac{\dfrac{(n_1-1)S_1^2}{\sigma_1^2} + \dfrac{(n_2-1)S_2^2}{\sigma_2^2}}{n_1+n_2-2}}} \rightsquigarrow t(n_1+n_2-2)
            \end{equation*}
    \end{enumerate}
    \end{proof}
\end{coro}

Gracias a estos resultados, podemos inferir sobre $\mu_1 - \mu_2$ y $\nicefrac{\sigma_1^2}{\sigma_2^2}$ a partir de dos m.a.s. independientes de poblaciones normales. En particular, tenemos los siguientes resultados:
\begin{enumerate}
    \item Si queremos inferir sobre $\mu_1 - \mu_2$ (comparación de medias):
    \begin{enumerate}
        \item Si conocemos las varianzas poblacionales $\sigma_{10}^2$ y $\sigma_{20}^2$, entonces:
        \begin{equation*}
            \dfrac{(\ol{X} - \ol{Y}) - (\mu_1 - \mu_2)}{\sqrt{\dfrac{\sigma_{10}^2}{n_1} + \dfrac{\sigma_{20}^2}{n_2}}} \rightsquigarrow \cc{N}(0,1)
        \end{equation*}

        \item Si las varianzas son desconocidas pero iguales ($\sigma_1^2 = \sigma_2^2 = \sigma^2$), entonces:
        \begin{equation*}
            \dfrac{(\ol{X} - \ol{Y}) - (\mu_1 - \mu_2)}{\sqrt{(n_1 - 1)S_1^2 + (n_2 - 1)S_2^2}} \sqrt{\dfrac{n_1 + n_2 - 2}{\dfrac{1}{n_1} + \dfrac{1}{n_2}}} \rightsquigarrow t(n_1 + n_2 - 2)
        \end{equation*}
    \end{enumerate}

    \item Si queremos inferir sobre $\nicefrac{\sigma_1^2}{\sigma_2^2}$ (comparación de varianzas):
    \begin{enumerate}
        \item Si conocemos las medias $\mu_{10}$ y $\mu_{20}$, entonces:
        \begin{equation*}
            \dfrac{n_1 \sigma_1^2 \sum\limits_{i=1}^{n_2} {(Y_i - \mu_{20})}^{2}}{n_2 \sigma_2^2 \sum\limits_{i=1}^{n_1} {(X_i - \mu_{10})}^{2}} \rightsquigarrow F(n_2, n_1)
        \end{equation*}

        \item Si no conocemos las medias $\mu_1$ y $\mu_2$, entonces:
        \begin{equation*}
            \dfrac{\sigma_1^2 S_2^2}{\sigma_2^2 S_1^2} \rightsquigarrow F(n_2 - 1, n_1 - 1)
        \end{equation*}
    \end{enumerate}
\end{enumerate}

    \section{Suficiencia y completitud}
\begin{ejemplo}
    Sean $X_1, X_2$ variables aleatorias independientes tal que $X_1,X_2\rightsquigarrow \cc{P}(\lm)$. Probar que $X_1+X_2$ es suficiente para $\lm$ pero $X_1+2X_2$ no lo es.
    \begin{proof}
        Notamos por $P_\lm$ a la función masa de probabilidad, que dependerá de $\lm$. Queremos probar que, condicionada a un valor de $X_1+X_2$, no depende de $\lm$. Fijado $k\in\bb{N}$, tenemos que:
        \begin{align*}
            P_\lm[X_1=x_1, X_2=x_2 \mid X_1+X_2=k] &= \dfrac{P_\lm[X_1=x_1, X_2=x_2, X_1+X_2=k]}{P_\lm[X_1+X_2=k]}
            =\\&= \begin{cases}
                \dfrac{P_\lm[X_1=x_1, X_2=x_2]}{P_\lm[X_1+X_2=k]} & \text{si } x_1+x_2=k \\
                0 & \text{si } x_1+x_2\neq k
            \end{cases}
        \end{align*}

        Puesto que $X_1$ y $X_2$ son independientes, para $x_1+x_2=k$ tenemos que:
        \begin{align*}
            P_\lm[X_1=x_1, X_2=x_2 \mid X_1+X_2=k] &= \dfrac{P_\lm[X_1=x_1] P_\lm[X_2=x_2]}{P_\lm[X_1+X_2=k]}
            \AstIg \dfrac{e^{-\lm}\frac{\lm^{x_1}}{x_1!} \cdot e^{-\lm}\frac{\lm^{x_2}}{x_2!}}{e^{-2\lm}\frac{(2\lm)^k}{k!}}
            = \\&= \dfrac{k!}{2^k x_1! x_2!}
        \end{align*}
        donde en $(\ast)$ hemos usado la reproductividad de la distribución Poisson, que nos dice que $X_1+X_2\rightsquigarrow \cc{P}(2\lm)$. Notemos que el resultado no depende de $\lm$, por lo que $X_1+X_2$ es suficiente para $\lm$.\\

        Por otro lado, queremos ver que $X_1+2X_2$ no es suficiente para $\lm$. Para ello, no es necesario probarlo para todos los valores posibles, sino que basta con encontrar un contraejemplo. Tomemos $k=2$ y la muestra $(x_1,x_2)=(0,1)$, de forma que $x_1+2x_2=2$. Tenemos que:
        \begin{align*}
            P_\lm[X_1=0, X_2=1 \mid X_1+2X_2=2] &= \dfrac{P_\lm[X_1=0, X_2=1, X_1+2X_2=2]}{P_\lm[X_1+2X_2=2]} =\\&= \dfrac{P_\lm[X_1=0, X_2=1]}{P_\lm[X_1+2X_2=2]}
        \end{align*}

        Por el Teorema de Cambio de Variable, tenemos que:
        \begin{align*}
            P_\lm[X_1+2X_2=2]
            &= \sum\limits_{\substack{x_1,x_2\in\bb{N}\\x_1+2x_2=2}} P_\lm[X_1=x_1, X_2=x_2]
            =\\&= P_\lm[X_1=0, X_2=1] + P_\lm[X_1=2, X_2=0]
        \end{align*}

        Por tanto, tenemos que:
        \begin{align*}
            P_\lm[X_1=0, X_2=1 \mid X_1+2X_2=2] &= \dfrac{P_\lm[X_1=0]P_\lm[X_2=1]}{P_\lm[X_1=0]P_\lm[X_2=1] + P_\lm[X_1=2]P_\lm[X_2=0]} =\\&= \dfrac{e^{-\lm}\frac{\lm^0}{0!} \cdot e^{-\lm}\frac{\lm^1}{1!}}{e^{-\lm}\frac{\lm^0}{0!} \cdot e^{-\lm}\frac{\lm^1}{1!} + e^{-\lm}\frac{\lm^2}{2!} \cdot e^{-\lm}\frac{\lm^0}{0!}} 
            = \dfrac{\lm}{\lm + \frac{\lm^2}{2}} = \dfrac{2}{2+\lm}
        \end{align*}
        Notemos que el resultado depende de $\lm$, por lo que $X_1+2X_2$ no es suficiente para $\lm$.
    \end{proof}
\end{ejemplo}

\begin{ejemplo}
    En $n$ lanzamientos de una moneda, el número de caras es suficiente para la probabilidad de cara (notada $p$).
    \begin{proof}
        Sea $X$ la v.a. que representa el número de caras en un lanzamiento de la moneda, de forma que $X\rightsquigarrow \cc{B}(1,p)$. Sea $(X_1, \ldots, X_n)$ una m.a.s. de $n$ lanzamientos de la moneda, y sea $T = \sum\limits_{i=1}^n X_i$ el número total de caras en los $n$ lanzamientos. Queremos ver que $Y$ es suficiente para $p$. Para ello, fijado un valor $k\in\{0, \ldots, n\}$, tenemos que:
        \begin{align*}
            P_p[X_1=x_1, \ldots, X_n=x_n \mid T=k] &= \dfrac{P_p[X_1=x_1, \ldots, X_n=x_n, T=k]}{P_p[T=k]} =\\&= \begin{cases}
                \dfrac{P_p[X_1=x_1, \ldots, X_n=x_n]}{P_p[T=k]} & \text{si } x_1+\ldots+x_n=k \\
                0 & \text{si } x_1+\ldots+x_n\neq k
            \end{cases}
        \end{align*}

        Por otro lado, como $X_1, \ldots, X_n$ son independientes, en el caso de que se tenga que $x_1+\ldots+x_n=k$, tenemos que:
        \begin{align*}
            P_p[X_1=x_1, \ldots, X_n=x_n \mid T=k] &= \dfrac{P_p[X_1=x_1]\cdots P_p[X_n=x_n]}{P_p[T=k]}
            \AstIg\\&\AstIg \dfrac{\prod\limits_{i=1}^n p^{x_i}(1-p)^{1-x_i}}{\binom{n}{k}p^k(1-p)^{n-k}}
            = \dfrac{p^{\sum\limits_{i=1}^n x_i}(1-p)^{\sum\limits_{i=1}^n 1-x_i}}{\binom{n}{k}p^k(1-p)^{n-k}}
            =\\&= \dfrac{p^k(1-p)^{n-k}}{\binom{n}{k}p^k(1-p)^{n-k}} = \dfrac{1}{\binom{n}{k}}
        \end{align*}
        donde en $(\ast)$ hemos usado que $T\rightsquigarrow \cc{B}(n,p)$. Notemos que el resultado no depende de $p$, por lo que $T$ es suficiente para $p$.
    \end{proof}
\end{ejemplo}

\begin{ejemplo}
    Se supone que se ha lanzado una moneda 100 veces, y se sabe que se han obtenido 60 caras pero se han perdido los datos originales de si salió cara o cruz en cada tirada. Reconstruir la muestra.\\

    Sea $X$ la v.a. que representa el número de caras en un lanzamiento de la moneda, de forma que $X\rightsquigarrow \cc{B}(1,p)$. Sea $(X_1, \ldots, X_{100})$ una m.a.s. de 100 lanzamientos de la moneda, y sea $T = \sum\limits_{i=1}^{100} X_i$ el número total de caras en los 100 lanzamientos, que sabemos que es suficiente para $p$, luego la información que nos da la muestra original se encuentra en $T$. Por tanto, podemos reconstruir la muestra original muestreando de la distribución condicional de $(X_1, \ldots, X_{100})$ dado que $T=60$, que en el ejemplo anterior hemos visto que es:
    \begin{equation*}
        P[X_1=x_1, \ldots, X_{100}=x_{100} \mid T=60] = \begin{cases}
            \dfrac{1}{\binom{100}{60}} & \text{si } x_1+\ldots+x_{100}=60 \\
            0 & \text{si } x_1+\ldots+x_{100}\neq 60
        \end{cases}
    \end{equation*}

    Por tanto, si $(X_1^*, \ldots, X_{100}^*)$ es una m.a.s. de la distribución condicionada al valor de $T=60$, tenemos que $(X_1^*, \ldots, X_{100}^*)$ sigue una distribución uniforme discreta sobre el conjunto de todas las muestras posibles de 100 lanzamientos de la moneda que contienen exactamente 60 caras; es decir, $(X_1^*, \ldots, X_{100}^*)\rightsquigarrow \cc{U}(\cc{S})$, donde:
    \begin{equation*}
        \cc{S} = \left\{(x_1, \ldots, x_{100})\in\{0,1\}^{100} : x_1+\ldots+x_{100}=60\right\}
    \end{equation*}
\end{ejemplo}

\begin{teo}[de Factorización de Neyman-Fisher (Discreto)]
    Sea una v. a. $X\rightsquigarrow F\in \{F_\theta : \theta\in\Theta\}$, y sea $(X_1, \ldots, X_n)$ una m.a.s. de $X$. Sea $f_\theta$ la función masa de probabilidad de $F_\theta$. Sea $T = T(X_1, \ldots, X_n)$ un estadístico. Entonces, $T$ es suficiente para $\theta$ si y solo si existen funciones $g_\theta$ y $h$ tales que:
    \begin{equation*}
        f_\theta(x_1, \ldots, x_n) = h(x_1, \ldots, x_n) g_\theta(T(x_1, \ldots, x_n)) \qquad \forall (x_1, \ldots, x_n)\in\cc{X}^n,\ \forall \theta\in\Theta
    \end{equation*}
    donde $h$ no depende de $\theta$.
    \begin{proof} Demostramos cada una de las implicaciones:
    \begin{description}
        \item[$\Longrightarrow$)] Supongamos que $T$ es suficiente para $\theta$. Entonces, para todo $(x_1, \ldots, x_n)\in\cc{X}^n$, $\theta\in\Theta$, y para cualquier $t\in T(\cc{X}^n)$ fijo, tenemos que:
        \begin{align*}
            P[X_1=x_1, \ldots, X_n=x_n\mid T=t] &= \dfrac{P_\theta[X_1=x_1, \ldots, X_n=x_n, T=t]}{P_\theta[T=t]}
        \end{align*}

        Evaluando en $t=T(x_1, \ldots, x_n)$, tenemos que:
        \begin{align*}
            P[X_1=x_1, \ldots, X_n=x_n\mid T=t] &= \dfrac{P_\theta[X_1=x_1, \ldots, X_n=x_n]}{P_\theta[T=t]}
        \end{align*}

        Despejando, tenemos que:
        \begin{align*}
            P_\theta[X_1=x_1, \ldots, X_n=x_n] &= P_\theta[T=t] P[X_1=x_1, \ldots, X_n=x_n\mid T=t]
        \end{align*}

        Por tanto, identificando
        $$g_\theta(t) = P_\theta[T=t]\qquad \qquad h(x_1, \ldots, x_n) = P[X_1=x_1, \ldots, X_n=x_n\mid T=T(x_1, \ldots, x_n)]$$
        tenemos que:
        \begin{align*}
            P_\theta[X_1=x_1, \ldots, X_n=x_n] = h(x_1, \ldots, x_n) g_\theta(T(x_1, \ldots, x_n))
        \end{align*}

        \item[$\Longleftarrow$)] Supongamos que existen funciones $g_\theta$ y $h$ tales que, para todo $(x_1, \ldots, x_n)\in\cc{X}^n$ y $\theta\in\Theta$, se cumple que:
        \begin{align*}
            P_\theta[X_1=x_1, \ldots, X_n=x_n] = h(x_1, \ldots, x_n) g_\theta(T(x_1, \ldots, x_n))
        \end{align*}

        Entonces, para cualquier $t\in T(\cc{X}^n)$, tenemos que:
        \begin{align*}
            P_\theta[X_1=x_1, \ldots, X_n=x_n\mid T=t] &= \dfrac{P_\theta[X_1=x_1, \ldots, X_n=x_n, T=t]}{P_\theta[T=t]} =\\&= \begin{cases}
                \dfrac{P_\theta[X_1=x_1, \ldots, X_n=x_n]}{P_\theta[T=t]} & \text{si } T(x_1, \ldots, x_n)=t \\
                0 & \text{si } T(x_1, \ldots, x_n)\neq t
            \end{cases}
        \end{align*}

        En el caso de que $T(x_1, \ldots, x_n)=t$, tenemos que:
        \begin{align*}
            P_\theta[X_1=x_1, \ldots, X_n=x_n\mid T=t] &= \dfrac{h(x_1, \ldots, x_n) g_\theta(T(x_1, \ldots, x_n))}{\sum\limits_{\substack{(x_1, \ldots, x_n)\in\cc{X}^n \\ T(x_1, \ldots, x_n)=t}} P_\theta[X_1=x_1, \ldots, X_n=x_n]} =\\&= \dfrac{h(x_1, \ldots, x_n) g_\theta(t)}{\sum\limits_{\substack{(x_1, \ldots, x_n)\in\cc{X}^n \\ T(x_1, \ldots, x_n)=t}} h(x_1, \ldots, x_n) g_\theta(t)} =\\&= \dfrac{h(x_1, \ldots, x_n)}{\sum\limits_{\substack{(x_1, \ldots, x_n)\in\cc{X}^n \\ T(x_1, \ldots, x_n)=t}} h(x_1, \ldots, x_n)}
        \end{align*}
        Notemos que el resultado no depende de $\theta$, por lo que $T$ es suficiente para $\theta$.\qedhere
    \end{description}
    \end{proof}
\end{teo}

\begin{ejemplo}
    Sea $X$ una v.a. tal que $X\rightsquigarrow F\in \{\cc{N}(\mu, \sigma^2) : \mu\in\bb{R}, \sigma^2>0\}$. Estudiar la suficiencia en los casos de:
    \begin{enumerate}
        \item $\sigma_0^2$ conocida.
        
        Consideramos una m.a.s. $(X_1, \ldots, X_n)$ de $X$. Sea $T=T(X_1, \ldots, X_n)$ un estadístico. Por el Teorema de Factorización de Neyman-Fisher, $T$ es suficiente para $\mu$ si y solo si existen funciones $g_\mu$ y $h$ tales que:
        \begin{equation*}
            f_\mu(x_1, \ldots, x_n) = h(x_1, \ldots, x_n) g_\mu(T(x_1, \ldots, x_n)) \qquad \forall (x_1, \ldots, x_n)\in\bb{R}^n,\ \forall \mu\in\bb{R}
        \end{equation*}

        Tenemos que:
        \begin{align*}
            f_\mu(x_1, \ldots, x_n) 
            &\stackrel{\text{i.i.d.}}{=}
            \prod_{i=1}^n f_\mu(x_i)
            = \prod_{i=1}^n \dfrac{1}{\sqrt{2\pi\sigma_0^2}} \exp\left(-\dfrac{(x_i-\mu)^2}{2\sigma_0^2}\right)
            =\\&= \dfrac{1}{(2\pi\sigma_0^2)^{\nicefrac{n}{2}}} \exp\left(-\dfrac{\sum\limits_{i=1}^n (x_i-\mu)^2}{2\sigma_0^2}\right)
            =\\&= \dfrac{1}{(2\pi\sigma_0^2)^{\nicefrac{n}{2}}} \exp\left(-\dfrac{\sum\limits_{i=1}^n (x_i^2 - 2\mu x_i + \mu^2)}{2\sigma_0^2}\right)
            =\\&= \dfrac{1}{(2\pi\sigma_0^2)^{\nicefrac{n}{2}}} \exp\left(-\dfrac{\sum\limits_{i=1}^n x_i^2 - 2\mu \sum\limits_{i=1}^n x_i + n\mu^2}{2\sigma_0^2}\right)
            =\\&= \dfrac{1}{(2\pi\sigma_0^2)^{\nicefrac{n}{2}}} \exp\left(-\dfrac{\sum\limits_{i=1}^n x_i^2}{2\sigma_0^2}\right) \exp\left(-\dfrac{- 2\mu \sum\limits_{i=1}^n x_i + n\mu^2}{2\sigma_0^2}\right)
        \end{align*}

        Por tanto, identificando términos, tenemos que:
        \begin{align*}
            h(x_1, \ldots, x_n) &= \dfrac{1}{(2\pi\sigma_0^2)^{\nicefrac{n}{2}}} \exp\left(-\dfrac{\sum\limits_{i=1}^n x_i^2}{2\sigma_0^2}\right) \\
            g_\mu(T(x_1, \ldots, x_n)) &= \exp\left(-\dfrac{- 2\mu T(x_1, \ldots, x_n) + n\mu^2}{2\sigma_0^2}\right)
        \end{align*}

        Por tanto, por el Teorema de Factorización de Neyman-Fisher, tenemos que $T(X_1, \ldots, X_n) = \sum\limits_{i=1}^n X_i$ es suficiente para $\mu$. En particular, como $\ol{X}= \nicefrac{1}{n} \cdot T$, tenemos que $\ol{X}$ es suficiente para $\mu$.
        \item $\mu_0$ conocida.
        
        En este caso, identificando términos, tenemos que:
        \begin{align*}
            h(x_1, \ldots, x_n) &= 1 \\
            g_{\sigma^2}(T(x_1, \ldots, x_n)) &= \dfrac{1}{(2\pi\sigma^2)^{\nicefrac{n}{2}}} \exp\left(-\dfrac{T(x_1, \ldots, x_n)}{2\sigma^2}\right)
        \end{align*}

        Por tanto, por el Teorema de Factorización de Neyman-Fisher, tenemos que $T(X_1, \ldots, X_n) = \sum\limits_{i=1}^n (X_i-\mu_0)^2$ es suficiente para $\sigma^2$. Como $\mu_0=E[X]=E[\ol{X}]$, tenemos que $S^2 = \nicefrac{1}{n-1} \cdot T$ es suficiente para $\sigma^2$.
        \item $\mu$ y $\sigma^2$ desconocidas.
        
        En este caso, identificando términos, tenemos de forma directa que:
        \begin{align*}
            h(x_1, \ldots, x_n) &= 1
        \end{align*}

        Puesto que el parámetro es bidimensional, el estadístico suficiente debe ser, al menos, bidimensional. Por tanto, podemos tomar como estadístico suficiente $(T_1, T_2)$, donde:
        \begin{equation*}
            T_1(X_1, \ldots, X_n) = \sum\limits_{i=1}^n X_i \qquad \qquad T_2(X_1, \ldots, X_n) = \sum\limits_{i=1}^n (X_i)^2
        \end{equation*}

        De esta forma, tenemos que:
        \begin{multline*}
            g_{\mu, \sigma^2}(T_1(x_1, \ldots, x_n), T_2(x_1, \ldots, x_n)) =\\= \dfrac{1}{(2\pi\sigma^2)^{\nicefrac{n}{2}}} \exp\left(-\dfrac{T_2(x_1, \ldots, x_n) - 2\mu T_1(x_1, \ldots, x_n) + n\mu^2}{2\sigma^2}\right)
        \end{multline*}
        
        De esta forma, por el Teorema de Factorización de Neyman-Fisher, tenemos que $(T_1, T_2)$ es suficiente para $(\mu, \sigma^2)$.
    \end{enumerate}
\end{ejemplo}



\begin{ejemplo}
    Sea $(X_1, \ldots, X_n)$ una m.a.s de la v.a. $X \rightsquigarrow \cc{B}(1, p)$, con $p \in \left]0, 1\right[$. Encontrar un estadístico suficiente y completo asociado a la muestra.\\

    Buscamos un estadístico suficiente. Por el Teorema de Factorización de Neyman-Fisher, $T$ es suficiente para $p$ si y solo si existen funciones $g_p$ y $h$ tales que:
    \begin{equation*}
        P_p[X_1=x_1, \ldots, X_n=x_n] = h(x_1, \ldots, x_n) g_p(T(x_1, \ldots, x_n)) \qquad \forall (x_1, \ldots, x_n)\in\{0, 1\}^n,\ \forall p\in\left]0, 1\right[
    \end{equation*}

    Tenemos que:
    \begin{align*}
        P_p[X_1=x_1, \ldots, X_n=x_n] &\stackrel{\text{i.i.d.}}{=} \prod_{i=1}^n P_p[X_i=x_i] = \prod_{i=1}^n p^{x_i}(1-p)^{1-x_i} = p^{\sum\limits_{i=1}^n x_i}(1-p)^{\sum\limits_{i=1}^n 1-x_i} =\\&= p^{\sum\limits_{i=1}^n x_i}(1-p)^{n-\sum\limits_{i=1}^n x_i}
    \end{align*}

    Por tanto, identificando términos, tenemos que:
    \begin{align*}
        h(x_1, \ldots, x_n) &= 1 \\
        g_p(T(x_1, \ldots, x_n)) &= p^{T(x_1, \ldots, x_n)}(1-p)^{n-T(x_1, \ldots, x_n)}
    \end{align*}

    Por tanto, por el Teorema de Factorización de Neyman-Fisher, tenemos que $T(X_1, \ldots, X_n) = \sum\limits_{i=1}^n X_i$ es suficiente para $p$. Veamos ahora si $T$ es completo. Sea $g$ una función medible tal que $E_p[g(T)] = 0$ para todo $p\in\left]0, 1\right[$. Probemos que $P_p[g(T)=0] = 1$ para todo $p\in\left]0, 1\right[$. Por reproductividad, se tiene $T \rightsquigarrow \cc{B}(n, p)$. Por tanto, para todo $p\in\left]0, 1\right[$, tenemos que:
    \begin{align*}
        E_p[g(T)] &= \sum\limits_{k=0}^n g(k) P_p[T=k] = \sum\limits_{k=0}^n g(k) \binom{n}{k} p^k (1-p)^{n-k}
        =\\&= (1-p)^n \sum\limits_{k=0}^n g(k) \binom{n}{k} \left(\dfrac{p}{1-p}\right)^k = 0
    \end{align*}

    Sea $q = \nicefrac{p}{1-p}$. Entonces, para todo $q\in\left]0, +\infty\right[$, tenemos que:
    \begin{align*}
        \sum\limits_{k=0}^n g(k) \binom{n}{k} q^k = 0
    \end{align*}

    Notemos que el polinomio $\sum\limits_{k=0}^n g(k) \binom{n}{k} q^k$ tiene grado a lo sumo $n$. Por tanto, como tiene infinitas raíces, debe ser el polinomio nulo. Por tanto, tenemos que:
    \begin{align*}
        g(k) \binom{n}{k} = 0 \qquad \forall k\in\{0, \ldots, n\}
    \end{align*}
    Por tanto, tenemos que $g(k) = 0$ para todo $k\in\{0, \ldots, n\}$. Por tanto, hemos visto que:
    \begin{equation*}
        \{0, \ldots, n\} \subset \{k\in\bb{N} : g(k) = 0\}
    \end{equation*}
    Por tanto:
    \begin{equation*}
        1=P_p[T\in \{0, \ldots, n\}] \leq P_p[g(T)=0]\leq 1
    \end{equation*}
    Por tanto, $P_p[g(T)=0] = 1$ para todo $p\in\left]0, 1\right[$; luego $T$ es completo para $p$.
\end{ejemplo}

\begin{ejemplo}
    Sea $(X_1, \ldots, X_n)$ una m.a.s. de la v.a. $X \rightsquigarrow 
    \cc{U}(0, \theta)$, con $\theta > 0$. Encontrar un estadístico suficiente y completo asociado a la muestra.\\
    
    
    Buscamos un estadístico suficiente. Por el Teorema de Factorización de Neyman-Fisher, $T$ es suficiente para $\theta$ si y solo si existen funciones $g_\theta$ y $h$ tales que:
    \begin{equation*}
        f_\theta(x_1, \ldots, x_n) = h(x_1, \ldots, x_n) g_\theta(T(x_1, \ldots, x_n)) \qquad \forall (x_1, \ldots, x_n)\in\bb{R}^n,\ \forall \theta > 0
    \end{equation*}
    Tenemos que:
    \begin{align*}
        f_\theta(x_1, \ldots, x_n) &\stackrel{\text{i.i.d.}}{=} \prod_{i=1}^n f_\theta(x_i) = \prod_{i=1}^n \dfrac{1}{\theta} = \dfrac{1}{\theta^n}\qquad \text{si } (x_1, \ldots, x_n)\in [0, \theta]^n
    \end{align*}

    Notemos que la función de densidad no depende de $(x_1, \ldots, x_n)$, salvo por la condición de que $(x_1, \ldots, x_n)\in [0, \theta]^n$. Por tanto, hemos de emplear la función indicadora para expresar la dependencia de $(x_1, \ldots, x_n)$. Notemos que:
    \begin{align*}
        f_\theta(x_1, \ldots, x_n) & = \dfrac{1}{\theta^n} \ind_{[0,\infty]}(x_{(1)}) \ind_{\left]-\infty, 0\right]}(x_{(n)}-\theta)\qquad \forall (x_1, \ldots, x_n)\in\bb{R}^n,\ \forall \theta > 0
    \end{align*}
    Por tanto, identificando términos, tenemos que:
    \begin{align*}
        h(x_1, \ldots, x_n) &= \ind_{[0,\infty]}(x_{(1)}) \\
        g_\theta(T(x_1, \ldots, x_n)) &= \dfrac{1}{\theta^n} \ind_{\left]-\infty, 0\right]}(T(x_1, \ldots, x_n)-\theta)
    \end{align*} 

    Por tanto, por el Teorema de Factorización de Neyman-Fisher, tenemos que $T(X_1, \ldots, X_n) = X_{(n)}$ es suficiente para $\theta$. Para estudiar su completitud, hemos de calcular en primer lugar la función de densidad de $T$. Tenemos que:
    \begin{align*}
        F_{X_{(n)}}(x) &= P[X_{(n)} \leq x] = P[X_1\leq x, \ldots, X_n\leq x] = \prod_{i=1}^n P[X_i\leq x] = \prod_{i=1}^n F_X(x) = F_X(x)^n
    \end{align*}
    Por tanto, tenemos que:
    \begin{align*}
        f_{X_{(n)}}(x) &= \dfrac{d}{dx} F_{X_{(n)}}(x) = n F_X(x)^{n-1} f_X(x) = n \left(\dfrac{x}{\theta}\right)^{n-1} \dfrac{1}{\theta} = \dfrac{n x^{n-1}}{\theta^n} \qquad \text{si } x\in [0, \theta]
    \end{align*}

    Sea $g$ una función medible tal que $E_\theta[g(T)] = 0$ para todo $\theta > 0$. Probemos que $P_\theta[g(T)=0] = 1$ para todo $\theta > 0$. Tenemos que:
    \begin{align*}
        E_\theta[g(T)] &= \int_0^\theta g(x) f_{X_{(n)}}(x) dx = \int_0^\theta g(x) \dfrac{n x^{n-1}}{\theta^n} dx = \dfrac{n}{\theta^n} \int_0^\theta g(x) x^{n-1} dx = 0\qquad \forall \theta > 0
    \end{align*}

    Por tanto, como el primer término no se puede anular, se tiene:
    \begin{align*}
        \int_0^\theta g(x) x^{n-1} dx = 0\qquad \forall \theta > 0
    \end{align*}
    % // TODO: Por qué es continua??
    Notemos que la función $g(x) x^{n-1}$ es continua, y por tanto integrable. Derivando con respecto a $\theta$, tenemos que:
    \begin{align*}
        g(\theta) \theta^{n-1} = 0\qquad \forall \theta > 0
    \end{align*}
    Esto implica que $g(\theta) = 0$ para todo $\theta > 0$, luego:
    \begin{equation*}
        \bb{R}^+ \subset \{x\in\bb{R} : g(x) = 0\}
    \end{equation*}


    Por tanto:
    \begin{equation*}
        1 = P_\theta[T\in \bb{R}^+ ] \leq P_\theta[g(T)=0]\leq 1
    \end{equation*}
    Por tanto, $P_\theta[g(T)=0] = 1$ para todo $\theta > 0$; luego $T$ es completo para $\theta$.
\end{ejemplo}

\begin{ejemplo}
    Demostrar que las siguientes familias de distribuciones son una familia exponencial uniparamétrica:
    \begin{enumerate}
        \item $\cc{P}(\lm)$, con $\lm > 0$
        
        En primer lugar, tenemos que $\lm\in \Theta = \bb{R}^+$. Además, tenemos que:
        \begin{align*}
            \chi = \bb{N} \qquad \qquad \forall \lm\in\Theta
        \end{align*}

        Por último, para cada $\lm\in\Theta$ y $x\in\chi$, tenemos que:
        \begin{align*}
            P_\lm[X=x] = e^{-\lm}\dfrac{\lm^x}{x!}
            = \exp\left(-\lm + \ln\left(\dfrac{\lm^x}{x!}\right)\right)
            = \exp\left(\ln(\lm)x - \lm - \ln(x!)\right)
        \end{align*}

        Por tanto, identificando términos, tenemos que:
        \begin{align*}
            Q(\lm) &= \ln(\lm) \\
            T(x) &= x \\
            D(\lm) &=- \lm \\
            S(x) &=- \ln(x!)
        \end{align*}

        Por tanto, $\cc{P}(\lm)$ es una familia exponencial uniparamétrica.
        \item $\cc{B}(k_0, p)$, con $p\in\left]0, 1\right[$ y $k_0\in\bb{N}$ fijo.
        
        En primer lugar, tenemos que $p\in \Theta = \left]0, 1\right[$. Además, tenemos que:
        \begin{align*}
            \chi = \{0, 1, \ldots, k_0\}\qquad \qquad \forall p\in\Theta
        \end{align*}

        Por último, tenemos que:
        \begin{align*}
            P_p[X=x] &= \binom{k_0}{x} p^x (1-p)^{k_0-x}
            = \exp\left(\ln\left(\binom{k_0}{x} p^x (1-p)^{k_0-x}\right)\right)
            =\\&= \exp\left(\ln\left(\binom{k_0}{x}\right) + x\ln(p) + (k_0-x)\ln(1-p)\right)
            =\\&= \exp\left(\ln\left(\binom{k_0}{x}\right) + x\ln\left(\dfrac{p}{1-p}\right) + k_0\ln(1-p)\right)
            \qquad \forall x\in\chi,\ \forall p\in\Theta
        \end{align*}

        Por tanto, identificando términos, tenemos que:
        \begin{align*}
            Q(p) &= \ln\left(\dfrac{p}{1-p}\right) \\
            T(x) &= x \\
            D(p) &= k_0\ln(1-p) \\
            S(x) &= \ln\left(\binom{k_0}{x}\right)
        \end{align*}
    \end{enumerate}
    
\end{ejemplo}


\begin{ejemplo}
    La $\cc{N}(\mu, \sigma^2)$ es una familia exponencial uniparamétrica en cada parámetro y bi-paramétrica en los dos parámetros.\\

    Veamos primero que $\cc{N}(\mu, \sigma^2)$ es una familia exponencial uniparamétrica en $\mu$ para $\sigma^2$ conocida. En este caso, tenemos que:
    \begin{align*}
        \Theta &= \bb{R} \\
        \chi &= \bb{R} \\
        f_\mu(x) &= \dfrac{1}{\sqrt{2\pi\sigma^2}} \exp\left(-\dfrac{(x-\mu)^2}{2\sigma^2}\right) = \exp\left(-\dfrac{(x-\mu)^2}{2\sigma^2} - \dfrac{1}{2}\ln(2\pi\sigma^2)\right)
        =\\&= \exp\left(-\dfrac{x^2 - 2\mu x + \mu^2}{2\sigma^2} - \dfrac{1}{2}\ln(2\pi\sigma^2)\right)
        = \exp\left(\dfrac{\mu x}{\sigma^2} - \dfrac{\mu^2}{2\sigma^2} - \dfrac{x^2}{2\sigma^2} - \dfrac{1}{2}\ln(2\pi\sigma^2)\right)
    \end{align*}

    Por tanto, identificando términos, tenemos que:
    \begin{align*}
        Q(\mu) &= \dfrac{\mu}{\sigma^2} \\
        T(x) &= x \\
        D(\mu) &= -\dfrac{\mu^2}{2\sigma^2} \\
        S(x) &= -\dfrac{x^2}{2\sigma^2} - \dfrac{1}{2}\ln(2\pi\sigma^2)
    \end{align*}

    Vemos por tanto que $\cc{N}(\mu, \sigma^2)$ es una familia exponencial uniparamétrica en $\mu$ para $\sigma^2$ conocida. Veamos ahora que $\cc{N}(\mu, \sigma^2)$ es una familia exponencial uniparamétrica en $\sigma^2$ para $\mu$ conocida. En este caso, tenemos que:
    \begin{align*}
        \Theta &= \bb{R}^+ \\
        \chi &= \bb{R} \\
        f_{\sigma^2}(x) &= \exp\left(-\dfrac{(x-\mu)^2}{2\sigma^2} - \dfrac{1}{2}\ln(2\pi\sigma^2)\right)
    \end{align*}

    Por tanto, identificando términos, tenemos que:
    \begin{align*}
        Q(\sigma^2) &= -\dfrac{1}{2\sigma^2} \\
        T(x) &= (x-\mu)^2 \\
        D(\sigma^2) &= -\dfrac{1}{2}\ln(2\pi\sigma^2) \\
        S(x) &= 0
    \end{align*}

    Por tanto, $\cc{N}(\mu, \sigma^2)$ es una familia exponencial uniparamétrica en $\sigma^2$ para $\mu$ conocida. Veamos por último que $\cc{N}(\mu, \sigma^2)$ es una familia exponencial bi-paramétrica en $(\mu, \sigma^2)$. En este caso, tenemos que:
    \begin{align*}
        \Theta &= \bb{R}\times\bb{R}^+ \\
        \chi &= \bb{R} \\
        f_{\mu, \sigma^2}(x) &= \exp\left(\dfrac{\mu x}{\sigma^2} - \dfrac{\mu^2}{2\sigma^2} - \dfrac{x^2}{2\sigma^2} - \dfrac{1}{2}\ln(2\pi\sigma^2)\right)
    \end{align*}

    Por tanto, identificando términos, tenemos que:
    \begin{align*}
        Q_1(\mu, \sigma^2) &= \dfrac{\mu}{\sigma^2} \\
        Q_2(\mu, \sigma^2) &= -\dfrac{1}{2\sigma^2} \\
        T_1(x) &= x \\
        T_2(x) &= x^2 \\
        D(\mu, \sigma^2) &= -\dfrac{\mu^2}{2\sigma^2} - \dfrac{1}{2}\ln(2\pi\sigma^2) \\
        S(x) &= 0
    \end{align*}

    Por tanto, $\cc{N}(\mu, \sigma^2)$ es una familia exponencial bi-paramétrica en $(\mu, \sigma^2)$.
\end{ejemplo}

\begin{teo}
    Si se tiene una v.a. $X \rightsquigarrow \{F_\theta, \theta\in\Theta\}$ con familia de funciones asociadas, $\{f_\theta(x), \theta\in\Theta\}$, siendo f.d.d. o f.m.p. según el caso, donde la familia de distribuciones es exponencial k-paramétrica, entonces la familia de distribuciones asociadas a $(X_1, \ldots, X_n)$ (que es una m.a.s. de $X$) es también exponencial k-paramétrica:
    \begin{equation*}
        f_\theta^n(x_1, \ldots, x_n) = \exp\left[\sum\limits_{h=1}^k Q_h(\theta) \sum\limits_{i=1}^n T_h(x_i) + \sum\limits_{i=1}^n S(x_i) + n D(\theta)\right], \qquad (x_1, \ldots, x_n)\in\chi^n
    \end{equation*}
    y se tiene:
    \begin{enumerate}
        \item El estadístico $\left(\sum\limits_{i=1}^n T_1(X_i), \ldots, \sum\limits_{i=1}^n T_k(X_i)\right)$ es suficiente para $\theta$.
        \item Si $k\leq n$ y el conjunto imagen de la función $Q(\theta) = (Q_1(\theta), \ldots, Q_k(\theta))$ contiene a un abierto de $\bb{R}^k$, el estadístico $\left(\sum\limits_{i=1}^n T_1(X_i), \ldots, \sum\limits_{i=1}^n T_k(X_i)\right)$ es también completo.
    \end{enumerate}
    \begin{proof}
        Comprobemos que $(X_1, \ldots, X_n)$ es una familia exponencial k-paramétrica. Por ser la familia de distribuciones asociadas a $X$ una familia exponencial k-paramétrica, tenemos que $\Theta$ es un intervalo de $\bb{R}^k$. Además, tenemos que el espacio muestral de $(X_1, \ldots, X_n)$ es independiente de $\theta$, y es $\chi^n$. Por último, tenemos que:
        \begin{align*}
            f_\theta^n(x_1, \ldots, x_n) &= \prod_{i=1}^n f_\theta(x_i) = \prod_{i=1}^n \exp\left[\sum\limits_{h=1}^k Q_h(\theta) T_h(x_i) + S(x_i) + D(\theta)\right]
            =\\&= \exp\left[\sum\limits_{i=1}^n \sum\limits_{h=1}^k Q_h(\theta) T_h(x_i) + \sum\limits_{i=1}^n S(x_i) + \sum\limits_{i=1}^n D(\theta)\right]
            =\\&= \exp\left[\sum\limits_{h=1}^k Q_h(\theta) \sum\limits_{i=1}^n T_h(x_i) + \sum\limits_{i=1}^n S(x_i) + n D(\theta)\right]
        \end{align*}

        Por tanto, identificando términos, tenemos que:
        \begin{align*}
            Q_h^n(\theta) &= Q_h(\theta) \\
            T_h^n(x_1, \ldots, x_n) &= \sum\limits_{i=1}^n T_h(x_i) \\
            S^n(x_1, \ldots, x_n) &= \sum\limits_{i=1}^n S(x_i) \\
            D^n(\theta) &= n D(\theta)
        \end{align*}

        Por tanto, $(X_1, \ldots, X_n)$ es una familia exponencial k-paramétrica. Veamos ahora que $\left(\sum\limits_{i=1}^n T_1(X_i), \ldots, \sum\limits_{i=1}^n T_k(X_i)\right)$ es suficiente para $\theta$.
        \begin{align*}
            f_\theta^n(x_1, \ldots, x_n) &= \exp\left[\sum\limits_{h=1}^k Q_h(\theta) \sum\limits_{i=1}^n T_h(x_i) + \sum\limits_{i=1}^n S(x_i) + n D(\theta)\right] =\\&= \underbrace{\exp\left[\sum\limits_{i=1}^n S(x_i)\right]}_{h(x_1, \ldots, x_n)} \underbrace{\exp\left[\sum\limits_{h=1}^k Q_h(\theta) \sum\limits_{i=1}^n T_h(x_i) + n D(\theta)\right]}_{g_\theta\left(\sum\limits_{i=1}^n T_1(x_i), \ldots, \sum\limits_{i=1}^n T_k(x_i)\right)}
        \end{align*}
        
        Por el Teorema de Factorización de Neyman-Fisher, tenemos que el estadístico $\left(\sum\limits_{i=1}^n T_1(X_i), \ldots, \sum\limits_{i=1}^n T_k(X_i)\right)$ es suficiente para $\theta$. Respecto a la complitud, al ser más complicado, no se incluye la demostración.
    \end{proof}
\end{teo}

\begin{ejemplo}
    Sea $X$ una v.a. con distribución $\cc{N}(\mu, \sigma^2)$ con $\mu \in \mathbb{R}$ y $\sigma^2 > 0$. Buscar un estadístico suficiente y completo basado en una muestra de tamaño $n$.\\

    Como hemos visto en el ejemplo anterior, $\cc{N}(\mu, \sigma^2)$ es una familia exponencial bi-paramétrica en $(\mu, \sigma^2)$. Por tanto, por el Teorema anterior, tenemos que la familia de distribuciones asociadas a $(X_1, \ldots, X_n)$ es también una familia exponencial bi-paramétrica. Por tanto, el estadístico:
    \begin{equation*}
        \left(\sum\limits_{i=1}^n X_i, \sum\limits_{i=1}^n X_i^2\right)
    \end{equation*}
    es suficiente para $(\mu, \sigma^2)$. Además, veamos que la imagen de $Q(\mu, \sigma^2)$ contiene un abierto de $\bb{R}^2$. Consideremos la aplicación:
    \Func{Q}{\Theta=\bb{R}\times\bb{R}^+}{\bb{R}^2}{(\mu, \sigma^2)}{\left(\dfrac{\mu}{\sigma^2}, -\dfrac{1}{2\sigma^2}\right)}

    Observemos que al variar $\sigma^2 \in \bb{R}^+$, la segunda componente $Q_2 = \nicefrac{-1}{2\sigma^2}$ recorre todo el intervalo $\bb{R}^-$. Para cualquier valor fijado de $Q_2 \in \bb{R}^-$, al variar $\mu \in \bb{R}$, la primera componente $Q_1 = -2\mu Q_2$ recorre todo $\bb{R}$. Por consiguiente, el conjunto imagen de $Q$ es:
    \begin{equation*}
        Q(\Theta) = \bb{R} \times \bb{R}^-
    \end{equation*}
    que es un abierto de $\bb{R}^2$. Por tanto, por el Teorema anterior, si $n \geq 2$, el estadístico:
    \begin{equation*}
        \left(\sum_{i=1}^n X_i, \sum_{i=1}^n X_i^2\right)
    \end{equation*}
    es también completo para $(\mu, \sigma^2)$.
\end{ejemplo}
    % \section{Estimación puntual. Insesgadez y mínima varianza}

\begin{ejemplo}
    Sea $(X_1, \ldots, X_n)$ una muestra aleatoria simple de una variable aleatoria $X$ con distribución $\cc{P}(\lm)$, $\lm > 0$. La media muestral $\ol{X}$ es un estimador de $\lm$. Además, cualquier función medible de la muestra que sea independiente del parámetro $\lm$ y tome valores positivos es un estimador del parámetro.
    \begin{equation*}
        \ol{X} = \frac{1}{n} \sum_{i=1}^{n} X_i \geq 0
    \end{equation*}
    donde la positividad de $\ol{X}$ se debe a que $X_i \geq 0$ para todo $i = 1, \ldots, n$.

    Por lo tanto, $\ol{X}$ es un estimador de $\lm$ y cualquier función medible de la muestra que sea independiente del parámetro $\lm$ y tome valores positivos es un estimador del parámetro.
\end{ejemplo}

\begin{ejemplo}
    Sea $(X_1, \ldots, X_n)$ una muestra aleatoria simple de una variable aleatoria $X$ con distribución $\cc{B}(1, p)$, $p \in (0, 1)$. La media muestral $\ol{X}$ es un estimador de $p$. Además, cualquier función medible de la muestra que sea independiente del parámetro $p$ y tome valores en el intervalo $(0, 1)$ es un estimador del parámetro.
    \begin{equation*}
        \ol{X} = \frac{1}{n} \sum_{i=1}^{n} X_i \in \left[0, 1\right]
    \end{equation*}
    donde la pertenencia al intervalo $\left[0, 1\right]$ se debe a que $X_i \in \{0, 1\}$ para todo $i = 1, \ldots, n$. Por lo tanto, $\ol{X}$ es un estimador de $p$ y cualquier función medible de la muestra que sea independiente del parámetro $p$ y tome valores en el intervalo $\left[0, 1\right]$ es un estimador del parámetro.
\end{ejemplo}

\begin{prop}
    Sea $X$ una v.a. con función de distribución en una familia de distribuciones paramétricas, $F \in \{F_\theta : \theta \in \Theta\}$ y $(X_1, \ldots, X_n)$ una m.a.s. de $X$. Sea $T = T(X_1, \ldots, X_n)$ un estimador de $\theta$. El error cuadrático medio (ECM) de $T$ se define como:
    \begin{equation*}
        ECM_T(\theta) := \bb{E}_\theta\left[(T - \theta)^2\right]
    \end{equation*}

    El ECM puede descomponerse en términos de la varianza y una función denominada sesgo:
    \begin{equation*}
        ECM_T(\theta) = Var_\theta(T) + B^2_T(\theta)\qquad \text{donde } B_T(\theta) := \bb{E}_\theta[T] - \theta
    \end{equation*}
    \begin{proof}
        Tenemos que:
        \begin{align*}
            ECM_T(\theta) &= \bb{E}_\theta\left[(T - \theta)^2\right] \\
            &= \bb{E}_\theta\left[(T - \bb{E}_\theta[T] + \bb{E}_\theta[T] - \theta)^2\right] \\
            &= \bb{E}_\theta\left[(T - \bb{E}_\theta[T])^2 + 2(T - \bb{E}_\theta[T])(\bb{E}_\theta[T] - \theta) + (\bb{E}_\theta[T] - \theta)^2\right] \\
            &= Var_\theta(T) + 2(\bb{E}_\theta[T] - \theta)\bb{E}_\theta[T - \bb{E}_\theta[T]] + (\bb{E}_\theta[T] - \theta)^2 \\
            &= Var_\theta(T) + 2(\bb{E}_\theta[T] - \theta)\cancel{(\bb{E}_\theta[T] - \bb{E}_\theta[T])} + (\bb{E}_\theta[T] - \theta)^2 \\
            &= Var_\theta(T) + (\bb{E}_\theta[T] - \theta)^2 \\
            &= Var_\theta(T) + B^2_T(\theta)
        \end{align*}
    \end{proof}
\end{prop}

\begin{ejemplo}
    Sea $(X_1, \ldots, X_n)$ una muestra aleatoria simple de alguna población. Probar que, en general, la cuasidesviación típica no es un estimador insesgado de la desviación típica poblacional.
    \begin{proof}
        Sea $(X_1, \ldots, X_n)$ una muestra aleatoria simple de $X\rightsquigarrow \cc{N}(\mu, \sigma^2)$, $\mu \in \bb{R}$ y $\sigma^2 > 0$. Sabemos que $E[S^2] = \sigma^2$. Queremos ver si $E[S] = \sigma$. Sabemos que:
        \begin{align*}
            Y = \frac{(n-1)S^2}{\sigma^2} \sim \chi^2(n-1)
        \end{align*}
        Por tanto:
        \begin{align*}
            E\left[\sqrt{Y}\right] &= \int_0^\infty \sqrt{y} f_Y(y) dy = \int_0^\infty \sqrt{y}\dfrac{1}{2^{(n-1)/2}\Gamma((n-1)/2)} y^{(n-3)/2} e^{-y/2} dy
            =\\&= \dfrac{1}{2^{(n-1)/2}\Gamma((n-1)/2)} \int_0^\infty y^{n/2-1} e^{-y/2} dy
            \AstIg\\&\AstIg \dfrac{1}{2^{(n-1)/2}\Gamma((n-1)/2)} \cdot \dfrac{\Gamma(n/2)}{(1/2)^{n/2}}
        \end{align*}
        donde en $(\ast)$ hemos aplicado que la función de densidad de la distribución $\chi^2(n)$ integra $1$. Por tanto:
        \begin{align*}
            E[\sqrt{Y}] &= E\left[\sqrt{\frac{(n-1)S^2}{\sigma^2}}\right] = \dfrac{\sqrt{n-1}E[S]}{\sigma}
        \end{align*}

        Por tanto:
        \begin{align*}
            E[S] = \dfrac{\sigma}{\sqrt{n-1}} E[\sqrt{Y}] = \dfrac{\sigma}{\sqrt{n-1}} \cdot \dfrac{1}{2^{(n-1)/2}\Gamma((n-1)/2)} \cdot \dfrac{\Gamma(n/2)}{(1/2)^{n/2}}
            \neq \sigma
        \end{align*}
        Por tanto, $S$ no es un estimador insesgado de $\sigma$.
    \end{proof}
\end{ejemplo}

\begin{ejemplo}
    Sea $(X_1, \ldots, X_n)$ una m.a.s. de una distribución binomial $\cc{B}(1, p)$. Probar que
    \begin{equation*}
        \dfrac{T^2 - T}{n^2 - n}
    \end{equation*}
    con $T = \sum\limits_{i=1}^{n} X_i$ es un estimador insesgado para la función paramétrica $g(p) = p^2$. Probar, además, que para $n = 1$ no existe un estimador insesgado de $p^2$.
    \begin{proof}
        Es directo ver que se trata de un estimador. Veamos que es insesgado, para lo que calculamos su esperanza:
        \begin{align*}
            \bb{E}\left[\dfrac{T^2 - T}{n^2 - n}\right] &= \dfrac{1}{n^2 - n} \bb{E}[T^2 - T] = \dfrac{1}{n^2 - n} \left(\bb{E}[T^2] - \bb{E}[T]\right) =\\&= \dfrac{1}{n^2 - n} \left(Var[T] + \bb{E}[T]^2 - \bb{E}[T]\right) = \dfrac{1}{n^2 - n} \left(np(1-p) + (np)^2 - np\right) =\\&= \dfrac{1}{n^2 - n} \left(n^2p^2 - np^2\right) = p^2
        \end{align*}
        Por tanto, $\dfrac{T^2 - T}{n^2 - n}$ es un estimador insesgado de $p^2$. Veamos ahora que para el valo de $n = 1$ no existe un estimador insesgado de $p^2$. Supongamos que existe un estimador insesgado $T(X)$ de $p^2$ basado en una muestra de tamaño $1$. Entonces, por definición de insesgadez, se cumple que:
        \begin{equation*}
            p^2 = \bb{E}_p[T(X)] = T(0)(1-p) + T(1)p \Longrightarrow p^2 = T(0) - T(0)p + T(1)p
        \end{equation*}
        La expresión de la derecha es una ecuación polinómica de primer grado en $p$, mientras que $p^2$ es un polinomio de segundo grado en $p$. Por tanto, no existe un estimador insesgado de $p^2$ basado en una muestra de tamaño $1$.
    \end{proof}
\end{ejemplo}

\begin{ejemplo}
    Sea $X$ una v.a. con distribución $\cc{P}(\lm)$, ¿existe algún estimador insesgado de $\nicefrac{1}{\lm}$ basado en una muestra de tamaño $1$?\\

    Supongamos que existe un estimador insesgado $T(X)$ de $\nicefrac{1}{\lm}$ basado en una muestra de tamaño $1$. Entonces, por definición de insesgadez, se cumple que:
    \begin{equation*}
        \dfrac{1}{\lm} = \bb{E}_\lm[T(X)] = \sum_{x=0}^{\infty} T(x) \dfrac{e^{-\lm}\lm^x}{x!}
        \Longrightarrow
        e^{\lm} = \sum_{x=0}^{\infty} T(x) \dfrac{\lm^{x+1}}{x!}
    \end{equation*}

    Empleando la serie de Taylor de $e^{\lm}$, obtenemos:
    \begin{equation*}
        \sum_{x=0}^{\infty} \dfrac{\lm^x}{x!} = \sum_{x=0}^{\infty} T(x) \dfrac{\lm^{x+1}}{x!}
    \end{equation*}

    La primera sumatoria tiene un término independiente de $\lm$, mientras que la segunda sumatoria no. Como ha de ser cierta la igualdad para todo $\lm > 0$, se llega a una contradicción. Por tanto, no existe un estimador insesgado de $\nicefrac{1}{\lm}$ basado en una muestra de tamaño $1$.
\end{ejemplo}
    % \include{Ejercicios5}
    % \include{Ejercicios6}
    % \include{Ejercicios7}
    % \include{Ejercicios8}
    % \include{Ejercicios9}
    \chapter{Relaciones de Ejercicios}
    \section{Estadísticos muestrales}

\begin{ejercicio}
    Sea $(X_1, \ldots, X_n)$ una muestra aleatoria simple de una variable aleatoria $X$. Dar el espacio muestral y calcular la función masa de probabilidad de $(X_1, \ldots, X_n)$ en cada uno de los siguientes casos:
    \begin{enumerate}[label=\alph*)]
        \item $X\rightsquigarrow \{\cc{B}(k_0,p) : p\in (0,1)\}$ Binomial.

            El espacio muestral en este caso es $\cc{X}^n$, donde:
            \begin{equation*}
                \cc{X} = \{0, 1, ..., k_0\}
            \end{equation*}
            Recordamos que si $X\rightsquigarrow \cc{B}(k_0,p)$, entonces:
            \begin{equation*}
                P[X=x] = \binom{k_0}{x} p^x {(1-p)}^{k_0-x} \qquad \forall x\in \cc{X}
            \end{equation*}
            Por tanto, para nuestra m.a.s. tendremos la función masa de probabilidad:
            \begin{align*}
                P[X_1 &= x_1, \ldots, X_n = x_n] \stackrel{\text{indep.}}{=} \prod_{i=1}^{n}P[X_i = x_i]\stackrel{\text{id. d.}}{=} \prod_{i=1}^{n} P[X=x_i] \\
                      &= \prod_{i=1}^{n} \binom{k_0}{x_i} p^{x_i} {(1-p)}^{k_0-x_i} = p^{\sum\limits_{i=1}^{n}x_i} {(1-p)}^{nk_0 - \sum\limits_{i=1}^{n}x_i} \prod_{i=1}^{n}\binom{k_0}{x_i} \\
                      & \qquad \forall (x_1, \ldots, x_n) \in \cc{X}^n
            \end{align*}
        \item $X\rightsquigarrow\{\cc{P}(\lm) : \lm \in \mathbb{R}^+\}$ Poisson.

            El espacio muestral en este caso es $\cc{X}^n$, donde:
            \begin{equation*}
                \cc{X} = \mathbb{N} \cup \{0\}
            \end{equation*} 
            Recordamos que si $X\rightsquigarrow \cc{P}(\lm)$, entonces:
            \begin{equation*}
                P[X=x] = e^{-\lm} \dfrac{\lm^x}{x!} \qquad \forall x\in \cc{X}
            \end{equation*}
            Por tanto:
            \begin{align*}
                P[X_1 &= x_1, \ldots, X_n = x_n] \stackrel{\text{indep.}}{=} \prod_{i=1}^{n}P[X_i = x_i]\stackrel{\text{id. d.}}{=} \prod_{i=1}^{n} P[X=x_i] \\
                      &= \prod_{i=1}^{n} e^{-\lm} \dfrac{\lm^{x_i}}{x_i!} = e^{-n\lm} \prod_{i=1}^{n} \dfrac{\lm^{x_i}}{x_i!} = e^{-n\lm} \cdot \dfrac{\lm^{\sum\limits_{i=1}^n x_i}}{\prod\limits_{i=1}^{n}x_i!}  \qquad \forall (x_1, \ldots, x_n) \in \cc{X}^n
            \end{align*}
        \item $X\rightsquigarrow\{BN(k_0,p) : p\in (0,1)\}$ Binomial Negativa.

            El espacio muestral en este caso es $\cc{X}^n$, donde:
            \begin{equation*}
                \cc{X} = \mathbb{N}\cup \{0\}
            \end{equation*}
            Recordamos que si $X\rightsquigarrow BN(k_0,p)$, entonces:
            \begin{equation*}
                P[X=x] = \binom{x+k_0-1}{x} {(1-p)}^{x}p^{k_0} \qquad \forall x\in \cc{X}
            \end{equation*}
            Por tanto:
            \begin{align*}
                P[X_1 &= x_1, \ldots, X_n = x_n] \stackrel{\text{indep.}}{=} \prod_{i=1}^{n}P[X_i = x_i]\stackrel{\text{id. d.}}{=} \prod_{i=1}^{n} P[X=x_i] \\
                      &= \prod_{i=1}^{n} \binom{x_i+k_0-1}{x_i}{(1-p)}^{x_i}p^{k_0} = p^{nk_0}{(1-p)}^{\sum\limits_{i=1}^n x_i} \prod_{i=1}^{n}\binom{x_i+k_0-1}{x_i} \\
                      &\forall (x_1,\ldots,x_n)\in \cc{X}^n
            \end{align*}
        \item $X\rightsquigarrow\{G(p) : p\in (0,1)\}$ Geométrica.

            El espacio muestral en este caso es $\cc{X}^n$, donde:
            \begin{equation*}
                \cc{X} = \mathbb{N}\cup \{0\}
            \end{equation*}
            Recordamos que $G(p)\equiv BN(1,p)$, por lo que si sustituimos en la fórmula obtenida en la Binomial Negativa $k_0 = 1$:
            \begin{equation*}
                P[X_1=x_1, \ldots, X_n = x_n] = p^{n}{(1-p)}^{\sum\limits_{i=1}^n x_i} \qquad \forall (x_1,\ldots,x_n)\in \cc{X}^n
            \end{equation*}
        \item $X\rightsquigarrow\{P_N : N\in \mathbb{N}\}$, $\quad P_N[X=x] = \dfrac{1}{N}$, $\quad x=1,\ldots,N$.
         
            El espacio muestral en este caso es $\cc{X}^n$, donde (tal y como nos dan en el enunciado):
            \begin{equation*}
                \cc{X} = \{1, \ldots, N\}
            \end{equation*}
            
            Calculemos la masa de probabilidad:
            \begin{align*}
                P[X_1 &= x_1, \ldots, X_n = x_n] \stackrel{\text{indep.}}{=} \prod_{i=1}^{n}P[X_i = x_i]\stackrel{\text{id. d.}}{=} \prod_{i=1}^{n} P[X=x_i] \\
                      &= \prod_{i=1}^{n} \dfrac{1}{N} = {\left(\dfrac{1}{N}\right)}^{n} \qquad \forall (x_1,\ldots,x_n) \in \cc{X}^n
            \end{align*}
    \end{enumerate}
\end{ejercicio}

\begin{ejercicio}
    Sea $(X_1, \ldots, X_n)$ una muestra aleatoria simple de una variable aleatoria $X$. Dar el espacio muestral y calcular la función de densidad de $(X_1, \ldots, X_n)$ en cada uno de los siguientes casos:
    \begin{enumerate}[label=\alph*)]
        \item $X\rightsquigarrow\{U(a,b) : a,b\in \mathbb{R}, a<b\}$ Uniforme.

            El espacio muestral en este caso es $\cc{X}^n$, donde:
            \begin{equation*}
                \cc{X} = [a,b]
            \end{equation*}
            Recordamos que si $X\rightsquigarrow U(a,b)$, entonces:
            \begin{equation*}
                f_X(x) = \dfrac{1}{b-a} \qquad \forall x\in [a,b]
            \end{equation*}
            Por lo que:
            \begin{align*}
                f_{(X_1, \ldots, X_n)}(x_1, \ldots, x_n) &\stackrel{\text{indep.}}{=} \prod_{i=1}^{n} f_{X_i}(x_i) \stackrel{\text{id. d.}}{=} \prod_{i=1}^{n} f_X(x_i) = \prod_{i=1}^{n} \dfrac{1}{b-a} \\ &= {\left(\dfrac{1}{b-a}\right)}^{n} \qquad \forall (x_1, \ldots, x_n) \in \cc{X}^n
            \end{align*}

        \item $X\rightsquigarrow\{\cc{N}(\mu, \sigma^2) : \mu \in \mathbb{R}, \sigma^2 \in \mathbb{R}^+\}$ Normal.

            El espacio muestral en este caso es $\cc{X}^n$, donde $\cc{X} = \mathbb{R}$. Recordamos que si $X\rightsquigarrow \cc{N}(\mu, \sigma^2)$, entonces:
            \begin{equation*}
                f_X(x) = \dfrac{1}{\sqrt{2\pi} \sigma} e^{-\dfrac{{(x-\mu)}^{2}}{2\sigma^2}} \qquad \forall x\in \mathbb{R}
            \end{equation*}
            Por lo que:
            \begin{align*}
                f_{(X_1, \ldots, X_n)}(x_1, \ldots, x_n) &\stackrel{\text{indep.}}{=} \prod_{i=1}^{n} f_{X_i}(x_i) \stackrel{\text{id. d.}}{=} \prod_{i=1}^{n} f_X(x_i) 
                = \prod_{i=1}^{n} \dfrac{1}{\sqrt{2\pi} \sigma} e^{-\dfrac{{(x_i-\mu)}^{2}}{2\sigma^2}}  \\
                &= {\left(\dfrac{1}{\sqrt{2\pi}\sigma}\right)}^{n} \prod_{i=1}^{n} e^{-\dfrac{{(x_i-\mu)}^{2}}{2\sigma^2}}  = {\left(\dfrac{1}{\sqrt{2\pi}\sigma}\right)}^{n} e^{-\sum\limits_{i=1}^n\dfrac{{(x_i-\mu)}^{2}}{2\sigma^2}}   \\
                &= {\left(\dfrac{1}{\sqrt{2\pi}\sigma}\right)}^{n} e^{\frac{-1}{2\sigma^2}\sum\limits_{i=1}^n {(x_i-\mu)}^{2}} \qquad \forall (x_1,\ldots,x_n) \in \mathbb{R}^n
            \end{align*}
        \item $X\rightsquigarrow\{\Gamma(p,a) : p,a\in \mathbb{R}^+\}$ Gamma.

            El espacio muestral en este caso es $\cc{X}^n$, donde  $\cc{X}=\mathbb{R}^+_0$. Recordamos que si $X\rightsquigarrow \Gamma(p,a)$, entonces:
            \begin{equation*}
                f_X(x) = \dfrac{a^p}{\Gamma(p)} x^{p-1} e^{-ax} \qquad \forall x\in \mathbb{R}^+_0
            \end{equation*}
            Por lo que:
            \begin{align*}
                f_{(X_1, \ldots, X_n)}(x_1, \ldots, x_n) &\stackrel{\text{indep.}}{=} \prod_{i=1}^{n} f_{X_i}(x_i) \stackrel{\text{id. d.}}{=} \prod_{i=1}^{n} f_X(x_i) 
                = \prod_{i=1}^{n} \dfrac{a^p}{\Gamma(p)} x_i^{p-1} e^{-ax_i} \\
                                                         &= {\left(\dfrac{a^p}{\Gamma(p)}\right)}^{n} \cdot e^{-a\sum\limits_{i=1}^n x_i} \cdot \prod_{i=1}^{n} x_i^{p-1} \qquad \forall (x_1,\ldots,x_n)\in \cc{X}^n
            \end{align*}
        \item $X\rightsquigarrow\{\beta(p,q) : p,q\in \mathbb{R}^+\}$ Beta.

            El espacio muestral en este caso es $\cc{X}^n$, donde $\cc{X} = [0,1]$. Recordamos que si $X\rightsquigarrow \beta(p,q)$, entonces:
            \begin{equation*}
                f_X(x) = \dfrac{1}{\beta(p,q)} x^{p-1} {(1-x)}^{q-1} \qquad \forall x\in [0,1]
            \end{equation*}
            Donde:
            \begin{equation*}
                \beta(p,q) = \dfrac{\Gamma(p)\Gamma(q)}{\Gamma(p+q)}
            \end{equation*}
            Por tanto:
            \begin{align*}
                f_{(X_1, \ldots, X_n)}(x_1, \ldots, x_n) &\stackrel{\text{indep.}}{=} \prod_{i=1}^{n} f_{X_i}(x_i) \stackrel{\text{id. d.}}{=} \prod_{i=1}^{n} f_X(x_i) 
                = \prod_{i=1}^{n} \dfrac{1}{\beta(p,q)} x_i^{p-1} {(1-x_i)}^{q-1} \\
                                                         &= \dfrac{1}{{\beta(p,q)}^{n}} \prod_{i=1}^{n}x_i^{p-1} {(1-x_i)}^{q-1} \qquad \forall (x_1,\ldots,x_n) \in \cc{X}^n
            \end{align*}
        \item $X\rightsquigarrow\{P_\theta : \theta \in \mathbb{R}^+\}$, $\quad f_\theta(x) = \dfrac{1}{2\sqrt{x\theta}}$, $\quad 0<x<\theta$.
        
            El espacio muestral en este caso es $\cc{X}^n$, donde (tal y como nos dan en el enunciado):
            \begin{equation*}
                \cc{X} = \left]0,\theta\right[
            \end{equation*}
            
            Calculamos la función de densidad conjunta:
            \begin{align*}
                f_{(X_1, \ldots, X_n)}(x_1, \ldots, x_n) &\stackrel{\text{indep.}}{=} \prod_{i=1}^{n} f_{X_i}(x_i) \stackrel{\text{id. d.}}{=} \prod_{i=1}^{n} f_X(x_i) 
                = \prod_{i=1}^{n} \dfrac{1}{2\sqrt{x_i \theta}} \\
                                                         &= \dfrac{1}{{\left(2\sqrt{\theta}\right)}^{n}} \prod_{i=1}^{n} \dfrac{1}{\sqrt{x_i}} \qquad \forall (x_1, \ldots, x_n) \in  \cc{X}^n
            \end{align*}
    \end{enumerate}
\end{ejercicio}

\begin{ejercicio}
    Se miden los tiempos de sedimentación de una muestra de partículas flotando en un líquido. Los tiempos observados son: 
    \begin{gather*}
        11.5; 1.8; 7.3; 12.1; 1.8; 21.3; 7.3; 15.2; 7.3; 12.1; 15.2;\\ 7.3; 12.1; 1.8; 10.5; 15.2; 21.3; 10.5; 15.2; 11.5
    \end{gather*}
    \begin{enumerate}
        \item Construir la función de distribución muestral asociada a a dichas observaciones.
        
            En primer lugar, ordenamos los datos de menor a mayor, puesto que nos será de utilidad para construir la función de distribución muestral:
            \begin{gather*}
                1.8; 1.8; 1.8; 7.3; 7.3; 7.3; 7.3; 10.5; 10.5; 11.5;\\ 11.5; 12.1; 12.1; 12.1; 15.2; 15.2; 15.2; 15.2; 21.3; 21.3
            \end{gather*}

            Si aplicamos la definición de función de distribución muestral obtenemos que esta viene dada por la siguiente ecuación, y está mostrada en la Figura~\ref{fig:funcion_distribucion_muestral}:
            \begin{equation*}
                F_n^\ast(x) =
                \begin{cases} 
                0 & \text{si\ } x < 1.8 \\[6pt]
                \nicefrac{3}{20} &\text{si\ } 1.8 \leq x < 7.3 \\[6pt]
                \nicefrac{7}{20} &\text{si\ } 7.3 \leq x < 10.5 \\[6pt]
                \nicefrac{9}{20} &\text{si\ } 10.5 \leq x < 11.5 \\[6pt]
                \nicefrac{11}{20} &\text{si\ } 11.5 \leq x < 12.1 \\[6pt]
                \nicefrac{14}{20} &\text{si\ } 12.1 \leq x < 15.2 \\[6pt]
                \nicefrac{18}{20} &\text{si\ } 15.2 \leq x < 21.3 \\[6pt]
                \nicefrac{20}{20} &\text{si\ } x \geq 21.3
                \end{cases}
            \end{equation*}

            \begin{figure}[H]
                \centering
                \begin{tikzpicture}
                \begin{axis}[
                    width=12cm, height=7cm, xlabel={$x$}, ylabel={$F_n^\ast(x)$}, ymin=0, ymax=1.1,
                    xmin=0, xmax=23, xtick={0,1.8,7.3,10.5,12.1,15.2,21.3,23},
                    ytick={0,0.1,0.2,0.3,0.4,0.5,0.6,0.7,0.8,0.9,1},
                    grid=both, domain=0:23, samples=200,
                ]

                % Graficamos la función escalonada
                \addplot[
                    thick, blue
                ] coordinates {
                    (0,0) (1.8,0) (1.8,3/20) (7.3,3/20) (7.3,7/20) (10.5,7/20)
                    (10.5,9/20) (11.5,9/20) (11.5,11/20) (12.1,11/20) (12.1,14/20)
                    (15.2,14/20) (15.2,18/20) (21.3,18/20) (21.3,20/20) (23,20/20)
                };
                \end{axis}
                \end{tikzpicture}
                \caption{Gráfica de la función de distribución muestral.}
                \label{fig:funcion_distribucion_muestral}
            \end{figure}
        \item Hallar los valores de los tres primeros momentos muestrales respecto al origen y respecto a la media.

            Calculamos primero los tres primeros momentos respecto al origen para luego calcular los centrados respecto a la media a partir de ellos:
            \begin{align*}
                a_1 &= \sum_{i=1}^{n} f_i x_i = 10.915 \qquad 
                a_2 = \sum_{i=1}^{n} f_i x_i^2 = 148.9325 \\
                a_3 &= \sum_{i=1}^{n} f_i x_i^3 = 2280.98365 \\
                b_1 &= \sum_{i=1}^{n} f_i {(x_i - \overline{x})}= 0\qquad  \\
                b_2 &= \frac{1}{n}\sum_{i=1}^{n}{(x_i-\overline{x})}^{2} = \frac{1}{n}\sum_{i=1}^{n}(x_i^2 - 2x_i \overline{x}+\overline{x}^2)  \\ &= \frac{1}{n}\sum_{i=1}^{n}x_i^2 - \frac{2\overline{x}}{n}\sum_{i=1}^{n}x_i + \overline{x}^2 = a_2 -2a_1^2 + a_1^2 = a_2 - a_1^2 \\
                    &= 148.9325 - 10.915^2  = 29.795275\\
                b_3 &= \frac{1}{n}\sum_{i=1}^{n} {(x_i-\overline{x})}^{3} = \frac{1}{n}\sum_{i=1}^{n}\left(x_i^3 - 3x_i^2 \overline{x} + 3x_i\overline{x}^2 -\overline{x}^3\right) \\
                    &= \frac{1}{n}\sum_{i=1}^{n}x_i^3 - \frac{3\overline{x}}{n}\sum_{i=1}^{n}x_i^2 + \frac{3\overline{x}^2}{n}\sum_{i=1}^{n}x_i - \overline{x}^3 = a_3 - 3a_1a_2 + 3a_1^3 - a_1^3 \\
                    &= a_3 - 3a_1a_2 + 2a_1^3 = 4.95455925
            \end{align*}
        \item Determinar los valores de los cuartiles muestrales y el percentil 70.

            Para obtener los cuartiles, primero ordenamos los datos de menor a mayor y los agrupamos en grupos de $\nicefrac{20}{4} = 5$ en 5:
            \begin{gather*}
                1.8;\ 1.8;\ 1.8;\ 7.3;\ 7.3;\ \red{7.3;\ 7.3;\ 10.5;\ 10.5;\ 11.5};\ 11.5;\ 12.1;\ 12.1;\\ 12.1;\ 15.2;\ \red{15.2;\ 15.2;\ 15.2;\ 21.3;\ 21.3}
            \end{gather*}
            Como en los cambios de agrupaciones de números estos se repiten, hemos obtenido el valor de los cuartiles:
            \begin{equation*}
                q_1 = 7.3 \qquad q_2 = 11.5 \qquad q_3 = 15.2 \qquad q_4 = 21.3
            \end{equation*}
            Para el percentil $70$, calculamos:
            \begin{equation*}
                0.7\cdot 20 = 14
            \end{equation*}
            Como hemos obtenido un número entero, el percentil 70 será:
            \begin{equation*}
                c_{70} = \dfrac{X_{(14)} + X_{(15)}}{2} = \dfrac{12.1 + 15.2}{2} = 13.65
            \end{equation*}
            En el caso de haber obtenido un número no entero (por ejemplo, $14.2$), sería $X_{(15)}$.
    \end{enumerate}
\end{ejercicio}

\begin{ejercicio}
    Se dispone de una muestra aleatoria simple de tamaño 40 de una distribución exponencial de media 3, ¿cuál es la probabilidad de que los valores de la función de distribución muestral y la teórica, en $x=1$, difieran menos de $0.01$? Aproximadamente, ¿cuál debe ser el tamaño muestral para que dicha probabilidad sea como mínimo $0.98$?\\

    \noindent
    Como dice el enunciado, tenemos una m.a.s. $(X_1, \ldots, X_n)$ con $n=40$, todas ellas idénticamente distribuidas a $X\rightsquigarrow \exp(\lm)$. Sabemos de la asignatura de Probabilidad que:
    \begin{equation*}
        E[X] = \dfrac{1}{\lm} = 3 \Longrightarrow \lm = \dfrac{1}{3}
    \end{equation*}
    Por lo que $X\rightsquigarrow \exp\left(\nicefrac{1}{3}\right)$. Denotaremos por comodidad:
    \begin{equation*}
        F_n^\ast(x) = F_{(X_1, \ldots, X_n)}^\ast(x) 
    \end{equation*}
    Y el enunciado nos pregunta por:
    \begin{equation*}
        P[|F_n^\ast(1) - F_X(1)| < 0.01]
    \end{equation*}
    Para ello, primero calculamos $F_X(1)$, que por comodidad denotaremos por $\alpha$:
    \begin{equation*}
        F_X(1) = 1 - e^{-\lm\cdot  1} = 1-e^{-\lm} = 1-e^{\nicefrac{-1}{3}} =: \alpha
    \end{equation*}
    Por lo que nos disponemos ya a calcular la probabilidad:
    \begin{align*}
        P[|F_n^\ast(1) - F_X(1)| < 0.01] &= P[|F_n^\ast(1) - \alpha| < 0.01] = P[-0.01 < F_n^\ast(1) - \alpha < 0.01] \\
                                         &= P[-0.01+\alpha < F_n^\ast(1) < 0.01+\alpha] \\
                                         &= P[40(-0.01+\alpha) < 40F_n^\ast(1) < 40(0.01+\alpha)]
    \end{align*}
    Y como sabemos que $Y = 40F_n^\ast(1) \rightsquigarrow \cc{B}(40, F_X(1)) \equiv \cc{B}(40, \alpha)$:
    \begin{equation*}
        P[|F_n^\ast(1) - F_X(1)| < 0.01] = P[40(-0.01+\alpha) < Y< 40(0.01+\alpha)] 
    \end{equation*}
    Si ahora tomamos:
    \begin{equation*}
        \alpha = 1-e^{\nicefrac{-1}{3}}\approx 0.283469
    \end{equation*}
    Entonces:
    \begin{align*}
        40(0.01 + \alpha) &\approx 40(0.01+0.283469)= 11.73876 \\
        40(-0.01 + \alpha) &\approx 40(-0.01 + 0.283469)  = 10.93876
    \end{align*}
    Por lo que, usando que $Y$ es una variable aleatoria discreta, tenemos que:
    \begin{align*}
        P[|F_n^\ast(1) - F_X(1)| < 0.01] &\approx P[10.93876 < Y < 11.73876] = P[Y=11]
    \end{align*}
    De donde usando la masa de probabilidad de la Binomial:
    \begin{equation*}
        P[Y=11] = \binom{40}{11} {(0.283469)}^{11} {(1-0.283469)}^{40-11} \approx 0.139
    \end{equation*}

    Por tanto, la probabilidad de que los valores de la función de distribución muestral y la teórica, en $x=1$, difieran menos de $0.01$ es aproximadamente $0.139$, luego la probabilidad es baja. Para aumentar dicha probabilidad, necesitamos un tamaño muestral mayor, que es lo que nos piden en el segundo apartado.\\

    \noindent
    Para el segundo apartado, como para $n=40$ obtenemos una probabilidad de $0.139$, podemos intuir que para que dicha probabilidad sea como mínimo $0.98$, nos es necesario un valor de $n$ grande, por lo que podemos suponer que:
    \begin{equation*}
        F_n^\ast(1) \rightsquigarrow \cc{N}\left(\alpha, \dfrac{\alpha(1-\alpha)}{n}\right)
    \end{equation*}
    De donde:
    \begin{equation*}
        Z = \dfrac{\sqrt{n}(F_n^\ast(1)-\alpha)}{\sqrt{\alpha(1-\alpha)}} \rightsquigarrow \cc{N}(0,1)
    \end{equation*}
    Buscamos el valor de $n$ que verifica:
    \begin{equation*}
        0.98 \leq P[|F_n^\ast(1) - F_X(1)| < 0.01] = P\left[|Z| < \dfrac{\sqrt{n}0.01}{\sqrt{\alpha(1-\alpha)}}\right]
    \end{equation*}
    Si aplicamos propiedades conocidas de la función de distribución de una variable aleatoria, si $a\in \mathbb{R}$, entonces:
    \begin{equation*}
        P[|Z| < a] = P[-a < Z < a] = P[Z<a] - P[Z\leq -a]
        = P[Z<a] - P[Z< -a]
    \end{equation*}
    Pero, empleando las particularidades de la distribución normal, sabemos que:
    \begin{equation*}
        P[Z< -a] = P[Z>a] = 1-P[Z<a]
    \end{equation*}
    Por lo que:
    \begin{equation*}
        P[|Z| < a] = P[Z<a] - P[Z < -a] = 2P[Z<a] - 1
    \end{equation*}
    Volviendo al caso que nos interesa:
    \begin{equation*}
        0.98 \leq P\left[|Z| < \dfrac{\sqrt{n}0.01}{\sqrt{\alpha(1-\alpha)}}\right] = 2P\left[Z<\dfrac{\sqrt{n}0.01}{\sqrt{\alpha(1-\alpha)}}\right] - 1
    \end{equation*}
    Luego:
    \begin{equation*}
        0.99 = \dfrac{0.98+1}{2} \leq P\left[Z < \dfrac{\sqrt{n}0.01}{\sqrt{\alpha(1-\alpha)}}\right]
    \end{equation*}
    Si consultamos la tabla de la normal $\cc{N}(0,1)$, observamos que el primer valor que supera la probbailidad de $0.99$ es $2.33$, por lo que:
    \begin{equation*}
        2.33 = \dfrac{\sqrt{n}0.01}{\sqrt{\alpha(1-\alpha)}} = \dfrac{\sqrt{n}0.01}{\sqrt{0.283469(1-0.283469)}} \approx 0.0221886 \sqrt{n}
    \end{equation*}
    De donde:
    \begin{equation*}
        5.4289 = {(2.33)}^{2} = {(0.0221886\sqrt{n})}^{2} = 0.00049233 n \Longrightarrow n = \dfrac{5.4289}{0.00049233} = 11026.95347
    \end{equation*}
    Por lo que para $n \geq 11027$ podemos asegurar que la probabilidad es como mínimo $0.98$.
\end{ejercicio}

\begin{ejercicio}
    Se dispone de una muestra aleatoria simple de tamaño 50 de una distribución de Poisson de media 2, ¿cuál es la probabilidad de que los valores de la función de distribución muestral y la teórica, en $x=2$, difieran menos de $0.02$? Aproximadamente, ¿qué tamaño muestral hay que tomar para que dicha probabilidad sea como mínimo $0.99$?\\

    \noindent
    Tenemos una m.a.s. $(X_1, \ldots, X_n)$ con $n=50$ idénticamente distribuidas a $X\rightsquigarrow~\cc{P}(\lm)$. Como $E[X] = \lm = 2$, tenemos que $X\rightsquigarrow \cc{P}(2)$. Notamos por comodidad:
    \begin{equation*}
        F_{(X_1, \ldots, X_n)}^\ast(x) = F_n^\ast(x)
    \end{equation*}
    Nos preguntan por:
    \begin{equation*}
        P[|F_n^\ast(2) - F_X(2)| < 0.02]
    \end{equation*}
    Para ello primero calculamos:
    \begin{equation*}
        F_X(2) = \sum_{k=0}^{2} e^{-2}\dfrac{2^k}{k!} =  e^{-2}\left(\dfrac{2^0}{0!} + \dfrac{2^1}{1!} + \dfrac{2^2}{2!}\right) = e^{-2} (1+2+2) \approx 0.6767
    \end{equation*}
    Por lo que:
    \begin{equation*}
        P[|F_n^\ast(2) - 0.6767| < 0.02] = P[-0.02 < F_n^\ast(2) - 0.6767 < 0.02] = P[0.6567 < F_n^\ast(2) < 0.6967]
    \end{equation*}
    Como sabemos por lo visto en teoría que:
    \begin{equation*}
        Y = 50F_n^\ast(2) \rightsquigarrow \cc{B}(50, F_X(2)) \equiv \cc{B}(50,\ \ 0.6767)
    \end{equation*}
    Multiplicamos por $50$ la última expresión:
    \begin{align*}
        P[|F_n^\ast(2) - 0.6767| < 0.02] = P[0.6567 < F_n^\ast(2) < 0.6967] &= P[32.835 < Y< 34.835] \\
                                         &\AstIg P[Y = 33] + P[Y=34]
    \end{align*}
    donde en $(\ast)$ hemos empleado la naturaleza discreta de la variable aleatoria $Y$. Calculamos estas dos probabilidades:
    \begin{align*}
        P[Y=33] &= \binom{50}{33} {(0.6767)}^{33} {(1-0.6767)}^{50-33} \approx 0.114734\\
        P[Y=34] &= \binom{50}{34} {(0.6767)}^{34} {(1-0.6767)}^{50-34} \approx 0.120075
    \end{align*}
    Por lo que:
    \begin{equation*}
        P[|F_n^\ast(2) - 0.6767| < 0.02] \approx 0.114734 + 0.120075 = 0.234809
    \end{equation*}

    \noindent
    Para el segundo apartado, como para $n=50$ obtenemos una probabilidad de $0.234809$, podemos intuir que para que dicha probabilidad sea como mínimo $0.99$, nos es necesario un valor de $n$ grande, por lo que podemos suponer que:
    \begin{equation*}
        F_n^\ast(2) \rightsquigarrow \cc{N}\left(0.6767, \dfrac{0.6767(1-0.6767)}{n}\right) \equiv \cc{N}\left(0.6767, \dfrac{0.218777}{n}\right)
    \end{equation*}
    Por lo que:
    \begin{equation*}
        Z = \dfrac{\sqrt{n}(F_n^\ast(2) - 0.6767)}{\sqrt{0.218777}} \rightsquigarrow\cc{N}(0,1)
    \end{equation*}
    En dicho caso, buscamos $n$ de forma que:
    \begin{equation*}
        0.99 \leq P[|F_n^\ast(2) - F_X(2)| < 0.02] = P\left[|Z| < \dfrac{\sqrt{n}0.02}{\sqrt{0.218777}}\right]
    \end{equation*}
    De forma análoga al ejercicio anterior:
    \begin{equation*}
        P\left[|Z| < \dfrac{\sqrt{n}0.02}{\sqrt{0.218777}}\right] = 2P\left[Z < \dfrac{\sqrt{n}0.02}{\sqrt{0.218777}}\right] - 1
    \end{equation*}
    Luego:
    \begin{equation*}
        0.995 = \dfrac{0.99 + 1}{2} \leq P\left[Z < \dfrac{\sqrt{n}0.02}{\sqrt{0.218777}}\right] 
    \end{equation*}
    Y si miramos la tabla de la Normal observamos que el primer valor que supera la probabilidad de $0.995$ es $2.58$, luego:
    \begin{equation*}
        2.58 = \dfrac{\sqrt{n}0.02}{\sqrt{0.218777}} = 0.042759 \sqrt{n}
    \end{equation*}
    Por lo que:
    \begin{equation*}
        6.6564 = {(2.58)}^{2} = {(0.042759 \sqrt{n})}^{2} = 0.00182833 n
    \end{equation*}
    Luego:
    \begin{equation*}
        n = \dfrac{6.6564}{0.00182833} \approx 3640.7
    \end{equation*}
    Por lo que para $n\geq 3641$ podemos asegurar que la probabilidad es como mínimo $0.99$.
\end{ejercicio}

\begin{ejercicio}
   Sea $X\rightsquigarrow \cc{B}(1,p)$ y $(X_1, X_2, X_3)$ una muestra aleatoria simple de $X$. Calcular la función masa de probabilidad de los estadísticos $\overline{X}$, $S^2$, $\min X_i$ y $\max X_i$.\\

   \noindent
   Para resolver este ejercicio, como $X$ sigue una distribución discreta, buscamos aplicar el teorema de cambio de variable de discreta a discreta. Para ello, la forma más cómoda será analizar cada uno de los valores que puede tomar la muestra aleatoria simple $(X_1, X_2, X_3)$ y determinar en consecuencia cada uno de los valores que toman $\overline{X}$, $S^2$, $\min X_i$ y $\max X_i$. Acompañaremos la tabla junto con la probabilidad de que la muestra tome dicho valor, es decir, en la fila correpondiente a $(x_1, x_2, x_3)$ incluiremos $P[X_1 = x_1, X_2 = x_2, X_3 = x_3]$:
   \begin{equation*}
   \begin{array}{c|c|c|c|c|c}
       P & (X_1, X_2, X_3) & \overline{X} & S^2 & \min X_i & \max X_i  \\
       \hline
       {(1-p)}^{3} & (0,0,0) & 0 & 0 & 0 & 0 \\
       p{(1-p)}^{2} & (0,0,1) & \nicefrac{1}{3} & \nicefrac{1}{3} & 0 & 1 \\
       p{(1-p)}^{2} & (0,1,0) & \nicefrac{1}{3} & \nicefrac{1}{3} & 0 & 1 \\
       p^2(1-p) & (0,1,1) & \nicefrac{2}{3} & \nicefrac{1}{3} & 0 & 1 \\
       p{(1-p)}^{2} & (1,0,0) & \nicefrac{1}{3} & \nicefrac{1}{3} & 0 & 1 \\
       p^2(1-p) & (1,0,1) & \nicefrac{2}{3} & \nicefrac{1}{3} & 0 & 1 \\
       p^2(1-p) & (1,1,0) & \nicefrac{2}{3} & \nicefrac{1}{3} & 0 & 1 \\
       p^3 & (1,1,1) & 1 & 0 & 1 & 1 
   \end{array}
   \end{equation*}
   Podemos ya calcular la función masa de probabilidad de cada uno de los estadísticos, simplemente sumando las probabilidades de la tabla que corresponden a cada valor del espacio muestral de cada estadístico:
   \begin{itemize}
       \item Para $\overline{X}$:
           \begin{align*}
               P[\overline{X} = 0] &= P[X_1 = 0, X_2 = 0, X_3=0] = {(1-p)}^{3}  \\
               P[\overline{X} = \nicefrac{1}{3}] &= \sum_{i=1}^{3}p{(1-p)}^{2} = 3p{(1-p)}^{2} \\
               P[\overline{X} = \nicefrac{2}{3}] &= \sum_{i=1}^{3}p^2(1-p) = 3p^2(1-p) \\
               P[\overline{X} = 1] &= P[X_1 = 1, X_2 = 1, X_3 = 1] = p^3
           \end{align*}
       \item Para $S^2$:
           \begin{align*}
               P[S^2 = 0] &= P[X_1 = 0, X_2 = 0, X_3 = 0] + P[X_1 = 1, X_2 = 1, X_3 = 1] = p^3 + {(1-p)}^{3} \\
               P[S^2 = \nicefrac{1}{3}] &= \sum_{i=1}^{3} p{(1-p)}^{2} + \sum_{i=1}^{3}p^2(1-p) = 3p(1-p)(p+1-p) = 3p(1-p)
           \end{align*}
       \item Para $\min X_i$:
           \begin{align*}
               P[\min X_i = 1] &= P[X_1 = 1, X_2 = 1, X_3 = 1] = p^3 \\
               P[\min X_i = 0] &= 1-P[\min X_i = 1] = 1-p^3
           \end{align*}
       \item Para $\max X_i$:
           \begin{align*}
               P[\max X_i = 0] &= P[X_1 = 0, X_2 = 0, X_3 = 0] = {(1-p)}^{3}\\
               P[\max X_i = 1] &= 1-P[\max X_i = 0] = 1-{(1-p)}^{3}
           \end{align*}
   \end{itemize}
\end{ejercicio}

\begin{ejercicio}\label{ej:funcion_masa_densidad}
    Obtener la función masa de probabilidad o función de densidad de $\overline{X}$ en el muestreo de una variable de Bernoulli, de una Poisson y de una exponencial.\\

    \noindent
    Calculamos la masa de probabilidad o función de densidad en cada caso, suponiendo que tenemos $(X_1, \ldots, X_n)$ una muestra aleatoria simple con variables aleatorias idénticamente distribuidas a $X$, que sigue una distribución distinta en cada caso y estaremos interesados en calcular la masa de:
    \begin{equation*}
        \overline{X} = \frac{1}{n}\sum_{i=1}^{n} X_i
    \end{equation*}
    \begin{description}
        \item [Bernoulli.] Veamos dos opciones:
        \begin{description}
            \item[Opción 1] Supuesto que $X\rightsquigarrow \cc{B}(1,p)$ para cierto $p\in \left]0,1\right[$, si tomamos:
            \begin{equation*}
                Y = \sum_{i=1}^{n}X_i
            \end{equation*}
            Por la propiedad reproductiva de la Bernoulli, tenemos que $Y\rightsquigarrow \cc{B}(n,p)$. En dicho caso:
            \begin{equation*}
                P[Y=k] = \binom{n}{k} p^k {(1-p)}^{n-k} \qquad \forall k\in \{0,\ldots,n\}
            \end{equation*}

            Por tanto, tendremos que:
            \begin{equation*}
                P[\overline{X} = x] = P\left[\frac{1}{n}\sum_{i=1}^{n}X_i = x\right] = P[Y = nx]
            \end{equation*}
            Como $Y$ es una variable aleatoria discreta, $nx$ debe ser un número entero entre $0$ y $n$, por lo que $x$ debe ser de la forma $\nicefrac{k}{n}$ para cierto $k\in \{0,\ldots,n\}$. Por lo que:
            \begin{equation*}
                P\left[\overline{X} = \frac{k}{n}\right] = P[Y=k] = \binom{n}{k} p^k {(1-p)}^{n-k} \qquad \forall k\in \{0,\ldots,n\}
            \end{equation*}
            
            \item[Opción 2] Calculemos la función generatriz de momentos de $\overline{X}$, que por la igualdad $(\ast)$ vista en teoría, tenemos que:
            \begin{align*}
                M_{\overline{X}}(t) &\AstIg {(M_X(\nicefrac{t}{n}))}^{n} = {\left(p e^{\nicefrac{t}{n}} + (1-p)\right)}^{n} = {\left(p e^{\nicefrac{t}{n}} + (1-p)\right)}^{n}
            \end{align*}
            Observamos que obtenemos una función generatriz de momentos para $\overline{X}$ igual que para una variable aleatoria de distribución $\cc{B}(n,p)$, en la que los valores posibles de $\overline{X}$ son $\nicefrac{k}{n}$ para cierto $k\in \{0,\ldots,n\}$.
        \end{description}
        \item [Poisson.] Supuesto que $X\rightsquigarrow \cc{P}(\lm)$ para cierto $\lm\in \mathbb{R}^+$, si tomamos:
            \begin{equation*}
                Y = \sum_{i=1}^{n}X_i
            \end{equation*}
            Por la propiedad reproductiva de la Poisson, tendremos que:
            \begin{equation*}
                Y\rightsquigarrow\cc{P}\left(\sum_{i=1}^{n}\lm\right) \equiv \cc{P}(n\lm)
            \end{equation*}
            En dicho caso:
            \begin{equation*}
                P[Y=x] = e^{-n\lm} \dfrac{{(n\lm)}^{x}}{x!} \qquad \forall x\in \mathbb{N}
            \end{equation*}
            Por lo que, empleando la misma idea que en el caso de la Bernoulli, tendremos que:
            \begin{equation*}
                P[\overline{X}=\nicefrac{x}{n}] = P[Y=x] = e^{-n\lm} \dfrac{{(n\lm)}^{x}}{x!} \qquad \forall x\in \mathbb{N}
            \end{equation*}
        \item [Exponencial.] Supuesto ahora que $X\rightsquigarrow \exp(\lm)$ para cierto $\lm\in \mathbb{R}^+$, tendremos entonces que:
            \begin{equation*}
                M_X(t) = \dfrac{\lm}{\lm - t} \qquad t<\lm
            \end{equation*}
            Si aplicamos la igualdad $(\ast)$ vista en teoría:
            \begin{align*}
                M_{\overline{X}}(t) &\AstIg {(M_X(\nicefrac{t}{n}))}^{n} = {\left(\dfrac{\lm}{\lm - \nicefrac{t}{n}}\right)}^{n} = {\left(\dfrac{n\lm}{n\lm - t}\right)}^{n}
            \end{align*}
            Observamos que obtenemos una función generatriz de momentos para $\overline{X}$ igual que para una variable aleatoria de distribución $\Gamma(n,n\lm)$. Como la función generatriz de momentos de una variable aleatoria caracteriza su distribución, concluimos que $\overline{X}\rightsquigarrow\Gamma(n,n\lm)$. Por tanto, la función de densidad de $\overline{X}$ es:
            \begin{equation*}
                f_{\overline{X}}(x) = \dfrac{{(n\lm)}^{n}}{\Gamma(n)} x^{n-1} e^{-n\lm x} \qquad \forall x>0
            \end{equation*}
    \end{description}
\end{ejercicio}

\begin{ejercicio}
    Calcular las funciones de densidad de los estadísticos $\max X_i$ y $\min X_i$ en el muestreo de una variable $X$ con funcion de densidad:
    \begin{equation*}
        f_\theta(x) = e^{\theta-x}, \qquad x>\theta.
    \end{equation*}

    \noindent
    Calculamos primero la función de distribución, para calcular con mayor comodidad las funciones de distribución de $X_{(n)}$ y $X_{(1)}$:
    \begin{equation*}
        F_\theta(x) = \int_{\theta}^{x} f_\theta(t)~dt = \int_{\theta}^{x} e^{\theta-t}~dt  = \left[-e^{\theta-t}\right]_\theta^x = 1 - e^{\theta - x} \qquad \forall x>\theta
    \end{equation*}
    Supuesto ahora que disponemos de una m.a.s. $(X_1, \ldots, X_n)$ idénticamente distribuidas a $X$ cuya función de densidad es la anteriormente dicha, podemos aplicar las fórmulas obtenidas en teoría para calcular las funciones de distribución del mínimo y del máximo. Para el máximo:
    \begin{equation*}
        F_{X_{(n)}}(x) = {(F_X(x))}^{n} = {(1-e^{\theta-x})}^{n} \Longrightarrow f_{X_{(n)}}(x) = n{(1-e^{\theta - x})}^{n-1}e^{\theta -x} \qquad \forall x>\theta
    \end{equation*}

    Para el mínimo:
    \begin{equation*}
        F_{X_{(1)}}(x) = 1 - {(1-F_X(n))}^{n} = 1-{(1-1+e^{\theta-x})}^{n} = 1-e^{n(\theta-x)} \qquad \forall x>\theta
    \end{equation*}
    de donde:
    \begin{equation*}
        f_{X_{(1)}}(x) = ne^{n(\theta-x)} \qquad \forall x>\theta
    \end{equation*}
\end{ejercicio}

\begin{ejercicio}
    El número de pacientes que visitan diariamente una determinada consulta médica es una variable aleatoria con varianza de 16 personas. Se supone que el número de visitas de cada día es independiente de cualquier otro. Si se observa el número de visitas diarias durante 64 días, calcular aproximadamente la probabilidad de que la media muestral no difiera en más de una persona del valor medio verdadero de visitas diarias.\\

    \noindent
    Sea $X$ una variable aleatoria que indica el número de pacientes que visitan diariamente dicha consulta médica. Por cómo nos definen $X$ sabemos que $X\rightsquigarrow\cc{P}(\lm)$. Como además nos dicen que la varianza de dicha variable aleatoria es $16$, tenemos que $Var(X) = \lm = 16$. Si tenemos ahora una muestra aleatoria simple $(X_1, \ldots, X_n)$ con $n=64$, nos preguntan por:
    \begin{equation*}
        P[|\overline{X} - E[X]| \leq 1]
    \end{equation*}
    Notemos que se incluye el valor $1$ en la desigualdad.

    Donde $E[X] = \lm = 16$, ya que $X\rightsquigarrow\cc{P}(16)$. Calculamos:
    \begin{align*}
        P[|\overline{X} - E[X]| \leq 1] &= P[-1 \leq \overline{X}-16 \leq 1] = P[15 \leq \overline{X} \leq 17]
    \end{align*}
    Aplicamos ahora lo visto en el Ejercicio~\ref{ej:funcion_masa_densidad}, ya que si $X\rightsquigarrow\cc{P}(\lm)$, entonces tendremos que $n\overline{X}\rightsquigarrow \cc{P}(n\lm)$, gracias a la propiedad reproductiva de la Poisson:
    \begin{align*}
        P[|\overline{X} - E[X]| \leq 1] &= P[15 \leq \overline{X} \leq 17] = P[64\cdot 15 \leq 64\overline{X}\leq 64\cdot 17] \\
                                     &= P[960 \leq 64\overline{X} \leq 1088]
    \end{align*}
    Donde $64\overline{X}\rightsquigarrow\cc{P}(64\cdot 16) \equiv \cc{P}(1024)$. Para calcular dicha probabilidad, aproximaremos la Poisson a una distribución normal empleando el Teorema Central del Límite, ya que $n=64$ es un tamaño muestral suficientemente grande. Por lo que:
    \begin{equation*}
        \cc{P}(1024) \approx \cc{N}(1024, 1024)
    \end{equation*}
    Por lo que:
    \begin{align*}
        P[|\overline{X} - E[X]| \leq 1] &= P[960 \leq 64\overline{X} \leq 1088] \approx P\left[\dfrac{960-1024}{\sqrt{1024}} \leq Z \leq \dfrac{1088 - 1024}{\sqrt{1024}}\right] \\
                                     &= P[-2 \leq Z \leq 2] = 2P[Z\leq 2] -1  \\
                                     &= 2\cdot 0.97725 - 1 = 0.9545
    \end{align*}
\end{ejercicio}

\begin{ejercicio}
    Una máquina de refrescos está arreglada para que la cantidad de bebida que sirve sea una variable aleatoria con media \unit[200]{ml}. y desviación típica \unit[15]{ml}. Calcular de forma aproximada la probabilidad de que la cantidad media servida en una muestra aleatoria de tamaño 36 sea al menos \unit[204]{ml}.\\

    \noindent
    Sea $(X_1, \ldots, X_n)$ una muestra aleatoria simple idénticamente distribuida a la variable aleatoria $X$ de tamaño $n=36$. De la información que nos dan, sabemos que $E[X] = 200$ y $Var(X) = 15^2$, pero no conocemos la distribución de $X$.

    Una primera aproximación sería emplear el Teorema Central del Límite para aproximar la distribución $X$ a una distribución normal, y de ahí intentar obtener la distribución de $\overline{X}$. No obstante, un resultado más directo es directamente calcular $E[\overline{X}]$ y $Var(\overline{X})$.
    \begin{align*}
        E[\overline{X}] &= E\left[\frac{1}{n}\sum_{i=1}^{n}X_i\right] = \frac{1}{n}\sum_{i=1}^{n}E[X_i] = \frac{1}{n}\sum_{i=1}^{n}E[X] = E[X] = 200 \\
        Var[\overline{X}] &= Var\left[\frac{1}{n}\sum_{i=1}^{n}X_i\right] = \frac{1}{n^2}\sum_{i=1}^{n}Var[X_i] = \frac{1}{n^2}\sum_{i=1}^{n}Var[X] = \frac{Var[X]}{n} = \frac{15^2}{36}
    \end{align*}

    Por tanto, aplicando ahora el Teorema Central del Límite, obtenemos que $\overline{X}\rightsquigarrow\cc{N}\left(200, \frac{15^2}{36}\right)$. Calculamos la probabilidad:
    \begin{align*}
        P[\overline{X}\geq 204] \stackrel{\text{tipificamos}}{\approx} P \left[Z \geq \dfrac{\sqrt{36}(204-200)}{15}\right] &= P[Z\geq 1.6] = 1 - P[Z< 1.6] \\
        &= 1 - 0.9452 = 0.0548
    \end{align*}
\end{ejercicio}

\subsection*{Ejercicios Adicionales}

\begin{ejercicio}[Reto]
    Sea $(X_1, \ldots, X_n)$ una muestra aleatoria simple de una variable aleatoria $X$ con función de densidad:
    \begin{equation*}
        f_\theta(x) = \theta x^{\theta-1}, \qquad 0<x<1, \quad \theta>0
    \end{equation*}

    Obtener la función de densidad de $X_{(n)}$ y de $X_{(1)}$.\\

    Obtenemos primero la función de distribución de $X$:
    \begin{equation*}
        F_\theta(x) = \int_{0}^{x} f_\theta(t)~dt = \int_{0}^{x} \theta t^{\theta-1}~dt = \left[t^\theta\right]_0^x = x^\theta \qquad \forall x\in \left]0,1\right[
    \end{equation*}

    Pasamos ahora a calcular la función de distribución de $X_{(n)}$:
    \begin{equation*}
        F_{X_{(n)}}(x) = P[X_{(n)} \leq x] = P[X_1 \leq x, \ldots, X_n \leq x] \stackrel{\text{idd}}{=} \prod_{i=1}^{n} P[X \leq x] = {(F_\theta(x))}^{n} = x^{n\theta} \qquad \forall x\in \left]0,1\right[
    \end{equation*}
    Derivando, obtenemos la función de densidad de $X_{(n)}$:
    \begin{equation*}
        f_{X_{(n)}}(x) = n\theta x^{n\theta-1} \qquad \forall x\in \left]0,1\right[
    \end{equation*}

    Pasamos ahora a calcular la función de distribución de $X_{(1)}$:
    \begin{align*}
        F_{X_{(1)}}(x) = P[X_{(1)} \leq x] &= 1 - P[X_{(1)} > x] = 1 - P[X_1 > x, \ldots, X_n > x] \stackrel{\text{idd}}{=} 
        \\&\stackrel{\text{idd}}{=} 1 - \prod_{i=1}^{n} P[X > x] = 1 - {(1-F_\theta(x))}^{n} = 1 - {(1-x^\theta)}^{n} \qquad \forall x\in \left]0,1\right[
    \end{align*}
    Derivando, obtenemos la función de densidad de $X_{(1)}$:
    \begin{equation*}
        f_{X_{(1)}}(x) = n\theta x^{\theta-1} {(1-x^\theta)}^{n-1} \qquad \forall x\in \left]0,1\right[
    \end{equation*}
\end{ejercicio}



    \newpage
\section{Distribuciones en el muestreo de poblaciones normales}

\begin{ejercicio}
    Se toma una muestra aleatoria simple de tamaño 5 de una variable aleatoria con distribución $\cc{N}(2.5,\ 36)$. Calcular:
    \begin{enumerate}[label=\alph*)]
        \item Probabilidad de que la cuasivarianza muestral esté comprendida entre $1.863$ y $2.674$.

            Tenemos una m.a.s. $(X_1, X_2, X_3, X_4, X_5)$, todas ellas idénticamente distribuidas a $X\rightsquigarrow \cc{N}(2.5,\ 36)$. Tomamos la cuasivarianza de dichos datos:
            \begin{equation*}
                S^2 = \dfrac{1}{n-1} \sum_{i=1}^n {(X_i - \ol{X})}^{2}
            \end{equation*}
            Y queremos calcular la probabilidad:
            \begin{equation*}
                P\left[1.863 < S^2 < 2.674\right]
            \end{equation*}
            Puesto que no es fácil obtener la distribución de $S^2$, vamos a usar la distribución $\chi^2$ para calcular dicha probabilidad. Para ello, en teoría hemos visto que como estamos muestreando de una población normal, se cumple que:
            \begin{equation*}
                \dfrac{(n-1)S^2}{\sigma^2}= \dfrac{4S^2}{36} \rightsquigarrow\chi^2(4) \equiv \chi^2(n-1)
            \end{equation*}
            Por tanto:
            \begin{align*}
                P[1.863 < S^2 < 2.674] &= P\left[\dfrac{4\cdot 1.863}{36}<\dfrac{4S^2}{36} < \dfrac{4\cdot 2.674}{36}\right] \\
                                       &= P\left[0.207 < \dfrac{4S^2}{36} < 0.2971\right]  \\
                                       &= P\left[\dfrac{4S^2}{36} > 0.207\right] - P\left[\dfrac{4S^2}{36}>0.2971\right]
            \end{align*}
            Si consultamos la tabla de la $\chi^2(4)$:
            \begin{equation*}
                P\left[\dfrac{4S^2}{36} > 0.207\right] \approx 0.995, \qquad 
                P\left[\dfrac{4S^2}{36}>0.2971\right] \approx 0.99
            \end{equation*}
            Por lo que:
            \begin{align*}
                P[1.863 < S^2 < 2.674] &= P\left[\dfrac{4S^2}{36} > 0.207\right] - P\left[\dfrac{4S^2}{36}>0.2971\right] \approx 0.995 - 0.99 \\
                                       &= 0.005
            \end{align*}
        \item Probabilidad de que la media muestral esté comprendida entre $1.3$ y $3.5$, supuesto que la cuasivarianza muestral está entre $30$ y $40$. 

            En la situación del apartado anterior, ahora tenemos también en consideración la media muestral:
            \begin{equation*}
                \overline{X} = \dfrac{1}{n}\sum_{i=1}^n X_i
            \end{equation*} 
            Bajo estas condiciones, el Lema de Fisher nos dice que $S^2$ y $\overline{X}$ son independientes, por lo que la probabilidad de que la media muestral esté comprendida entre $1.3$ y $3.5$ suponiendo que la cuasivarianza muestral está comprendida entre $30$ y $40$ es igual a la probabilidad de que la media muestral esté comprendida entre $1.3$ y $3.5$:
            \begin{equation*}
                P[1.3<\overline{X}<3.5\ |\ 30 < S^2 < 40] = P[1.3<\overline{X}<3.5]
            \end{equation*}

            Sabemos que $\overline{X}\rightsquigarrow\cc{N}\left(\mu, \frac{\sigma^2}{n}\right)$, nos disponemos a calcular dicha probabilidad:
            \begin{align*}
                P[1.3<\overline{X}<3.5] &\stackrel{\text{tipificamos}}{=} P\left[\dfrac{\sqrt{5}(1.3-2.5)}{\sqrt{36}}< Z < \dfrac{\sqrt{5}(3.5-2.5)}{\sqrt{36}}\right] \\
                                        &= P[-0.447214 < Z < 0.372678] \\ 
                                        &= P[Z<0.372678] - P[Z>0.447214] \\
                                        &= P[Z<0.372678] -1 + P[Z<0.447214] \\
                                        &\approx 0.64431 - 1 + 0.67003 = 0.31434
            \end{align*}
    \end{enumerate}
\end{ejercicio}

\begin{ejercicio}
    La longitud craneal en una determinada población humana es una variable aleatoria que sigue una distribución normal con media $\unit[185.6]{mm}$ y desviación típica $\unit[12.78]{mm}$. ¿Cuál es la probabilidad de que una muestra aleatoria simple de tamaño 20 de esa población tenga media mayor que $\unit[190]{mm}$?\\

    \noindent
    Tenemos una población que sigue una variable aleatoria $X\rightsquigarrow \cc{N}\left(\mu, \sigma^2\right)$ con parámetros $\mu = 185.6$ y $\sigma=12.78$. Tomamos una m.a.s. $(X_1, \ldots, X_n)$ con $n=20$ y tomamos la media muestral:
    \begin{equation*}
        \overline{X} = \dfrac{1}{n}\sum_{i=1}^{n}{(X_i - \mu)}^{2}
    \end{equation*}
    Queremos calcular:
    \begin{equation*}
        P\left[\overline{X}>190\right]
    \end{equation*}
    Para ello, sabemos por lo visto en teoría que $\overline{X}\rightsquigarrow\cc{N}\left(\mu, \frac{\sigma^2}{n}\right)$, por lo que:
    \begin{align*}
        P\left[\overline{X}>190\right] &\stackrel{\text{tipificamos}}{=} P\left[Z > \dfrac{\sqrt{20}(190-185.6)}{12.78}\right] = P[Z > 1.5397] \\ 
                                                                     &= 1-P[Z<1.5397] \approx 1-0.93822 = 0.06178
    \end{align*}
\end{ejercicio}

\begin{ejercicio}
    ¿De qué tamaño mínimo habría que seleccionar una muestra de una variable con distribución normal $\cc{N}(\mu, 4)$ para poder afirmar, con probabilidad mayor que $0.9$, que la media muestral diferirá de la poblacional menos de $0.1$?\\

    \noindent
    Supuesto que tenemos una m.a.s. $(X_1, \ldots, X_n)$ de tamaño $n\in \mathbb{N}$ de variables aleatorias idénticamente distribuidas a $X\rightsquigarrow\cc{N}(\mu, 4)$, queremos buscar el menor $n$ de forma que:
    \begin{equation*}
        0.9 < P\left[|\overline{X}-\mu| < 0.1\right]
    \end{equation*}
    Para ello, usaremos que hemos visto en teoría que $\overline{X}\rightsquigarrow\cc{N}\left(\mu, \frac{\sigma^2}{n}\right)$:
    \begin{align*}
        0.9 < P[|\overline{X}-\mu| < 0.1] &= P[-0.1 < \overline{X}-\mu < 0.1] = P\left[\dfrac{-\sqrt{n}0.1}{2} < Z < \dfrac{\sqrt{n}0.1}{2}\right] \\ 
                                             &= P\left[-\sqrt{n}0.05 < Z<\sqrt{n}0.05\right] = 2P\left[Z<\sqrt{n}0.05\right] - 1 
    \end{align*}
    De donde:
    \begin{equation*}
        0.95 = \dfrac{0.9+1}{2} < P[Z<\sqrt{n}0.05]
    \end{equation*}
    De donde mirando la tabla de la normal $\cc{N}(0,1)$:
    \begin{equation*}
        \sqrt{n}0.05 \geq 1.65 \Longrightarrow n \geq {\left(\dfrac{1.65}{0.05}\right)}^{2} = 1089
    \end{equation*}
    Por lo que habría que seleccionar como mínimo una muestra de tamaño $1089$ para poder afirmar con probabilidad mayor que $0.9$ que la media muestral difiere de la poblacional menos de $0.1$.
\end{ejercicio}

\begin{ejercicio}
    Sea $(X_1, \ldots, X_n)$ una muestra aleatoria simple de una variable con distribución normal.  Calcular la probabilidad de que la cuasivarianza muestral sea menor que un $50\%$ de la varianza poblacional para $n = 16$ y para $n = 1000$.\\

    \noindent
    Dicha muestra aleatoria simple es de variables idénticamente distribuidas a $X\rightsquigarrow\cc{N}(\mu, \sigma^2)$. Bajo estas condiciones, sabemos por lo visto en teoría que:
    \begin{equation*}
        \dfrac{(n-1)S^2}{\sigma^2} \rightsquigarrow\chi^2(n-1)
    \end{equation*}
    Y queremos calcular:
    \begin{align*}
        P[S^2 < 0.5\sigma^2] = P\left[\dfrac{(n-1)S^2}{\sigma^2} < \dfrac{(n-1)0.5\cancel{\sigma^2}}{\cancel{\sigma^2}}\right] = 1-P\left[\dfrac{(n-1)S^2}{\sigma^2} > (n-1)0.5\right]
    \end{align*}
    \begin{itemize}
        \item Para $n=16$, mirando en la tabla de $\chi^2(15)$ tenemos que:
            \begin{equation*}
                P[S^2 < 0.5\sigma^2] = 1-P\left[\dfrac{(n-1)S^2}{\sigma^2} > 7.5\right] \approx 1-0.95 = 0.05
            \end{equation*}
        \item Para $n=1000$, no disponemos de tabla para $\chi^2(999)$, por lo que aproximaremos $\chi^2(999)$ por $\cc{N}(999, 2\cdot 999)\equiv \cc{N}(999, 1998)$. De esta forma, si tomamos:
            \begin{equation*}
                \dfrac{(n-1)S^2}{\sigma^2}\rightsquigarrow \chi^2(999) \approx \cc{N}(999, 1998)
            \end{equation*}
            Tenemos entonces que:
            \begin{align*}
                P[S^2 < 0.5\sigma^2] &= P\left[\dfrac{(n-1)S^2}{\sigma^2} < \dfrac{(n-1)0.5\cancel{\sigma^2}}{\cancel{\sigma^2}}\right] = P\left[\dfrac{(n-1)S^2}{\sigma^2} < 499.5\right]
                \approx \\& \stackrel{\text{tipificamos}}{\approx}
                P\left[Z < \dfrac{499.5-999}{\sqrt{1998}}\right] = P[Z < -11.17]
                =\\&= 1-P[Z<11.17] \approx 1-1 = 0
            \end{align*}
    \end{itemize}
\end{ejercicio}

\begin{ejercicio}
    Sean $S_1^2$ y $S_2^2$ las cuasivarianzas muestrales de dos muestras independientes de tamaños $n_1 = 5$ y $n_2 = 4$ de dos poblaciones normales con la misma varianza. Calcular la probabilidad de que $\nicefrac{S_1^2}{S_2^2}$ sea menor que $5.34$ o mayor que $9.12$.\\

    \noindent
    Bajo estas hipótesis, sabemos por lo visto en teoría que (como $\sigma_1=\sigma_2$):
    \begin{equation*}
        \dfrac{S_1^2}{S_2^2} \rightsquigarrow F(n_1-1, n_2-1) \equiv F(4,3)
    \end{equation*}
    Por tanto, notando $Y = \frac{S_1^2}{S_2^2}$ por comodidad, calculamos (son sucesos disjuntos):
    \begin{equation*}
        P[Y < 5.34 \text{\ o\ } Y>9.12] = P[Y<5.34] + P[Y>9.12]
    \end{equation*}
    Observamos la tabla de $F(4,3)$:
    \begin{align*}
        P[Y<5.34] &\approx 0.9 \\
        P[Y>9.12] &\approx 1-P[Y<9.12] = 1-0.95 = 0.05
    \end{align*}
    En definitiva:
    \begin{equation*}
        P[Y < 5.34 \text{\ o\ } Y>9.12] = P[Y<5.34] + P[Y>9.12] \approx 0.95
    \end{equation*}
\end{ejercicio}

\begin{ejercicio}
    Se consideran dos poblaciones de bombillas cuyas longitudes de vida siguen una ley normal con la misma media y desviaciones típicas $425$ y $375$ horas, respectivamente.  Con objeto de realizar un estudio comparativo de ambas poblaciones, se considera una muestra aleatoria simple de $10$ bombillas en la primera población y una de tamaño 6 en la segunda. ¿Cuál es la probabilidad de que la media muestral del primer grupo menos la del segundo sea menor que la observada en dos realizaciones muestrales que dieron $1325$ horas y $1215$ horas, respectivamente?\\

    \noindent
    Como describe el enunciado, tenemos una m.a.s. $(X_1, \ldots, X_n)$ de tamaño $n=10$ de variables aleatorias idénticamente distribuidas a $X\rightsquigarrow\cc{N}\left(\mu, {425}^{2}\right)$; y otra m.a.s. $(Y_1, \ldots, Y_m)$ de tamaño $m=6$ de variables aleatorias idénticamente distribuidas a $Y\rightsquigarrow\cc{N}\left(\mu, {375}^{2}\right)$. Nos piden calcular:
    \begin{equation*}
        P[\overline{X}-\overline{Y}<1325-1215]
    \end{equation*}
    
    
    Bajo estas condiciones, sabemos por lo visto en teoría que:
    \begin{equation*}
        \dfrac{\overline{X}-\overline{Y}}{\sqrt{\dfrac{\sigma_1^2}{n_1}+\dfrac{\sigma_2^2}{n_2}}} = \dfrac{\overline{X}-\overline{Y}-(\mu - \mu)}{\sqrt{\dfrac{\sigma_1^2}{n_1}+\dfrac{\sigma_2^2}{n_2}}}  \rightsquigarrow \cc{N}(0,1)
    \end{equation*}

    De esta forma, si notamos a esta última variable como $Z$ (pues es una normal estándar), tenemos que:
    \begin{align*}
        P[\overline{X}-\overline{Y}<1325-1215] &= P\left[Z < \dfrac{1325-1215}{\sqrt{\dfrac{\sigma_1^2}{n_1}+\dfrac{\sigma_2^2}{n_2}}}\right] = P\left[Z <  \dfrac{1325-1215}{\sqrt{\dfrac{{425}^{2}}{10}+\dfrac{{375}^{2}}{6}}}\right] \\
                                               &= P[Z<0.5399] \approx 0.7054
    \end{align*}
\end{ejercicio}

\begin{ejercicio}
    Sean $X_1, \ldots, X_n, X_{n+1}$ variables aleatorias independientes e idénticamente distribuidas según una $\cc{N}(\mu, \sigma^2)$, y sean $\overline{X}$ y $S^2$ la media y la cuasivarianza muestral de $(X_1, \ldots, X_n)$. Calcular la distribución de
    \begin{equation*}
        \dfrac{X_{n+1}- \overline{X}}{S}\sqrt{\dfrac{n}{n+1}}
    \end{equation*}

    \noindent
    Bajo dichas hipótesis, tenemos:
    \begin{equation*}
        \left.\begin{array}{r}
            X_{n+1} \rightsquigarrow \cc{N}(\mu, \sigma^2) \\
            \overline{X} \rightsquigarrow \cc{N}\left(\mu, \frac{\sigma^2}{n}\right)
    \end{array}\right\}\Longrightarrow X_{n+1}-\overline{X} \rightsquigarrow\cc{N}\left(\mu-\mu, \sigma^2 + \frac{\sigma^2}{n}\right) \equiv \cc{N}\left(0, \sigma^2\left(1+\frac{1}{n}\right)\right)
    \end{equation*}
    de donde, si tipificamos la variable:
    \begin{equation*}
        \dfrac{X_{n+1}-\overline{X}}{\sqrt{\sigma^2\left(1+\frac{1}{n}\right)}} = \dfrac{X_{n+1}-\overline{X}}{\sigma\sqrt{\dfrac{n+1}{n}}} = \dfrac{X_{n+1}-\overline{X}}{\sigma}\sqrt{\dfrac{n}{n+1}}\rightsquigarrow \cc{N}(0,1)
    \end{equation*}

    Notemos que esta distribución es similar a la que buscamos, salvo que en el denominador tenemos $\sigma$ en lugar de $S$. Podemos emplear el siguiente resultado, visto en teoría, que relaciona la cuasivarianza muestral con la varianza poblacional:
    \begin{equation*}
        \dfrac{(n-1)S^2}{\sigma^2} \rightsquigarrow\chi^2(n-1)
    \end{equation*}

    En vistas de aplicar la cosntrucción de la distribución $t$ de Student, necesitamos que las variables aleatorias con las que trabajamos sean independientes. La primera es función de $\ol{X}$ y $X_{n+1}$, mientras que la segunda es función de $S^2$. Por tanto, nos basta con ver que $(\overline{X}, X_{n+1})$ es independiente de $S^2$. Esto se cumple, ya que $\overline{X}$ y $S^2$ son independientes por el Lema de Fisher, y $X_{n+1}$ es independiente de $(X_1, \ldots, X_n)$, por lo que $(\overline{X}, X_{n+1})$ es independiente de $S^2$. Podemos por tanto aplicar la construccion de la distribución $t$ de Student:
    \begin{equation*}
        \left.\begin{array}{l}
                U\rightsquigarrow \cc{N}(0,1) \\
                V\rightsquigarrow\chi^2(n)\\
                U \text{ y } V \text{ independientes}
        \end{array}\right\} \Longrightarrow \dfrac{U}{\sqrt{\nicefrac{V}{n}}} \rightsquigarrow t(n)
    \end{equation*}
    Si lo aplicamos a nuestras variables:
    \begin{equation*}
        \dfrac{\dfrac{X_{n+1}-\overline{X}}{\sigma}\sqrt{\dfrac{n}{n+1}}}{\sqrt{\dfrac{(n-1)S^2}{\sigma^2(n-1)}}} = \dfrac{\dfrac{X_{n+1}-\overline{X}}{\sigma}\sqrt{\dfrac{n}{n+1}}}{\dfrac{S}{\sigma}} = \dfrac{X_{n+1}- \overline{X}}{S}\sqrt{\dfrac{n}{n+1}} \rightsquigarrow t(n-1)
    \end{equation*}
\end{ejercicio}

\begin{ejercicio}
    Sean $(X_1, \ldots, X_n)$, $(Y_1, \ldots, Y_m)$ muestras aleatorias simples independientes de poblaciones $\cc{N}(\mu_1,\sigma^2)$ y $\cc{N}(\mu_2, \sigma^2)$, respectivamente. Sean $\alpha,\beta\in \mathbb{R}$ y $\overline{X},\overline{Y},S_1^2, S_2^2$ las medias y cuasivarianzas de las dos muestras. Calcular la distribución de
    \begin{equation*}
        \dfrac{\alpha\left(\overline{X}-\mu_1\right)+\beta\left(\overline{Y}-\mu_2\right)}{\sqrt{\dfrac{(n-1)S_1^2 + (m-1)S_2^2}{n+m-2}} \sqrt{\dfrac{\alpha^2}{n} + \dfrac{\beta^2}{m}}}
    \end{equation*}

    \noindent
    Sabemos que:
    \begin{equation*}
        \left.\begin{array}{r}
                \overline{X} \rightsquigarrow\cc{N}\left(\mu_1, \frac{\sigma^2}{n}\right) \\
                \overline{Y} \rightsquigarrow\cc{N}\left(\mu_2, \frac{\sigma^2}{m}\right) 
        \end{array}\right\} \Longrightarrow 
        \left.\begin{array}{r}
                \overline{X}-\mu_1 \rightsquigarrow\cc{N}\left(0, \frac{\sigma^2}{n}\right) \\
                \overline{Y}-\mu_2 \rightsquigarrow\cc{N}\left(0, \frac{\sigma^2}{m}\right) 
        \end{array}\right\} \Longrightarrow 
        \left.\begin{array}{r}
                \alpha(\overline{X}-\mu_1 )\rightsquigarrow\cc{N}\left(0, \frac{\alpha^2\sigma^2}{n}\right) \\
                \beta(\overline{Y}-\mu_2 )\rightsquigarrow\cc{N}\left(0, \frac{\beta^2\sigma^2}{m}\right) 
        \end{array}\right\} 
    \end{equation*}
    Como ambas son independientes, podemos aplicar la propeidad reproductiva de la normal:
    \begin{equation*}
        \alpha(\overline{X}-\mu_1 ) +\beta(\overline{Y}-\mu_2 ) \rightsquigarrow \cc{N}\left(0, \dfrac{\alpha^2\sigma^2}{n} + \dfrac{\beta^2\sigma^2}{m}\right) \equiv \cc{N}\left(0, \sigma^2\left(\frac{\alpha^2}{n}+\frac{\beta^2}{m}\right) \right) 
    \end{equation*}
    Si tipificamos la variable:
    \begin{equation*}
        \dfrac{\alpha(\overline{X}-\mu_1 ) +\beta(\overline{Y}-\mu_2) }{\sigma\sqrt{\frac{\alpha^2}{n}+\frac{\beta^2}{m}}} \rightsquigarrow\cc{N}(0,1)
    \end{equation*}
    Por otra parte, tenemos que:
    \begin{equation*}
        \dfrac{(n-1)S_1^2}{\sigma^2} \rightsquigarrow\chi^2(n-1), \qquad 
        \dfrac{(m-1)S_2^2}{\sigma^2} \rightsquigarrow\chi^2(m-1) 
    \end{equation*}
    Como ambas son independientes por ser $S_1^2$ y $S_2^2$ independientes (ya que las muestras aleatorias simples eran independientes), podemos aplicar la propiedad reproductiva de $\chi^2$, obteniendo que:
    \begin{equation*}
        \dfrac{(n-1)S_1^2}{\sigma^2} + \dfrac{(m-1)S_2^2}{\sigma^2} = \dfrac{(n-1)S_1^2+(m-1)S_2^2}{\sigma^2}\rightsquigarrow\chi^2(n+m-2) 
    \end{equation*}
    Finalmente, la extensión del Lema de Fisher nos dice que $(\overline{X},\overline{Y})$ es independiente de $(S_1^2, S_2^2)$, por lo que las dos variables aleatorias con las que trabajamos son independientes, lo que nos permite aplicar la construcción de la distribución $t$ de Student:
    \begin{equation*}
        \dfrac{\dfrac{\alpha(\overline{X}-\mu_1 ) +\beta(\overline{Y}-\mu_2) }{\cancel{\sigma}\sqrt{\frac{\alpha^2}{n}+\frac{\beta^2}{m}}}}{\sqrt{\dfrac{\dfrac{(n-1)S_1^2+(m-1)S_2^2}{\cancel{\sigma^2}}}{n+m-2}}} = \dfrac{\alpha\left(\overline{X}-\mu_1\right)+\beta\left(\overline{Y}-\mu_2\right)}{\sqrt{\dfrac{(n-1)S_1^2 + (m-1)S_2^2}{n+m-2}} \sqrt{\dfrac{\alpha^2}{n} + \dfrac{\beta^2}{m}}} \rightsquigarrow t(n+m-2)
    \end{equation*}
\end{ejercicio}


\begin{ejercicio}
    Sean $X_1, \ldots, X_{10}$ una muestra aleatoria simple de una variable con distribución normal $\cc{N}(5, \sigma^2)$ y $Y_1, \ldots, Y_{10}$ una muestra aleatoria simple de una variable con distribución normal $\cc{N}(\mu_i, \sigma^2)$, siendo ambas muestras independientes. Calcular:
    \begin{equation*}
        P\left[\dfrac{\sum\limits_{i=1}^{10}{(X_i-5)}^{2}}{S_Y^2} > 31.4\right]
    \end{equation*}

    \noindent
    Puesto que $E[\ol{X}]=E[X]=5$, podemos escribir:
    \begin{equation*}
        P\left[\dfrac{\sum\limits_{i=1}^{10}{(X_i-5)}^{2}}{S_Y^2} > 31.4\right]
        = P\left[\dfrac{\sum\limits_{i=1}^{10}{(X_i-5)}^{2}}{9\cdot S_Y^2} > \dfrac{31.4}{9}\right]
        = P\left[\dfrac{S_X^2}{S_Y^2} > \dfrac{31.4}{9}\right]
    \end{equation*}

    Como sabemos que $\dfrac{S_X^2}{S_Y^2} \rightsquigarrow F(9,9)$, tenemos que:
    \begin{equation*}
        P\left[\dfrac{S_X^2}{S_Y^2} > \dfrac{31.4}{9}\right] = 1-P\left[\dfrac{S_X^2}{S_Y^2} < 3.49\right] \approx 1-0.95 = 0.05
    \end{equation*}
\end{ejercicio}



    \newpage
\section{Suficiencia y completitud}

\begin{ejercicio}
    Sea $(X_1, \ldots, X_n)$ una muestra aleatoria simple de una variable $X\rightsquigarrow\{B(k,p) : p\in \left]0,1\right[\}$ y sea $T(X_1, \ldots, X_n) = \sum\limits_{i=1}^n X_i$. Probar que $T$ es suficiente para $p$:
    \begin{enumerate}[label=\alph*)]
        \item usando la definición
        \item aplicando el teorema de factorización
    \end{enumerate}

    \begin{enumerate}[label=\alph*)]
        \item Si notamos para abreviar $T=T(X_1, \ldots, X_n)$, tenemos que probar que la distribución de la muestra condicionada a cualquier valor del estadístico no depende del parámetro $p$ para probar que $T$ es suficiente para $p$. Para ello:
            \begin{align*}
                P_p[X_1 = x_1, \ldots, X_n = x_n \ &|\ T=t] = \dfrac{P_p[X_1=x_1, \ldots, X_n = x_n, T=t]}{P_p[T=t]} \\
                                                           &= \left\{\begin{array}{ll}
                                                               0 & \text{si\ } T(x_1,\ldots,x_n)\neq t \\
                                                               \dfrac{P_p[X_1=x_1, \ldots, X_n=x_n]}{P_p[T=t]}& \text{si\ } T(x_1,\ldots,x_n) = t
                                                           \end{array}\right. 
            \end{align*}
            Como $0$ obviamente no depende de $p$, el caso $T(x_1,\ldots,x_2)\neq t$ se encuentra ya estudiado, por lo que nos centramos en el caso $T(x_1,\ldots,x_2)=t$:
            \begin{align*}
                \dfrac{P_p[X_1=x_1, \ldots, X_n = x_n]}{P_p[T=t]} &\stackrel{\text{iid.}}{=} \dfrac{\prod\limits_{i=1}^{n}P_p[X=x_i]}{P_p[T=t]} \AstIg \dfrac{\prod\limits_{i=1}^n \binom{k}{x_i} p^{x_i} {(1-p)}^{k-x_i}}{\binom{nk}{t}p^t{(1-p)}^{nk-t}} \\
                                                                  &= \dfrac{p^{\sum\limits_{i=1}^n x_i} {(1-p)}^{nk - \sum\limits_{i=1}^n x_i}\prod\limits_{i=1}^n \binom{k}{x_i}}{\binom{nk}{t}p^t{(1-p)}^{nk-t}} \stackrel{(\ast\ast)}{=} \dfrac{\prod\limits_{i=1}^{n}\binom{k}{x_i}}{\binom{nk}{t}}
            \end{align*}
            Donde en $(\ast)$ hemos usado que $(X_1, \ldots, X_n)$ es una m.a.s. (variables independientes) y la reproductividad de la Binomial, por lo que $T\rightsquigarrow B(nk,p)$; y en $(\ast\ast)$ hemos usado que $t=T(x_1, \ldots, x_n) = \sum\limits_{i=1}^{n}x_i$. En definitiva, hemos obtenido que la distribución de la muestra condicionada a cualquier valor del estadístico no dependende del parámetro $p$, por lo que $T$ es suficiente para $p$.
        \item Si podemos usar el Teorema de factorización, escribimos la función masa de probabilidad de la distribución conjunta de la muestra aleatoria simple:
            \begin{align*}
                P[X_1=x_1, \ldots, X_n = x_n] &\stackrel{\text{iid.}}{=} \prod_{i=1}^{n} P[X=x_i] = \prod_{i=1}^{n} \binom{k}{x_i} p^{x_i} {(1-p)}^{k-x_i} \\
                                              &= p^{\sum\limits_{i=1}^n x_i} {(1-p)}^{nk - \sum\limits_{i=1}^n x_i} \prod_{i=1}^{n} \binom{k}{x_i} 
            \end{align*}
            Si tomamos:
            \begin{equation*}
                h(x_1, \ldots, x_n) = \prod_{i=1}^{n}\binom{k}{x_i}, \qquad 
                T(X_1, \ldots, X_n) = \sum_{i=1}^{n} X_i, \qquad 
                g_p(t) = p^t {(1-p)}^{nk - t}
            \end{equation*}
            Podemos aplicar el Teorema de Factorización de Neymann-Fisher, obteniendo que $T$ es un estadístico suficiente para $p$.
    \end{enumerate}
\end{ejercicio}

\begin{ejercicio}
    Sea $(X_1, \ldots, X_n)$ una muestra aleatoria simple de una variable $X\rightsquigarrow\{\cc{P}(\lm) : \lm\in \mathbb{R}^+\}$ y sea $T(X_1, \ldots, X_n) = \sum\limits_{i=1}^n X_i$. Probar que $T$ es suficiente para $\lm$:
    \begin{enumerate}[label=\alph*)]
        \item usando la definición.
        \item aplicando el teorema de factorización.
    \end{enumerate}
    

    \begin{enumerate}[label=\alph*)]
        \item Notando $T=T(X_1, \ldots, X_n)$, seguimos los mismos pasos que en el ejercicio anterior:
            \begin{align*}
                P_\lm[X_1 = x_1, \ldots, X_n = x_n \ &|\ T=t] = \dfrac{P_\lm[X_1=x_1, \ldots, X_n = x_n, T=t]}{P_\lm[T=t]} \\
                                                           &= \left\{\begin{array}{ll}
                                                               0 & \text{si\ } T(x_1,\ldots,x_n)\neq t \\
                                                               \dfrac{P_\lm[X_1=x_1, \ldots, X_n=x_n]}{P_\lm[T=t]}& \text{si\ } T(x_1,\ldots,x_n) = t
                                                           \end{array}\right. 
            \end{align*}
            Y ahora nos interesamos por el segundo término, que es el que puede depender de $\lm$:
            \begin{align*}
                \dfrac{P_\lm[X_1 = x_1, \ldots, X_n = x_n]}{P_\lm[T=t]} &\stackrel{\text{iid.}}{=} \dfrac{\prod\limits_{i=1}^{n} P_\lm[X=x_i]}{P_\lm[T=t]} \AstIg \dfrac{\prod\limits_{i=1}^n e^{-\lm} \dfrac{\lm^{x_i}}{x_i!}}{e^{-n\lm} \dfrac{{(n\lm)}^{t}}{t!}} = \dfrac{e^{-n\lm} \cdot \dfrac{\lm^{\sum\limits_{i=1}^n x_i}}{\prod\limits_{i=1}^n x_i!}}{e^{-n\lm}\cdot  \dfrac{n^t \lm^t}{t!}} \\
                                                                        &\stackrel{(\ast\ast)}{=} \dfrac{t!}{n^t \prod\limits_{i=1}^{n}x_i!}
            \end{align*}
            Donde en $(\ast)$ usamos que $T\rightsquigarrow \cc{P}(n\lm)$, por la reproductividad de la Poisson y en $(\ast\ast)$ usamos que $t = T(x_1, \ldots, x_n) = \sum\limits_{i=1}^{n} x_i$. Obtenemos una cantidad que no depende de $\lm$, por lo que $T$ es suficiente para $\lm$.
        \item Si podemos aplicar el Teorema de factorización, escribimos:
            \begin{equation*}
                P_\lm[X_1=x_1, \ldots, X_n = x_n] \stackrel{\text{iid.}}{=} \prod_{i=1}^{n} P_\lm[X=x_i] = \prod_{i=1}^{n}e^{-\lm}\dfrac{\lm^{x_i}}{x_i!} = e^{-n\lm} \cdot \dfrac{\lm^{\sum\limits_{i=1}^{n} x_i}}{\prod\limits_{i=1}^{n}x_i!}
            \end{equation*}
            Si tomamos:
            \begin{equation*}
                h(x_1, \ldots, x_n) = \dfrac{1}{\prod\limits_{i=1}^{n}x_i!}, \qquad T(X_1, \ldots, X_n) = \sum_{i=1}^{n}X_i, \qquad g_\lm(t) = e^{-n\lm} \cdot \lm^t
            \end{equation*}
            Por el Teorema de Factorización de Neymann-Fisher obtenemos que $T$ es suficiente para $\lm$.
    \end{enumerate}
\end{ejercicio}

\begin{ejercicio}
    Sea $(X_1, X_2, X_3)$ una muestra aleatoria simple de una variable $X\rightsquigarrow\{B(1,p) : p \in \left]0,1\right[\}$. Probar que el estadístico $X_1 + 2X_2 + 3X_3$ no es suficiente.\\

    \noindent
    Consideramos el estadístico $T(X_1, X_2, X_3) = X_1 + 2X_2 + 3X_3$. El espacio muestral de $X$ es $\cc{X}=\{0,1\}$, por lo que el espacio muestral de $T$ es $\cc{T} = \{0, 1, 2, 3, 4, 5, 6\}$. Sabemos por un ejemplo visto en teoría que el ``truco'' para demostrar que $T(X_1, X_2,X_3)$ no es suficiente para $p$ es buscar un valor del espacio muestral $\cc{T}$ que provenga de varias combinaciones de estados del espacio muestral $\cc{X}^3$. Como $0,1,2$ solo provienen de una combinación del espacio muestral $\cc{X}^3$ (son $(0,0,0),\ (1,0,0)$ y $(0,1,0)$ respectivamente), probamos buscar el contraejemplo con $t=3$, que proviene de considerar las observaciones de la muestra $(1,1,0)$ y $(0,0,1)$.

    Una vez explicado el procedimiento para buscar cuál es el valor de $t$ que funciona, procedemos a probar que la distribución de la muestra condicionada a dicho valor de $t$ depende del parámetro $p$. Para ello:
            \begin{align*}
                P_p[X_1 = x_1, X_2=x_2, X_3 = x_3 \ &|\ T=3] = \dfrac{P_p[X_1=x_1, X_2=x_2, X_3 = x_3, T=3]}{P_p[T=3]} \\
                                                           &= \left\{\begin{array}{ll}
                                                               0 & \text{si\ } T(x_1,x_2,x_2)\neq 3 \\
                                                               \dfrac{P_p[X_1=x_1, X_2=x_2, X_3=x_3]}{P_p[T=3]}& \text{si\ } T(x_1,x_2,x_3) = 3
                                                           \end{array}\right. 
            \end{align*}
            Vemos que el primer caso no puede depender nunca de $p$, por lo que buscamos probar que el segundo caso sí que depende de $p$. Para ello:
            \begin{align*}
                \dfrac{P_p[X_1=x_1, X_2=x_2, X_3=x_3]}{P_p[T=3]} &\stackrel{\text{iid.}}{=} \dfrac{P_p[X=x_1]P_p[X=x_2]P_p[X=x_3]}{P_p[T=3]} \\
                                                                 &\AstIg \dfrac{p^{x_1}{(1-p)}^{1-x_1} p^{x_2}{(1-p)}^{1-x_2}p^{x_3}{(1-p)}^{1-x_3}}{P_p[X_1=1,X_2=1,X_3=0]+P_p[X_1=0,X_2=0,X_3=1]} \\
                                                                 &= \dfrac{p^{x_1}{(1-p)}^{1-x_1} p^{x_2}{(1-p)}^{1-x_2}p^{x_3}{(1-p)}^{1-x_3}}{p\cdot p\cdot (1-p) + (1-p)(1-p)p} \\
            \end{align*}
            Donde en $(\ast)$ se ha usado la propiedad reproductiva de la Binomial, por lo que $T\rightsquigarrow B(3,p)$, así como que la condición $t=3$ provenía de los valores $(1,1,0)$ y $(0,0,1)$.
            Ahora, si tomamos $(x_1,x_2,x_3) = (1,1,0)$, tenemos que:
            \begin{align*}
                \dfrac{P_p[X_1=1, X_2=1, X_3=0]}{P_p[T=3]} &= \dfrac{p^{1}{(1-p)}^{0} p^{1}{(1-p)}^{0}p^{0}{(1-p)}^{1}}{p\cdot p\cdot (1-p) + (1-p)(1-p)p} \\
                                                                 &= \dfrac{pp(1-p)}{pp(1-p)+(1-p)(1-p)p} \\
                                                                 &= \dfrac{p^2(1-p)}{p(1-p)\cancelto{1}{(p+1-p)}} = p
            \end{align*}
            Que claramente depende de $p$, por lo que $T$ no es suficiente para $p$.
\end{ejercicio}

\begin{ejercicio}
    Aplicando el teorema de factorización, y basándose en una muestra de tamaño arbitrario, encontrar un estadístico suficiente para cada una de las siguientes familias de distribuciones (en las familias biparamétricas, suponer los casos de sólo un parámetro desconocido y de los dos desconocidos).
    \begin{enumerate}[label=\alph*)]
        \item $X\rightsquigarrow\{\cc{U}(\nicefrac{-\theta}{2},\nicefrac{\theta}{2}) : \theta>0\}$
        \item $X\rightsquigarrow\{\Gamma(p,a) : p,a > 0\}$
        \item $X\rightsquigarrow\{\beta(p,q) : p,q>0\}$
        \item $X\rightsquigarrow\{P_{N_1,N_2} : N_1, N_2 \in \mathbb{N}, N_1 \leq N_2\}$ y la masa de probabilidad viene dada por:
            \begin{equation*}
                P_{N_1, N_2}[X=x] = \dfrac{1}{N_2 - N_1 + 1} \qquad x\in \{N_1, \ldots, N_2\}
            \end{equation*}
    \end{enumerate}

    \noindent
    En lo que sigue, supondremos que tenemos una m.a.s. $(X_1, \ldots, X_n)$ de variables idénticamente distribuidas a la respectiva variable $X$:
    \begin{enumerate}[label=\alph*)]
        \item Si $X\rightsquigarrow \cc{U}(\nicefrac{-\theta}{2},\nicefrac{\theta}{2})$ con $\theta>0$:
            \begin{equation*}
                f(x_1, \ldots, x_n) \stackrel{\text{iid.}}{=}\prod_{i=1}^{n}f(x_i) = \prod_{i=1}^{n} \dfrac{1}{\theta} = \dfrac{1}{\theta^n} \qquad x_i \in \left]\nicefrac{-\theta}{2},\nicefrac{\theta}{2}\right[, \quad \forall i \in \{1,\ldots,n\}
            \end{equation*}
            Por lo que tendremos $\nicefrac{-\theta}{2}<X_{(1)}\leq X_{(n)} < \nicefrac{\theta}{2}$:
            \begin{equation*}
                f(x_1, \ldots, x_n) = \dfrac{1}{\theta^n} \ind_{\left]0,+\infty\right[}(X_{(1)}+\nicefrac{\theta}{2}) \ind_{\left]-\infty,0\right[}(X_{(n)}-\nicefrac{\theta}{2})
            \end{equation*}
            Si tomamos:
            \begin{gather*}
                h(x_1,\ldots,x_n)=1,\qquad T(X_1, \ldots, X_n) = (X_{(1)}, X_{(n)}) \\ g_{\theta}(t_1,t_2) = \dfrac{1}{\theta^n} \ind_{\left]0,+\infty\right[}(t_1 + \nicefrac{\theta}{2}) \ind_{\left]-\infty,0\right[}(t_2-\nicefrac{\theta}{2})
            \end{gather*}
            Por el Teorema de factorización, tenemos que $T(X_1, \ldots, X_n)$ es un estadístico suficiente para $\theta$.

            % // TODO: Unidimensional?
        \item Si $X\rightsquigarrow\Gamma(p,a)$ con $p,a>0$:
            \begin{align*}
                f(x_1, \ldots, x_n) &\stackrel{\text{iid.}}{=} \prod_{i=1}^{n}f(x_i) = \prod_{i=1}^{n} \dfrac{a^p}{\Gamma(p)} x_i^{p-1} e^{-ax_i} = {\left(\dfrac{a^p}{\Gamma(p)}\right)}^{n} e^{-a\sum\limits_{i=1}^{n}x_i} \prod_{i=1}^{n}x_i^{p-1} \\ 
                                    &= {\left(\dfrac{a^p}{\Gamma(p)}\right)}^{n} e^{-a\sum\limits_{i=1}^{n}x_i} {\left(\prod_{i=1}^{n}x_i\right)}^{p-1} \qquad x_i \geq 0 \qquad \forall i \in \{1,\ldots,n\}
            \end{align*}
            \begin{itemize}
                \item Suponiendo que $p$ es conocida, podemos tomar:
                    \begin{gather*}
                        h(x_1, \ldots, x_n) = {\left(\prod_{i=1}^{n}x_i\right)}^{p-1}, \qquad  T(X_1, \ldots, X_n) = \sum_{i=1}^{n} X_i \\
                        g_a(t) = {\left(\dfrac{a^p}{\Gamma(p)}\right)}^{n}e^{-at}
                    \end{gather*}
                \item Suponiendo ahora que $a$ es conocida:
                    \begin{gather*}
                        h(x_1, \ldots, x_n) = e^{-a\sum\limits_{i=1}^n x_i}, \qquad T(X_1, \ldots, X_n) = \prod_{i=1}^{n}x_i \\
                        g_p(t) = {\left(\dfrac{a^p}{\Gamma(p)}\right)}^{n} t^{p-1}
                    \end{gather*}
                \item Si ahora tanto $p$ como $a$ son desconocidas, podemos tomar:
                    \begin{gather*}
                        h(x_1, \ldots, x_n) = 1, \qquad T(X_1, \ldots, X_n) = \left(\sum_{i=1}^{n}X_i, \prod_{i=1}^{n} X_i\right) \\
                        g_{(a,p)}(t_1,t_2) = {\left(\dfrac{a^p}{\Gamma(p)}\right)}^{n}e^{-at_1}t_2^{p-1}
                    \end{gather*}
            \end{itemize}
        \item Si $X\rightsquigarrow\beta(p,q)$ con $p,q>0$:
            \begin{align*}
                f(x_1, \ldots, x_n) &\stackrel{\text{iid.}}{=} \prod_{i=1}^{n}f(x_i) = \prod_{i=1}^{n} \dfrac{1}{\beta(p,q)} x_i^{p-1} {(1-x_i)}^{q-1} \\
                &= \dfrac{1}{{\beta(p,q)}^{n}} {\left(\prod_{i=1}^{n}x_i\right)}^{p-1} {\left(\prod_{i=1}^{n}(1-x_i)\right)}^{q-1} \quad x_i \in [0,1], \forall i \in \{1,\ldots,n\}
            \end{align*}
            \begin{itemize}
                \item Si $p$ es conocida, tomamos:
                    \begin{gather*}
                        h(x_1, \ldots, x_n) = {\left(\prod_{i=1}^{n}x_i\right)}^{p-1}, \qquad T(X_1, \ldots, X_n) = \prod_{i=1}^{n}(1-X_i) \\
                        g_q(t) = \dfrac{1}{{\beta(p,q)}^{n}} t^{q-1}
                    \end{gather*}
                \item Si $q$ es conocida:
                    \begin{gather*}
                        h(x_1, \ldots, x_n) = {\left(\prod_{i=1}^{n}(1-x_i)\right)}^{q-1}, \qquad T(X_1, \ldots, X_n) = \prod_{i=1}^{n}X_i \\
                        g_p(t) = \dfrac{1}{{\beta(p,q)}^{n}} t^{p-1}
                    \end{gather*}
                \item Si tanto $p$ como $q$ son parámetros:
                    \begin{gather*}
                        h(x_1, \ldots, x_n) = 1, \qquad T(X_1, \ldots, X_n) = \left(\prod_{i=1}^{n}X_i, \prod_{i=1}^{n}(1-X_i)\right) \\
                        g_{(p,q)}(t_1, t_2) = \dfrac{1}{{\beta(p,q)}^{n}}t_1^{p-1}t_2^{q-1}
                    \end{gather*}
            \end{itemize}
        \item Si $X\rightsquigarrow P_{N_1,N_2}$ con $N_1,N_2\in \mathbb{N}$, $N_1\leq N_2$, entonces:
            \begin{align*}
                P[X_1=x_1, \ldots, X_n=x_n]&\stackrel{\text{iid.}}{=} \prod_{i=1}^{n} P[X=x_i] =\prod_{i=1}^{n} \dfrac{1}{N_2-N_1+1} \\
                                           &= \dfrac{1}{{(N_2-N_1+1)}^{n}}, \quad x_i \in \{N_1, \ldots N_2\} 
            \end{align*}
            Como se tiene $N_1\leq X_{(1)}\leq X_{(n)}\leq N_2$, entonces:
            \begin{equation*}
                P[X_1=x_1, \ldots, X_n=x_n] = \dfrac{\ind_{\bb{N}_0}(X_{(1)}-N_1)\ind_{-\bb{N}_0}(X_{(n)}-N_2)}{{(N_2-N_1+1)}^{n}}
            \end{equation*}
            \begin{itemize}
                \item Si $N_1$ es conocido:
                    \begin{gather*}
                        h(x_1, \ldots, x_n) = \ind_{\bb{N}_0}(X_{(1)}-N_1), \qquad T(X_1, \ldots, X_n) = X_{(n)} \\
                        g_{N_2}(t) = \dfrac{\ind_{-\bb{N}_0}(t-N_2)}{{(N_2-N_1+1)}^{n}}
                    \end{gather*}
                \item Si $N_2$ es conocida:
                    \begin{gather*}
                        h(x_1, \ldots, x_n) = \ind_{-\bb{N}_0}(X_{(n)}-N_2), \qquad T(X_1, \ldots, X_n) = X_{(1)} \\
                        g_{N_1}(t) = \dfrac{\ind_{\bb{N}_0}(t-N_1)}{{(N_2-N_1+1)}^{n}}
                    \end{gather*}
                \item Si tanto $N_1$ como $N_2$ son parámetros:
                    \begin{gather*}
                        h(x_1,\ldots,x_n) = 1, \qquad T(X_1, \ldots, X_n) = \left(X_{(1)}, X_{(n)}\right) \\
                        g_{(N_1,N_2)}(t_1,t_2) = \dfrac{\ind_{\bb{N}_0}(t_1-N_1)\ind_{-\bb{N}_0}(t_2-N_2)}{{(N_2-N_1+1)}^{n}}
                    \end{gather*}
            \end{itemize}
            Por lo que en cualquier caso obtenemos un estadístico suficiente, por el Teorema de Factorización.
    \end{enumerate}
\end{ejercicio}

\begin{ejercicio}\label{ej:3.5}
    Sea $X\rightsquigarrow\{P_N : N\in \mathbb{N}\}$, siendo $P_N$ la distribución uniforme en los puntos $\{1,\ldots,N\}$, y sea $(X_1, \ldots, X_n)$ una muestra aleatoria simple de $X$. Probar que $\max(X_1, \ldots, X_n)$ es un estadístico suficiente y completo.\\

    \noindent
    Buscamos aplicar el Teorema de factorización de Neymann-Fisher:
    \begin{equation*}
        P[X_1 = x_1, \ldots, X_n = x_n] \stackrel{\text{iid.}}{=} \prod_{i=1}^{n}P[X=x_i] =  \prod_{i=1}^{n} \dfrac{1}{N}, \quad x_i \in \{1,\ldots, N\}, \forall i \in \{1,\ldots,n\}
    \end{equation*}
    Por lo que $1\leq X_{(1)} \leq X_{(n)} \leq N$:
    \begin{equation*}
        P[X_1 = x_1, \ldots, X_n=x_n] = \dfrac{\ind_{\bb{N}_0}(X_{(1)}-1)\ind_{-\bb{N}_0}(X_{(n)}-N)}{N^n}
    \end{equation*}
    De donde podemos tomar:
    \begin{gather*}
        h(x_1, \ldots, x_n) = \ind_{\bb{N}_0}(X_{(1)}-1), \qquad T(X_1, \ldots, X_n) = X_{(n)} \\ 
        g_N(t) = \dfrac{\ind_{-\bb{N}_0}(t-N)}{N^n}
    \end{gather*}
    Por el Teorema de factorización de Neymann-Fisher, tenemos que el estadístico $T(X_1, \ldots, X_n) = X_{(n)}$ es suficiente. Comprobamos que también es completo: sea $g$ cualquier función medible, supongamos que (abreviaremos $T(X_1, \ldots, X_n) = T$):

    \begin{equation*}
        0 = E[g(T)] = \sum_{t=1}^{N}g(t)P[T=t] \qquad \forall N\in \mathbb{N}
    \end{equation*}
    Calculamos la función masa de probabilidad de $T$:
    \begin{equation*}
        F_T(t) = {(F_X(t))}^{n} \Longrightarrow P[T=t] = P[T\leq t] - P[T\leq t-1] = {(F_X(t))}^{n} - {(F_X(t-1))}^{n}
    \end{equation*}
    Como $F_X(t) = \frac{t}{N}$, tenemos entonces que:
    \begin{equation*}
        P[T=t] = {(F_X(t))}^{n} - {(F_X(t-1))}^{n} = \dfrac{t^n}{N^n} - \dfrac{{(t-1)}^{n}}{N^n} = \dfrac{t^n - {(t-1)}^{n}}{N^n}
    \end{equation*}
    Por lo que:
    \begin{equation*}
        0 = E[g(T)] = \sum_{t=1}^{N} g(t) \dfrac{t^n -{(t-1)}^{n}}{N^n} = \dfrac{1}{N^n} \sum_{t=1}^{N} g(t) (t^n - {(t-1)}^{n}) \qquad \forall N\in \mathbb{N}
    \end{equation*}

    de donde:
    \begin{equation*}
        \sum_{t=1}^{N}g(t)(t^n -{(t-1)}^{n}) = 0 \qquad \forall N\in \mathbb{N}
    \end{equation*}

    Puesto que la potencia $n-$ésima es una función estrictamente creciente en el intervalo $\left[0,+\infty\right[$, tenemos que $t^n - {(t-1)}^{n} > 0$ para todo $t\in \mathbb{N}$. Para que una suma de términos positivos sea cero, todos los sumandos han de ser cero, por lo que $g(t) = 0$ para todo $t\in \mathbb{N}$. En definitiva, tenemos que:
    \begin{equation*}
        \mathbb{N} \subseteq \{t:g(t) = 0\}
    \end{equation*}
    Por lo que:
    \begin{equation*}
        1 \geq P[g(T) = 0] \geq P[T \in \mathbb{N}] = 1 \Longrightarrow P[g(T)= 0] = 1
    \end{equation*}
    Lo que demuestra que $X_{(n)}$ es completo.

    \begin{comment}
    Probemos por inducción sobre $N$ que $g(t) = 0\quad \forall t\in \mathbb{N}$:
    \begin{itemize}
        \item \underline{Para $N=1$:} tenemos que:
            \begin{equation*}
                g(1) = g(1)(1^n-{(1-1)}^{n}) = \sum_{t=1}^{1}g(t) (t^n - {(t-1)}^{n}) = 0
            \end{equation*}
        \item \underline{Supuesto que $g(t) = 0$ para $t<N$:}
            \begin{equation*}
                g(N)(N^n - {(N-1)}^{n}) = \sum_{t=1}^{N}g(t)(t^n-{(t-1)}^{n}) = 0
            \end{equation*}
            Por lo que:
            \begin{itemize}
                \item $g(N) = 0$.
                \item $N^n - {(N-1)}^{n}=0$, que es imposible, puesto que la potencia $n-$ésima es una función estrictamente creciente en el intervalo $\left[0,+\infty\right[$.
            \end{itemize}
            
    \end{itemize}
    \end{comment}

    
\end{ejercicio}

\begin{ejercicio}
    Basándose en una muestra de tamaño arbitrario, obtener un estadístico suficiente y completo para la familia de distribuciones definidas por todas las densidades de la forma
    \begin{equation*}
        f_\theta(x) = e^{\theta-x}, \qquad x>\theta
    \end{equation*}

    \noindent
    Sea $(X_1, \ldots, X_n)$ una m.a.s. de tamaño $n\in \mathbb{N}$: 
    \begin{equation*}
        f_{\theta}(x_1, \ldots, x_n)  \stackrel{\text{iid.}}{=} \prod_{i=1}^{n}f_\theta(x_i) = \prod_{i=1}^{n}e^{\theta - x_i} = e^{n\theta - \sum\limits_{i=1}^{n}x_i} = \dfrac{e^{n\theta}}{e^{\sum\limits_{i=1}^n x_i}}, \quad x_i > \theta\ \forall i \in \{1,\ldots,n\}
    \end{equation*}
    Por lo que ha de ser $\theta < X_{(1)}$:
    \begin{equation*}
        f_{\theta}(x_1, \ldots, x_n) = \dfrac{e^{n\theta}\cdot \ind_{\left]0, \infty\right[}(X_{(1)}-\theta)}{e^{\sum\limits_{i=1}^n x_i}}\qquad x_i\in \bb{R}, \forall i \in \{1,\ldots,n\}
    \end{equation*}
    Si tomamos:
    \begin{gather*}
        h(x_1, \ldots, x_n) = e^{-\sum\limits_{i=1}^n x_i}, \qquad T(X_1, \ldots, X_n) = X_{(1)} \\
        g_\theta(t) = e^{n\theta}\cdot \ind_{\left]0, \infty\right[}(t-\theta)
    \end{gather*}
    Por el Teorema de factorización de Neymann-Fisher tenemos que el estadístico $X_{(1)}$ es suficiente para $\theta$. Comprobemos si es completo: sea $g$ cualquier función medible, suponemos que (y escribimos $T=T(X_1, \ldots, X_n)$ para abreviar):
    \begin{equation*}
        0 = E[g(T)] = \int_{\theta}^{+\infty} g(t)f_T(t) ~dt  \qquad \forall \theta \in \mathbb{R}
    \end{equation*}
    Como $T=X_{(1)}$, tenemos que:
    \begin{equation*}
        F_T(t) = 1-{(1-F_\theta(t))}^{n} \Longrightarrow f_T(t) = n{(1-F_\theta(t))}^{n-1}f_\theta(t)
    \end{equation*}

    donde:
    \begin{equation*}
        F_\theta(t) = \int_{\theta}^{t} e^{\theta-x} ~dx  = \left[-e^{\theta-x}\right]_{\theta}^t = 1-e^{\theta-t}, \qquad t>\theta
    \end{equation*}

    por lo que:
    \begin{equation*}
        f_T(t) = n{\left(e^{\theta-t}\right)}^{n-1}e^{\theta-t} = n{\left(e^{\theta-t}\right)}^{n}, \qquad t>\theta
    \end{equation*}

    volviendo al caso que nos interesa:
    \begin{equation*}
        0 = E[g(T)] = \int_{\theta}^{+\infty} g(t)n{\left(e^{\theta-t}\right)}^{n} ~dt  = ne^{n\theta} \int_{\theta}^{+\infty} g(t)e^{-nt} ~dt  \qquad \forall \theta \in \mathbb{R}
    \end{equation*}

    por lo que:
    \begin{equation*}
        \int_{\theta}^{+\infty} g(t)e^{-nt}~dt  = 0 \qquad \forall \theta\in \mathbb{R}
    \end{equation*}
    Por simplicidad de cálculos, \underline{supondremos que $g$ es continua}, con lo que si $G(t)$ es una primitiva de $g(t)e^{-nt}$, entonces:
    \begin{equation*}
        \int_{\theta}^{+\infty} g(t)e^{-nt}~dt = \lim_{n\to\infty}\int_{\theta}^{n}g(t)e^{-nt} ~dt  = \lim_{n\to\infty}G(n)-G(\theta) =  0 \qquad \forall \theta\in \mathbb{R}
    \end{equation*}
    Si derivamos ahora respecto a $\theta$, tenemos que:
    \begin{equation*}
        -g(\theta)e^{-n\theta} = 0\qquad \forall \theta\in \mathbb{R}
    \end{equation*}

    con lo que $g(\theta) = 0 \quad \forall \theta\in \mathbb{R}$. Es decir:
    \begin{equation*}
        \mathbb{R}\subseteq \{t:g(t) = 0\}
    \end{equation*}

    luego:
    \begin{equation*}
        1\geq P[g(T)=0] \geq P[T\in \mathbb{R}] = 1 \Longrightarrow P[g(T)=0]=1
    \end{equation*}

    por lo que $T$ es completo.
\end{ejercicio}

\begin{ejercicio}
    Comprobar que las siguientes familias de distribuciones son exponenciales uniparamétricas y, considerando una muestra aleatoria simple de una variable con distribución en dicha familia, obtener, si existe, un estadístico suficiente y completo.
    \begin{enumerate}[label=\alph*)]
        \item $\{B(k_0,p) : 0<p<1\}$

            Comprobamos todas las condiciones:
            \begin{enumerate}[label=\arabic*.]
                \item El espacio paramétrico es: $\left]0,1\right[\subseteq \mathbb{R}$.
                \item El espacio muestral es $\cc{X} = \{0,\ldots,k_0\}$, que no depende de $p$.
                \item Para la tercera condición:
                    \begin{align*}
                        P[X=x]_p &= \binom{k_0}{x} p^{x}{(1-p)}^{k_0-x} = \exp\left[\ln\left(\binom{k_0}{x}p^{x}{(1-p)}^{k_0-x} \right)\right]\\
                               &= \exp\left[\ln\binom{k_0}{x} + x\ln p + (k_0-x)\ln(1-p)\right] \\
                               &= \exp\left[\ln\binom{k_0}{x} + x\ln p + k_0 \ln (1-p) - x\ln(1-p)\right] \\
                               &= \exp\left[\ln\binom{k_0}{x} + k_0 \ln (1-p) + x\ln\left(\frac{p}{1-p}\right)\right] 
                    \end{align*}
                    Tomando:
                    \begin{align*}
                        T(x) = x, &\qquad S(x) = \ln\binom{k_0}{x} \\
                        Q(p) = \ln\left(\dfrac{p}{1-p}\right), &\qquad D(p) = k_0\ln(1-p)
                    \end{align*}
                    obtenemos la tercera condición. 
            \end{enumerate}
            Sea $(X_1, \ldots, X_n)$ una m.a.s. de variables aleatorias idénticamente distribuidas a $X\rightsquigarrow B(k_0,p)$ con $p\in \left]0,1\right[$, el Teorema visto en teoría para las familias de distribuciones exponenciales nos dice que el estadístico:
            \begin{equation*}
                T = T(X_1, \ldots, X_n) = \sum_{i=1}^{n}T(X_i) = \sum_{i=1}^{n}X_i
            \end{equation*}
            es suficiente para $p$. Para ver que $T$ es también completo, hemos de ver que $Im Q$ contiene un abierto de $\mathbb{R}$:
            \begin{description}
                \item [Opción 1.] Como $Q:\left]0,1\right[\to \mathbb{R}$ es continua, no constante y definida sobre un intervalo, por el Teorema del Valor Intermedio su imagen ha de ser un intervalo, que es un abierto de $\mathbb{R}$, por lo que $T$ es completo.
                \item [Opción 2.] 
                \begin{align*}
                    Im Q = \left\{\ln\left(\frac{p}{1-p}\right) : p \in \left]0,1\right[\right\}
                \end{align*}
                Definimos la siguiente función:
                \Func{f}{\left]0,1\right[}{\mathbb{R}^+}{t}{\dfrac{t}{1-t}}
                \begin{itemize}
                    \item Está bien definida, puesto que si $t\in \left]0,1\right[$, entonces $1-t>0$, con lo que $f(t)\in \mathbb{R}^+$.
                    \item Sea $y\in \mathbb{R}^+$, tenemos que:
                        \begin{equation*}
                            \dfrac{t}{1-t} = y \Longleftrightarrow t = y(1-t) \Longleftrightarrow t=y-yt \Longleftrightarrow t(1+y) = y \Longleftrightarrow t = \dfrac{y}{y+1}<1
                        \end{equation*}
                        Por lo que $f(t) = y$ con $t\in \left]0,1\right[$, con lo que $f$ es sobreyectiva.
                \end{itemize}
                Por tanto, $Im f = \mathbb{R}^+$, de donde deducimos que:
                \begin{equation*}
                    Im Q = \left\{\ln\left(\frac{p}{1-p}\right) : p \in \left]0,1\right[\right\} = \{\ln(f(p)) : p\in \left]0,1\right[\} = \ln(Im f) = \ln(\mathbb{R}^+) = \mathbb{R}
                \end{equation*}
                Como obviamente $\mathbb{R}$ contiene algún abierto de $\mathbb{R}$, deducimos que $T$ era un estadístico completo.
            \end{description}
        \item $\{\cc{P}(\lm) : \lm>0\}$

            Comprobamos las condiciones:
            \begin{enumerate}[label=\arabic*.]
                \item El espacio paramétrico es $\mathbb{R}^+ \subseteq \mathbb{R}$.
                \item El espacio muestral es $\cc{X} = \mathbb{N}\cup \{0\}$, que no depende de $\lm$.
                \item Para la tercera condición:
                    \begin{align*}
                        P_\lm[X=x] &= e^{-\lm} \dfrac{\lm^x}{x!} = \exp\left[\ln\left(e^{-\lm} \dfrac{\lm^x}{x!}\right)\right] = \exp[-\lm + x\ln \lm - \ln(x!)]
                    \end{align*}
                    Tomando:
                    \begin{align*}
                        T(x) = x, &\qquad S(x) = -\ln(x!) \\
                        Q(\lm) = \ln\lm, &\qquad D(\lm) = -\lm
                    \end{align*}
                    obtenemos la tercera condición.
            \end{enumerate}
            Sea $(X_1, \ldots, X_n)$ una m.a.s. de variables aleatorias idénticamente distribuidas a $X\rightsquigarrow \cc{P}(\lm)$ con $\lm\in \mathbb{R}^+$, el Teorema visto en teoría para las familias de distribuciones exponenciales nos dice que el estadístico:
            \begin{equation*}
                T = T(X_1, \ldots, X_n) = \sum_{i=1}^{n}T(X_i) = \sum_{i=1}^{n}X_i
            \end{equation*}
            es suficiente para $\lm$. Para ver que $T$ es también completo, hemos de ver que $Im Q$ contiene un abierto de $\mathbb{R}$. Como:
            \begin{equation*}
                Im Q = \{\ln(\lm) : \lm \in \mathbb{R}^+\} = \ln(\mathbb{R}^+) = \mathbb{R}
            \end{equation*}
            Tenemos que $Im Q = \mathbb{R}$ claramente contiene un abierto de $\mathbb{R}$, por lo que $T$ es completo.
        \item $\{BN(k_0,p) : 0<p<1\}$

            Comprobamos las condiciones:
            \begin{enumerate}[label=\arabic*.]
                \item El espacio paramétrico es $\left]0,1\right[\subseteq \mathbb{R}$.
                \item El espacio muestral es $\cc{X} = \mathbb{N}\cup \{0\}$, que no depende de $p$.
                \item Para la tercera condición:
                    \begin{align*}
                        P_p[X=x] &= \binom{x+k_0-1}{x} {(1-p)}^{x}p^{k_0 }= \exp\left[\ln\left(\binom{x+k_0-1}{x} {(1-p)}^{x}p^{k_0} \right)\right] \\
                                 &= \exp\left[\ln\binom{x+k_0-1}{x} + x\ln(1-p) + k_0\ln p\right]
                    \end{align*}
                    Tomando:
                    \begin{align*}
                        T(x) = x, &\qquad S(x) = \ln\binom{x+k_0-1}{x} \\
                        Q(p) = \ln(1-p), &\qquad D(p) = k_0\ln p
                    \end{align*}
                    Obtenemos la tercera condición.
            \end{enumerate}
            Sea $(X_1, \ldots, X_n)$ una m.a.s. de variables aleatorias idénticamente distribuidas a $X\rightsquigarrow BN(k_0,p)$ con $p\in \left]0,1\right[$, el Teorema visto en teoría para las familias de distribuciones exponenciales nos dice que el estadístico:
            \begin{equation*}
                T = T(X_1, \ldots, X_n) = \sum_{i=1}^{n}T(X_i) = \sum_{i=1}^{n}X_i
            \end{equation*}
            es suficiente para $p$. Para ver que $T$ es también completo, hemos de ver que $Im Q$ contiene un abierto de $\mathbb{R}$. Como $Q:\left]0,1\right[\to\mathbb{R}$ es una función continua, no constante y definida en un intervalo, tenemos que su imagen es un intervalo, por lo que contiene abiertos de $\mathbb{R}$, de donde $T$ es completo.
        \item $\{\exp(\lm) : \lm>0\}$

            Comprobamos las condiciones:
            \begin{enumerate}[label=\arabic*.]
                \item El espacio paramétrico es $\mathbb{R}^+ \subseteq \mathbb{R}$.
                \item El espacio muestral es $\cc{X} = \mathbb{R}^+$, que no depende de $\lm$.
                \item Para la tercera condición:
                    \begin{equation*}
                        f_\lm(x) = \lm e^{-\lm x} = \exp\left[\ln\left(\lm e^{-\lm x} \right)\right] = \exp\left[\ln(\lm) -\lm x\right]
                    \end{equation*}
                    Tomando:
                    \begin{align*}
                        T(x) = x, &\qquad S(x) = 0 \\
                        Q(\lm) = -\lm, &\qquad D(\lm) = \ln(\lm)
                    \end{align*}
                    tenemos la tercera condición.
            \end{enumerate}
            Sea $(X_1, \ldots, X_n)$ una m.a.s. de variables aleatorias idénticamente distribuidas a $X\rightsquigarrow \exp(\lm)$ con $\lm \in \mathbb{R}^+$, el Teorema visto en teoría para las familias de distribuciones exponenciales nos dice que el estadístico:
            \begin{equation*}
                T = T(X_1, \ldots, X_n) = \sum_{i=1}^{n}T(X_i) = \sum_{i=1}^{n}X_i
            \end{equation*}
            es suficiente para $\lm$. Para ver que $T$ es también completo, hemos de ver que $Im Q$ contiene un abierto de $\mathbb{R}$. Como $Q:\mathbb{R}^+\to\mathbb{R}$ es una función continua, no constante y definida en un intervalo, tenemos que su imagen es un intervalo, por lo que contiene abiertos de $\mathbb{R}$, de donde $T$ es completo.
    \end{enumerate}
\end{ejercicio}

\begin{ejercicio}
    Estudiar si las siguientes familias de distribuciones son exponenciales biparamétricas.  En caso afirmativo, considerando una muestra aleatoria simple de una variable con distribución en dicha familia, obtener, si existe, un estadístico suficiente y completo.
    \begin{enumerate}[label=\alph*)]
        \item $\{\Gamma(p,a) : p,a>0\}$

            Comprobamos las condiciones:
            \begin{enumerate}[label=\arabic*.]
                \item El espacio paramétrico es $\mathbb{R}^+\times \mathbb{R}^+\subseteq \mathbb{R}^2$.
                \item El espacio muestral es $\mathbb{R}^+$, que no depende de $p$ ni de $a$.
                \item Para la tercera condición:
                    \begin{align*}
                        f_{(p,a)}(x) &= \dfrac{a^p}{\Gamma(p)} x^{p-1}e^{-ax} = \exp\left[\ln\left(\dfrac{a^p}{\Gamma(p)} x^{p-1}e^{-ax} \right)\right] \\ &= \exp\left[\ln\left(\dfrac{a^p}{\Gamma(p)}\right) + (p-1)\ln x - ax\right]
                    \end{align*}
                Tomando:
                \begin{align*}
                    T_1(x) = \ln x, &\qquad  T_2(x) = x, \qquad \qquad   S(x) = 0 \\
                    Q_1(p,a) = (p-1), &\qquad Q_2(p,a) = -a, \qquad D(p,a) = \ln\left(\dfrac{a^p}{\Gamma(p)}\right) 
                \end{align*}
                Tenemos la tercera condición.
            \end{enumerate}
            Sea $(X_1, \ldots, X_n)$ una m.a.s. de variables aleatorias idénticamente distribuidas a $X\rightsquigarrow \Gamma(p,a)$ con $p,a\in \mathbb{R}^+$, el Teorema visto en teoría para las familias de distribuciones exponenciales multiparamétricas nos dice que el estadístico:
            \begin{equation*}
                T = T(X_1, \ldots, X_n) = \left(\sum_{i=1}^{n}T_1(X_i), \sum_{i=1}^{n}T_2(X_i)\right) = \left(\sum_{i=1}^{n}\ln(X_i), \sum_{i=1}^{n}X_i\right)
            \end{equation*}
            es suficiente para $(p,a)$. Para ver que $T$ es también completo, hemos de ver que $Im Q$ contiene un abierto de $\mathbb{R}^2$, donde $Q:{(\mathbb{R}^+)}^{2}\to \mathbb{R}^2$, con $Q=(Q_1,Q_2)$. Para ello:
            \begin{align*}
                Im Q = \left\{(p-1,-a) : (p,a)\in {(\mathbb{R}^+)}^{2}\right\} &= \{(x-1,y) : x\in \mathbb{R}^+, y\in \mathbb{R}^-\} \\ &= \left]-1,+\infty\right[\times \mathbb{R}^-
            \end{align*}
            Como claramente $\left]-1,+\infty\right[\times\mathbb{R}^-$ contiene un abierto de $\mathbb{R}^2$, tenemos que $T$ es completo.

            Observemos que también podríamos haber tomado:
            \begin{equation*}
                T(X_1, \ldots, X_n) = \left(\prod_{i=1}^n X_i, \sum_{i=1}^n X_i\right)
            \end{equation*}
        \item $\{\beta(p,q) : p,q>0\}$

            Comprobamos las condiciones:
            \begin{enumerate}[label=\arabic*.]
                \item El espacio paramétrico es $\mathbb{R}^+\times \mathbb{R}^+\subseteq \mathbb{R}^2$.
                \item El espacio muestral es $[0,1]$, que no depende de $p$ ni de $q$.
                \item Para la tercera condición:
                    \begin{align*}
                        f_{(p,q)}(x) &= \dfrac{1}{\beta(p,q)}x^{p-1}{(1-x)}^{q-1} = \exp\left[\ln\left(\dfrac{1}{\beta(p,q)}x^{p-1}{(1-x)}^{q-1} \right)\right] \\
                                     &= \exp\left[\ln\left(\dfrac{1}{\beta(p,q)}\right) + (p-1)\ln x + (q-1)\ln(1-x)\right]
                    \end{align*}
                Tomando:
                \begin{align*}
                    T_1(x) = \ln x, &\qquad T_2(x) = \ln(1-x), \qquad S(x) = 0 \\
                    Q_1(p,q) = p-1, &\qquad Q_2(p,q) = q-1, \qquad D(p,q) = \ln\left(\dfrac{1}{\beta(p,q)}\right) 
                \end{align*}
                Tenemos la tercera condición.
            \end{enumerate}
            Sea $(X_1, \ldots, X_n)$ una m.a.s. de variables aleatorias idénticamente distribuidas a $X\rightsquigarrow \beta(p,q)$ con $p,q\in \mathbb{R}^+$, el Teorema visto en teoría para las familias de distribuciones exponenciales multiparamétricas nos dice que el estadístico:
            \begin{equation*}
                T = T(X_1, \ldots, X_n) = \left(\sum_{i=1}^{n}T_1(X_i), \sum_{i=1}^{n}T_2(X_i)\right) = \left(\sum_{i=1}^{n}\ln(X_i), \sum_{i=1}^{n}\ln(1-X_i)\right)
            \end{equation*}
            es suficiente para $(p,q)$. Para ver que $T$ es también completo, hemos de ver que $Im Q$ contiene un abierto de $\mathbb{R}^2$, donde $Q:{(\mathbb{R}^+)}^{2}\to \mathbb{R}^2$, con $Q=(Q_1,Q_2)$. Para ello:
            \begin{equation*}
                Im Q = \{(p-1,q-1):p,q\in \mathbb{R}^+\} = \left]-1,+\infty\right[\times \left]-1,+\infty\right[
            \end{equation*}
            Como claramente este conjunto contiene un abierto de $\mathbb{R}^2$, tenemos que $T$ es completo.
    \end{enumerate}
\end{ejercicio}

    % \section{Estimación puntual. Insesgadez y mínima varianza}

\begin{ejercicio}
    Sea $(X_1, \ldots, X_n)$ una muestra de una variable $X\rightsquigarrow\cc{N}(\mu, \sigma^2)$ con $\mu\in \mathbb{R}$, $\sigma\in \mathbb{R}^+$. Probar que
    \begin{equation*}
        T(X_1, \ldots, X_n) = \left\{\begin{array}{ll}
            1 & \text{si\ } \overline{X}\leq 0 \\
            0 & \text{si\ } \overline{X} > 0
        \end{array}\right. 
    \end{equation*}
    es un estimador insesgado de la función paramétrica $\Phi\left(\frac{-\mu \sqrt{n}}{\sigma}\right)$, siendo $\Phi$ la función de distribución de la $\cc{N}(0,1)$.\\

    \noindent
    Tenemos $T(X_1, \ldots, X_n) = I_{\left]-\infty,0\right]}(\overline{X})$. Como $X\rightsquigarrow \cc{N}(\mu, \sigma^2)$, sabemos por lo visto en el Tema 1 que entonces:
    \begin{equation*}
        \overline{X} \rightsquigarrow \cc{N}\left(\mu, \frac{\sigma^2}{n}\right)
    \end{equation*}

    de donde (escribiendo $T=T(X_1,\ldots,X_n)$):
    \begin{equation*}
    T = I_{\left]-\infty,0\right]}(\overline{X}) \rightsquigarrow B(1,P[\overline{X}\leq 0])
    \end{equation*}

    estamos ya en condiciones de ver que $T$ es insesgado para dicha función:
    \begin{equation*}
        E[T] \AstIg P[\overline{X}\leq 0] \stackrel{\text{tipif.}}{=} P\left[Z \leq \dfrac{-\mu\sqrt{n}}{\sigma}\right] = \Phi\left(\dfrac{-\mu\sqrt{n}}{\sigma}\right) 
    \end{equation*}
    donde en $(\ast)$ usamos que conocemos bien la esperanza de una distribución Bernoulli.
\end{ejercicio}

\begin{ejercicio}
    Sea $(X_1, \ldots, X_n)$ una muestra aleatoria simple de $X\rightsquigarrow B(1,p)$ con $p\in \left]0,1\right[$ y sea $T=\sum\limits_{i=1}^{n}X_i$.
    \begin{enumerate}[label=\alph*)]
        \item Probar que si $k\in \mathbb{N}$ y $k\leq n$, el estadístico
            \begin{equation*}
                \dfrac{T(T-1)\cdot \ldots\cdot (T-k+1)}{n(n-1)\cdot \ldots\cdot (n-k+1)}
            \end{equation*}
            es un estimador insesgado de $p^k$. ¿Es este estimador el UMVUE?.
        \item Probar que si $k>n$, no existe ningún estimador insesgado para $p^k$.
        \item ¿Puede afirmarse que $\frac{T}{n}{\left(1-\frac{T}{n}\right)}^{2}$ es insesgado para $p{(1-p)}^{2}$? 
    \end{enumerate}~\\

    Veamos cada apartado:
    \begin{enumerate}[label=\alph*)]
        \item Sea $k\in \mathbb{N}$ con $k\leq n$, definimos:
            \begin{equation*}
                h(T) = \dfrac{T(T-1)\cdot \ldots\cdot (T-k+1)}{n(n-1)\cdot \ldots\cdot (n-k+1)}
            \end{equation*}
            Veamos que $h(T)$ es insesgado para $p^k$. En primer lugar, observemos que por la reproductividad de la binomial $T\rightsquigarrow B(n,p)$:
            \begin{equation*}
                E[h(T)] = \dfrac{E[T(T-1)\cdot \ldots\cdot (T-k+1)]}{n(n-1)\cdot \ldots\cdot (n-k+1)} = \dfrac{(n-k)!}{n!}E[T(T-1)\cdot \ldots\cdot (T-k+1)]
            \end{equation*}
            Calculamos ahora la esperanza:
            \begin{align*}
                E[T(T-1)\cdot \ldots\cdot (T-k+1)] &= \sum_{t=0}^{n}t(t-1)\cdot \ldots\cdot (t-k+1)P[T=t] \\ 
                                                   &= \sum_{t=0}^{n}t(t-1)\cdot \ldots\cdot (t-k+1)\binom{n}{t}p^t{(1-p)}^{n-t} \\
                                                   &= \sum_{t=k}^{n}t(t-1)\cdot \ldots\cdot(t-k+1) \binom{n}{t}p^t {(1-p)}^{n-t} \\
                                                   &= \sum_{t=k}^{n}\dfrac{t!}{(t-k)!}\dfrac{n!}{t!(n-t)!}p^t {(1-p)}^{n-t} \\
                                                   &= \sum_{t=k}^{n}\dfrac{n!}{(t-k)!(n-t)!}p^t {(1-p)}^{n-t}
            \end{align*}
            si desarrollamos ahora los primeros términos, observamos que:
            \begin{align*}
                &E[T(T-1)\cdot \ldots\cdot (T-k+1)] = \sum_{t=k}^{n}\dfrac{n!}{(t-k)!(n-t)!}p^t {(1-p)}^{n-t} \\
               &= \dfrac{n!}{(n-k)!}p^k {(1-p)}^{n-k} + \dfrac{n!}{(n-k-1)!}p^{k+1}{(1-p)}^{n-k-1} + \ldots
            \end{align*}
            donde podemos ver que podemos sacar factor común de la sumatoria ciertos términos:
            \begin{align*}
                &E[T(T-1)\cdot \ldots\cdot (T-k+1)] = \sum_{t=k}^{n}\dfrac{n!}{(t-k)!(n-t)!}p^t {(1-p)}^{n-t} \\
               &= \dfrac{n!\cdot p^k }{(n-k)!}\sum_{t=k}^{n}p^{t-k}{(1-p)}^{n-t}\dfrac{(n-k)!}{(t-k)!(n-t)!} \\
               &= \dfrac{n!\cdot p^k}{(n-k)!}\sum_{t=0}^{n-k} p^t {(1-p)}^{n-k-t}\dfrac{(n-k)!}{t!(n-k-t)!} \\
               &= \dfrac{n!\cdot p^k}{(n-k)!}\sum_{t=0}^{n-k}\binom{n-k}{t} p^t {(1-p)}^{n-k-t} \\
               &= \dfrac{n!\cdot p^k}{(n-k)!}\sum_{t=0}^{n-k}P[S=t] = \dfrac{n!\cdot p^k}{(n-k)!}P[0\leq S \leq n-k] \AstIg \dfrac{n!\cdot p^k}{(n-k)!}
            \end{align*}
            para cierta variable aleatoria $S\rightsquigarrow B(n-k,p)$, donde en $(\ast)$ usamos que $P[0\leq S \leq n-k]=1$. Ahora, vemos que:
            \begin{equation*}
                E[h(T)] = \dfrac{(n-k)!}{n!}\dfrac{n!}{(n-k)!}p^k = p^k
            \end{equation*}
            por lo que $h(T)$ es insesgado para $p^k$. Veamos ahora que $h(T)$ es un estimador. Para ello, observamos primero que:
            \begin{equation*}
                h(T) = \dfrac{T}{n}\cdot \dfrac{T-1}{n-1}\cdot \ldots \cdot \dfrac{T-k+1}{n-k+1} = \prod_{j=0}^{k-1}\dfrac{T-j}{n-j}
            \end{equation*}
            Para cada $j\in \{0,\ldots,k-1\}$, observemos que $T\in \{0,\ldots,n\}$, por lo que $T-j\in \{-j, \ldots, n-j\}$ con $j<k\leq n$, de donde deducimos que:
            \begin{equation*}
                \dfrac{T-j}{n-j} \in [-1, 1] \quad \Longrightarrow \quad  \prod_{j=0}^{k-1}\dfrac{T-j}{n-j} \in [-1,1] 
            \end{equation*}
            Tenemos que ver finalmente que dicho producto es positivo, con lo que habremos probado que $h(T)$ es un estimador. Para ello, si el prodcuto no fuera positivo es por la existencia de $j\in \{0,\ldots,k-1\}$ de forma que $T-j<0$, es decir, tenemos entonces que $T\in \{0,\ldots,j-1\}$ supuesto que $T=l \in \{0,\ldots,j-1\}$, tendremos entonces que:
            \begin{equation*}
                h(T) = \prod_{j=0}^{k-1}\dfrac{T-j}{n-j} = \left(\prod_{j=0}^{l-1}\dfrac{T-j}{n-j} \right) \dfrac{T-l}{n-l} \left(\prod_{j=l+1}^{k-1}\dfrac{T-j}{n-j}\right) = 0
            \end{equation*}
            es decir, siempre que un término del producto sea negativo el producto entero se anula, por lo que siempre el producto es positivo, de donde $h(T)\in [0,1]$, por lo que $h(T)$ es un estimador.

            Finalmente, como:
            \begin{equation*}
                E\left[{(h(T))}^{2}\right] = \sum_{t=0}^{n}{(h(t))}^{2}P[T=t] < \infty
            \end{equation*}
            y teníamos que $T$ era un estimador suficiente y completo (cuando vimos que $\{B(1,p) : p\in \left]0,1\right[\}$ era una familia exponencial), tenemos entonces que $E[h(T)/T] = h(T)$ es el UMVUE de $p^k$.
        \item Sea ahora $k>n$, veamos que no puede existir ningun estimador insesgado para $p^k$. Para ello, por reducción al absurdo, supongamos que $h(T)$ es un estimador insesgado para $p^k$, con lo que:
            \begin{equation*}
                p^k = E[h(T)] = \sum_{t=0}^{n}h(t)P[T=t] = \sum_{t=0}^{n}h(t)\binom{n}{t}p^t {(1-p)}^{n-t}
            \end{equation*}

            de donde:
            \begin{equation*}
                1 = \sum_{t=0}^{n}h(t)\binom{n}{t}p^{t-k} {(1-p)}^{n-t} \qquad \forall p \in \left]0,1\right[
            \end{equation*}
            en particular, tomando $p\to 0$, como $t-k < 0$ para todo $t \in \{0,\ldots,k\}$, tenemos que:
            \begin{equation*}
                1 = \lim_{p\to0} \sum_{t=0}^{n}h(t)\binom{n}{t}p^{t-k} {(1-p)}^{n-t} = \infty
            \end{equation*}
            \underline{contradicción}, con lo que para $k>n$ no puede existir un estimador insesgado para $p^k$. También podríamos haber justificado que son dos polinomios de distinto grado, para llegar a contradicción.
        \item Buscamos ahora comprobar si $h(T) = \frac{T}{n}{\left(1-\frac{T}{n}\right)}^{2}$ es insesgado para $p{(1-p)}^{2}$.

            Resulta que esto no puede afirmarse para todos los valores de $n$ y $p$. Como un contraejemplo, para $n=3$ resulta que tenemos:
            \begin{align*}
                E[h(T)] &= \sum_{t=0}^{3}\frac{t}{3}{\left(1-\frac{t}{3}\right)}^{2}\binom{3}{t}p^t{(1-p)}^{3-t} \AstIg \sum_{t=1}^{2}\frac{t}{3}{\left(1-\frac{t}{3}\right)}^{2}\binom{3}{t}p^t{(1-p)}^{3-t} \\
                        &= \frac{1}{3}{\left(1-\frac{1}{3}\right)}^{2}\binom{3}{1}p{(1-p)}^{2}+\frac{2}{3}{\left(1-\dfrac{2}{3}\right)}^{2}\binom{3}{2}p^2(1-p) \\
                        &= p(1-p)\left(\frac{1}{9}(1-p)+\frac{2}{9}p\right) = p(1-p)\frac{(1+p)}{9} \qquad \forall p\in \left]0,1\right[
            \end{align*}
            donde en $(\ast)$ hemos usado que la función $t\mapsto \frac{t}{n}{\left(1-\frac{t}{n}\right)}^{2}$ evaluada en 0 y 3 es igual a cero. Observamos que:
            \begin{equation*}
                E[h(T)] = p(1-p)\frac{(1+p)}{9} = p{(1-p)}^{2} \Longleftrightarrow \dfrac{1+p}{9} = 1-p
            \end{equation*}
            resultado que no es cierto para todo $p\in \left]0,1\right[$ (basta considerar $p=\nicefrac{1}{2}$), por lo que dicho estimardor no es insesgado para $p{(1-p)}^{2}$.
    \end{enumerate}
\end{ejercicio}

\begin{ejercicio}
    Sea $(X_1, \ldots, X_n)$ una muestra aleatoria simple de una variable $X\rightsquigarrow \cc{P}(\lm)$ con $\lm\in \mathbb{R}^+$. Encontrar, si existe, el UMVUE para $\lm^s$, siendo $s\in \mathbb{N}$ arbitrario.\\
    
    \noindent
    Veamos que $T(X_1, \ldots, X_n) = \sum\limits_{i=1}^{n}X_i$ es un estadístico suficiente y completo. Para ello, recordemos que $\{\cc{P}(\lm) : \lm>0\}$ es una familia exponencial:
    \begin{enumerate}
        \item El espacio paramétrico es $\mathbb{R}^+ \subseteq \mathbb{R}$.
        \item El espacio muestral es $\cc{X}=\mathbb{N}\cup \{0\}$, que no depende de $\lm$.
        \item Observamos que:
            \begin{align*}
                P_\lm[X=x] &= e^{-\lm}\dfrac{\lm^x}{x!} = exp\left[\ln\left(e^{-\lm}\dfrac{\lm^x}{x!}\right)\right] = exp\left(-\lm + x\ln\lm -\ln(x!)\right)
            \end{align*}
            por lo que basta tomar:
            \begin{equation*}
                Q(\lm) = \ln\lm, \qquad T(x) = x \qquad D(\lm) = -\lm, \qquad S(x) = -\ln(x!)
            \end{equation*}
    \end{enumerate}
    En consecuencia, por un Teorema visto en teoría, tenemos que el estadístico:
    \begin{equation*}
        T(X_1, \ldots, X_n) = \sum_{i=1}^{n}T(X_i) = \sum_{i=1}^{n}X_i
    \end{equation*}
    es suficiente y completo para $\lm$. Observemos que por la reproductividad de la Poisson tenemos que (notando $T=T(X_1, \ldots, X_n)$): 
    \begin{equation*}
        T\rightsquigarrow \cc{P}\left(\sum_{i=1}^{n}\lm\right) \equiv \cc{P}(n\lm)
    \end{equation*}
    Ahora, para buscar el UMVUE, buscamos una función $h$ medible de forma que:
    \begin{equation*}
        \lm^s = E[h(T)] = \sum_{t\in \mathbb{N}\cup \{0\}}h(t)P[T=t] = \sum_{t\in \mathbb{N}\cup \{0\}} h(t) e^{-n\lm} \dfrac{{(n\lm)}^{t}}{t!} 
    \end{equation*}

    por lo que:
    \begin{equation*}
        \lm^s e^{n\lm} = \sum_{t\in \mathbb{N}\cup \{0\}} h(t)\dfrac{{(n\lm)}^{t}}{t!}
    \end{equation*}

    y si aplicamos el desarrollo en serie de la exponencial, obtenemos:
    \begin{equation*}
        \lm^s \sum_{t\in \mathbb{N}\cup \{0\}} \dfrac{{(n\lm)}^{t}}{t!} = \lm^s e^{n\lm} = \sum_{t\in \mathbb{N}\cup\{0\}} h(t) \dfrac{{(n\lm)}^{t}}{t!} 
    \end{equation*}

    si desarrollamos cada uno de los términos:
    \begin{equation*}
        \lm^s + \lm^{s+1}n + \frac{\lm^{s+2}n^2}{2!} + \ldots = h(0) + h(1)(n\lm) + \ldots + h(s)\dfrac{{(n\lm)}^{s}}{s!} + \ldots
    \end{equation*}

    observamos que tomando:
    \begin{gather*}
        h(0) = \ldots = h(s-1) = 0,\quad  h(s) = \dfrac{s!}{n^s}\\h(s+1) = \dfrac{(s+1)!}{n^s}, \quad \ldots \quad  h(s+k) = \dfrac{(s+k)!}{n^s k!}
    \end{gather*}

    es decir:
    \begin{equation*}
        h(T) = \left\{\begin{array}{ll}
            0 & \text{si\ } T<s \\
            \dfrac{T!}{n^s(T-s)!}& \text{si\ } T\geq s
        \end{array}\right. 
    \end{equation*}
    tenemos que $h(T)$ es insesgado para $\lm^s$. Es claro además que $h(t)\in \mathbb{R}^+$ para cualquier valor de $t$, con lo que $h(T)$ es un estimador de $\lm^s$. Finalmente, observemos que:
    \begin{align*}
        E\left[{(h(T))}^{2}\right] &= \sum_{t\in \mathbb{N}\cup \{0\}} {(h(t))}^{2}P[T=t] = \sum_{t\geq s} {\left(\dfrac{t!}{n^s(t-s)!}\right)}^{2} e^{-n\lm} \dfrac{{(n\lm)}^{t}}{t!} \\ &= \dfrac{1}{n^s e^{n\lm}} \sum_{t\geq s} \dfrac{{(n\lm)}^{t}t!}{{((t-s)!)}^{2}}
    \end{align*}

    como:
    \begin{equation*}
        \dfrac{\dfrac{{(n\lm)}^{t+1}(t+1)!}{{((t+1-s)!)}^{2}}}{\dfrac{{(n\lm)}^{t}t!}{{((t-s)!)}^{2}}} = \dfrac{{(n\lm)}^{t+1}(t+1)!{((t-s)!)}^{2}}{{(n\lm)}^{t}t!{((t+1-s)!)}^{2}} = \dfrac{n\lm(t+1)}{{(t+1-s)}^{2}} \to 0 < 1
    \end{equation*}
    por el Criterio del cociente, tenemos que:
    \begin{equation*}
        \sum_{t\geq s} \dfrac{{(n\lm)}^{t}t!}{{((t-s)!)}^{2}} < \infty \Longrightarrow E\left[{(h(T))}^{2}\right]  =\dfrac{1}{n^s e^{n\lm}}\sum_{t\geq s} \dfrac{{(n\lm)}^{t}t!}{{((t-s)!)}^{2}}   < \infty
    \end{equation*}
    en consecuencia, tenemos que $h(T)$ es un estimador insesgado para $\lm^s$ y de momento de segundo orden finito y es función de un estadístico suficiente y completo, con lo que el Teorema de Lehmann-Scheffé nos dice que:
    \begin{equation*}
        E[h(T)/T] = h(T)
    \end{equation*}
    es un UMVUE para $\lm^s$.
\end{ejercicio}

\begin{ejercicio}
    Sea $(X_1, \ldots, X_n)$ una muestra aleatoria simple de una variable con distribución uniforme discreta en los puntos $\{1,\ldots,N\}$, siendo $N$ un número natural arbitrario. Encontrar el UMVUE para $N$.\\

    \noindent
    En el Ejercicio~\ref{ej:3.5} vimos que $T(X_1, \ldots, X_n) = X_{(n)}$ era un estadístico suficiente y completo. Si notamos $T = T(X_1,\ldots, X_n)$, tenemos que:
    \begin{equation*}
        F_T(t) = {(F_X(t))}^{n} \Longrightarrow P[T=t] = P[T\leq t] - P[T\leq t-1] = {(F_X(t))}^{n}-{(F_X(t-1))}^{n}
    \end{equation*}
    como $F_X(t) = \frac{t}{N}$ para $t\in \{1,\ldots,N\}$, tenemos:
    \begin{equation*}
        P[T=t] = {(F_X(t))}^{n}-{(F_X(t-1))}^{n} = \dfrac{t^n - {(t-1)}^{n}}{N^n}
    \end{equation*} 
    Buscamos ahora una función $h$ medible de forma que:
    \begin{equation*}
        N = E[h(T)] = \sum_{t=1}^{N}h(t) P[T=t] = \sum_{t=1}^{N}h(t)\dfrac{t^n-{(t-1)}^{n}}{N^n}
    \end{equation*}
    Si desarrollamos la suma para $h(t) = 1$, observamos un comportamiento telescópico:
    \begin{align*}
        \dfrac{1}{N^n}\sum_{t=1}^{N}(t^n - {(t-1)}^{n}) &= \dfrac{1}{N^n}(N^n - {(N-1)}^{n} + {(N-1)}^{n} - {(N-2)}^{n} + \ldots + 2^n - 1 + 1 - 0) \\
                                                        &= \dfrac{N^n}{N^n} = 1
    \end{align*}
    No hemos obteniedo lo que queríamos, puesto que queríamos que el resultado de la suma fuera $N^{n+1}$. Si tomamos sin embargo:
    \begin{equation*}
        h(t) = \dfrac{t^{n+1}-{(t-1)}^{n+1}}{t^n - {(t-1)}^{n}}
    \end{equation*}
    tenemos entonces que:
    \begin{align*}
        E[h(T)] &= \dfrac{1}{N^n}\sum_{t=1}^{N}h(t) (t^n-{(t-1)}^{n}) = \dfrac{1}{N^n}\sum_{t=1}^{N}t^{n+1}-{(t-1)}^{n+1} \\
                &= \dfrac{1}{N^n}(N^{n+1}-{(N-1)}^{n+1} + {(N-1)}^{n+1} + \ldots + 2^{n+1} - 1 + 1 -0) \\ 
                &= \dfrac{N^{n+1}}{N^n} = N
    \end{align*}
    Por lo que el estadístico $h(T)$ es insesgado para $N$. Nos falta comprobar si es un estimador y si tiene momento de segundo orden finito:
    \begin{itemize}
        \item No es estimador, puesto que $h(\mathbb{N})\nsubseteq \mathbb{N}$. Sin embargo, podemos realizar aproximaciones y quedarnos con el natural más próximo. Es decir, no va a haber un UMVUE pero podemos quedarnos con este estadístico, que es aquello que podemos conseguir más próximo a un UMVUE. % Literalmente dijo esto la profe.
        \item Para el momento de segundo orden, observamos que:
            \begin{equation*}
                E[{(h(T))}^{2}] = \dfrac{1}{N^n}\sum_{t=1}^{N}{(h(t))}^{2}(t^n-{(t-1)}^{n}) < \infty
            \end{equation*}
            al ser una suma finita.
    \end{itemize}
\end{ejercicio}

\begin{ejercicio}
    Sea $(X_1, \ldots, X_n)$ una muestra aleatoria simple de una variable aleatoria $X$ cuya función de densidad es de la forma
    \begin{equation*}
        f_\theta(x) = \dfrac{1}{2\sqrt{x\theta}}, \quad 0<x<\theta
    \end{equation*}
    Calcular, si existe, el UMVUE para $\theta$.\\

    \noindent
    En primer lugar, buscamos un estadístico suficiente y completo:
    \begin{equation*}
        f_\theta(x_1, \ldots, x_n) \stackrel{\text{iid.}}{=} \prod_{i=1}^{n}f_\theta(x_i) = \prod_{i=1}^{n}\dfrac{1}{2\sqrt{x_i\theta}} = \dfrac{1}{{\left(2\sqrt{\theta}\right)}^{n}} \prod_{i=1}^{n} \frac{1}{\sqrt{x_i}} \qquad x_i \in \left]0,\theta\right[
    \end{equation*}
    Si consideramos $0< X_{(1)} \leq X_{(n)} < \theta $, podemos escribir:
    \begin{equation*}
        f_\theta(x_1, \ldots, x_n) = \dfrac{I_{\left]0,+\infty\right[}(X_{(1)}-0) \cdot I_{\left]-\infty,0\right[}(X_{(n)}-\theta)}{{\left(2\sqrt{\theta}\right)}^{n}} \prod_{i=1}^{n}\dfrac{1}{\sqrt{x_i}}
    \end{equation*}
    Tomando:
    \begin{gather*}
        h(x_1, \ldots, x_n) = \prod_{i=1}^{n}\dfrac{1}{\sqrt{x_i}}\cdot I_{\left]0,+\infty\right[}(X_{(1)}-0)  , \qquad T(X_1, \ldots, X_n) = X_{(n)} \\
        g_\theta(t) = \dfrac{I_{\left]-\infty,0\right[}(t-\theta)}{{\left(2\sqrt{\theta}\right)}^{n}} 
    \end{gather*}
    por el Teorema de factorización de Neymann-Fisher tenemos que el estadístico $X_{(n)}$ es suficiente para $\theta$. Notando $T=X_{(n)}$ para abreviar, calculamos la distribución de $T$:
    \begin{equation*}
        F_T(t) = {(F_X(t))}^{n} \quad \Longrightarrow \quad  f_T(t) = n{(F_X(t))}^{n-1}f_\theta(t)
    \end{equation*}
    Calculamos la distribución de $X$:
    \begin{equation*}
        F_X(t) = \int_{0}^{t} \dfrac{1}{2\sqrt{x\theta}}~dx = \dfrac{1}{2\sqrt{\theta}}\int_{0}^{t} \dfrac{1}{\sqrt{x}}~dx  = \sqrt{\frac{t}{\theta}} \qquad t \in \left]0,\theta\right[
    \end{equation*}
    Por lo que:
    \begin{equation*}
        f_T(t) = n{(F_X(t))}^{n-1}f_\theta(t) = n{\left(\sqrt{\frac{t}{\theta}}\right)}^{n-1}\frac{1}{2\sqrt{t\theta}} = \dfrac{n{\left(\sqrt{t}\right)}^{n-2}}{2{\left(\sqrt{\theta}\right)}^{n}}\qquad t\in \left]0,\theta\right[
    \end{equation*}
    Sea ahora $h$ una función medible de forma que:
    \begin{equation*}
        0 = E[h(T)] = \int_{0}^{\theta} h(t) f_T(t)~dt = \frac{n}{2{\left(\sqrt{\theta}\right)}^{n}}\int_{0}^{\theta} h(t){\left(\sqrt{t}\right)}^{n-2}~dt  \qquad \forall \theta\in \mathbb{R}^+
    \end{equation*}

    tenemos entonces que:
    \begin{equation*}
        \int_{0}^{\theta} h(t){\left(\sqrt{t}\right)}^{n-2}~dt = 0 \qquad \forall \theta\in \mathbb{R}^+
    \end{equation*}
    Sea ahora $H(t)$ una primitiva de $h(t){\left(\sqrt{t}\right)}^{n-2}$, tenemos entonces que:
    \begin{equation*}
        H(\theta) - H(0) = 0 \qquad \forall \theta\in \mathbb{R}^+
    \end{equation*}

    por lo que derivando respecto a $\theta$:
    \begin{equation*}
        h(\theta) {\left(\sqrt{\theta}\right)}^{n-2} = 0 \qquad \forall \theta\in \mathbb{R}^+ \quad \Longrightarrow \quad  h(\theta) = 0 \qquad \forall \theta\in \mathbb{R}^+
    \end{equation*}
    En conclusión, tenemos que:
    \begin{equation*}
        \mathbb{R}^+ \subseteq \{t:h(t) = 0\}
    \end{equation*}

    por lo que:
    \begin{equation*}
        1 \geq P[h(T) = 0] \geq P[T\in \mathbb{R}^+] = 1 \quad \Longrightarrow \quad P[h(T)=0] = 1
    \end{equation*}
    lo que demustra que $T$ es un estadístico completo. Ahora, tratamos de buscar un estimador insesgado para $\theta$, que será nuestro candidato a UMVUE. Buscamos una función medible $h$ de forma que:
    \begin{equation*}
        \theta = E[h(T)] = \dfrac{n}{2{\left(\sqrt{\theta}\right)}^{n}}\int_{0}^{\theta} h(t){\left(\sqrt{t}\right)}^{n-2}~dt 
    \end{equation*} 
    Observemos que:
    \begin{equation*}
        \int_{0}^{\theta}{\left(\sqrt{t}\right)}^{n} ~dt  = \dfrac{2}{n+2}\left[{\left(\sqrt{t}\right)}^{n+2}\right]_0^\theta = \dfrac{2}{n+2}{\left(\sqrt{\theta}\right)}^{n+2} = \dfrac{2}{n+2}{\left(\sqrt{\theta}\right)}^{n}\theta
    \end{equation*}
    Por lo que si tomamos:
    \begin{equation*}
        h(t) = \dfrac{n+2}{n} \cdot t
    \end{equation*}
    tenemos que:
    \begin{align*}
        E[h(T)] &= \int_{0}^{\theta} \dfrac{n+2}{n}t\frac{n}{2}\frac{{\left(\sqrt{t}\right)}^{n-2}}{{\left(\sqrt{\theta}\right)}^{n}}~dt  = \dfrac{n+2}{2{\left(\sqrt{\theta}\right)}^{n}} \int_{0}^{\theta}t{\left(\sqrt{t}\right)}^{n-2} ~dt  \\
                &= \dfrac{n+2}{2{\left(\sqrt{\theta}\right)}^{n}} \int_{0}^{\theta}{\left(\sqrt{t}\right)}^{n} ~dt   = \dfrac{n+2}{2{\left(\sqrt{\theta}\right)}^{n}} \dfrac{2}{n+2}{\left(\sqrt{\theta}\right)}^{n}\theta  = \theta \qquad \forall \theta\in \mathbb{R}^+
    \end{align*}
    Por lo que $h(T)$ es insesgado para $\theta$. Ahora, si $t\in \mathbb{R}^+$, tendremos entonces que:
    \begin{equation*}
        h(t) = \dfrac{n+2}{n}\cdot t \geq t \in \mathbb{R}^+ \quad  \Longrightarrow \quad  h(t) \in \mathbb{R}^+
    \end{equation*}
    por lo que $h(T)$ es además un estimador. Si calculamos ahora:
    \begin{align*}
        E\left[{(h(T))}^{2}\right] = \dfrac{n}{2{\left(\sqrt{\theta}\right)}^{n}} \int_{0}^{\theta} {(h(t))}^{2}{(\sqrt{t})}^{n-2}~dt 
    \end{align*}
    Tenemos la integral en un compacto de una función continua, por lo que sabemos que $E[{(h(T))}^{2}]<\infty$. En conclusión, tenemos por el Teorema de Lehmann-Scheffé que: $E[h(T)/T] = h(T)$ es el UMVUE.
\end{ejercicio}

\begin{ejercicio}
    Sea $(X_1, \ldots, X_n)$ una muestra aleatoria simple de una variable aleatoria $X$ con función de densidad
    \begin{equation*}
        f_\theta(x) = \dfrac{\theta}{x^2}, \quad x>\theta > 0
    \end{equation*}
    Calcular, si existen, los UMVUE para $\theta$ y para $\nicefrac{1}{\theta}$.\\

    \noindent
    Calculamos en primer lugar un estadístico suficiente y completo para $\theta$:
    \begin{equation*}
        f_\theta(x_1,\ldots, x_n) \stackrel{\text{iid.}}{=} \prod_{i=1}^{n}f_\theta(x_i) =\prod_{i=1}^{n} \dfrac{\theta\cdot I_{\left]0,+\infty\right[}(X_{(1)}-\theta)}{x_i^2} = \theta^n \cdot I_{\left]0,+\infty\right[}(X_{(1)}-\theta)\prod_{i=1}^{n} \dfrac{1}{x_i^2}
    \end{equation*}
    Tomando:
    \begin{gather*}
        h(x_1, \ldots, x_n) = \prod_{i=1}^{n}\dfrac{1}{x_i^2}, \qquad T(X_1, \ldots, X_n) = X_{(1)} \\
        g_\theta(t) = \theta^n \cdot I_{\left]0,+\infty\right[}(t-\theta)
    \end{gather*}
    Podemos aplicar el Teorema de factorización de Neymann-Fisher, obteniendo que $T$ es suficiente para $\theta$. Notando $T = X_{(1)}$ para abreviar, calculamos la distribución de $T$:
    \begin{equation*}
        F_T(t) = 1-{(1-F_X(t))}^{n} \quad \Longrightarrow \quad f_T(t) = n{(1-F_X(t))}^{n-1}f_\theta(t)
    \end{equation*}
    Si calculamos:
    \begin{equation*}
        F_X(t) = \int_{\theta}^{t} \dfrac{\theta}{x^2}~dx  = \theta \int_{\theta}^{t} \dfrac{1}{x^2}~dx = \theta \left[\dfrac{-1}{x}\right]_\theta^t = \theta\left(\dfrac{-1}{t}+\dfrac{1}{\theta}\right) = \dfrac{t-\theta}{t} \qquad  t > \theta
    \end{equation*}
    Tenemos que:
    \begin{equation*}
        f_T(t) = n{\left(1-\dfrac{t-\theta}{t}\right)}^{n-1}\dfrac{\theta}{t^2} = n\cdot \dfrac{\theta^{n}}{t^{n+1}}
    \end{equation*}
    Sea $h$ una función medible de forma que:
    \begin{equation*}
        0 = E[h(T)] = \int_{\theta}^{+\infty} h(t)f_T(t)~dt  = \lim_{n\to\infty}\int_{\theta}^{n} \underbrace{h(t)\cdot n\cdot \frac{\theta^n}{t^{n+1}}}_{(\ast)}~dt  \qquad \forall \theta\in \mathbb{R}^+
    \end{equation*}
    Sea $H(t)$ una primitiva de $(\ast)$, tenemos entonces que:
    \begin{equation*}
         0 = \lim_{n\to\infty}H(n) - H(\theta) \qquad \forall \theta\in \mathbb{R}^+
    \end{equation*}

    con lo que derivando respecto $\theta$:
    \begin{equation*}
        0 = h(\theta) \cdot n \cdot \dfrac{\theta^n}{\theta^{n+1}} = h(\theta)\cdot \dfrac{n}{\theta} \quad \Longrightarrow \quad  h(\theta) = 0 \qquad \forall \theta\in \mathbb{R}^+
    \end{equation*}
    Por lo que:
    \begin{equation*}
        \mathbb{R}^+ \subseteq \{t:h(t) = 0\}
    \end{equation*}

    de donde:
    \begin{equation*}
        1 \geq P[h(T) = 0] \geq P[T\in \mathbb{R}^+] = 1 \quad \Longrightarrow \quad P[h(T)=0] = 1
    \end{equation*}
    Luego tenemos que $T$ es completo. Si calculamos la esperanza de $T$:
    \begin{align*}
        E[T] &= \int_{\theta}^{\infty} t\cdot n\cdot \dfrac{\theta^n}{t^{n+1}}~dt  = \int_{\theta}^{+\infty} \dfrac{n\theta^n}{t^n}~dt  = n\theta^n \int_{\theta}^{+\infty} \dfrac{1}{t^n}~dt  = \dfrac{-n\theta^n}{n-1} \left[\frac{1}{t^{n-1}}\right]_\theta^\infty \\
             &= \dfrac{n\theta^n}{(n-1)\theta^{n-1}} = \dfrac{n\theta}{n-1}
    \end{align*}
    Por lo que tomando:
    \begin{equation*}
        h(t) = \dfrac{n-1}{n}\cdot t
    \end{equation*}

    tenemos que:
    \begin{equation*}
        E[h(T)] = \dfrac{n-1}{n}E[T] = \theta
    \end{equation*}
    Es decir, $h(T)$ es un estadístico insesgado para $\theta$. Veamos que es estimador y que tiene momento de segundo orden:
    \begin{itemize}
        \item Si $t\in \mathbb{R}^+$, tenemos entonces que $h(t) \in \mathbb{R}^+_0$ para $n>1$, por lo que $h(T)$ es estimador.
        \item Calculamos para $n>2$:
            \begin{align*}
                E[{(h(T))}^{2}] &= \int_{\theta}^{+\infty} {(h(t))}^{2}\dfrac{n\theta^n}{t^{n+1}}~dt  = {\left(\dfrac{n-1}{n}\right)}^{2}\int_{\theta}^{+\infty} \dfrac{n\theta^n t^2}{t^{n+1}}~dt  \\ 
                                &= n\theta^n {\left(\dfrac{n-1}{n}\right)}^{2}\int_{\theta}^{+\infty} \dfrac{1}{t^{n-1}}~dt  = n\theta^n {\left(\dfrac{n-1}{n}\right)}^{2} \left[\frac{1}{(-n+2)t^{n-2}}\right]_\theta^{+\infty} \\
                                &= n\theta^n {\left(\dfrac{n-1}{n}\right)}^{2}\dfrac{1}{(n-2)\theta^{n-2}} = {\left(\dfrac{n-1}{n}\right)}^{2} \dfrac{n}{(n-2)}\theta^2 < \infty
            \end{align*}
    \end{itemize}
    Por lo que $h(T) = \frac{n-1}{n}\cdot t$ es UMVUE para $\theta$. Si buscamos ahora un UMVUE para $\nicefrac{1}{\theta}$, observamos que tomando $h(t) = \frac{n+1}{n}\cdot \frac{1}{t}$, tenemos que:
    \begin{align*}
        E[h(T)] &= \int_{\theta}^{+\infty} h(t) \dfrac{n\theta^n}{t^{n+1}}~dt  = (n+1)\theta^n \int_{\theta}^{+\infty} \frac{1}{t^{n+2}}~dt  = (n+1)\theta^n \left[\frac{1}{-(n+1)t^{n+1}}\right]_\theta^{+\infty} \\
                &= \dfrac{\theta^n}{\theta^{n+1}} = \dfrac{1}{\theta}
    \end{align*}
    por lo que $h(T)$ es insesgado para $T$. Observamos además que:
    \begin{itemize}
        \item Si $t\in \mathbb{R}^+$, entonces $h(t)\in \mathbb{R}^+$, por lo que $h(T)$ es un estimador.
        \item Calculamos si $n+3>1$:
            \begin{align*}
                E[{(h(T))}^{2}] &= \int_{\theta}^{+\infty} {(h(t))}^{2}\dfrac{n\theta^n}{t^{n+1}}~dt  = \dfrac{\theta^n{(n+1)}^{2}}{n}\int_{\theta}^{+\infty} \dfrac{1}{t^{n+3}}~dt  \\
                &= \dfrac{\theta^n{(n+1)}^{2}}{n} \left[\frac{1}{-(n+2)t^{n+2}}\right]_\theta^{+\infty} = \dfrac{{(n+1)}^{2}}{n(n+2)\theta^2} < \infty
            \end{align*}
    \end{itemize}
    Aplicando el Teorema de Lehmann-Scheffé, tenemos que $E[h(T)/T] = h(T)$ es UMVUE para $\nicefrac{1}{\theta}$.
\end{ejercicio}

\begin{ejercicio}\label{ej:7_rel4}
    Sea $X\rightsquigarrow P_\theta$ siendo $P_\theta$ una distribución con función de densidad
    \begin{equation*}
        f_\theta(x) = e^{\theta-x}, \quad x\geq \theta
    \end{equation*}
    Dada una muestra aleatoria simple de tamaño arbitrario, encontrar los UMVUE de $\theta$ y de $e^{\theta}$.\\

    \noindent
    Calculamos en primer lugar un estadístico suficiente y completo para $\theta$:
    \begin{equation*}
        f_\theta(x_1, \ldots, x_n) \stackrel{\text{iid.}}{=} \prod_{i=1}^{n}f_\theta(x_i) = \prod_{i=1}^{n} e^{\theta-x_i} \cdot I_{\left[0,+\infty\right[}(X_{(1)}-\theta)=I_{\left[0,+\infty\right[}(X_{(1)}-\theta) e^{n\theta}\prod_{i=1}^{n} e^{-x_i}
    \end{equation*}
    Tomando:
    \begin{gather*}
        h(x_1, \ldots, x_n) = \prod_{i=1}^{n}e^{-x_i}, \qquad T(X_1, \ldots, X_n) = X_{(1)} \\
        g_\theta(t) = I_{\left[0,+\infty\right[}(t-\theta) \cdot e^{n\theta}
    \end{gather*}
    Tenemos por el Teorema de factorización de Neymann-Fisher que $T=X_{(1)}$ es suficiente para $\theta$. Calculamos la distribución de $T$:
    \begin{equation*}
        F_T(t) = 1-{(1-F_X(t))}^{n} \quad \Longrightarrow \quad f_T(t) = n{(1-F_X(t))}^{n-1}f_\theta(t)
    \end{equation*}
    Si calculamos:
    \begin{equation*}
        F_X(t) = \int_{\theta}^{t} e^{\theta-x}~dx  = e^{\theta}\int_{\theta}^{t} e^{-x}~dx  = e^{\theta}[-e^{-x}]_\theta^t = e^{\theta}\left(e^{-\theta} - e^{-t}\right) = 1- e^{\theta-t}, \quad t\geq \theta
    \end{equation*}
    Tenemos entonces que:
    \begin{equation*}
        f_T(t) = n{\left(e^{\theta-t}\right)}^{n-1} e^{\theta-t} = n{\left(e^{\theta-t}\right)}^{n}
    \end{equation*}
    Sea $h$ una función medible de forma que:
    \begin{equation*}
        0 = E[h(T)] = \int_{\theta}^{+\infty} h(t)n{\left(e^{\theta-t}\right)}^{n}~dt  = ne^{n\theta} \int_{\theta}^{+\infty}h(t)e^{-nt} ~dt \qquad \forall \theta\in \mathbb{R}
    \end{equation*}
    Tenemos entonces que:
    \begin{equation*}
        0 = \lim_{n\to\infty}\int_{\theta}^{n} h(t)e^{-nt}~dt  \qquad \forall \theta\in \mathbb{R}
    \end{equation*}
    Sea $H(t)$ una primiriva de $h(t)e^{-nt}$, tenemos que:
    \begin{equation*}
        0 = \lim_{n\to\infty}H(n)-H(\theta)\qquad \forall \theta\in \mathbb{R}
    \end{equation*}

    derivando respecto $\theta$ obtenemos:
    \begin{equation*}
        0 = -h(\theta)e^{-n\theta} \quad \Longrightarrow \quad  h(\theta) = 0 \qquad \forall \theta\in \mathbb{R}
    \end{equation*}
    Por lo que:
    \begin{equation*}
        \mathbb{R}\subseteq \{t:h(t) = 0\}
    \end{equation*}

    de donde deducimos que $T$ es completo:
    \begin{equation*}
        1 \geq P[h(t) = 0] \geq P[T\in \mathbb{R}] = 1 \quad \Longrightarrow \quad P[h(t)=0] = 1
    \end{equation*}
    Buscamos ahora un estadístico insesgado para $\theta$, es decir, buscamos una función $h$ medible de forma que:
    \begin{equation*}
        \theta = E[h(T)] = \int_{\theta}^{+\infty} h(t) n{(e^{\theta-t})}^{n}~dt  = ne^{n\theta}\int_{\theta}^{+\infty} h(t)e^{-nt}~dt 
    \end{equation*}
    Sea $H(t)$ una primitiva de $h(t)e^{-nt}$, tenemos entonces que:
    \begin{equation*}
        \theta = ne^{n\theta}\left(\lim_{n\to\infty}H(n)-H(\theta )\right) \Longrightarrow \frac{\theta}{ne^{n\theta}} = \left(\lim_{n\to\infty}H(n)-H(\theta )\right)
    \end{equation*}
    y si derivamos a ambos lados:
    \begin{equation*}
        \dfrac{e^{-n\theta}}{n} - \theta e^{-n\theta} = -h(\theta)e^{-n\theta} \quad \Longrightarrow \quad  h(\theta) = \theta-\frac{1}{n}
    \end{equation*}
    Por lo que si tomamos:
    \begin{equation*}
        h(T) = T - \dfrac{1}{n}
    \end{equation*}
    Tenemos que $E[T] = \theta$, con lo que $h(T)$ es insesgado para $\theta$. Además, se trata de un estimador (puesto que el espacio paramétrico es $\mathbb{R}$). Finalmente, vemos que:
    \begin{equation*}
        E[{(h(T))}^{2}] = \int_{\theta}^{+\infty} {(h(t))}^{2}n{\left(e^{\theta-t}\right)}^{n}~dt  = ne^{n\theta} \int_{\theta}^{+\infty} \left(t^2 + \frac{1}{n^2}-\frac{2t}{n}\right)e^{-nt}~dt  
    \end{equation*}
    calculamos la integral a parte:
    \begin{align*}
        \int_{\theta}^{+\infty} \left(t^2 + \frac{1}{n^2}-\frac{2t}{n}\right)e^{-nt}~dt &= \left\{\begin{array}{ll}
                                u = {(t-\nicefrac{1}{n})}^{2} & du = 2t -\nicefrac{2}{n} \\
                                dv = e^{-nt} & v = \nicefrac{-e^{-nt}}{n}
                        \end{array}\right\} \\
                                                                                        &= \left[\frac{-{(t-\nicefrac{1}{n})}^{2}e^{-nt}}{n}\right]_\theta^{+\infty} +\frac{1}{n}\int_{\theta}^{+\infty} (2t-\nicefrac{2}{n})e^{-nt}~dt 
    \end{align*}

    tenemos que:
    \begin{equation*}
        \left[\frac{-{(t-\nicefrac{1}{n})}^{2}e^{-nt}}{n}\right]_\theta^{+\infty} < \infty
    \end{equation*}

    y volvemos a calcular a parte:
    \begin{align*}
        \int_{\theta}^{+\infty} (2t-\nicefrac{2}{n})e^{-nt}~dt  &= \left\{\begin{array}{ll}
                u = 2t-\nicefrac{2}{n} & du = 2 \\
                dv = e^{-nt} & v = \nicefrac{-e^{-nt}}{n}
        \end{array}\right\} \\
                                                                &= \left[\frac{-(2t-\nicefrac{2}{n})e^{-nt}}{n}\right]_\theta^{+\infty} +\frac{2}{n} \int_{\theta}^{+\infty} e^{-nt}~dt 
    \end{align*}

    tenemos que:
    \begin{equation*}
        \left[\frac{-(2t-\nicefrac{2}{n})e^{-nt}}{n}\right]_\theta^{+\infty}  < \infty
    \end{equation*}

    y también que:
    \begin{equation*}
        \frac{2}{n}\int_{\theta}^{+\infty}e^{-nt} ~dt = \frac{2}{n}\left[\frac{-e^{-nt}}{n}\right]_\theta^{+\infty} < \infty
    \end{equation*}
    En definitiva, $E[{(h(T))}^{2}] < \infty$, por lo que aplicando el Teorema de Lehmann-Scheffé, obtenemos que $E[h(T)/T] = h(T)$ es UMVUE.\\ 

    \noindent
    Si buscamos ahora un estadístico insesgado para $e^\theta$, sea $h$ una función medible de forma que:
    \begin{equation*}
        e^\theta = E[h(T)] = \int_{\theta}^{+\infty} h(t)n{(e^{\theta-t})}^{n}~dt = ne^{n\theta}\int_{\theta}^{+\infty}h(t)e^{-nt} ~dt
    \end{equation*}
    Sea $H(t)$ una primitiva de $h(t)e^{-nt}$ tenemos entonces que:
    \begin{equation*}
        e^\theta = ne^{n\theta} \left(\lim_{n\to\infty}H(n)-H(\theta) \right) \Longrightarrow \frac{e^{\theta(1-n)}}{n}=\lim_{n\to\infty}H(n)-H(\theta) 
    \end{equation*}
    y si derivamos a ambos lados:
    \begin{equation*}
        \frac{(1-n)e^{\theta(1-n)}}{n} = -h(\theta) e^{-n\theta} \Longrightarrow h(t) = \frac{n-1}{n}e^t
    \end{equation*}
    Obtenemos que claramente $h(T)$ es estimador, así como que:
    \begin{align*}
        E[{(h(T))}^{2}] &= \int_{\theta}^{+\infty} {\left(\frac{n-1}{n}\right)}^{2} e^{2t}ne^{n\theta-nt}~dt  = {\left(\frac{n-1}{n}\right)}^{2}ne^{n\theta} \int_{\theta}^{+\infty} e^{(2-n)t}~dt  \\
                        &= \frac{1}{2-n}\left[e^{(2-n)t}\right]_\theta^\infty = \frac{e^{(2-n)\theta}}{n-2} < \infty \qquad \text{si\ } n>2
    \end{align*}
    Por lo que aplicando el Teorema de Lehmann-Scheffé, obtenemos que $E[h(T)/T] = h(T)$ es UMVUE para $e^{\theta}$.
\end{ejercicio}

\subsubsection{Ejercicios de estimadores eficientes}

\begin{ejercicio}
    Sea $X$ la variable que describe el número de fracasos antes del primer éxito en una sucesión de pruebas de Bernoulli con probabilidad de éxito $\theta\in \left]0,1\right[$, y sea $(X_1, \ldots, X_n)$ una muestra aleatoria simple de $X$.
    \begin{enumerate}[label=\alph*)]
        \item Probar que la familia de distribuciones de $X$ es regular y calcular la función de información asociada a la muestra.

            Según el enunciado, ``$X$ modela el número de fracasos antes del primer éxito en una sucesión de pruebas de Bernoulli con probabilidad de éxito $\theta\in \left]0,1\right[$'', por lo que $X$ sigue una distribución geométrica con probabilidad de éxito $\theta$, $X\rightsquigarrow G(\theta)$, $\theta \in \left]0,1\right[$. Para ver que la familia de distribuciones de $X$ es regular:
            \begin{enumerate}
                \item[$i)$] El espacio paramétrico es $\Theta = \left]0,1\right[$, un intervalo abierto de $\mathbb{R}$.
                \item[$ii)$] El espacio muestral es $\cc{X}=\mathbb{N}\cup \{0\}$, que no depende de $\theta$.
                \item[$iii)$] Para comprobar la tercera condición, recordamos que la función masa de probabilidad viene dada por:
                    \begin{equation*}
                        P_\theta[X=x] = {(1-\theta)}^{x}\theta
                    \end{equation*}
                    que es derivable respecto $\theta$. Calculamos la derivada del logaritmo de la función masa de probabilidad:
                    \begin{align*}
                        \ln P_\theta[X=x] &= \ln\left({(1-\theta)}^{x}\theta\right) = x\ln(1-\theta) + \ln\theta \\
                        \dfrac{\partial \ln P_\theta[X=x]}{\partial \theta} &= \dfrac{-x}{1-\theta} + \dfrac{1}{\theta} = \dfrac{1-\theta - x\theta}{\theta(1-\theta)} 
                    \end{align*}
                    Y si calculamos su esperanza:
                    \begin{equation*} 
                        E\left[\dfrac{1-\theta - X\theta}{\theta(1-\theta)} \right] = \dfrac{1-\theta - \theta E[X]}{\theta(1-\theta)} \AstIg \dfrac{1-\theta - (1-\theta)}{\theta(1-\theta)} = 0
                    \end{equation*}
                    donde en $(\ast)$ usamos que $E[X] = \frac{1-\theta}{\theta}$. En definitiva, se cumple también la tercera condición de las familias exponenciales.
            \end{enumerate}
            Calculamos ahora la función de información asociada a la muestra:
            \begin{align*}
                I_X(\theta) &= Var\left(\dfrac{1-\theta - X\theta}{\theta(1-\theta)}\right) = \dfrac{\theta^2}{\theta^2{(1-\theta)}^{2}}Var(X) \\
                            &\AstIg \dfrac{1-\theta}{{(1-\theta)}^{2}\theta^2}  = \dfrac{1}{\theta^2(1-\theta)}
            \end{align*}
            donde en $(\ast)$ hemos usado que $Var(X) = \frac{1-\theta}{\theta^2}$.
        \item Especificar la clase de funciones paramétricas que admiten estimadores eficientes y los correspondientes estimadores.

            Una vez que sabemos que la familia es exponencial y que $0<I_X(\theta)<\infty$, lo que hacemos es buscar las funciones $a(\theta)$ y $g(\theta)$ que se usan en el enunciado del Teorema de caracterización de los estimadores eficientes. Para ello:
            \begin{align*}
                \dfrac{\partial \ln P_\theta[X_1=x_1, \ldots, X_n = x_n]}{\partial \theta} &= \sum_{i=1}^{n} \dfrac{\partial  \ln P_\theta[X=x_i]}{\partial  \theta} = \sum_{i=1}^{n} \left(\dfrac{1-\theta-x_i\theta}{\theta(1-\theta)}\right) \\
                                                                                           &= \dfrac{n(1-\theta) - \theta\sum\limits_{i=1}^n x_i}{\theta(1-\theta)} = \dfrac{\sum\limits_{i=1}^n x_i - \frac{n(1-\theta)}{\theta}}{\theta-1}
            \end{align*}
            Por lo que tomando:
            \begin{equation*}
                T(X_1, \ldots, X_n) = \sum_{i=1}^{n}X_i, \qquad g(\theta) = \dfrac{n(1-\theta)}{\theta}, \qquad a(\theta) = \frac{1}{\theta-1}
            \end{equation*}
            Tenemos que $a(\theta)\neq 0 \quad \forall \theta\in \left]0,1\right[$, que $g$ es derivable, con:
            \begin{equation*}
                g'(\theta) = \frac{-n}{\theta^2} \neq 0 \qquad \forall \theta\in \left]0,1\right[
            \end{equation*}
            Además, observamos que:
            \begin{equation*}
                a(\theta)g'(\theta) = \frac{-n}{\theta^2(\theta-1)} = \frac{n}{\theta^2(1-\theta)} = nI_X(\theta) = I_{(X_1, \ldots, X_n)}(\theta) \qquad \forall \theta\in \left]0,1\right[
            \end{equation*}
            Así como que $T$ es un estimador, pues: 
            \begin{equation*}
                T(\cc{X}^n) = [0,n]\cap \mathbb{N} \subset g(\left]0,1\right[)
            \end{equation*}
            Finalmente, por un Corolario visto en teoría sabemos que las únicas funciones paramétricas que admiten estimadores eficientes son de la forma:
            \begin{equation*}
                a\cdot \frac{n(1-\theta)}{\theta} + b, \qquad a,b\in \mathbb{R}, \quad a\neq 0
            \end{equation*}
            y que sus estimadores eficientes son de la forma:
            \begin{equation*}
                a\cdot \sum_{i=1}^{n}X_i + b
            \end{equation*}
        \item Calcular la varianza de cada estimador eficiente y comprobar que coincide con las correspondiente cota de Fréchet-Cramér-Rao.

            Si calculamos ahora la varianza de $T=a\cdot \sum\limits_{i=1}^n X_i$ y la cota de Fréchet-Cramér-Rao:
            \begin{align*}
                Var(T) &= Var\left(a\cdot \sum_{i=1}^{n}X_i+b\right) \stackrel{\text{indep.}}{=} a^2\sum_{i=1}^{n}Var(X_i) = a^2\sum_{i=1}^{n} \left(\frac{1-\theta}{\theta^2}\right) \\ &= \frac{a^2n(1-\theta)}{\theta^2} \\
                \frac{{(ag'(\theta))}^{2}}{I_{(X_1,\ldots,X_n)}(\theta)} &= \frac{{\left(\frac{-an}{\theta^2}\right)}^{2}}{\frac{n}{\theta^2(1-\theta)}} = \frac{a^2n(1-\theta)}{\theta^2}
            \end{align*}
        \item Calcular, si existen, los UMVUE para $P_\theta[X=0]$ y para $E_\theta[X]$ y decir si son eficientes. 
            \begin{itemize}
                \item Para $P_\theta[X=0]$:
                    \begin{equation*}
                        P_\theta[X=0] = {(1-\theta)}^{0}\theta = \theta
                    \end{equation*}
                    Como ya hemos demostrado que $T = \sum\limits_{i=1}^{n}X_i$ es un estimador eficiente para $g(\theta)=\frac{n(1-\theta)}{\theta}$, tenemos automáticamente que $T$ es un estadístico suficiente. Tomando ahora $Q(\theta)$ como una primitiva de $a(\theta)$, por ejemplo $Q(\theta) = \ln(|\theta-1|)$, como:
                    \begin{equation*}
                        Q(\left]0,1\right[) = \left]-\infty,0\right[
                    \end{equation*}
                    claramente contiene un abierto de $\mathbb{R}$ concluimos que $T$ es completo. Buscamos ahora una función $h$ medible de forma que $h(T)$ sea insesgado para $\theta$, sabiendo que $T\rightsquigarrow BN(n,\theta)$, con función masa de probabilidad:
                    \begin{equation*}
                        P_\theta[T = t] = \binom{t+n-1}{t} {(1-\theta)}^{t}\theta^n
                    \end{equation*}
                    Buscamos pues una función $h$ medible de forma que:
                    \begin{equation*}
                        E[h(T)] = \sum_{t=0}^{\infty} h(t) \dfrac{(t+n-1)!}{t!(n-1)!}{(1-\theta)}^{t}\theta^n = \theta
                    \end{equation*}
                    Es decir, que cumpla (dividiendo todo entre $\theta$):
                    \begin{equation*}
                        E[h(T)] = \sum_{t=0}^{\infty} h(t) \dfrac{(t+n-1)!}{t!(n-1)!}{(1-\theta)}^{t}\theta^{n-1} = 1
                    \end{equation*}
                    Notemos que si $Y\rightsquigarrow B(n-1,\theta)$, tendremos entonces que:
                    \begin{equation*}
                        P[Y=t] = \binom{t+n-2}{t}{(1-\theta)}^{t}\theta^{n-1} = \dfrac{(t+n-2)!}{t!(n-2)!}{(1-\theta)}^{t}\theta^{n-1}
                    \end{equation*}
                    Por lo que ha de ser (por ser una función masa de probabilidad):
                    \begin{equation*}
                        1 = \sum_{t=0}^{\infty}\dfrac{(t+n-2)!}{t!(n-2)!}{(1-\theta)}^{t}\theta^{n-1}
                    \end{equation*}
                    Por tanto, si tomamos:
                    \begin{equation*}
                        h(t) = \dfrac{n-1}{t+n-1}
                    \end{equation*}
                    obtenemos lo buscado. Observamos ahora que:
                    \begin{itemize}
                        \item $h(T)$ es un estimador, puesto que $h(T)\in \left]0,1\right[$.
                        \item Calculamos su momento de segundo orden:
                            \begin{align*}
                                E[{(h(t))}^{2}] &=\sum_{t=0}^{\infty}{(h(t))}^{2}\dfrac{(t+n-1)!}{t!(n-1)!}{(1-\theta)}^{t}\theta^n \\
                                                &=\sum_{t=0}^{\infty}{\left(\frac{n-1}{t+n-1}\right)}^{2}\dfrac{(t+n-1)!}{t!(n-1)!}{(1-\theta)}^{t}\theta^n
                            \end{align*}
                            Como tenemos que $h(t)\in \left]0,1\right[$ para todo $t\in \mathbb{N}$, tenemos que ${(h(t))}^{2}\in \left]0,1\right[$, por lo que cada sumando de esta nueva serie es menor o igual que cada sumando de la serie $E[h(t)]$, que era convergente, por lo que por el criterio de comparación esta nueva serie también será convergente, $E[{(h(t))}^{2}]<\infty$.
                    \end{itemize}
                    En definitiva, el UMVUE para $P_\theta[X=0]=\theta$ es $h(T)$, por el Teorema de Lehmann-Scheffé.
                \item Para $E_\theta[X]$:
                    \begin{equation*}
                        E[X] = \frac{1-\theta}{\theta}
                    \end{equation*}
                    Por el apartado anterior tenemos ya que $T=\sum\limits_{i=1}^{n}X_i$ es un estadístico suficiente y completo. Buscamos ahora una función $h$ medible tal que $h(T)$ sea insesgado para $\frac{1-\theta}{\theta}$, es decir:
                    \begin{equation*}
                        E[h(T)] = \sum_{t=0}^{\infty}h(t) \frac{(t+n-1)!}{t!(n-1)!}{(1-\theta)}^{t}\theta^n = \frac{1-\theta}{\theta}
                    \end{equation*}
                    Escribiendo los términos en el mismo lado para igualar a 1:
                    \begin{equation}\label{eq:despejar_ej8}
                        \sum_{t=0}^{\infty}h(t) \frac{(t+n-1)!}{t!(n-1)!}{(1-\theta)}^{t-1}\theta^{n+1} = 1
                    \end{equation}
                    Si $Y\rightsquigarrow BN(n+1,\theta)$, entonces:
                    \begin{equation*}
                        P[Y=t] = \binom{n+t}{t}{(1-\theta)}^{t}\theta^{n+1} = \dfrac{(t+n)!}{t!n!}{(1-\theta)}^{t}\theta^{n+1}
                    \end{equation*}
                    De donde sabemos que:
                    \begin{equation*}
                        \sum_{t=0}^{\infty}\dfrac{(t+n)!}{t!n!}{(1-\theta)}^{t}\theta^{n+1} = 1
                    \end{equation*}
                    Si sacamos fuera de la suma de~(\ref{eq:despejar_ej8}) el primer sumando, tenemos:
                    \begin{equation*}
                        h(0)\frac{\theta}{1-\theta} + \sum_{t=1}^{\infty}h(t)\frac{(t+n-1)!}{t!(n-1)!}{(1-\theta)}^{t-1}\theta^{n+1}
                    \end{equation*}
                    y podemos comenzar a numerar la serie desde 0:
                    \begin{equation*}
                        h(0)\frac{\theta}{1-\theta} + \sum_{t=0}^{\infty}h(t+1)\frac{(t+n)!}{(t+1)!(n-1)!}{(1-\theta)}^{t}\theta^{n+1}
                    \end{equation*}
                    De esta forma, si tomamos:
                    \begin{equation*}
                        h(t) = \frac{t}{n}
                    \end{equation*}
                    tendremos en particular que $h(0) = 0$, con lo que el primer sumando se anula y el segundo coincide con la suma en $\mathbb{N}$ de la función masa de probablidad de una variable aleatoria con distribución binomial negativa, que ha de sumar 1, como queríamos en la igualdad~(\ref{eq:despejar_ej8}). Observemos ahora que:
                    \begin{itemize}
                        \item $h(t)\in \left[0,1\right]$ para todo $t\in \{0, \ldots, n\}$, por lo que $h(T)$ es un estimador.
                        \item Para el momento de segundo orden:
                            \begin{equation*}
                                E[{(h(t))}^{2}] = \sum_{t=0}^{\infty}\frac{t^2}{n^2}\frac{(t+n-1)!}{t!(n-1)!}{(1-\theta)}^{t}\theta^n 
                            \end{equation*}
                            Aplicamos el criterio del cociente:
                            \begin{align*}
                                &\dfrac{\dfrac{{(t+1)}^{2}(t+n)!}{\cancel{n^2}(t+1)!\cancel{(n-1)!}}{(1-\theta)}^{t+1}\cancel{\theta^n}}{\dfrac{t^2(t+n-1)!}{\cancel{n^2}t!\cancel{(n-1)!}}{(1-\theta)}^{t}\cancel{\theta^n}} = \dfrac{{(t+1)}^{2}(t+n)!t!(1-\theta)}{(t+1)!t^2(t+n-1)!} \\ &= \dfrac{{(t+1)}^{2}(t+n)(1-\theta)}{t^2(t+1)} = \dfrac{(t+1)(t+n)(1-\theta)}{t^2} \rightarrow 1-\theta
                            \end{align*}
                            Con $\theta\in \left]0,1\right[$, por lo que $1-\theta\in \left]0,1\right[$, de donde la serie es convergente.
                    \end{itemize}
                    En definitiva, el UMVUE para $E_\theta[X]=\frac{1-\theta}{\theta}$ es $h(T)$, por el Teorema de Lehmann-Scheffé.
            \end{itemize}
            En los apartados anteriores observamos que la familia de funciones paramétricas que admitían estimadores eficientes eran las de la forma:
            \begin{equation*}
                a\cdot \frac{n(1-\theta)}{\theta} + b \qquad a,b\in \mathbb{R}, \quad a\neq 0
            \end{equation*}
            Y sus estimadores eficientes eran:
            \begin{equation*}
                a\cdot \sum_{i=1}^{n}X_i +b
            \end{equation*}
            Por lo que:
            \begin{itemize}
                \item Como $\theta$ no está en dicha familia de funciones paramétricas, el UMVUE obtenido para $\theta$ no puede ser eficiente.
                \item Si tomamos $a = \frac{1}{n}$ y $b=0$, tenemos que l UMVUE obtenido para $\frac{1-\theta}{\theta}$ coincide con el estimador eficiente en dicho caso.
            \end{itemize}
    \end{enumerate}
\end{ejercicio}

\begin{ejercicio}\label{ej:9_rel4} 
    Sea $(X_1, \ldots, X_n)$ una muestra aleatoria simple de una variable aleatoria $X$ con distribución exponencial.
    \begin{enumerate}[label=\alph*)]
        \item Probar que la familia de distribuciones de $X$ es regular.

            Supuesto que $X\rightsquigarrow exp(\lm)$ con $\lm\in \mathbb{R}^+$, la función de densidad de $X$ viene dada por:
            \begin{equation*}
                f_\lm(x) = \lm e^{-\lm x} \qquad x\geq 0
            \end{equation*}
            Y se cumple:
            \begin{equation*}
                E[X] = \frac{1}{\lm}, \qquad Var(X) = \frac{1}{\lm^2}
            \end{equation*}
            Comprobamos cada una de las propiedades:
            \begin{enumerate}
                \item[$i)$] El espacio paramétrico es $\Theta=\mathbb{R}^+$, intervalo abierto de $\mathbb{R}$.
                \item[$ii)$] El espacio muestral es $\mathbb{R}^+_0$, que no depende de $\lm$.
                \item[$iii)$] Para la tercera, observamos que $f_\lm$ es derivable respecto $\lm$, con:
                    \begin{equation*}
                        \ln f_\lm(X) = \ln\left(\lm e^{-\lm X}\right) = \ln\lm - \lm X \quad \Longrightarrow \quad  \dfrac{\partial \ln f_\lm(X)}{\partial \lm} = \frac{1}{\lm} - X
                    \end{equation*}
                    Y si calculamos:
                    \begin{equation*}
                        E\left[\dfrac{\partial \ln f_\lm(X)}{\partial \lm}\right] = E\left[\frac{1}{\lm}-X\right] = \frac{1}{\lm}-E[X] = \frac{1}{\lm}-\frac{1}{\lm} = 0
                    \end{equation*}
            \end{enumerate}
        \item Encontrar la clase de funciones paramétricas que admiten estimador eficiente y el estimador correspondiente. Calcular la varianza de estos estimadores.

            Calculamos primero la función de información de Fisher:
            \begin{equation*}
                I_X(\lm) = Var\left(\dfrac{\partial \ln f_\lm(X)}{\partial \lm}\right) = Var\left(\frac{1}{\lm}-X\right) = Var(X) = \frac{1}{\lm^2}
            \end{equation*}
            que como vemos verifica $0<I_X(\lm) <\infty$. Buscamos ahora las funciones $a$ y $g$:
            \begin{equation*}
                \dfrac{\partial \ln f_\lm^n(x_1, \ldots, x_n)}{\partial \lm} = \sum_{i=1}^{n}\dfrac{\partial \ln f_\lm(x_i)}{\partial \lm} = \sum_{i=1}^{n} \left(\frac{1}{\lm}-x_i\right) = \frac{n}{\lm} - \sum_{i=1}^{n}x_i = -\left(\sum_{i=1}^{n}x_i - \frac{n}{\lm}\right)
            \end{equation*}
            Por lo que tomando:
            \begin{equation*}
                T(X_1, \ldots, X_n) = \sum_{i=1}^{n} X_i, \qquad g(\lm) = \frac{n}{\lm}, \qquad a(\lm) = -1
            \end{equation*}
            Tenemos que $a(\lm) \neq 0 \quad \forall \lm\in \mathbb{R}^+$, que $g$ es derivable con:
            \begin{equation*}
                g'(\lm) = \frac{-n}{\lm^2} \neq 0 \qquad \forall \lm\in \mathbb{R}^+
            \end{equation*}
            y que $T$ es un estimador, pues $T(\mathbb{R}^+) = \mathbb{R}^+ = g(\mathbb{R}^+)$. Finalmente, falta comprobar que:
            \begin{equation*}
                a(\lm)g'(\lm) = \frac{n}{\lm^2} = nI_X(\lm) = I_{(X_1, \ldots, X_n)}(\lm)
            \end{equation*}
            Por lo que podemos deducir ya por el Teorema de caracterización de los estimadores eficientes que $T$ es un estimador eficiente. Más aún, por un Corolario del mismo, tenemos que las únicas funciones paramétricas de dicha familia que admiten estimadores eficientes son de la forma:
            \begin{equation*}
                a\cdot \frac{n}{\lm}+b, \qquad a,b\in \mathbb{R}\quad a\neq 0
            \end{equation*}
            y que dichos estimadores son $a\cdot T+b$. Nos disponemos a calcular la varianza de estos estimadores:
            \begin{equation*}
                Var\left(a\cdot T+b\right) = a^2 Var\left(\sum_{i=1}^{n}X_i\right) = a^2\sum_{i=1}^{n}Var(X) = a^2 \sum_{i=1}^{n} \left(\frac{1}{\lm^2}\right) = \frac{na^2}{\lm^2}
            \end{equation*}
            Veamos que coincide con la cota de Fréchet-Cramér-Rao, para comprobar nuestro trabajo calculando la familia de estimadores eficientes:
            \begin{equation*}
                \dfrac{{(ag'(\lm))}^{2}}{I_{(X_1,\ldots,X_n)}(\lm)} = \dfrac{{\left(\frac{-an}{\lm^2}\right)}^{2}}{\frac{n}{\lm^2}} = \frac{na^2}{\lm^2} = Var(a\cdot T+b)
            \end{equation*}
        \item Basándose en el apartado anterior, encontrar el UMVUE para la media de $X$.

            Como ya sabemos que $T=\sum\limits_{i=1}^{n}X_i$ es un estimador eficiente para $g(\lm)$, tenemos automáticamente que $T$ es suficiente. Para ver que $T$ es completo, sea $Q$ una primitiva de $a$, podemos tomar $Q(\lm) = -\lm$, tenemos que claramente $Q(\mathbb{R}^+)$ contiene un abierto de $\mathbb{R}$, por lo que $T$ es completo. 

            Como en el apartado anterior vimos que $T$ es eficiente para $g(\lm) = \frac{n}{\lm}$, tenemos entonces que $\frac{T}{n}$ es eficiente para $\tilde{g}(\lm) = \frac{g(\lm)}{n} = \frac{1}{\lm}$. Como la correspondencia $t \mapsto \frac{t}{n}$ es biunívoca, $\frac{T}{n}$ también será suficiente y completo (para ver que es completo no hace falta que la corresponencia sea unívoca). Por un Corolario del Teorema de caracterización de estimadores eficientes tenemos que $T$ es UMVUE para $\frac{1}{\lm} = E[X]$.

        \item Dar la cota de Fréchet-Cramér-Rao para la varianza de estimadores insesgados y regulares de $\lm^3$. ¿Es alcanzable dicha cota?

            Si tomamos $g(\lm) = \lm^3$, la cota de Fréchet-Cramér-Rao tiene la forma:
            \begin{equation*}
                \frac{{(g'(\lm))}^{2}}{I_{(X_1,\ldots,X_n)}(\lm)} = \frac{{(3\lm^2)}^{2}}{\frac{n}{\lm^2}} = \dfrac{9\lm^6}{n}
            \end{equation*}
            y como $g(\lm)$ no es de la forma $a\cdot \frac{n}{\lm}+b$, para ciertos $a,b\in \mathbb{R}$, sabemos que esta cota no es alcanzable.
    \end{enumerate}
\end{ejercicio}

\begin{ejercicio}\label{ej:10_rel4} 
    Sea $X$ una variable aleatoria con función de densidad de la forma
    \begin{equation*}
        f_\theta(x) = \theta x^{\theta-1}, \quad 0<x<1
    \end{equation*}
    \begin{enumerate}[label=\alph*)]
        \item Sabiendo que $E_\theta[\ln X] = -\frac{1}{\theta}$ y $Var_\theta[\ln X] = \frac{1}{\theta^2}$, comprobar que esta familia de distribuciones es regular.

            Veamos las propiedades, suponiendo que $\theta > 0$.
            \begin{enumerate}
                \item[$i)$] El espacio paramétrico es $\Theta = \mathbb{R}^+$, intervalo abierto de $\mathbb{R}$.
                \item[$ii)$] El espacio muestral es $\cc{X} = \left]0,1\right[$, que no depende de $\theta$.
                \item[$iii)$] La función de densidad es derivable respecto $\theta$, calculamos la derivada del logaritmo de la función de densidad:
                    \begin{equation*}
                        \ln f_\theta(X) = \ln(\theta X^{\theta-1}) = \ln\theta + (\theta-1)\ln X \quad \Longrightarrow \quad \dfrac{\partial \ln f_\theta(X)}{\partial \theta} = \frac{1}{\theta} + \ln X
                    \end{equation*}
                    Con lo que:
                    \begin{equation*}
                        E\left[\dfrac{\partial \ln f_\theta(X)}{\partial \theta}\right] = E\left[\frac{1}{\theta} + \ln X\right] = \frac{1}{\theta} + E[\ln X] = \frac{1}{\theta} - \frac{1}{\theta} = 0
                    \end{equation*}
            \end{enumerate}
        \item Basándose en una muestra aleatoria simple de $X$, dar la clase de funciones paramétricas con estimador eficiente, los estimadores y su varianza.

            Sea $(X_1, \ldots, X_n)$ una muestra aleatoria simple de $X$, en primer lugar calculamos la función de información de Fisher:
            \begin{equation*}
                I_X(\theta) = Var\left(\dfrac{\partial \ln f_\theta(X)}{\partial \theta}\right) = Var\left(\frac{1}{\theta} + \ln X\right) = Var\left(\ln X\right) = \frac{1}{\theta^2}
            \end{equation*}
            y observamos que $0<I_X(\theta) < \infty$. Ahora, buscamos las funciones $g$ y $a$:
            \begin{equation*}
                \dfrac{\partial \ln f_\theta^n(x_1, \ldots, x_n)}{\partial \theta} = \sum_{i=1}^{n} \dfrac{\partial  \ln f_\theta(x_i)}{\partial \theta} = \sum_{i=1}^{n} \left(\frac{1}{\theta}+\ln x_i\right) = \frac{n}{\theta} + \sum_{i=1}^{n}\ln x_i 
            \end{equation*}
            Tomando:
            \begin{equation*}
                T(X_1, \ldots, X_n) = \sum_{i=1}^{n}\ln X_i, \qquad g(\theta) = \frac{-n}{\theta}, \qquad a(\theta) = 1
            \end{equation*}

            tenemos que $a(\theta)\neq 0$ para todo $\theta\in \mathbb{R}^+$, que $g$ es derivable con:
            \begin{equation*}
                g'(\theta) = \frac{n}{\theta^2} \neq 0 \qquad \forall \theta\in \mathbb{R}^+
            \end{equation*}

            además de que $T$ es un estimador, puesto que:
            \begin{equation*}
                T({(\left]0,1\right[)}^{n}) \subseteq \left]-\infty,0\right[ = g(\mathbb{R}^+)
            \end{equation*}
            Comprobemos que:
            \begin{equation*}
                a(\theta)g'(\theta) = \frac{n}{\theta^2} = nI_X(\theta) = I_{(X_1, \ldots, X_n)}(\theta)
            \end{equation*}
            Por lo que $T$ es un estimador eficiente para $g(\theta) = \frac{-n}{\theta}$, luego las únicas funciones paramétricas que admiten estimadores eficientes son las de la forma:
            \begin{equation*}
                a\cdot \frac{-n}{\theta} + b, \qquad a,b\in \mathbb{R}, \quad a\neq 0
            \end{equation*}
            cuyos estimadores eficientes son:
            \begin{equation*}
                a\cdot T + b
            \end{equation*}
            Calculemos sus varianzas y veamos que coinciden con la cota de Fréchet-Cramér-Rao:
            \begin{align*}
                Var(a\cdot T+b) &= a^2 Var\left(\sum_{i=1}^{n}\ln X_i\right) \stackrel{\text{indep.}}{=} a^2 \sum_{i=1}^{n}Var(\ln X_i) = a^2 \sum_{i=1}^{n} \frac{1}{\theta^2} = \frac{a^2n}{\theta^2} \\
                \frac{{(ag'(\theta))}^{2}}{I_{(X_1,\ldots, X_n)}(\theta)} &= \frac{{(\frac{an}{\theta^2})}^{2}}{\frac{n}{\theta^2}} = \frac{a^2n}{\theta^2}
            \end{align*}
    \end{enumerate}
\end{ejercicio}

    % \section{Estimación de máxima verosimilitud y otros métodos}
\begin{ejercicio}
    Sea $X\rightsquigarrow P_\theta$ con $\theta\in \mathbb{R}$ siendo $P_\theta$ una distribución con función de densidad
    \begin{equation*}
        f_\theta(x) = e^{\theta-x}, \qquad x\geq \theta
    \end{equation*}
    Dada una muestra aleatoria simple de tamaño $n$, encontrar los estimadores máximo verosímiles de $\theta$ y de $e^\theta$. Basándose en los resultados del Ejercicio~\ref{ej:7_rel4}, decir si estos estimadores son insesgados.\\

    \noindent
    Sea $(X_1, \ldots, X_n)$ una m.a.s. de $X\rightsquigarrow  P_\theta$, tenemos que:
    \begin{equation*}
        f^n_\theta(x_1, \ldots, x_n) \stackrel{\text{iid.}}{=} \prod_{i=1}^{n}f_\theta(x_i) = \prod_{i=1}^{n} e^{\theta-x} = e^{n\theta - \sum\limits_{i=1}^n x_i} \qquad \forall x_i \geq \theta
    \end{equation*}
    Por lo que si consideramos $x_{(1)} = \min\limits_{i \in \{1,\ldots,n\}}\{x_i\}$:
    \begin{equation*}
        L_{x_1,\ldots, x_n}(\theta) = f^n_\theta(x_1,\ldots,x_n) = \left\{\begin{array}{ll}
                e^{n\theta - \sum\limits_{i=1}^n x_i}  & \text{si\ } \theta \leq x_{(1)}  \\
             0 & \text{en otro caso\ } 
        \end{array}\right. 
    \end{equation*}
    Observamos que $L_{x_1,\ldots,x_n}(\theta)$ es una función creciente en $\left]-\infty,x_{(1)}\right]$, por lo que alcanza su máximo en $x_{(1)}$. Como el espacio paramétrico es $\Theta = \mathbb{R}$, tenemos que $T(X_1, \ldots, X_n) = X_{(1)}$ es un estimador de $\theta$, con lo que es el EMV, al maximizar la función de verosimilitud. Además, el Teorema de Zehna nos dice que:
    \begin{equation*}
        \widehat{e^\theta} = e^{\hat{\theta}} = e^{X_{(1)}}
    \end{equation*}
    Por lo que $e^{X_{(1)}}$ es EMV para $e^{\theta}$.\\ 

    \noindent
    Buscamos ahora razonar si los estimadores obtenidos son o no insesgados, a partir de los resultados del Ejercicio~\ref{ej:7_rel4}. Si suponemos que los EMVs que hemos obtenido son insesgados, vimos anteriormente en dicho ejercicio que el mínimo era suficiente y completo, por lo que tendríamos una función del suficiente y completo que es insesgada. Si comprobamos que también tiene momento de segundo orden finito y que es estimador (hágase), tendríamos que es un UMVUE distinto al anterior, pero la unicidad del UMVUE nos dice que estos han de ser iguales, contradicción que viene de suponer que los EMVs son insesgados.
\end{ejercicio}

\begin{ejercicio}
    Sea $(X_1, \ldots, X_n)$ una muestra aleatoria simple de una variable aleatoria $X$ con distribución exponencial. Basándose en los resultados del Ejercicio~\ref{ej:9_rel4}, encontrar los estimadores máximo verosímiles de la media y de la varianza de $X$.\\

    \noindent
    Sea $X\rightsquigarrow exp(\lm)$ con $\lm\in \mathbb{R}^+$, nos piden encontrar EMVs para:
    \begin{equation*}
        E[X] = \frac{1}{\lm}, \qquad Var(X) = \frac{1}{\lm^2}
    \end{equation*}
    En el Ejercicio~\ref{ej:9_rel4} obtuvimos que las únicas funciones paramétricas de dicha familia de distribuciones que admiten estimadores eficientes son las de la forma:
    \begin{equation*}
        a\cdot \frac{n}{\lm} +b \qquad a,b\in \mathbb{R}, \quad a\neq 0
    \end{equation*}
    y sus estimadores eficientes son:
    \begin{equation*}
        a\cdot \sum_{i=1}^{n}X_i + b
    \end{equation*}
    Tomando $a=\frac{1}{n}, b=0$, obtenemos que el estimador $T=\frac{\sum\limits_{i=1}^nX_i}{n}$ es eficiente para $E[X] = \nicefrac{1}{\lm}$, por lo que un Teorema visto en teoría nos asegura que $T$ es el único EMV que podemos considerar para $\nicefrac{1}{\lm}$. Para $\nicefrac{1}{\lm^2}$ podemos aplicar el Teorema de Zehna, obteniendo:
    \begin{equation*}
          \boxed{ \widehat{\frac{1}{\lambda^2}} = \widehat{\frac{1}{\lambda}}^2 = \left(\frac{\sum_l{i=1}^n X_i}{n}\right)^2 } 
    \end{equation*}
    que es estimador, luego es EMV.
\end{ejercicio}

\begin{ejercicio}
    Sea $X$ una variable aleatoria con función de densidad de la forma
    \begin{equation*}
        f_\theta(x) = \theta x^{\theta -1}, \qquad 0<x<1
    \end{equation*}
    \begin{enumerate}[label=\alph*)]
        \item Calcular un estimador máximo verosímil para $\theta$.

            Maximizamos la función de verosimilitud de $\theta$:
            \begin{equation*}
                L_{x_1,\ldots,x_n}(\theta) = f^n_\theta(x_1, \ldots, x_n) \stackrel{\text{iid.}}{=}\prod_{i=1}^{n}f_\theta(x_i) = \prod_{i=1}^{n} \theta x_i^{\theta-1} = \theta^n \prod_{i=1}^{n} x_i^{\theta -1}
            \end{equation*}
            Si aplicamos logaritmo no cambia la abscisa a maximizar:
            \begin{equation*}
                \ln L_{x_1,\ldots,x_n}(\theta) = n\ln \theta + (\theta -1)\sum_{i=1}^{n} \ln x_i
            \end{equation*}
            hayamos los puntos críticos:
            \begin{equation*}
                \dfrac{\partial \ln L_{x_1,\ldots,x_n}(\theta)}{\partial  \theta} = \frac{n}{\theta} + \sum_{i=1}^{n}\ln x_i = \frac{n + \theta \sum\limits_{k=1}^{n}\ln x_i}{\theta} = 0 \Longleftrightarrow \theta = \frac{-n}{\sum\limits_{k=1}^{n}\ln x_i}
            \end{equation*}
            Por lo que tomando:
            \begin{equation*}
                \hat{\theta}(X_1, \ldots, X_n) = \frac{-n}{\sum\limits_{i=1}^{n} \ln X_i}
            \end{equation*}
            tenemos que $\hat{\theta}$ es un estimador de $\theta$, luego es EMV para $\theta$.
        \item Deducir dicho estimador a partir de los resultados del Ejercicio~\ref{ej:10_rel4}.

            En el Ejercicio~\ref{ej:10_rel4} vimos que las únicas funciones paramétricas que admiten un estimador eficiente son las de la forma:
            \begin{equation*}
                a\cdot \frac{-n}{\theta} +b\qquad a,b\in \mathbb{R}, \quad a\neq 0
            \end{equation*}
            y sus estimadores eficientes son:
            \begin{equation*}
                a\cdot \sum_{i=1}^{n}\ln X_i + b
            \end{equation*}
            Tomando $a=\frac{-1}{n},b=0$, tenemos que el estimador:
            \begin{equation*}
                T = \frac{\sum\limits_{i=1}^{n}X_i}{-n}
            \end{equation*}
            es eficiente para $\frac{1}{\theta}$, por lo que $T$ es el EMV de $\frac{1}{\theta}$. Si aplicamos el Teorema de Zehna, tenemos que:
            \begin{equation*}
                \widehat{\left(\frac{1}{\theta}\right)} = \frac{1}{\hat{\theta}} \Longrightarrow \hat{\theta} = \frac{1}{\widehat{\left(\frac{1}{\theta}\right)} } = \frac{-n}{\sum\limits_{i=1}^{n}\ln X_i}
            \end{equation*}
            Observamos que obtenemos el mismo resultado que en el primer apartado.
    \end{enumerate}
\end{ejercicio}

\begin{ejercicio}
    Sea $(X_1, \ldots, X_n)$ una muestra de una variable $X\rightsquigarrow B(k_0,p)$ para cierto $k_0\in \mathbb{N}$ y $p\in \left]0,1\right[$. Estimar, por máxima verosimilitud y por el método de los momentos, el parámetro $p$ y la varianza de $X$.\\

    \noindent
    \textit{Aplicación}: Se lanza $10$ veces un dado cargado y se cuenta el número de veces que sale un $4$. Este experimento se realiza $100$ veces de forma independiente, obteniéndose los siguientes resultados:
    \begin{table}[H]
    \centering
    \begin{tabular}{l|cccc}
        nº de 4 & 0 & 1 & 2 & 3 \\
        \hline
        frecuencia & 84 & 15 & 1 & 0
    \end{tabular}
    \end{table}
    \noindent
    Estimar, a partir de estos datos, la probabilidad de salir un cuatro.\\

    \begin{description}
        \item [Por máxima verosimilitud.] 
            Tenemos:
            \begin{align*}
                P[X_1 =x_1, \ldots, X_n = x_n] &\stackrel{\text{iid.}}{=}\prod_{i=1}^{n} P[X=x_i] = \prod_{i=1}^{n} \binom{k_0}{x_i} p^{x_i}{(1-p)}^{k_0-x_i} \\ &= p^{\sum\limits_{i=1}^{n}x_i}{(1-p)}^{nk_0-\sum\limits_{i=1}^{n}x_i} \prod_{i=1}^{n}\binom{k_0}{n}
            \end{align*}
            Por lo que:
            \begin{equation*}
                \ln L_{x_1,\ldots,x_n}(p) = \ln p \sum_{i=1}^{n}x_i + \ln(1-p)\left(nk_0 - \sum_{i=1}^{n}x_i\right) + \sum_{k=1}^{n}\ln\binom{k_0}{n}
            \end{equation*}
            Luego:
            \begin{align*}
                \dfrac{\partial \ln L_{x_1,\ldots,x_n}(p)}{\partial p} &= \frac{1}{p}\sum_{i=1}^{n}x_i - \frac{1}{1-p}\left(nk_0 - \sum_{i=1}^{n}x_i\right) = \frac{(1-p)\sum\limits_{i=1}^{n}x_i - pnk_0 + p\sum\limits_{i=1}^{n}x_i}{p(1-p)} \\
                                                                       &= \frac{\sum\limits_{i=1}^{n}x_i -pnk_0}{p(1-p)} = 0 \Longleftrightarrow p = \frac{\sum\limits_{i=1}^{n}x_i}{nk_0}
            \end{align*}
            Como $\sum\limits_{i=1}^{n}x_i\leq nk_0$, tenemos entonces que $p\in [0,1]$, por lo que esta fórmula nos dará una estimador, con lo que el EMV es:
            \begin{equation*}
                \hat{p} = \frac{1}{nk_0}\sum_{i=1}^{n}X_i = \frac{1}{k_0}\overline{X}
            \end{equation*}
            Para $Var(X) = k_0p(1-p)$, si aplicamos el Teorema de Zehna tenemos que:
            \begin{equation*}
                \widehat{Var(X)} = \widehat{k_0p(1-p)} = k_0\hat{p}(1-\hat{p}) = \frac{1}{n}\sum_{i=1}^{n}X_i\left(1-\frac{1}{nk_0}\sum_{i=1}^{n}X_i\right) = \overline{X}\left(1-\frac{1}{k_0}\overline{X}\right)
            \end{equation*}
        \item [Por el método de los momentos.] 
            Por el método de los momentos:
            \begin{equation*}
                \left\{\begin{array}{l}
                        k_0p = E[X] = \overline{X} = \frac{1}{n}\sum\limits_{i=1}^{n}X_i \\
                        k_0p(1-p) = Var(X) = \frac{1}{n}\left(\sum\limits_{k=1}^{n}X_i^2 -\overline{X}^2\right)
                \end{array}\right.
            \end{equation*}
            Del primero de ducimos que:
            \begin{equation*}
                \hat{p} = \frac{1}{nk_0}\sum_{i=1}^{n}X_i = \frac{1}{k_0}\overline{X}
            \end{equation*}
            Y del segundo que:
            \begin{equation*}
                \widehat{Var(X)} = \frac{1}{n}\left(\sum\limits_{k=1}^{n}X_i^2 -\overline{X}^2\right)
            \end{equation*}
    \end{description}
    Para estimar a partir de los datos mencionados la probabilidad de salir un cuatro, lo que hacemos primero es considerar:
    \begin{equation*}
        Y \equiv \text{``Número de veces que sale un 4 en 10 tiradas del dado''} \rightsquigarrow B(10,p)
    \end{equation*}
    donde estimaremos el valor de $p$ mediante la estimación obtenida:
    \begin{equation*}
        p = \frac{1}{nk_0}\sum_{i=1}^{n}x_i= \frac{1}{nk_0}\sum_{i=1}^{n}x_i = \frac{1}{1000}(0\cdot 84 + 1\cdot 15 + 2\cdot 1 + 3\cdot 0) = \frac{15+2}{1000} = \frac{17}{1000}
    \end{equation*}
    Por lo que la probabilidad de obtener un 4 es $\frac{17}{1000} = 0.017$.
\end{ejercicio}

\begin{ejercicio}\label{ej:5_rel5}
    Se lanza un dado hasta que salga un $4$ y se anota el número de lanzamientos necesarios; este experimento se efectúa veinte veces de forma independiente. A partir de los resultados obtenidos, estimar la probabilidad de sacar un $4$ por máxima verosimilitud.\\

    \noindent
    Sea:
    \begin{equation*}
        X \equiv \text{``número de lanzamientos antes del lanzamiento en el que sale 4''}
    \end{equation*}
    tenemos que si $p$ es la probabilidad de que salga un 4, entonces $X$ sigue una distribución geométrica de parámetro $p$: $X\rightsquigarrow G(p)$. Como se realiza el experimento 20 veces de forma independientes, disponemos de una muestra $x_1, \ldots, x_n$ donde $n=20$. Estimamos $p$ por máxima verosimilitud:
    \begin{equation*}
        L_{x_1,\ldots,x_{n}}(p) = P_p[X_1 =x_1, \ldots, X_n=x_n] \stackrel{\text{indep.}}{=} \prod_{i=1}^{n} P_p[X=x_i] = \prod_{i=1}^{n} {(1-p)}^{x_i}p = p^{n}{(1-p)}^{\sum\limits_{i=1}^{n}x_i}
    \end{equation*}
    de donde al aplicar logaritmos:
    \begin{equation*}
        \ln L_{x_1,\ldots,x_n}(p) = n\cdot \ln p + \ln(1-p)\sum_{i=1}^{n}x_i
    \end{equation*}
    si derivamos:
    \begin{equation*}
        \dfrac{\partial \ln L_{x_1,\ldots,x_n}(p)}{\partial p} = \frac{n}{p} - \frac{\sum\limits_{i=1}^{n}x_i}{1-p} = \frac{n - np - p\left(\sum\limits_{i=1}^{n}x_i\right)}{p(1-p)} = \frac{n - p\left(\sum\limits_{i=1}^{n}x_i +n\right)}{p(1-p)}
    \end{equation*}
    Obtenemos que:
    \begin{equation*}
        \dfrac{\partial \ln L_{x_1,\ldots,x_n}(p)}{\partial  p} = 0 \Longleftrightarrow p = \frac{n}{\sum\limits_{i=1}^{n}x_i + n}
    \end{equation*}
    Como siempre tenemos que:
    \begin{equation*}
        \frac{n}{\sum\limits_{i=1}^{n}x_i + n} \in [0,1] 
    \end{equation*}
    Tenemos entonces que el EMV es:
    \begin{equation*}
        \hat{p} = \frac{n}{\sum\limits_{i=1}^{n}X_i+n} = \frac{20}{\sum\limits_{i=1}^{20}X_i+20}
    \end{equation*}

    o bien:
    \begin{equation*}
        \hat{p} = \frac{1}{1+\overline{X}}
    \end{equation*}
\end{ejercicio}

\begin{ejercicio}
    En $20$ días muy fríos, una granjera pudo arrancar su tractor en el primer, tercer, quinto, primer, segundo, primer, tercer, séptimo, segundo, cuarto, cuarto, octavo, primer, tercer, sexto, quinto, segundo, primer, sexto y segundo intento. Suponiendo que la probabilidad de arrancar en cada intento es constante, y que las observaciones se han obtenido de forma independiente, dar la estimación más verosímil de la probabilidad de que el tractor arranque en el segundo intento.\\

    \noindent
    Si reunimos en una tabla la información del enunciado:
    \begin{table}[H]
    \centering
    \begin{tabular}{cccccccc}
        Primer & Segundo & Tercero & Cuarto & Quinto & Sexto & Séptimo & Octavo \\
        \hline
        5 & 4 & 3 & 2 & 2 & 2 & 1 & 1
    \end{tabular}
    \caption{Veces que arrancó en cada intento.}
    \end{table}

    Si consideramos:
    \begin{equation*}
        X \equiv \text{``Número de intentos fallidos de arrancar el tractor''}
    \end{equation*}
    tenemos que $X\rightsquigarrow G(p)$, donde $p$ es la probabilidad de arrancar el tractor en un intento. tenemos a nuestra disposición de una muestra de $X$ $x_1,\ldots, x_n$ de tamaño $n=20$. Si observamos la similitud con el Ejercicio~\ref{ej:5_rel5}, estamos bajo las mismas hipótesis, por lo que al calcular el EMV para $p$ obtendremos:
    \begin{equation*}
        \hat{p} = \frac{n}{\sum\limits_{i=1}^{n}X_i+n} = \frac{20}{\sum\limits_{i=1}^{20}X_i+20}
    \end{equation*}
    Sin embargo, queremos calcularlo para la probabilidad de que el tractor arranque en el segundo intento, es decir, para tener exactamente un intento fallido de arranzar el tractor:
    \begin{equation*}
        P_p[X = 1] = (1-p)p
    \end{equation*}
    Si amplicamos el Teorema de Zehna, podemos obtener el EMV de $P_p[X=1]$ a partir de $\hat{p}$:
    \begin{equation*}
        \widehat{P_p[X=1]} = \widehat{(1-p)p} = (1-\hat{p})\hat{p} = \left(1-\frac{20}{\sum\limits_{i=1}^{20}X_i+20}\right) \frac{20}{\sum\limits_{i=1}^{20}X_i+20}
    \end{equation*}
    Si sustituimos ahora en los valores de la muestra, observamos que:
    \begin{equation*}
        \sum_{i=1}^{n}x_i = 0\cdot 5 + 1\cdot 4 + 2\cdot 3 + 3\cdot 2 + 4\cdot 2 + 5\cdot 2 + 6\cdot 1 + 7\cdot 1 = 47
    \end{equation*}
    Por lo que la estimación máximo verosímil de la probabilida de que el tractor arranque en el segundo intento es igual a:
    \begin{equation*}
        \left(1-\frac{20}{\sum\limits_{i=1}^{20}X_i+20}\right) \frac{20}{\sum\limits_{i=1}^{20}X_i+20} = \left(1-\frac{20}{47+20}\right) \frac{20}{47+20} = \frac{940}{4489} \approx 0.2094
    \end{equation*}
\end{ejercicio}

\begin{ejercicio} % Se podría hacer también con una multinomial
    Una variable aleatoria discreta toma los valores $0$, $1$ y $2$ con las siguientes probabilidades
    \begin{equation*}
        P_p[X=0] = p^2, \qquad P_p[X=1] = 2p(1-p), \qquad P_p[X=2] = {(1-p)}^{2}
    \end{equation*}
    siendo $p$ un parámetro desconocido. En una muestra aleatoria simple de tamaño $100$, se ha presentado $22$ veces el $0$, $53$ veces el $1$ y $25$ veces el $2$. Calcular la función de verosimilitud asociada a dicha muestra y dar la estimación más verosímil de $p$.\\

    \noindent
    Los datos obtenidos son:
    \begin{table}[H]
    \centering
    \begin{tabular}{c|ccc}
        Dato & 0 & 1 & 2 \\
        \hline
        Veces & 22 & 53 & 25
    \end{tabular}
    \end{table}
    Sea $X$ una variable aleatoria cuya función masa de probabilidad nos viene dada, tenemos una muestra $x_1, \ldots, x_n$ de tamaño $n=100$. La función de verosimilitud asociada a la muestra es:
    \begin{align*}
        L_{x_1,\ldots,x_n}(p) &= P_p[X_1=x_1, \ldots, X_n = x_n] \stackrel{\text{iid.}}{=}\prod_{i=1}^{n}P_p[X=x_i]\\ 
                              &= \prod_{i=1}^{22}P_p[X=0] \prod_{i=1}^{53}P_p[X=1]\prod_{i=1}^{25}P_p[X=2] \\
                              &= \prod_{i=1}^{22}p^2 \prod_{i=1}^{53}2p(1-p)\prod_{i=1}^{25}{(1-p)}^{2} \\
                              &= p^{44} \cdot 2^{53}p^{53}{(1-p)}^{53} \cdot {(1-p)}^{50} = 2^{53}p^{97}{(1-p)}^{103}
    \end{align*}
    Si aplicamos logaritos:
    \begin{equation*}
        \ln L_{x_1,\ldots,x_n}(p) = 53\ln 2 + 97 \ln p + 103 \ln(1-p)
    \end{equation*}
    y derivamos:
    \begin{equation*}
        \dfrac{\partial \ln L_{x_1,\ldots,x_n}(p)}{\partial p} = \frac{97}{p} - \frac{103}{1-p} = \frac{97(1-p)-103p}{p(1-p)} = \frac{97 - 200p}{p(1-p)} = 0 \Longleftrightarrow p = \frac{97}{200}
    \end{equation*}
    Por lo que la estimación por máxima verosimilitud es $\frac{97}{200} = 0.485$.\\

    \noindent
    Este ejercicio se podría haber resuelto también usando la multinomial.
\end{ejercicio}

\begin{ejercicio}
    En el muestreo de una variable aleatoria con distribución $\cc{N}(\mu,1)$, $\mu \in \mathbb{R}$, se observa que no se obtiene un valor menor que $-1$ hasta la quinta observación. Dar una estimación máximo verosímil de $\mu$.\\

    \noindent
    Sea $Y\rightsquigarrow\cc{N}(\mu, 1)$, si tomamos:
    \begin{equation*}
        X\equiv \text{``número de observaciones hasta obtener una menor que -1 (incluida)''}
    \end{equation*}

    tenemos que:
    \begin{equation*}
        P_p[X=x] = {(1-p)}^{x-1}p
    \end{equation*}
    donde $p= P[Y<-1]$. Tenemos una muestra de tamaño 1 ($x=5$) de $X$, con lo que la función de verosimilitud es:
    \begin{equation*}
        L_5(p) = P_p[X=5] = {(1-p)}^{5-1}p = p{(1-p)}^{4}
    \end{equation*}
    aplicando logaritmos:
    \begin{equation*}
        \ln L_5(p) = 4\ln(1-p) + \ln p
    \end{equation*}
    de donde:
    \begin{equation*}
        \dfrac{\partial \ln L_5(p)}{\partial p} = \frac{-4}{1-p}+\frac{1}{p} = \frac{-4p+1-p}{p(1-p)} = \frac{1-5p}{p(1-p)} = 0 \Longleftrightarrow p = \frac{1}{5}
    \end{equation*}
    Si notamos por $\Phi$ la función de distribución de una $\cc{N}(0,1)$, tenemos que:
    \begin{equation*}
        p = P[Y< -1] = P[Z<-1-\mu] = \Phi(-1-\mu) = 1-\Phi(1+\mu)
    \end{equation*}

    de donde (usando que $\Phi$ es biyectiva):
    \begin{equation*}
        \mu = \Phi^{-1}\left(1-p\right) - 1
    \end{equation*}
    por lo que si aplicamos el Teorema de Invarianza de Zehna:
    \begin{equation*}
        \hat{\mu} = \Phi^{-1}(1-\hat{p}) - 1 = \Phi^{-1}\left(1-\frac{1}{5}\right) - 1 = \Phi^{-1}(0.8) - 1
    \end{equation*}
    Y si miramos en la tabla de la normal $\cc{N}(0,1)$ la abscisa en la que se da la probabilidad $0.8$, obtenemos que se alcanza en la abscisa $0.85$, por lo que:
    \begin{equation*}
        \hat{\mu} = \Phi^{-1}(0.8) - 1 = 0.85 - 1 = -0.15
    \end{equation*}
    es la estimación máximo verosímil de $\mu$.
\end{ejercicio}

\begin{ejercicio}
    En la producción de filamentos eléctricos la medida de interés, $X$, es el tiempo de vida de cada filamento, que tiene una distribución exponencial de parámetro $\theta$. Se eligen $n$ de tales filamentos de forma aleatoria e independiente, pero, por razones de economía, no conviene esperar a que todos se quemen y la observación acaba en el tiempo $T$. Dar el estimador máximo verosímil para la media de $X$ a partir del número de filamentos quemados durante el tiempo de observación.\\

    \noindent
    Sea $X\rightsquigarrow exp(\theta)$, si consideramos: 
    \begin{align*}
        Y\equiv &\text{``número de filamentos con tiempo de vida menor que} \\ &\qquad \qquad \text{$T$ en una muestra de tamaño $n$''} 
    \end{align*}
    tenemos que $Y\rightsquigarrow B(n,p)$ donde $p$ es la probabilidad de que el tiempo de vida de un filamento sea menor que $T$, $p = P[X<T]$. Disponemos de una muestra de tamaño $1$ de $Y$, $y$. Tratamos de calcular el EMV de $p$:
    \begin{equation*}
        L_y(p) = P_p[Y=y] = \binom{n}{y}p^y{(1-p)}^{n-y}
    \end{equation*}
    si aplicamos logaritmo:
    \begin{equation*}
        \ln L_y(p) = \ln\binom{n}{y} + y\ln p + (n-y)\ln (1-p)
    \end{equation*}
    y ahora derivamos:
    \begin{equation*}
        \dfrac{\partial \ln L_y(p)}{\partial p} = \frac{y}{p} - \frac{n-y}{1-p} = \frac{y - yp - np + yp}{p(1-p)} = \frac{y-np}{p(1-p)} = 0 \Longleftrightarrow p = \frac{y}{n}
    \end{equation*}
    tenemos que el EMV de $p$ es:
    \begin{equation*}
        \hat{p} = \frac{Y}{n}
    \end{equation*}
    Si observamos la definición de $p$:
    \begin{equation*}
        p = P[X<T] = F_X(T) = 1 - e^{-\theta T} \Longrightarrow e^{-\theta T} = 1-p \Longrightarrow -\theta T = \ln (1-p)
    \end{equation*}

    de donde:
    \begin{equation*}
        \theta = \frac{-\ln(1-p)}{T}
    \end{equation*}
    Sin embargo, lo que nos interesa es calcuar el EMV de $E[X] = \frac{1}{\theta}$, por lo que:
    \begin{equation*}
        E[X] = \frac{1}{\theta} = \frac{-T}{\ln(1-p)}
    \end{equation*}
    Si aplicamos ahora el Teorema de Zehna obtenemos finalmente el EMV de $E[X]$:
    \begin{equation*}
        \widehat{E[X]} = \widehat{\frac{1}{\theta}} = \widehat{\frac{-T}{\ln(1-p)}} = \frac{-T}{\ln(1-\hat{p})} = \frac{-T}{\ln\left(1-\frac{Y}{n}\right)}
    \end{equation*}
\end{ejercicio}

\begin{ejercicio}
    Sean $X_1, \ldots, X_n$ observaciones independientes de una variable $X\rightsquigarrow\{\Gamma(p,a) : p,a>0\}$. Estimar ambos parámetros mediante el método de los momentos.\\

    \noindent
    \textit{Aplicación}: Ciertos neumáticos radiales tuvieron vidas útiles de $35200$, $41000$, $44700$, $38600$ y $41500$ kilómetros. Suponiendo que estos datos son observaciones independientes de una variable con distribución exponencial de parámetro $\theta$, dar una estimación de dicho parámetro por el método de los momentos.\\

    \noindent
    Planteamos el sistema a resolver, sabiendo que si $X\rightsquigarrow \Gamma(p,a)$, entonces:
    \begin{equation*}
        E[X] = \frac{p}{a}, \qquad Var(X) = \frac{p}{a^2}
    \end{equation*}

    por lo que:
    \begin{equation*}
        \left\{\begin{array}{c}
                \nicefrac{p}{a} = E[X] = \overline{X} \\
                \nicefrac{p}{a^2} = Var(X) = A_2 - {\left(\overline{X}\right)}^2
        \end{array}\right.
    \end{equation*}
    de la primera ecuación tenemos que:
    \begin{equation*}
        p = a\overline{X}
    \end{equation*}
    y si sustituimos en la segunda:
    \begin{equation*}
        \frac{\overline{X}}{a} = A_2 - {\left(\overline{X}\right)}^2 \Longrightarrow a = \frac{\overline{X}}{A_2-{\left(\overline{X}\right)}^2}
    \end{equation*}
    Sustituyendo ahora $a$ en la primera:
    \begin{equation*}
        p = \frac{{\left(\overline{X}\right)}^{2}}{A_2-{\left(\overline{X}\right)}^{2}}
    \end{equation*}
    Por lo que las estimaciones a considerar son:
    \begin{equation*}
        \hat{a} = \frac{\overline{X}}{A_2-{\left(\overline{X}\right)}^2}, \qquad \hat{p} = \frac{{\left(\overline{X}\right)}^{2}}{A_2-{\left(\overline{X}\right)}^{2}}
    \end{equation*}~\\

    \noindent
    Sea $X\rightsquigarrow exp(\theta)$, por el método de los momentos tenemos que:
    \begin{equation*}
        \frac{1}{\theta} = E[X] = \overline{X} \Longrightarrow \theta = \frac{1}{\overline{X}}
    \end{equation*}
    Por tanto, para estimar $\theta$ lo primero que hacemos es calcular la media de los valores:
    \begin{equation*}
        \frac{35200 + 41000 + 44700 + 38600 + 41500}{5} = 40200
    \end{equation*}

    y la estimación será:
    \begin{equation*}
        \hat{\theta} = \frac{1}{\overline{x}} = \frac{1}{40200} \approx 2.488\cdot 10^{-5}
    \end{equation*}
\end{ejercicio}

    % \section{Estimación por intervalos de confianza}

\begin{ejercicio}
    Sea $\overline{X}$ la media de una muestra aleatoria de tamaño $n$ de una población $\cc{N}(\mu, 16)$ Encontrar el menor valor de $n$ para que $\left]\overline{X}-1,\overline{X}+1\right[$ sea un intervalo de confianza para $\mu$ al nivel de confianza $0.9$.\\

    \noindent
    Si $\overline{X}$ es la media de una m.a.s. de tamaño $n$ de una población $\cc{N}(\mu,16)$, tenemos entonces que $\overline{X}\rightsquigarrow\cc{N}(\mu, \frac{16}{n})$. Calculamos $n$ para que el intervalo $\left]\overline{X}-1,\overline{X}+1\right[$ sea un intervalo de confianza para $\mu$ a nivel $0.9$, buscando:
    \begin{align*}
        P_\mu\left[\overline{X}-1 \leq \mu \leq \overline{X}+1\right] \geq 0.9 &\Longleftrightarrow P_\mu[\mu-1\leq \overline{X}\leq \mu+1] \geq 0.9 \\
                                                                    &\stackrel{\text{tipif.}}{\Longleftrightarrow} P_\mu\left[\frac{-\sqrt{n}}{4} \leq Z \leq \frac{\sqrt{n}}{4}\right] \geq 0.9 \\
                                                                    &\Longleftrightarrow 2P_\mu\left[Z\leq \frac{\sqrt{n}}{4}\right] - 1 \geq 0.9 \\
                                                                    &\Longleftrightarrow 2P_\mu\left[Z\leq \frac{\sqrt{n}}{4}\right] \geq 1.9 \\
                                                                    &\Longleftrightarrow P_\mu\left[Z\leq \frac{\sqrt{n}}{4}\right] \geq 0.95 
    \end{align*}
    Consultando la tabla de la normal $\cc{N}(0,1)$, tenemos que la primera abscisa en la que se alcanza una probabilidad superior a $0.95$ es $1.65$, por lo que:
    \begin{equation*}
        \frac{\sqrt{n}}{4} = 1.65 \Longleftrightarrow \sqrt{n} = 6.6 \Longleftrightarrow n = 43.56
    \end{equation*}
    Por tanto, el menor valor de $n$ para el cual el intervalo $\left]\overline{X}-1,\overline{X}+1\right[$ es un intervalo de confianza para $\mu$ a nivel de confianza $0.9$ es $44$.
\end{ejercicio}

\begin{ejercicio}
    La altura en cm. de los individuos varones de una población sigue una distribución $\cc{N}(\mu,\ 56.25)$. Si en una muestra aleatoria simple de tamaño $12$ de dicha población se obtiene una altura media de $175$ cm., determinar un intervalo de confianza para $\mu$ al nivel de confianza $0.95$. ¿Qué tamaño de muestra es necesario para que el intervalo de confianza a dicho nivel tenga longitud menor que 1 cm?\\

    \noindent
    Buscamos un intervalo de confianza para $\mu$ en una población normal con varianza $\sigma_0$ conocida a nivel de confianza $1-\alpha = 0.95$ (por lo que $\alpha= 0.05$). Tenemos pues una muestra aleatoria simple $(X_1, \ldots, X_n)$ de $X\rightsquigarrow\cc{N}(\mu,56.25)$. Usaremos el método de la cantidad pivotal, usando para ello la función pivote (donde $n=12$):
    \begin{equation*}
        T(X_1, \ldots, X_n; \mu) = \frac{\overline{X}-\mu}{\nicefrac{\sigma_0}{\sqrt{n}}} \rightsquigarrow\cc{N}(0,1)
    \end{equation*}
    Que:
    \begin{itemize}
        \item Es estrictamente decreciente respecto $\mu$.
        \item Si tomamos $\lm$ de forma que:
            \begin{equation*}
                \lm = \frac{\overline{X}-\mu}{\nicefrac{\sigma_0}{\sqrt{n}}} \Longrightarrow \mu = \overline{X} - \lm \frac{\sigma_0}{\sqrt{n}}
            \end{equation*}
    \end{itemize}
    Por lo que el intervalo de confianza que consideraremos viene dado por:
    \begin{equation*}
        \left]\overline{X}-\lm_2 \frac{\sigma_0}{\sqrt{n}}, \overline{X}-\lm_1 \frac{\sigma_0}{\sqrt{n}}\right[
    \end{equation*}
    donde $\lm_1$ y $\lm_2$ están sujetos a la restricción $P_\mu[\lm_1<T<\lm_2]=1-\alpha$. Obtenemos un intervalo de longitud esperada:
    \begin{equation*}
        E[L] = E\left[(\lm_2-\lm_1)\frac{\sigma_0}{\sqrt{n}}\right] = (\lm_2-\lm_1)\frac{\sigma_0}{\sqrt{n}}
    \end{equation*}
    Como $\nicefrac{\sigma_0}{\sqrt{n}}$ es una constante positiva, bastará minimizar la cantidad $\lm_2-\lm_1$ (sujeta a la restricción $P_\mu[\lm_1<T<\lm_2]=1-\alpha$) para minimizar la longitud del intervalo. La restricción mencionada puede reescribirse como: 
    \begin{equation*}
        1-\alpha = P_\mu[\lm_1<T<\lm_2] = F_Z(\lm_2) - F_Z(\lm_1)
    \end{equation*}
    Consideramos por el método de los multiplicadores de Lagrange:
    \begin{equation*}
        F(\lm_1,\lm_2) = (\lm_2-\lm_1) - \lm(F_Z(\lm_2) - F_Z(\lm_1) - (1-\alpha))
    \end{equation*}
    y buscamos los valores de $\lm_1$ y $\lm_2$ que minimicen la expresión, por lo que calculamos sus derivadas parciales:
    \begin{equation*}
        \dfrac{\partial F}{\partial \lm_1} = -1+\lm f_Z(\lm_1) \qquad \dfrac{\partial F}{\partial \lm_2} = 1-\lm f_Z(\lm_2)
    \end{equation*}
    Si igualamos ambas a cero y tratamos de despejar $\lm$:
    \begin{equation*}
        \left.\begin{array}{l}
            0 = \dfrac{\partial F}{\partial \lm_1} = -1+\lm f_Z(\lm_1) \\
            0 = \dfrac{\partial F}{\partial \lm_2} = 1-\lm f_Z(\lm_2)
        \end{array}\right\} \Longrightarrow \left.\begin{array}{l}
            \lm = \frac{1}{f_Z(\lm_1)} \\
            \lm = \frac{1}{f_Z(\lm_2)}
        \end{array}\right\} \Longrightarrow f_Z(\lm_1) = f_Z(\lm_2)
    \end{equation*}
    Como $f_Z$ es una función simétrica respecto al origen, las únicas posibilidades son bien $\lm_1 = \lm_2$ bien $\lm_1 = -\lm_2$. Como estos valores han de cumplir la restricción $P_\mu[\lm_1<T<\lm_2]=1-\alpha > 0$, la primera opción es imposible, por lo que tiene que ser $\lm_1 = -\lm_2$. Como han de dejar entre ellos un valor de $1-\alpha$, han de ser:
    \begin{equation*}
        \lm_2 = Z_{\nicefrac{\alpha}{2}}, \qquad \lm_1 = -\lm_2 = -Z_{\nicefrac{\alpha}{2}}
    \end{equation*}
    Por lo que el intervalo a considerar es:
    \begin{equation*}
        \left]\overline{X}-Z_{\nicefrac{\alpha}{2}}\frac{\sigma_0}{\sqrt{n}}, \overline{X}+Z_{\nicefrac{\alpha}{2}}\frac{\sigma_0}{\sqrt{n}}\right[
    \end{equation*}
    A partir de los datos proporcionados, tenemos que:
    \begin{itemize}
        \item $\alpha = 0.05$.
        \item $\sigma_0 = \sqrt{56.25} = 7.5$.
        \item $\sqrt{n} = \sqrt{12} = 3.4641$.
        \item $\overline{x} = 175$.
        \item $Z_{\nicefrac{\alpha}{2}}$ = $Z_{0.025} = 1.96$.
    \end{itemize}
    Por lo que el intervalo a considerar es:
    \begin{equation*}
        \left]175-1.96\cdot \frac{7.5}{3.4641},175+1.96\cdot \frac{7.5}{3.4641}\right[ = \left]170.7564755, 179.2435245\right[
    \end{equation*}
    Buscamos ahora el tamaño de la muestra para que la longitud del intervalo sea menor que $1$cm. Para ello, vemos que nuestro intervalo tiene longitud:
    \begin{equation*}
        2Z_{\nicefrac{\alpha}{2}}\frac{\sigma_0}{\sqrt{n}} = 2\cdot 1.96 \cdot \frac{7.5}{\sqrt{n}} = \frac{29.4}{\sqrt{n}}  \leq 1 \Longleftrightarrow 29.4 \leq \sqrt{n} \Longleftrightarrow 864.36 \leq n
    \end{equation*}
    Por lo que nicesitaremos al menos una nuestra de tamaño $865$ para que el intervalo de confianza a dicho nivel tenga longitud menor que 1 cm.
\end{ejercicio}

\begin{ejercicio}\label{ej:3_rel6}
    Una fábrica produce tornillos cuyo diámetro medio es $3$ mm. Se seleccionan aleatoria e independientemente 12 de estos tornillos y se miden sus diámetros, que resultan ser $3.01,\ 3.05,\ 2.99,\ 2.99,\ 3.00,\ 3.02,\ 2.98,\ 2.99,\ 2.97,\ 2.97,\ 3.02$ y $3.01$. Suponiendo que el diámetro es una variable aleatoria con distribución normal, determinar un intervalo de confianza para la varianza al nivel de confianza $0.99$, y una cota superior de confianza al mismo nivel. Interpretar los resultados en términos de la desviación típica del diámetro de los tornillos.\\

    \noindent
    Sea $(X_1, \ldots, X_n)$ una m.a.s. de $X\rightsquigarrow\cc{N}(3,\sigma^2)$, vamos a calcular un intervalo de confianza para $\sigma^2$ en una población normal con media $\mu_0$ conocida a nivel de confianza $1-\alpha$ mediante el método de la cantidad pivotal, usando la función pivote:
    \begin{equation*}
        T(X_1, \ldots, X_n;\sigma^2) = \frac{\sum\limits_{i=1}^n{(X_i-\mu_0)}^{2}}{\sigma^2} \rightsquigarrow\chi^2(n)
    \end{equation*}
    Que:
    \begin{itemize}
        \item Es estrictamente decreciente en $\sigma^2$.
        \item Si tenemos $\lm$ con:
            \begin{equation*}
                \lm = \frac{\sum\limits_{i=1}^n{(X_i-\mu_0)}^{2}}{\sigma^2} \rightsquigarrow\chi^2(n) \Longrightarrow \sigma^2 = \frac{\sum\limits_{i=1}^n{(X_i-\mu_0)}^{2}}{\lm}
            \end{equation*}
    \end{itemize}
    Por lo que el intervalo a considerar será:
    \begin{equation*}
        \left]\frac{1}{\lm_2}\sum\limits_{i=1}^n{(X_i-\mu_0)}^{2},\frac{1}{\lm_1}\sum\limits_{i=1}^n{(X_i-\mu_0)}^{2}\right[
    \end{equation*}
    donde $\lm_1$ y $\lm_2$ están sujetos a la restricción:
    \begin{equation*}
        1-\alpha \leq P_{\sigma^2}[\lm_1<T<\lm_2] = F_T(\lm_2) - F_T(\lm_1) \qquad (T\rightsquigarrow\chi^2(n))
    \end{equation*}
    Tenemos que la longitud esperada del intervalo es:
    \begin{equation*}
        E[L] = E\left[\left(\frac{1}{\lm_1}-\frac{1}{\lm_2}\right)\sum\limits_{i=1}^n{(X_i-\mu_0)}^{2}\right] = \left(\frac{1}{\lm_1}-\frac{1}{\lm_2}\right)E\left[\sum\limits_{i=1}^n{(X_i-\mu_0)}^{2}\right]
    \end{equation*}
    Donde la última esperanza es una cantidad constante y positiva, por lo que trataremos de minimizar simplemente $\nicefrac{1}{\lm_1}-\nicefrac{1}{\lm_2}$. Mediante el método de los multiplicadores de Lagrange:
    \begin{equation*}
        F(\lm_1,\lm_2) = \left(\frac{1}{\lm_1}-\frac{1}{\lm_2}\right) - \lm(F_T(\lm_2)-F_T(\lm_1)-(1-\alpha))
    \end{equation*}
    Calculando las derivadas parciales e igualando a cero llegamos a:
    \begin{equation*}
        \frac{f_T(\lm_1)}{f_T(\lm_2)} = \frac{\lm_2^2}{\lm_1^2}
    \end{equation*}
    Se trata de un despeje difícil, por lo que en la práctica se toman los valores:
    \begin{equation*}
        \lm_1 = \chi^2_{n,1-\nicefrac{\alpha}{2}}, \qquad \lm_2 = \chi^2_{n,\nicefrac{\alpha}{2}}
    \end{equation*}
    En definitiva, consideraremos el intervalo:
    \begin{equation*}
        \left]\frac{1}{\chi^2_{n,\nicefrac{\alpha}{2}}}\sum\limits_{i=1}^n{(X_i-\mu_0)}^{2},\frac{1}{\chi^2_{n,1-\nicefrac{\alpha}{2}}}\sum\limits_{i=1}^n{(X_i-\mu_0)}^{2} \right[
    \end{equation*}
    Para determinar una cota superior, lo que hacemos es tomar $\lm_2 = \infty$, por lo que:
    \begin{equation*}
        1-\alpha = F_T(\lm_2)-F_T(\lm_1) = 1-F_T(\lm_1) \Longrightarrow F_T(\lm_1) = \alpha
    \end{equation*}
    Luego podemos tomar $\lm_1 = \chi^2_{n,1-\alpha}$, obteniendo la cota superior:
    \begin{equation*}
        \frac{1}{\chi^2_{n,1-\alpha}}\sum\limits_{i=1}^n{(X_i-\mu_0)}^{2} 
    \end{equation*}
    Si situimos ahora en los parámetros que tenemos:
    \begin{itemize}
        \item $\mu_0 = 3$.
        \item $n = 12$.
        \item $\alpha = 0.01$.
        \item $\chi^2_{n,\nicefrac{\alpha}{2}} = \chi^2_{12,\ 0.005} = 28.2997$.
        \item $\chi^2_{n,1-\nicefrac{\alpha}{2}} = \chi^2_{12,\ 0.995} = 3.0738$.
        \item $\chi^2_{n,1-\alpha}= \chi^2_{12,0.99} = 3.5706$.
        \item Finalmente:
            \begin{equation*}
                \sum_{i=1}^{n}{(x_i - \mu_0)}^{2} = 0.0059
            \end{equation*}
    \end{itemize}
    El intervalo de confianza a nivel de confianza $0.99$ será:
    \begin{equation*}
        \left]\frac{1}{28.2997} 0.0059, \frac{1}{3.0738} 0.0059\right[ = \left]0.0002084,0.00191945\right[
    \end{equation*}
    Y la cota superior será:
    \begin{equation*}
        \frac{1}{3.5706} 0.0059 = 0.00165238
    \end{equation*}
    La desviación típica de la muestra dada es:
    \begin{equation*}
        \frac{1}{n}\sum_{i=1}^{n}{(x_i-\mu_0)}^{2} = \frac{1}{12}\cdot 0.0059 = 0.000491667
    \end{equation*}
    Que como veos pertence al intervalo y es menor que la cota dada.
\end{ejercicio}

\begin{ejercicio} % // TODO: ESTA MAL
    Las notas en cierta asignatura de 7 alumnos de una clase, elegidos de forma aleatoria e independiente son: $4.5$, $3, 6$, $7$, $1.5$, $5.2$ y $3.6$. Suponiendo que las notas tienen distribución normal, dar un intervalo de confianza para la varianza de las mismas al nivel de confianza $0.95$.\\

    \noindent
    Sea $(X_1, \ldots, X_n)$ una m.a.s. de $X\rightsquigarrow\cc{N}(\mu_0, \sigma^2)$. Tenemos que el intervalo de confianza para $\sigma^2$ a nivel de confianza $1-\alpha$ es:
    \begin{equation*}
        \left]\frac{(n-1)S^2}{\chi^2_{n-1,\nicefrac{\alpha}{2}}},\frac{(n-1)S^2}{\chi^2_{n-1,1-\nicefrac{\alpha}{2}}}\right[
    \end{equation*}
    Pero ahora tenemos los valores:
    \begin{itemize}
        \item $n = 7$.
        \item $\alpha = 0.05$.
        \item $\chi^2_{n,\nicefrac{\alpha}{2}} = \chi^2_{7,0.025} = 16.0128$.
        \item $\chi^2_{n,1-\nicefrac{\alpha}{2}} = \chi^2_{7,0.975} = 1.6899$.
    \end{itemize}
    Por lo que el intervalo será:
    \begin{equation*}
        \left]1.451936, 16.959525\right[
    \end{equation*}
\end{ejercicio}

\begin{ejercicio}
    Dos muestras independientes, cada una de tamaño 7, de poblaciones normales con igual varianza, producen medias $4.8$ y $5.4$ y cuasivarianzas muestrales $8.38$ y $7.62$, respectivamente. Encontrar un intervalo de confianza para la diferencia de medias al nivel de confianza $0.95$.\\

    \noindent
    Si tenemos $(X_1, \ldots, X_{n_1})$ m.a.s.  de $X\rightsquigarrow\cc{N}(\mu_1,\sigma^2_1)$ y $(Y_1, \ldots, Y_{n_2})$ m.a.s. de $Y\rightsquigarrow\cc{N}(\mu_2,\sigma_2^2)$. Si queremos calcular el intervalo de confianza de $\mu_1-\mu_2$ a nivel de confianza $1-\alpha$, hemos visto en teoría que se obtiene:
    \begin{equation*}
        \left]\overline{X}-\overline{Y}-t_{n_1+n_2-2,\nicefrac{\alpha}{2}}S_p\sqrt{\frac{1}{n_1}+\frac{1}{n_2}},\overline{X}-\overline{Y}+t_{n_1+n_2-2,\nicefrac{\alpha}{2}}S_p\sqrt{\frac{1}{n_1}+\frac{1}{n_2}}\right[
    \end{equation*}
    donde:
    \begin{equation*}
        S_p^2 = \frac{(n_1-1)S_1^2 + (n_2-1)S_2^2}{n_1+n_2-2}
    \end{equation*}
    Calculamos el intervalo, usando para ello que:
    \begin{itemize}
        \item $n_1 = 7 = n_2$.
        \item $\overline{x} = 4.8$, $\overline{y} = 5.4$.
        \item $s_1^2 = 8.38$, $s_2^2 = 7.62$.
        \item $\alpha = 0.05$.
    \end{itemize}
    Por lo que:
    \begin{equation*}
        S_p = \sqrt{\frac{6\cdot 8.38+6\cdot 7.62}{7 + 7 -2}} = 2.828427
    \end{equation*}
    De donde el intervalo será, usando que:
    \begin{equation*}
        t_{n_1+n_2-2,\nicefrac{\alpha}{2}} = t_{12,\ 0.025} = 2.1788
    \end{equation*}
    \begin{align*}
        &\left]4.8-5.4-2.1788\cdot 2.828427\cdot \sqrt{\frac{1}{7}+\frac{1}{7}}, 4.8-5.4+2.1788\cdot 2.828427\cdot \sqrt{\frac{1}{7}+\frac{1}{7}}\right[ \\
        &= \left]-3.894036, 2.694036\right[    
    \end{align*}
\end{ejercicio}

\begin{ejercicio}
    La siguiente tabla presenta los salarios anuales (en miles de euros) de dos grupos de recién graduados de dos carreras diferentes. Suponiendo normalidad en los salarios de ambos grupos, determinar un intervalo de confianza para el cociente de las varianzas al nivel de confianza $0.90$.
    \begin{table}[H]
    \centering
    \begin{tabular}{|c|cccccccccc|}
        \hline 
        GRUPO 1 & 16.3 & 18.2 & 17.5 & 16.1 & 15.9 & 15.4 & 15.8 & 17.3 & 14.9 & 15.1 \\
        \hline
        GRUPO 2 & 13.2 & 15.1 & 13.9 & 14.7 & 15.6 & 15.8 & 14.9 & 18.1 & 15.6 & 15.3 \\ 
                & 16.2 & 15.2 & 15.4 & 16.6 & & & & & &  \\
        \hline
    \end{tabular}
    \end{table}
    \noindent
    Si tenemos $(X_1, \ldots, X_{n_1})$ m.a.s.  de $X\rightsquigarrow\cc{N}(\mu_1,\sigma^2_1)$ (que se corresponderá con el grupo 1) y $(Y_1, \ldots, Y_{n_2})$ m.a.s. de $Y\rightsquigarrow\cc{N}(\mu_2,\sigma_2^2)$ (que se corresponderá con el grupo 2). Si queremos calcular el intervalo de confianza de $\frac{\sigma_1^2}{\sigma_2^2}$ a nivel de confianza $1-\alpha$, hemos visto en teoría que este será:
    \begin{equation*}
        \left]F_{n_1-1,n_2-1,1-\nicefrac{\alpha}{2}}\frac{S_1^2}{S_2^2},F_{n_1-1,n_2-1,\nicefrac{\alpha}{2}}\frac{S_1^2}{S_2^2} \right[
    \end{equation*}
    Con los datos:
    \begin{itemize}
        \item $\alpha = 0.1$.
        \item $n_1 = 10$.
        \item $n_2 = 14$.
        \item $S_1^2 = 1.187222$.
        \item $S_2^2 = 1.352308$.
    \end{itemize}
    Tenemos ($F_{9,13}$ no sale en la tabla, tomo los valores de $F_{9,12}$):
    \begin{itemize}
        \item $F_{n_1-1,n_2-1,1-\nicefrac{\alpha}{2}} = F_{9,13,0.95} = 0.325$.
        \item $F_{n_1-1,n_2-1,\nicefrac{\alpha}{2}} = F_{9,13,0.05} = 2.8$.
    \end{itemize}
    Por lo que el intervalo es:
    \begin{equation*}
        \left]0.325\cdot \frac{1.189222}{1.352308}, 2.8\cdot \frac{1.189222}{1.352308}\right[ = \left]0.285325, 2.458184\right[
    \end{equation*}
\end{ejercicio}

\begin{ejercicio}
    Con objeto de estudiar la efectividad de un agente diurético, se eligen al azar 11 pacientes, aplicando dicho fármaco a seis de ellos y un placebo a los cinco restantes. La variable observada en esta experiencia fue la concentración de sodio en la orina a las 24 horas, que se supone tiene una distribución normal en ambos casos. Los resultados observados fueron:
    \begin{itemize}
        \item Diurético: $20.4$, $62.5$, $61.3$, $44.2$, $11.1$, $23.7$.
        \item Placebo: $1.2$, $6.9$, $38.7$, $20.4$, $17.2$.
    \end{itemize}
    \begin{enumerate}[label=\alph*)]
        \item Calcular un intervalo de confianza para el cociente de las varianzas al nivel de confianza $0.95$.
        \item Suponiendo que las varianzas son iguales, calcular un intervalo de confianza para la diferencia de las medias al nivel de confianza $0.9$, y una cota inferior de confianza al mismo nivel. Interpretar los resultados.
    \end{enumerate} % // TODO: El ejercicio está hecho con alpha = 0.95, cambiarlo
    
    \noindent
    \textbf{Solución.} Tenemos:
    \begin{align*}
        X&\equiv \text{``Concentración de \texttt{Na} en la orina a las 24h tomando diurético''} \rightsquigarrow \cc{N}(\mu_1,\sigma_1^2) \\
        Y&\equiv \text{``Concentración de \texttt{Na} en la orina a las 24h tomando placebo''} \rightsquigarrow \cc{N}(\mu_2, \sigma_2^2)
    \end{align*}
    Y tenemos dos m.a.s. de ellas: $(X_1, \ldots, X_{n_1})$ de $X$ y $(Y_1, \ldots, Y_{n_2})$ de $Y$.
    \begin{enumerate}[label=\alph*)]
        \item Calcular un intervalo de confianza para el cociente de las varianzas al nivel de confianza $0.95$.

            Buscamos un intervalo de confianza para $\frac{\sigma_1^2}{\sigma_2^2}$ a nivel de confianza $1-\alpha = 0.95\Longrightarrow \alpha = 0.05$. La función pivote en este caso es:
            \begin{equation*}
                T\left(X_1, \ldots, X_{n_1},Y_1, \ldots, Y_{n_2};\frac{\sigma_1^2}{\sigma_2^2}\right) = \frac{S_2^2}{S_1^2} \frac{\sigma_1^2}{\sigma_2^2} \rightsquigarrow F(n_2-1,n_1-1)
            \end{equation*}
            Que:
            \begin{itemize}
                \item Sale creciente respecto el parámetro.
                \item Si tomamos $\lm$ con:
                    \begin{equation*}
                        \frac{S_2^2}{S_1^2} \frac{\sigma_1^2}{\sigma_2^2} \rightsquigarrow F(n_2-1,n_1-1) = \lm \Longrightarrow \frac{\sigma_1^2}{\sigma_2^2} = \frac{S_1^2}{S_1^2}\lm
                    \end{equation*}
            \end{itemize}
            Por lo que el intervalo general a considerar será:
            \begin{equation*}
                \left]\frac{S_1^2}{S_2^2}\lm_1, \frac{S_1^2}{S_2^2}\lm_2\right[
            \end{equation*}
            Por lo que vimos en teoría, tendremos:
            \begin{equation*}
                \lm_1 = F_{n_2-1,n_1-1,1-\nicefrac{\alpha}{2}}, \qquad \lm_2 = F_{n_2-1,n_1-1,\nicefrac{\alpha}{2}}
            \end{equation*}
            Por lo que el intervalo a considerar será:
            \begin{equation*}
                \left]\frac{S_1^2}{S_2^2}F_{n_2-1,n_1-1,1-\nicefrac{\alpha}{2}}, \frac{S_1^2}{S_2^2}F_{n_2-1,n_1-1,\nicefrac{\alpha}{2}}\right[
            \end{equation*}
            Calculamos las cuasivarianzas:
            \begin{equation*}
                S_1^2 = \frac{\sum\limits_{i=1}^n{(x_i-\overline{x})}^{2}}{n_1-1} = \frac{(n-1)}{n}Var(X_1, \ldots, X_n)
            \end{equation*}
            Usando la calculadora, obtenemos:
            \begin{equation*}
                S_1^2 = 483.12 \qquad S_2^2 = 208.517
            \end{equation*}
            Y consultando la tabla de la $F$ de Snedecor:
            \begin{equation*}
                F_{n_2-1,n_1-1,0.975}= 0.107, \qquad F_{n_2-1,n_1-1,0.025} = 7.39
            \end{equation*}
            Por lo que el intervalo obtenido es:
            \begin{equation*}
                \left]0.2479, 17.1219\right[
            \end{equation*}
            Si lo hubiéramos hecho para $\frac{\sigma_2^2}{\sigma_1^2}$ habríamos obtenido:
            \begin{equation*}
                \left]0.0583, 4.0398\right[
            \end{equation*}
        \item Suponiendo que las varianzas son iguales, calcular un intervalo de confianza para la diferencia de las medias al nivel de confianza $0.9$, y una cota inferior de confianza al mismo nivel. Interpretar los resultados.

            Como el valor $1$ para $\frac{\sigma_1^2}{\sigma_2^2}$ está en los intervalos, no podemos descartar que las varianzas sean iguales, con lo que podemos asumirlo. En teoría vimos que el intervalo a obtener es:
            \begin{equation*}
                \left]\overline{X}-\overline{Y}-t_{n_1+n_2-2,\nicefrac{\alpha}{2}}S_p\sqrt{\frac{1}{n_1}+\frac{1}{n_2}},\overline{X}-\overline{Y}+t_{n_1+n_2-2,\nicefrac{\alpha}{2}}S_p\sqrt{\frac{1}{n_1}+\frac{1}{n_2}}\right[
            \end{equation*}
            Y la cota inferior es:
            \begin{equation*}
                \overline{X}-\overline{Y}-t_{n_1+n_2-2,\alpha}S_p\sqrt{\frac{1}{n_1}+\frac{1}{n_2}}
            \end{equation*}
            donde:
            \begin{equation*}
                S_p^2 = \frac{(n_1-1)S_1^2 + (n_2-1)S_2^2}{n_1+n_2-2}
            \end{equation*}
            Como tenemos los datos:
            \begin{itemize}
                \item $n_1 = 6$.
                \item $n_2 = 5$.
                \item $\alpha = 0.05$.
                \item $S_1^2 = 483.12$.
                \item $S_2^2 = 208.517$.
                \item $\overline{x} = 37.2$.
                \item $\overline{y} = 16.88$.
            \end{itemize}
            Podemos también consultar:
            \begin{itemize}
                \item $t_{n_1+n_2-2,\nicefrac{\alpha}{2}} = t_{9,0.025} = 2.2622$.
                \item $t_{n_1+n_2-2,\alpha} = t_{9,0.05} = 1.8331$.
            \end{itemize}
            Por lo que:
            \begin{equation*}
                S_p = \sqrt{\frac{5\cdot 483.12+4\cdot 16.88}{6+5-2}} = 16.610305
            \end{equation*}
            y el intervalo y la cota inferior son:
            \begin{align*}
                &\left]37.2-16.88-2.2622\cdot 16.610305\cdot \sqrt{\frac{1}{6}+\frac{1}{5}},37.2-16.88+\right.\\&\qquad \qquad \qquad \qquad \qquad \qquad \qquad \qquad \qquad \qquad \left. +2.2622\cdot 16.610305\cdot \sqrt{\frac{1}{6}+\frac{1}{5}} \right[ \\
                &=\left]-2.433296, 43.073269\right[
            \end{align*}
            Vemos que $0$ está en dicho intervalo, por lo que no podemos descartar que las dos medias sean iguales.
    \end{enumerate}
\end{ejercicio}

\begin{ejercicio}
    Sea $(X_1, \ldots, X_n)$ una muestra aleatoria simple de una variable aleatoria con distribución $U(0,\theta)$. Dado un nivel de confianza arbitrario, calcular el intervalo de confianza para $\theta$ de menor longitud media uniformemente basado en un estadistico suficiente.\\

    \noindent
    Si la muestra es de $X\rightsquigarrow U(0,\theta)$ y consideramos un nivel de confianza $1-\alpha$ para $\alpha \in \left]0,1\right[$, calculamos un estadístico suficiente para $\theta$:
    \begin{equation*}
        f_\theta^n(x_1, \ldots, x_n) \stackrel{\text{iid.}}{=} \prod_{i=1}^{n}f_\theta(x_i) = \prod_{i=1}^{n}\frac{1}{\theta} = \frac{1}{\theta^n} \qquad x_i \in \left]0,\theta\right[
    \end{equation*}
    Por lo que para $0<X_{(1)}<X_{(n)}<\theta$ tenemos:
    \begin{equation*}
        f_\theta^n(x_1, \ldots, x_n) = I_{\mathbb{R}^+}(X_{(1)}) I_{\mathbb{R}^-}(X_{(n)}-\theta)\frac{1}{\theta^n}
    \end{equation*}
    Tomando:
    \begin{gather*}
        h(X_1, \ldots, X_n) = I_{\mathbb{R}^+}(X_{(1)}), \qquad T(X_1, \ldots, X_n) = X_{(n)} \\
        g_\theta(t) = I_{\mathbb{R}^-}(t-\theta)\frac{1}{\theta^n}
    \end{gather*}
    Tenemos por el Teorema de factorización de Neymann-Fisher que el estadístico $T(X_1, \ldots, X_n) = X_{(n)}$ es suficiente para $\theta$. Calculamos su función de distribución:
    \begin{equation*}
        F_T(t) = {(F_X(t))}^{n} = {\left(\frac{t}{\theta}\right)}^{n} \qquad t\in [0,\theta]
    \end{equation*}
    Por lo que una función pivote para aplicar el método de la cantidad pivotal es:
    \begin{equation*}
        T(X_1, \ldots, X_n;\theta) = F_T(T) = {\left(\frac{X_{(n)}}{\theta}\right)}^{n} \rightsquigarrow U(0,1)
    \end{equation*}
    Tenemos que:
    \begin{itemize}
        \item La función es estrictamente decreciente en $\theta$.
        \item Si tenemos $\lm$ de forma que:
            \begin{equation*}
                \lm = {\left(\frac{X_{(n)}}{\theta}\right)}^{n} \Longrightarrow \theta = \frac{X_{(n)}}{\sqrt[n]{\lm}}
            \end{equation*}
    \end{itemize}
    Por lo que el intervalo de confianza para $\theta$ a nivel de confianza $1-\alpha$ será:
    \begin{equation*}
        \left]\frac{X_{(n)}}{\sqrt[n]{\lm_2}},\frac{X_{(n)}}{\sqrt[n]{\lm_1}} \right[
    \end{equation*}
    donde $\lm_1$ y $\lm_2$ verifican:
    \begin{equation*} % // TODO: A PARTIR DE AQUI
        1-\alpha\leq P_\theta[\lm_1<T<\lm_2] = \lm_2 - \lm_1
    \end{equation*}
    La longitud esperada del intervalo es:
    \begin{equation*}
        E[L] = E\left[\left(\frac{1}{\sqrt[n]{\lm_1}} - \frac{1}{\sqrt[n]{\lm_2}}\right) X_{(n)}\right] = \left(\frac{1}{\sqrt[n]{\lm_1}} - \frac{1}{\sqrt[n]{\lm_2}}\right)E[X_{(n)}]
    \end{equation*}
    Con $E[X_{(n)}]\geq 0$, por lo que será suficiente con minimizar el primer término. Por el método de los multiplicadores de Lagrange:
    \begin{equation*}
        F(\lm_1,\lm_2) = \left(\frac{1}{\sqrt[n]{\lm_1}} - \frac{1}{\sqrt[n]{\lm_2}}\right) - \lm\left(\lm_2-\lm_1 - (1-\alpha)\right)
    \end{equation*}
    Si derivamos e igualamos a cero despejando $\lm$:
    \begin{align*}
        \dfrac{\partial F}{\partial \lm_1} = \frac{-1}{n\sqrt[n]{\lm_1^{n+1}}} + \lm = 0 \quad &\Longrightarrow \quad  \lm = \frac{1}{n\sqrt[n]{\lm_1^{n+1}}} = \frac{1}{n\lm_1\sqrt[n]{\lm_1}} \\
        \dfrac{\partial F}{\partial \lm_2} = \frac{1}{n\sqrt[n]{\lm_2^{n+1}}} -\lm= 0 \quad &\Longrightarrow \quad \lm = \frac{1}{n\sqrt[n]{\lm_2^{n+1}}} = \frac{1}{n\lm_2\sqrt[n]{\lm_2}}
    \end{align*} % // TODO: TERMINAR, pedir a Irina
\end{ejercicio}

\begin{ejercicio}
    Utilizando la desigualdad de Chebychev, dar un intervalo de confianza para $p$ a nivel de confianza arbitrario, basado en una muestra de tamaño arbitrario de una variable aleatoria con distribución $B(1,p)$.\\

    \noindent
    Sea $(X_1, \ldots, X_n)$ una m.a.s. de $X\rightsquigarrow B(1,p)$, buscamos por el método de Chebyshev un intervalo de confianza para $p$ a nivel de confianza $1-\alpha$. Para ello, es necesario tener un estimador insesgado de $p$. Si consideramos:
    \begin{equation*}
        T(X_1, \ldots, X_n) = \overline{X}
    \end{equation*}
    tenemos que $\overline{X}\in [0,1]$, así como que $\overline{X}$ es insesgado, puesto que:
    \begin{equation*}
        E[\overline{X}] = E\left[\frac{1}{n}\sum_{i=1}^{n}X_i\right] = \frac{1}{n}\sum_{i=1}^{n}E[X_i] = \frac{1}{n}\sum_{i=1}^{n}p = \frac{np}{n} = p
    \end{equation*}
    Usando la reproductividad de la binomial vemos que:
    \begin{equation*}
        Var(\overline{X}) = Var\left(\frac{1}{n}\sum_{i=1}^{n}X_i\right) = \frac{1}{n^2}Var\left(\sum_{i=1}^{n}X_i\right) = \frac{np(1-p)}{n^2} = \frac{p(1-p)}{n}
    \end{equation*}
    Y si tomamos $c$ una constante de forma que:
    \begin{equation*} % // TODO: tomar c = 1/n o c = 1/4n
        \frac{p(1-p)}{n} \leq c \qquad \forall p\in \left]0,1\right[
    \end{equation*}
    tenemos entonces que el intervalo:
    \begin{equation*}
        \left]\overline{X}-\sqrt{\frac{c}{\alpha}},\overline{X}+\sqrt{\frac{c}{\alpha}}\right[
    \end{equation*}
    Es un intervalo de confianza para $p$ a nivel de confianza $1-\alpha$. % // TODO: dar la constante c
\end{ejercicio}

\begin{ejercicio}
    Para una muestra de tamaño $n$ de una variable aleatoria con función de densidad
    \begin{equation*}
        f_\theta(x) = \frac{2x}{\theta^2}, \qquad 0<x<\theta
    \end{equation*}
    encontrar el intervalo de confianza para $\theta$ de menor longitud media uniformemente a nivel de confianza $1-\alpha$, basado en un estadístico suficiente.\\

    \noindent
    Buscamos un estadístico suficiente de $\theta$:
    \begin{equation*}
        f_\theta^n(X_1, \ldots, X_n) \stackrel{\text{iid.}}{=} \prod_{i=1}^{n}f_\theta(X_i) = \prod_{i=1}^{n}\frac{2X_i}{\theta^2} = \frac{2^n}{\theta^{2n}}\prod_{i=1}^{n}X_i \qquad 0<X_i<\theta
    \end{equation*}
    Por lo que tomando $0<X_{(1)}<X_{(n)}<\theta$ tenemos:
    \begin{equation*}
        f_\theta^n(X_1, \ldots, X_n) = \frac{2^n}{\theta^{2n}}I_{\mathbb{R}^+}(X_{(1)})I_{\mathbb{R} ^-}(X_{(n)}-\theta)\prod_{i=1}^n X_i
    \end{equation*}
    Si tomamos:
    \begin{gather*}
        h(X_1, \ldots, X_n) = I_{\mathbb{R}^+}(X_{(1)})\prod_{i=1}^n X_i, \qquad T(X_1, \ldots, X_n) = X_{(n)} \\
        g_\theta(t) = \frac{2^n}{\theta^{2n}}I_{\mathbb{R}^-}(t-\theta)
    \end{gather*}
    Tenemos por el Teorema de factorización de Neymann-Fisher que $X_{(n)}$ es un estadístico suficiente para $\theta$. Calculamos su función de distribución:
    \begin{equation*}
        F_T(t) = {(F_X(t))}^{n}
    \end{equation*}
    Y para ello vemos que  tenemos que calcular $F_X$:
    \begin{equation*}
        F_X(t) = \int_{0}^{t} \frac{2x}{\theta^2}~dx = \left[\frac{x^2}{\theta^2}\right]_0^t = \frac{t^2}{\theta^2}
    \end{equation*}
    Por lo que:
    \begin{equation*}
        F_T(t) = {(F_X(t))}^{n} = {\left(\frac{t}{\theta}\right)}^{2n}
    \end{equation*}
    La función pivote que consideramos para aplicar el método de la cantidad pivotal es:
    \begin{equation*}
        T(X_1, \ldots, X_n;\theta) = F_T(T) = {\left(\frac{X_{(n)}}{\theta}\right)}^{2n} \rightsquigarrow U(0,1)
    \end{equation*}
    Tenemos que:
    \begin{itemize}
        \item Es estrictamente decreciente respecto $\theta$.
        \item Si tenemos $\lm$ de forma que:
            \begin{equation*}
                \lm = {\left(\frac{X_{(n)}}{\theta}\right)}^{2n} \quad \Longrightarrow\quad  \theta = \frac{X_{(n)}}{\sqrt[2n]{\lm}}
            \end{equation*}
    \end{itemize}
    Tenemos por tanto que el intervalo es:
    \begin{equation*}
        \left]\frac{X_{(n)}}{\sqrt[2n]{\lm_2}}, \frac{X_{(n)}}{\sqrt[2n]{\lm_1}}\right[
    \end{equation*}
    donde $\lm_1$ y $\lm_2$ verifican:
    \begin{equation*}
        1-\alpha\leq P_\theta[\lm_1<T<\lm_2] = F_T(\lm_2) - F_T(\lm_1) = \lm_2 - \lm_1
    \end{equation*}
    La longitud esperada del intervalo es:
    \begin{equation*}
        E[L] = E\left[\left(\frac{1}{\sqrt[2n]{\lm_1}} - \frac{1}{\sqrt[2n]{\lm_2}}\right) X_{(n)}\right] = \left(\frac{1}{\sqrt[2n]{\lm_1}} - \frac{1}{\sqrt[2n]{\lm_2}}\right)E[X_{(n)}]
    \end{equation*}
    Con $E[X_{(n)}]\geq 0$, por lo que será suficiente con minimizar el primer término. Si despejamos $\lm_2$:
    \begin{equation*}
        \lm_2 = 1 - \alpha + \lm_1
    \end{equation*}
    Tenemos que la función a minimizar es:
    \begin{equation*}
        F(\lm_1) 0 \lm_1^\frac{-1}{2n} - {(1-\alpha+\lm_1)}^{\frac{-1}{2n}}
    \end{equation*}
    Calculamos su derivada e igualdando a cero:
    \begin{equation*}
        F'(\lm_1) = \frac{-1}{2n}\lm_1^{\frac{-1}{2n}-1} + \frac{1}{2n}{(1-\alpha+\lm_1)}^{\frac{-1}{2n}-1} = 0
    \end{equation*}
    Tenemos que $F'(\lm_1) = 0 \Longleftrightarrow \lm_1 = 1-\alpha+\lm_1 \Longleftrightarrow \alpha = 1$. Por lo que no podemos encontrar una solución, luego $F$ es estrictamente monótona:
    \begin{itemize}
        \item Supuesto que $F$ es estrictamente crecientes, tenemos entonces que $\lm_1 = 0$ y $\lm_2 = 1-\alpha$, por lo que el intervalo es:
            \begin{equation*}
                \left]\frac{X_{(n)}}{\sqrt[2n]{1-\alpha}},+\infty\right[
            \end{equation*}
        \item Supuesto que $F$ es estrictamente decreciente, tenemos que $\lm_2 = 1$ y $\lm_1 = \alpha$, por lo que el intervalo es:
            \begin{equation*}
                \left]X_{(n)}, \frac{X_{(n)}}{\sqrt[2n]{\alpha}}\right[
            \end{equation*}
    \end{itemize}
    Como en la segunda opción nos sale un intervalo acotado, este es el de menor longitud esperada, por lo que nos quedamos con:
    \begin{equation*}
        \left]X_{(n)}, \frac{X_{(n)}}{\sqrt[2n]{\alpha}}\right[
    \end{equation*}
\end{ejercicio}

\begin{ejercicio} 
    Para una muestra de tamaño $n$ de una variable aleatoria con función de densidad
    \begin{equation*}
        f_\theta(x) = \frac{\theta}{x^2}, \qquad x>\theta
    \end{equation*}
    encontrar el intervalo de confianza para $\theta$ de menor longitud media uniformemente a nivel de confianza $1 - \alpha$, basado en el estimador máximo verosímil de $\theta$.\\

    \noindent
    La función de verosimilitud es:
    \begin{equation*}
        L_{x_1,\ldots,x_n}(\theta) = \prod_{i=1}^{n} \frac{\theta}{x_i^2} \qquad \forall x_i > \theta > 0
    \end{equation*}
    Por lo que:
    \begin{equation*}
        L_{x_1,\ldots,x_n}(\theta) = \left\{\begin{array}{ll}
            \prod\limits_{i=1}^{n}\frac{\theta}{x_i^2} & \text{si\ } \theta < x_{(1)}  \\
             0 & \text{en otro caso} 
        \end{array}\right. 
    \end{equation*}
    Y vemos que $L_{x_1,\ldots,x_n}(\theta)$ es creciente, con lo que alcanza su máximo en $\theta = x_{(1)}$. En definitiva, el EMV es $\hat{\theta} = X_{(1)}$, por lo que buscamos un intervalo de confianza basado en $X_{(1)}$.\\

    \noindent
    La función de distribución del mínimo es:
    \begin{equation*}
        F_{X_{(1)}}(t) = 1-{(1-F_X(t))}^{n}
    \end{equation*}
    Por lo que calculamos:
    \begin{equation*}
        F_X(t) = \int_{\theta}^{t} \frac{\theta}{x^2}~dx  = \left[\frac{-\theta}{x}\right]_\theta^t = \frac{-\theta}{t} + 1 \qquad t>\theta
    \end{equation*}

    de donde:
    \begin{equation*}
        F_{X_{(1)}}(t) = 1-{(1-F_X(t))}^{n} = 1 - {\left(1-\left(1-\frac{\theta}{t}\right)\right)}^{n} = 1-{\left(\frac{\theta}{t}\right)}^{n} \qquad t> \theta
    \end{equation*}
    Tomamos por tanto como función pivote:
    \begin{equation*}
        T(X_1, \ldots, X_n;\theta) = 1-{\left(\frac{\theta}{X_{(1)}}\right)}^{n} \rightsquigarrow U(0,1)
    \end{equation*}
    Comprobamos las condiciones del método de la cantidad pivotal:
    \begin{itemize}
        \item Si derivamos:
            \begin{equation*}
                \dfrac{\partial T}{\partial \theta} = -n{\left(\frac{\theta}{X_{(1)}}\right)}^{n-1} \frac{1}{X_{(1)}} = \frac{-n\theta^{n-1}}{X_{(1)}^n} < 0 \qquad \forall \theta
            \end{equation*}
            Por lo que $T(X_1, \ldots, X_n;\theta)$ es estrictamente decreciente en función de $\theta$.
        \item Si tratamos de despejar el parámetro para cierto $\lm$:
            \begin{equation*}
                1-{\left(\frac{\theta}{X_{(1)}}\right)}^{n} = \lm \Longrightarrow \theta = \sqrt[n]{(1-\lm)X_{(1)}^n}  = X_{(1)}\sqrt[n]{1-\lm}
            \end{equation*}
    \end{itemize}
    Por lo que hemos obtenido como intervalo de confianza (donde los índices de $\lm$ son debidos a que es estrictamente decreciente):
    \begin{equation*}
        \left]X_{(1)}\sqrt[n]{1-\lm_2}, X_{(1)}\sqrt[n]{1-\lm_2}\right[
    \end{equation*}
    Tratamos ahora de minimizar la longitud del intervalo, que es:
    \begin{equation*}
        L = X_{(1)}\left(\sqrt[n]{1-\lm_1} - \sqrt[n]{1-\lm_2}\right)
    \end{equation*}

    de donde:
    \begin{equation*}
        E_\theta[L] = E_\theta[X_{(1)}]\left(\sqrt[n]{1-\lm_1} - \sqrt[n]{1-\lm_2}\right)
    \end{equation*}
    Y como $E_\theta(X_{(1)})$ es una constante positiva, podemos obviarla a la hora de minimizar, por lo que buscamos minimizar el segundo trozo, sujeto a la restricción:
    \begin{equation*}
        P_\theta[\lm_1 < T < \lm_2] = 1-\alpha
    \end{equation*}
    Y como tenemos:
    \begin{equation*}
        P[\lm_1 < T < \lm_2] = F_T(\lm_2) - F_T(\lm_1) = \lm_2 - \lm_1
    \end{equation*}
    Por lo que tenemos la restricción $\lm_2-\lm_1 = 1-\alpha$. Para meterla en la función a minimizar:
    \begin{itemize}
        \item Bien despejamos $\lm_1$ o $\lm_2$, sustituimos en la función a minimizar y obtenemos una función de una variable, que sabemos minimizar.
        \item Usamos los multiplicadores de Lagrange, sumamos a la función la restricción (igualdada a 0) multiplicada por un cierto parámetro, que llamaremos $\lm$, por lo que buscamos minimizar:
            \begin{equation*}
                {(1-\lm_1)}^{\nicefrac{1}{n}}- {(1-\lm_2)}^{\nicefrac{1}{n}} + \lm (\lm_2 - \lm_1 - 1+\alpha)
            \end{equation*}
    \end{itemize}
    Cuando las restricciones son tan sencillas el método de despejar y sustituir es más sencillo, por lo que optamos por dicho método:
    \begin{equation*}
        \lm_2 = 1-\alpha+\lm_1
    \end{equation*}
    Sustituyendo en la función a minimizar, buscamos la constante  $\lm_1$ que minimiza la expresión:
    \begin{equation*}
        f(\lm_1) = {(1-\lm_1)}^{\nicefrac{1}{n}}-{(1-(a-\alpha+\lm_1))}^{\nicefrac{1}{n}} = {(1-\lm_1)}^{\nicefrac{1}{n}} - {(\alpha-\lm_1)}^{\nicefrac{1}{n}}
    \end{equation*}
    Calculamos la derivada y la igualamos a $0$, para despejar $\lm_1$:
    \begin{equation*}
        f'(\lm_1) = \frac{1}{n}{(1-\lm_1)}^{\nicefrac{1}{n}-1}(-1) - \frac{1}{n}{(\alpha-\lm_1)}^{\nicefrac{1}{n}-1}(-1) = 0
    \end{equation*}
    Haciendo cálculos llegamos a que $f'(\lm_1)=0$ si y solo si $\alpha=1$, pero $\alpha\in \left]0,1\right[$, por lo que $f'(\lm_1)\neq 0$, con lo que $f$ es estrictamente monónotona. Sabemos también que $0<\lm_1<\lm_2<1$, buscamos razonar si $f$ es estrictamente creciente o estrictamente decreciente. Para ello, podemos estudiar el signo de la derivada o bien razonarlo de la siguiente forma:
    \begin{itemize}
        \item Si la función es estrictamente creciente, esta alcanza su mínimo en $\lm_1=0$, con lo que usando que $\lm_2-\lm_1 = 1-\alpha$ tenemos que $\lm_2 = 1-\alpha$. Sustituyendo, obtenemos el intervalo:
            \begin{equation*}
                \left]X_{(1)}\sqrt[n]{1-(1-\alpha)},X_{(1)}\sqrt[n]{1-0}\right[ = \left]X_{(1)}\sqrt[n]{\alpha},X_{(1)}\right[
            \end{equation*}
        \item Si la función es estrictamente decreciente, entonces esta alcana su mínimo cuando $\lm_2=1$, con lo que $\lm_1 = \alpha$, de donde obtenemos el intervalo:
            \begin{equation*}
                \left]X_{(1)}\sqrt[n]{0},X_{(1)}\sqrt[n]{1-\alpha}\right[ = \left]0,X_{(1)}\sqrt[n]{1-\alpha}\right[
            \end{equation*}
    \end{itemize}
    El segundo intervalo no está acotado por debajo (ya que el mínimo a considerar en el espacio paramétrico es $0$), pero en el primero sí que tenemos un valor acotado por debajo y otro por encima.  Preferimos elegir el primer intervalo, puesto que restringe más la longitud del intervalo (es decir, lo acota por ambos lados), aunque puede ocurrir que en función de los valores de $X_{(1)}$ y $\alpha$ el intervalo mínimo sea a veces el primero y otras veces el segundo.
\end{ejercicio}


    % \section{Contraste de hipótesis}

% // TODO: donde no diga nada de sigma = sigma, comprobar si son = y decir que la muestra no lo contradice
% // TODO: El 13 está mal, H0: mu1 <= mu2

\begin{ejercicio}
    Se toma una observación de una variable con distribución de Poisson para contrastar que la media vale 1 frente a que vale 2.
    \begin{enumerate}[label=\alph*)]
        \item Construir un test no aleatorizado con nivel de significación $0.05$ para el contraste planteado. Calcular las probabilidades de cometer error de tipo 1 y de tipo 2, el tamaño y la potencia del test frente a la hipótesis alternativa.
        \item ¿Cómo debe aleatorizarse el test para alcanzar el tamaño $0.05$? ¿Cuál es la potencia de este test?
    \end{enumerate}
    
    \textbf{Solución}
    \begin{enumerate}[label=\alph*)]
        \item Sea $X\rightsquigarrow \cc{P}(\lm)$, queremos resolver el contraste:
            \begin{equation*}
                \left\{\begin{array}{l}
                    H_0 : \lm = 1 \\
                    H_1 : \lm = 2
                \end{array}\right.
            \end{equation*}
            Trabajaremos por tanto con:
            \begin{equation*}
                \Theta_0 = \{1\}, \qquad \Theta_1 = \{2\}, \qquad \Theta = \{1,2\}
            \end{equation*}
            y escribiremos $\lm_0 = 1$, $\lm_1 = 2$. Pensamos la forma que ha de tener el test: el espacio muestral de $X$ es $\mathbb{R}^+_0$:
            \begin{figure}[H]
                \centering
                \begin{tikzpicture}
                % Dibuja la recta real
                \draw[-Stealth] (-1,0) -- (3,0) node[right] {};

                % Marca los puntos 1 y 2
                \foreach \x in {0,1,2} {
                    \draw (\x,0.1) -- (\x,-0.1) node[below] {$\x$};
                }
                \end{tikzpicture}
            \end{figure}
            Es claro que cuanto más a la izquierda nos encontremos con dicha observación $X$, no rechazaremos la hipótesis nula, así como que cuanto más a la derecha estemos la rechazaremos. Por tanto, establecemos un cierto punto $c\in \mathbb{R}^+_0$ que delimite la región crítica del test (que será la región $\left]c,+\infty\right[$). Para determinar dicho punto $c$ imponemos un nivel de significación $\alpha = 0.05$. Por tanto, el test será de la forma:
            \begin{equation*}
                \varphi(X) = \left\{\begin{array}{ll}
                    1 & \text{si\ } X> c \\
                    0 & \text{si\ } X \leq c
                \end{array}\right. 
            \end{equation*}
            Imponemos nivel de significancia $\alpha$:
            \begin{equation*}
                \alpha \geq \sup_{\lm \in \Theta_0}\beta_\varphi(\lm) = \sup_{\lm\in \{\lm_0\}} \beta_\varphi(\lm) = \beta_\varphi(\lm_0) = E_{\lm_0}[\varphi]  = P_{\lm_0}[X>c]  = 1-P_{\lm_0}[X\leq c]
            \end{equation*}
            Y lo que hacemos ahora es ir probando con distintos valores de $c$ hasta acercarnos lo máximo posible a $\alpha = 0.05$ pero sin pasarnos:
            \begin{itemize}
                \item Para $c=0$:
                    \begin{equation*}
                        P_1[X\leq 0] = P_1[X = 0] = 0.3679 \quad\Longrightarrow\quad 1-P_1[X\leq 0] = 0.6321 > \alpha
                    \end{equation*}
                \item Para $c=1$:
                    \begin{align*}
                        P_1[X\leq 1] &= P_1[X=0] + P_1[X=1] = 0.3679 + 0.3679 = 0.7358  \\
                                     &\Longrightarrow 1-P_1[X\leq 1] = 0.2642 > \alpha
                    \end{align*}
                \item Para $c=2$:
                    \begin{align*}
                        P_1[X\leq 2] &= P_1[X\leq 1] + P_1[X=2] = 0.7358 + 0.1839 = 0.9197 \\
                                     &\Longrightarrow 1-P_1[X\leq 2] = 0.0803 > \alpha
                    \end{align*}
                \item Para $c=3$:
                    \begin{align*}
                        P_1[X\leq 3] &= P_1[X\leq 2] + P_1[X=3] = 0.9197 + 0.0613 = 0.981 \\
                                     &\Longrightarrow 1-P_1[X\leq 3] = 0.019 \leq \alpha
                    \end{align*}
            \end{itemize}
            Por tanto, el valor a tomar es $c=3$, de donde el test sería:
            \begin{equation*}
                \varphi(X) = \left\{\begin{array}{ll}
                    1 & \text{si\ } X> 3 \\
                    0 & \text{si\ } X \leq 3
                \end{array}\right. 
            \end{equation*}
            Calculamos ahora:
            \begin{description}
                \item [Probabilidad de cometer error tipo 1.] 
                    \begin{equation*}
                        P_{\lm\in \Theta_0}[X\in C] = P_{\lm_0}[X>3] = P_1[X>3] = 0.019
                    \end{equation*}
                \item [Probabilidad de cometer error tipo 2.]
                    \begin{align*}
                        P_{\lm \in \Theta_1}\left[X\in \overline{C}\right] &= P_{\lm_1}[X\leq 3] = P_2[X\leq 3] = \sum_{k=0}^{3} P_2[X=k] \\
                                                                           &= (0.1353 + 0.2707 + 0.2707 + 0.1804) = 0.8571
                    \end{align*}
                \item [Tamaño del test.] El tamaño resulta ser igual a la probabilidad encontrada cuando buscábamos el valor de $c$:
                    \begin{equation*}
                        \sup_{\lm \in \Theta_0}\beta_\varphi(\lm) = \beta_\varphi(\lm_0) = E_{\lm_0}[\varphi] = P_{\lm_0}[X>3] = 0.019
                    \end{equation*}
                \item [Potencia frente a hipótesis alternativa.] La potencia del test es:
                    \begin{align*}
                        \sup_{\lm \in \Theta_1} \beta_\varphi(\lm) &= \beta_\varphi(\lm_1) = E_{\lm_1}[\varphi] = P_{\lm_1}[X>3] = P_2[X>3] = 1-P_2[X\leq 3] \\
                                                                   &= 1- 0.8571 = 0.1429
                    \end{align*}
            \end{description}
        \item Para alcanzar un tamaño de $0.05$ mediante un test aleatoriezado, lo que hacemos ahora es considerar un test del tipo:
            \begin{equation*}
                \varphi(X) = \left\{\begin{array}{ll}
                    1 & \text{si\ } X>3 \\
                    \gamma & \text{si\ } X=3 \\
                    0 & \text{si\ } X<3 
                \end{array}\right. 
            \end{equation*}
            Para cierto $\gamma\in [0,1]$, de forma que a la hora de imponer un tamaño menor o igual que $\alpha=0.05$ calculemos $\gamma$ para obtener la igualdad. De esta forma, el tamaño del test es:
            \begin{equation*}
                \sup_{\lm \in \Theta_0}\beta_\varphi(\lm) = \beta_\varphi(\lm_0) = E_{\lm_0}[\varphi] = P_{\lm_0}[X>3] + \gamma P_{\lm_0}[X=3] = 0.019 + \gamma 0.0613
            \end{equation*}
            Si imponemos que sea igual a $\alpha$:
            \begin{equation*}
                0.05 = \alpha = 0.019 + \gamma \cdot 0.0613 \quad\Longrightarrow\quad \gamma = \frac{0.05 - 0.019}{0.0613} \approx 0.5057096
            \end{equation*}
            Calculamos la potencia del test:
            \begin{equation*}
                \sup_{\lm \in \Theta_1}\beta_\varphi(\lm) = E_{\lm_1}[\varphi] = P_{\lm_1}[X>3] + \gamma P_{\lm_1}[X=3] = 0.1429 + \gamma 0.1804 \approx 0.23413
            \end{equation*}
    \end{enumerate}
\end{ejercicio}

\begin{ejercicio}
    Una urna contiene 10 bolas, blancas y negras. Para contrastar que el número de bolas blancas es 5 frente a que dicho número es 6 o 7, se extraen tres bolas con reemplazamiento y se rechaza $H_0$ sólo si se obtienen 2 o 3 bolas blancas. Calcular el tamaño de este test y la potencia frente a las alternativas.\\

    \noindent 
    Sea:
    \begin{equation*}
        X\equiv \text{``Número de bolas blancas en una extracción''} \rightsquigarrow B(1,p)
    \end{equation*}
    Se quiere resolver el contraste de hipótesis:
    \begin{equation*}
        \left\{\begin{array}{l}
            H_0 : p = p_0 =  \nicefrac{5}{10} \\
            H_1 : p \in \{\nicefrac{6}{10}, \nicefrac{7}{10}\}
        \end{array}\right.
    \end{equation*}
    Por lo que consideraremos:
    \begin{equation*}
        \Theta_0 = \{\nicefrac{5}{10}\}, \qquad \Theta_1 = \{\nicefrac{6}{10}, \nicefrac{7}{10}\}, \qquad \Theta = \{\nicefrac{5}{10}, \nicefrac{6}{10}, \nicefrac{7}{10}\}
    \end{equation*}
    Y tenemos $(X_1, X_2, X_3)$ una m.a.s. de $X$. Planteamos el test:
    \begin{equation*}
        \varphi(X_1, X_2, X_3) = \left\{\begin{array}{ll}
            1 & \text{si\ } X_1+X_2+X_3 \in \{2,3\} \\
            0 & \text{en otro caso\ } 
        \end{array}\right. 
    \end{equation*}
    Por la reproductividad de la binomial tenemos que $X_1+X_2+X_3\rightsquigarrow B(3,p)$. Calculamos: 
    \begin{description}
        \item [Tamaño del test.] 
            \begin{align*}
                \sup_{p\in \Theta_0} \beta_\varphi(p) &= \beta_\varphi(p_0) = E_{p_0}[\varphi] = P_{p_0}[X_1+X_2+X_3 = 2] + P_{p_0}[X_1 + X_2 + X_3 = 3] \\
                                                      &= P_{\nicefrac{1}{2}}[X_1+X_2+X_3 = 2] + P_{\nicefrac{1}{2}}[X_1 + X_2 + X_3 = 3]  \\
                                                      &= 0.375 + 0.125 = 0.5
            \end{align*}
        \item [Potencia frente a hipótesis alternativas.] 
            Calculamos cada una de ellas:
            \begin{itemize}
                \item Para $p = \nicefrac{6}{10}$:
                    \begin{align*}
                        \beta_\varphi(p) &= E_p[\varphi] = P_p[X_1 + X_2 + X_3 = 2] + P_p[X_1+X_2+X_3 = 3] \\
                                         &= P_{0.6}[X_1 + X_2 + X_3 =2] + P_{0.6}[X_1+X_2+X_3 =3]
                    \end{align*}
                    Si consideramos $Y \rightsquigarrow B(3,1-p) = B(3,\ 0.4)$ tenemos:
                    \begin{align*}
                        \beta_\varphi(p) &= P_{0.6}[X_1 + X_2 + X_3 =2] + P_{0.6}[X_1+X_2+X_3 =3] \\
                                         &= P[Y=3-2] + P[Y=3-3] = P[Y=1] + P[Y=0] \\
                                         &= 0.432 + 0.216 = 0.648
                    \end{align*}
                \item para $p=\nicefrac{7}{10}$:
                    \begin{align*}
                        \beta_\varphi(p) &= P_p[X_1 + X_2 + X_3 = 2] + P_p[X_1+X_2+X_3 = 3] \\
                                         &= P_{0.7}[X_1 + X_2 + X_3 =2] + P_{0.7}[X_1+X_2+X_3 =3]
                    \end{align*}
                    Si consideramos $Y\rightsquigarrow B(3,1-p) = B(3,\ 0.3)$ tenemos:
                    \begin{align*}
                        \beta_\varphi(p) &= P_{0.7}[X_1 + X_2 + X_3 =2] + P_{0.7}[X_1+X_2+X_3 =3] \\
                                         &= P[Y=1] + P[Y=0] = 0.441 + 0.343 = 0.784
                    \end{align*}
            \end{itemize}
    \end{description}
\end{ejercicio}

\begin{ejercicio}
    Sea $(X_1, \ldots, X_n)$ una muestra aleatoria simple de una variable aleatoria con distribución de Poisson de parámetro $\lm$. Encontrar el test más potente de tamaño $\alpha$ para resolver el problema de contraste
    \begin{align*}
        H_0 &: \lm = \lm_0 \\
        H_1 &: \lm = \lm_1
    \end{align*}
    \textit{Aplicación:} En una centralita telefónica el número de llamadas por minuto sigue una distribución de Poisson. Si en cinco minutos se han recibido 12 llamadas, ¿puede aceptarse que el número medio de llamadas por minuto es $1.5$, frente a que dicho número es 2, al nivel de significación $0.05$? Calcular la potencia del test obtenido.\\

    \noindent
    Sea $X\rightsquigarrow \cc{P}(\lm)$, en este caso tenemos:
    \begin{equation*}
        \Theta_0 = \{\lm_0\}, \qquad \Theta_1 = \{\lm_1\}, \qquad \Theta = \{\lm_0,\lm_1\}
    \end{equation*}
    Como nos encontramos ante un contraste simple frente a simple y queremos hayar el test más potente de tamaño $\alpha$, recurrimos a un test de Neymann-Pearson, que será de la forma:
    \begin{equation*}
        \varphi(X_1, \ldots, X_n) = \left\{\begin{array}{ll}
            1 & \text{si\ } \lm(X_1, \ldots, X_n) > k \\
             \gamma & \text{si\ } \lm(X_1, \ldots, X_n) = k \\
             0 & \text{si\ } \lm(X_1, \ldots, X_n) < k
        \end{array}\right. 
    \end{equation*}
    Para ciertos valores $\gamma\in [0,1]$, $k\in \mathbb{R}$ y una cierta función $\lm$ que nos disponemos a calcular. Primero observemos que:
    \begin{equation*}
        P_{\lm}[X_1 = x_1, \ldots, X_n = x_n] \stackrel{\text{iid.}}{=} \prod_{i=1}^{n} P_{\lm}[X=x_i] = \prod_{i=1}^{n} e^{-\lm} \frac{\lm^{x_i}}{x_i!} = e^{-n\lm} \cdot \frac{\lm^{\sum\limits_{k=1}^{n}x_i}}{\prod\limits_{i=1}^{n}x_i!} > 0
    \end{equation*}
    Por lo que:
    \begin{align*}
        \lm(x_1,\ldots, x_n) &= \frac{P_{\lm_1}[X_1=x_1,\ldots, X_n=x_n]}{P_{\lm_0}[X_1=x_1,\ldots, X_n=x_n]} = \frac{e^{-n\lm_1} \cdot \frac{\lm_1^{\sum\limits_{i=1}^{n}x_i}}{\prod\limits_{i=1}^{n}x_i!} }{e^{-n\lm_0} \cdot \frac{\lm_0^{\sum\limits_{i=1}^{n}x_i}}{\prod\limits_{i=1}^{n}x_i!} } = e^{n(\lm_0-\lm_1)} {\left(\frac{\lm_1}{\lm_0}\right)}^{\sum\limits_{i=1}^{n}x_i}
    \end{align*}
    Y para simplificar el test buscamos uno equivalente, si consideramos el estadístico $T(X_1, \ldots, X_n) = \sum\limits_{i=1}^{n}X_i\rightsquigarrow \cc{P}(n\lm)$, estudiamos el comportamiento de:
    \begin{equation*}
        g(t) = e^{n(\lm_0-\lm_1)}{\left(\frac{\lm_1}{\lm_0}\right)}^{t}
    \end{equation*}
    Para ello:
    \begin{gather*}
        \ln g(t) = n(\lm_0-\lm_1) + t\ln\left(\frac{\lm_1}{\lm_0}\right) \\
        g'(t) = \ln\left(\frac{\lm_1}{\lm_0}\right)
    \end{gather*}
    Y observamos que:
    \begin{equation*}
        g'(t) > 0 \Longleftrightarrow \ln\left(\frac{\lm_1}{\lm_0}\right) > 0 \Longleftrightarrow \frac{\lm_1}{\lm_0} > 1 \Longleftrightarrow \lm_1 > \lm_0
    \end{equation*}
    Por lo que distinguimos casos (sabemos que $\lm_0\neq \lm_1$):
    \begin{description}
        \item [Supuesto que $\lm_0<\lm_1$.] Tenemos entonces que $g$ es estrictamente creciente:

            \begin{figure}[H]
                \centering
                \begin{tikzpicture}
                    % Ejes
                    \draw[-Stealth] (-0.2,0) -- (3,0) node[right] {$T$};
                    \draw[-Stealth] (0,-0.2) -- (0,3) node[above] {$\lm$};

                    % Recta con pendiente negativa
                    \draw[thick] (0,0.65) -- (3,2.65);

                    % Punto c en el eje x
                    \coordinate (C) at (2,0);
                    \draw (C) node[below] {$c$};

                    % Proyección sobre la recta
                    \coordinate (K) at (2,2); % Ajustado para caer sobre la recta dibujada
                    \draw[dashed] (C) -- (K);
                    \draw[dashed] (K) -- (0,2);

                    % Etiqueta k en el eje y
                    \draw (0,2) node[left] {$k$};

                    % Punto en la recta
                    \fill (K) circle (2pt);
                \end{tikzpicture}
            \end{figure}
            Por tanto:
            \begin{equation*}
                \left.\begin{array}{l}
                    \lm > k \\
                    \lm = k \\
                    \lm < k
                \end{array}\right\} \Longleftrightarrow \left\{\begin{array}{l}
                    T > k \\
                    T = k \\
                    T < k 
                \end{array}\right.
            \end{equation*}
            De donde el test será:
            \begin{equation*}
                \varphi'(X_1, \ldots, X_n) = \left\{\begin{array}{ll}
                    1 & \text{si\ } \sum\limits_{i=1}^{n} X_i > c \\
                     \gamma& \text{si\ } \sum\limits_{i=1}^{n}X_i = c \\
                     0& \text{si\ } \sum\limits_{i=1}^{n}X_i < c
                \end{array}\right. 
            \end{equation*}
            Imponiendo tamaño $\alpha$:
            \begin{equation*}
                \alpha = \sup_{\lm \in \Theta_0}\beta_\varphi(\lm) = \beta_\varphi(\lm_0) = E_{\lm_0}[\varphi] = P_{\lm_0}[T > c] + \gamma P_{\lm_0}[T=c]
            \end{equation*}
            Y lo que haremos será para cierto $\alpha$ calcular los valores de $c$ y $\gamma$.
        \item [Supuesto que $\lm_1 < \lm_0$.] En este caso tendremos que $g$ es estrictamente decreciente:
            \begin{figure}[H]
                \centering
                \begin{tikzpicture}
                    % Ejes
                    \draw[-Stealth] (-0.2,0) -- (3,0) node[right] {$T$};
                    \draw[-Stealth] (0,-0.2) -- (0,3) node[above] {$\lm$};

                    % Recta con pendiente negativa
                    \draw[thick] (0,2.65) -- (3,0.65);

                    % Punto c en el eje x
                    \coordinate (C) at (1.5,0);
                    \draw (C) node[below] {$c$};

                    % Proyección sobre la recta
                    \coordinate (K) at (1.5,1.65); % Ajustado para caer sobre la recta dibujada
                    \draw[dashed] (C) -- (K);
                    \draw[dashed] (K) -- (0,1.65);

                    % Etiqueta k en el eje y
                    \draw (0,1.65) node[left] {$k$};

                    % Punto en la recta
                    \fill (K) circle (2pt);
                \end{tikzpicture}
            \end{figure}
            Por lo que obtendremos un test equivalente de la forma:
            \begin{equation*}
                \varphi'(X_1, \ldots, X_n) = \left\{\begin{array}{ll}
                    1 & \text{si\ } \sum\limits_{i=1}^{n}X_i<c \\
                     \gamma& \text{si\ } \sum\limits_{i=1}^{n}X_i = c \\
                     0 & \text{si\ } \sum\limits_{i=1}^{n}X_i > c
                \end{array}\right. 
            \end{equation*}
            Imponiendo tamaño $\alpha$:
            \begin{equation*}
                \alpha = \sup_{\lm \in \Theta_0}\beta_\varphi(\lm) = \beta_\varphi(\lm_0) = E_{\lm_0}[\varphi] = P_{\lm_0}[T < c] + \gamma P_{\lm_0}[T=c]
            \end{equation*}
            Y lo que haremos será para cierto $\alpha$ calcular los valores de $c$ y $\gamma$.
    \end{description}~\\

    \noindent
    \textit{Aplicación.} Sea:
    \begin{equation*}
        X \equiv \text{``Número de llamadas por minuto''}\rightsquigarrow\cc{P}(\lm)
    \end{equation*}
    En 5 observaciones $(x_1, \ldots, x_5)$ se ha obtenido:
    \begin{equation*}
        \sum_{i=1}^{n} x_i = 12
    \end{equation*}
    Y planteamos el contraste de hipótesis:
    \begin{equation*}
        \left\{\begin{array}{l}
            H_0 : \lm = 1.5 \\
            H_1 : \lm = 2
        \end{array}\right.
    \end{equation*}
    Sea $\alpha = 0.05$, como $\lm_0<\lm_1$ nos encontramos en el primer caso, por lo que usamos el test:
    \begin{equation*}
        \varphi'(X_1, \ldots, X_n) = \left\{\begin{array}{ll}
            1 & \text{si\ } \sum\limits_{i=1}^{n}X_i>c \\
             \gamma& \text{si\ } \sum\limits_{i=1}^{n}X_i = c \\
             0 & \text{si\ } \sum\limits_{i=1}^{n}X_i < c
        \end{array}\right. 
    \end{equation*}
    Imponiendo tamaño $\alpha$:
    \begin{equation*}
        0.05 = \alpha = \sup_{\lm \in \Theta_0}\beta_\varphi(\lm) = P_{\lm_0}[T > c] + \gamma P_{\lm_0}[T=c] = P_{1.5}[T>c] + \gamma P_{1.5}[T=c]
    \end{equation*}
    Recordamos que $T\rightsquigarrow \cc{P}(n\lm)$, y como nos encontramos bajo hipótesis nula, tenemos que $T\rightsquigarrow \cc{P}(5\cdot 1.5) \equiv \cc{P}(7.5)$. Buscamos el primer valor de $c$ de forma que se tenga $1-P[T\leq c] = P[T>c]\leq \alpha$:
    \begin{itemize}
        \item Para $c=12$, se tiene:
            \begin{equation*}
                P[T>c] = 0.0427 \leq 0.05 =  \alpha
            \end{equation*}
        \item Para $c=11$ se tiene que $P[T>c]$ es claramente mayor que $0.05$
    \end{itemize}
    Por lo que tomamos $c=12$. Calculamos ahora $\gamma$:
    \begin{equation*}
        0.05 = 0.0427 + \gamma \cdot 0.0366 \quad\Longrightarrow\quad \gamma = \frac{0.05-0.0427}{0.0366} \approx 0.199454
    \end{equation*}
    De esta forma, el test a usar es:
    \begin{equation*}
        \varphi'(X_1, \ldots, X_n) = \left\{\begin{array}{ll}
            1 & \text{si\ } \sum\limits_{i=1}^{n}X_i<12 \\
             0.199454 & \text{si\ } \sum\limits_{i=1}^{n}X_i = 12 \\
             0 & \text{si\ } \sum\limits_{i=1}^{n}X_i > 12
        \end{array}\right. 
    \end{equation*}
    Como $\sum\limits_{i=1}^{5} x_i = 12$, tenemos que:
    \begin{equation*}
        \varphi'(x_1, \ldots, x_5) = 0.199454
    \end{equation*}
    Por lo que rechazamos la hipótesis nula con dicha probabilidad. Lo mejor sería repetir el experimento para obtener una decisión clara. Calculamos ahora la potencia del test:
    \begin{equation*}
        \sup_{\lm \in \Theta_1}\beta_\varphi(\lm) = \beta_\varphi(\lm_1) = E_{\lm_1}[\varphi] = P_{\lm_1}[T>12] + 0.199454\cdot  P_{\lm_1}[T=12]
    \end{equation*}
    Recordemos que $T\rightsquigarrow \cc{P}(5\lm_1)\equiv \cc{P}(10)$:
    \begin{equation*}
        \sup_{\lm\in \Theta_1}\beta_\varphi(\lm) = P_2[T> 12] + 0.199454 \cdot P_2[T=12] = 0.227408
    \end{equation*}
\end{ejercicio}

\begin{ejercicio} % // TODO: HACER
    Dada una muestra de tamaño $n$ de una variable con distribución $\cc{N}(\mu, \sigma_0^2)$, deducir el test más potente de tamaño arbitrario para contrastar hipótesis simples sobre $\mu$.
\end{ejercicio}

\begin{ejercicio}% // TODO: HACER
    Dada una muestra de tamaño $n$ de una variable aleatoria con distribución $U(-\theta,\theta)$, deducir el test más potente de tamaño $\alpha$ para contrastar $H_0 : \theta = \theta_0$ frente a $H_1 : \theta = \theta_1$ y calcular su potencia. ¿Cuál es el test óptimo fijado un nivel de significación arbitrario?
\end{ejercicio}

\begin{ejercicio}% // TODO: HACER
    Deducir el test más potente de tamaño arbitrario para contrastar $h_0 : \theta = \theta_0$ frente a $H_1 : \theta = \theta_1$, basándose en una muestra de tamaño $n$ de una variable aleatoria con función de densidad
    \begin{equation*}
        f_\theta(x) = \frac{\theta}{x^2}, \qquad x>\theta
    \end{equation*}
    Deducir el test óptimo para un nivel de significación arbitrario.
\end{ejercicio}

\begin{ejercicio}% // TODO: HACER
    Construir el test de Neyman-Pearson de tamaño $\alpha$ para contrastar $H_0:\theta = \theta_0$ frente a $H_1 : \theta = \theta_1$, basándose en una muestra de tamaño $n$ de una variable aleatoria con función de densidad
    \begin{equation*}
        f_\theta(x) = \frac{1}{x\ln \theta}, \qquad 1<x<\theta
    \end{equation*}
    Deducir el test óptimo para un nivel de significación arbitrario.\\

    \noindent
    En este caso tenemos:
    \begin{equation*}
        \Theta_0 = \{\theta_0\}, \qquad \Theta_1 = \{\theta_1\}, \qquad \Theta = \{\theta_0,\theta_1\}
    \end{equation*}
    Sea $X$ una variable aleatoria cuya función de densidad es la enunciada y sea $(X_1, \ldots, X_n)$ una muestra aleatoria simple de la misma, el test de Neymann-Pearson de tamaño $\alpha$ para el contraste de hipótesis anunciadas será de la forma:
    \begin{equation*}
        \varphi(X_1, \ldots, X_n) = \left\{\begin{array}{ll}
            1 & \text{si\ } \lm(X_1, \ldots, X_n)>k \\
            \gamma & \text{si\ } \lm(X_1, \ldots, X_n) = k \\
            0 & \text{si\ } \lm(X_1, \ldots, X_n) < k 
        \end{array}\right. 
    \end{equation*}
    Para ciertas constantes $\gamma$ y $k$. Calculamos en primer lugar:
    \begin{align*}
        f_\theta^n(x_1, \ldots, x_n) &= \prod_{i=1}^{n} f_\theta(x_i) = \prod_{i=1}^{n} \frac{I_{\mathbb{R}^+}(x_{(1)}-1)I_{\mathbb{R}^-}(x_{(n)}-\theta)}{x_i \ln\theta} \\ &= {\left({(\ln\theta)}^{n} \prod\limits_{i=1}^{n}x_i\right)}^{-1}I_{\mathbb{R}^+}(x_{(1)}-1)I_{\mathbb{R}^-}(x_{(n)}-\theta)
    \end{align*}
    Y vemos que $f_\theta^n(x_1, \ldots, x_n) \neq 0 \Longleftrightarrow x_{(1)}>1 \land x_{(n)}<\theta$. Para dichas realizaciones muestrales calculamos:
    \begin{equation*}
        \lm(x_1, \ldots, x_n) = \frac{f_{\theta_0}^n(x_1, \ldots, x_n)}{f_{\theta_1}^n(x_1, \ldots, x_n)} = \frac{{(\ln\theta_{1})}^{n}\prod\limits_{i=1}^{n}x_i}{{(\ln\theta_0)}^{n}\prod\limits_{i=1}^{n}x_i} = {\left(\frac{\ln\theta_1}{\ln\theta_0}\right)}^{n}
    \end{equation*} % // TODO: TERMINAR
\end{ejercicio}

\begin{ejercicio}% // TODO: HACER
    Sea $X$ una observación de una variable aleatoria con función de densidad
    \begin{equation*}
        f_\theta(x) = \frac{1}{\theta} e^{\nicefrac{-x}{\theta}}, \qquad x>\theta
    \end{equation*}
    Construir el test de razón de verosimilitudes de tamaño $\alpha$ arbitrario para contrastar
    \begin{align*}
        H_0 &: \theta = \theta_0 \\
        H_1 &: \theta < \theta_1
    \end{align*}
\end{ejercicio}

\begin{ejercicio}
    En base a una observación de $X\rightsquigarrow \{B(n,p) : p\in \left]0,1\right[\}$, deducir el test de razón de verosimilitudes para contrastar la hipótesis de que el parámetro $p$ no supera un determinado valor, $p_0$.\\

    \noindent
    Planteamos el contraste de hipótesis:
    \begin{equation*}
        \left\{\begin{array}{l}
            H_0 : p\leq p_0 \\
            H_1 : p> p_0
        \end{array}\right.
    \end{equation*}
    Por lo que tendremos:
    \begin{equation*}
        \Theta_0 = \left]0,p_0\right], \qquad \Theta_1 = \left]p_0,1\right[, \qquad \Theta = \left]0,1\right[
    \end{equation*}
    Calculamos el EMV de $p$ para $x\in \cc{X}$:
    \begin{gather*}
        L_x(p) = P_p[X=x] = \binom{n}{x}p^x {(1-p)}^{n-x} \\ 
        \ln L_x(p) =\ln\binom{n}{x} + x\ln p + (n-x)\ln(1-p) \\
        \dfrac{\partial \ln L_x(p)}{\partial p} = \frac{x}{p} + \frac{x-n}{1-p} = \frac{x - xp + xp - np}{p(1-p)} = \frac{x-np}{p(1-p)} = 0 \Longleftrightarrow p = \frac{x}{n}
    \end{gather*}
    De esta forma:
    \begin{equation*}
        \sup_{p \in \Theta_0}L_x(p) = \left\{\begin{array}{ll}
            L_x(\nicefrac{x}{n}) & \text{si\ } p_0 \geq \nicefrac{x}{n}  \\
             L_x(p_0) & \text{si\ } p_0 < \nicefrac{x}{n}
        \end{array}\right. 
    \end{equation*}
    Por lo que:
    \begin{equation*}
        \lm(x) = \frac{\sup\limits_{p\in \Theta_0}L_x(p)}{\sup\limits_{p\in \Theta}L_x(p)} = \left\{\begin{array}{ll}
            1 & \text{si\ } p_0 \geq \nicefrac{x}{n}  \\
             \dfrac{L_x(p_0)}{L_x(\nicefrac{x}{n})}& \text{si\ } p_0<\nicefrac{x}{n}
        \end{array}\right. 
    \end{equation*}
    Y calculamos para el caso $p_0<\frac{x}{n}$ dicha fracción:
    \begin{equation*}
        \frac{L_x(p_0)}{L_x(\nicefrac{x}{n})} = \frac{\binom{n}{x}p_0^x {(1-p_0)}^{n-x}}{\binom{n}{x}{\left(\frac{x}{n}\right)}^{x}{\left(1-\frac{x}{n}\right)}^{n-x}} = {\left(\frac{np_0}{x}\right)}^{x} {\left(\frac{1-p_0}{1-\frac{x}{n}}\right)}^{n-x}
    \end{equation*}
    Y tratamos de ver si es creciente o decreciente en función de $x$, para lo que consideramos:
    \begin{gather*}
        g(x) = {\left(\frac{np_0}{x}\right)}^{x} {\left(\frac{1-p_0}{1-\frac{x}{n}}\right)}^{n-x}  \\
        \ln g(x) = x \ln (np_0) - x\ln x + (n-x)\ln (1-p_0) - (n-x)\ln \left(1-\frac{x}{n}\right) \\
        \dfrac{\partial \ln g(x)}{\partial  x} = \ln(np_0) - \ln x \cancel{- 1} - \ln(1-p_0) + \ln\left(1-\frac{x}{n}\right) \cancel{+1} = \ln\left(\frac{np_0}{x}\right) + \ln\left(\frac{1-\frac{x}{n}}{1-p_0}\right)
    \end{gather*}
    Para el primer sumando:
    \begin{equation*}
        p_0 < \frac{x}{n} \Longleftrightarrow np_0 < x \Longleftrightarrow \frac{np_0}{x}<1 \Longleftrightarrow \ln\left(\frac{np_0}{x}\right) < 0
    \end{equation*}
    Para el segundo:
    \begin{equation*}
        p_0 < \frac{x}{n}\Longleftrightarrow -p_0 > -\frac{x}{n} \Longleftrightarrow 1-p_0>1-\frac{x}{n} \Longleftrightarrow \frac{1-\frac{x}{n}}{1-p_0} < 1 \Longleftrightarrow \ln\left(\frac{1-\frac{x}{n}}{1-p_0}\right) < 0
    \end{equation*}
    Por tanto tenemos que $g$ es estrictamente decreciente, de donde si consideramos $\lm$ como función de $x$ ($p_0 \geq \nicefrac{x}{n} \Longleftrightarrow x \leq np_0$):
    \begin{figure}[H]
        \centering
        \begin{tikzpicture}
            % Ejes
            \draw[-Stealth] (-0.2,0) -- (6,0) node[right] {$x$};
            \draw[-Stealth] (0,-0.2) -- (0,3.5) node[left] {$\lm$};

            % Recta con pendiente negativa
            \draw[thick] (3,2.65) -- (6,0.65);
            \draw[thick] (0,2.65) -- (3,2.65);

            % Punto c en el eje x
            \coordinate (C) at (4.5,0);
            \draw (C) node[below] {$k$};

            % Proyección sobre la recta
            \coordinate (K) at (4.5,1.65); % Ajustado para caer sobre la recta dibujada
            \draw[dashed] (C) -- (K);
            \draw[dashed] (K) -- (0,1.65);

            \draw[dashed] (3,0) -- (3,2.65);
            \draw (3,0) node[below] {$np_0$};

            % Etiqueta k en el eje y
            \draw (0,1.65) node[left] {$c$};
            \draw (0,2.65) node[left] {$1$};

            % Punto en la recta
            \fill (K) circle (2pt);
        \end{tikzpicture}
    \end{figure}
    \noindent
    Tendremos que:
    \begin{equation*}
        \left.\begin{array}{l}
            \lm < c \\
            \lm \geq c
        \end{array}\right\} \Longleftrightarrow \left\{\begin{array}{l}
            x > k \\
            x \leq k
        \end{array}\right. \qquad \text{con\ } k\geq np_0
    \end{equation*}
    Por lo que el test a considerar será:
    \begin{equation*}
        \varphi(X) = \left\{\begin{array}{ll}
            1 & \text{si\ } X>k \\
             0 & \text{si\ } X\leq k
        \end{array}\right. 
    \end{equation*}
    Y determinaremos el valor de $k$ imponiendo un nivel de significación $\alpha$:
    \begin{equation*}
        \alpha \geq \sup_{p\in \Theta_0}\beta_\varphi(p) = \sup_{p\in \Theta_0} E_p[\varphi] = \sup_{p\in \Theta_0} P_p[X>k] = \sup_{p\in \Theta_0} \left(\sum_{t=\lfloor k \rfloor+1}^{n} P_p[X=t] \right)
    \end{equation*} 
    Veamos ahora que fijado $t\in \{\lfloor k\rfloor+1,\ldots,n\}$ se tiene que $P_p[X=t]$ es creciente en función de $p\in \Theta_0$. Para ello, si tomamos:
    \begin{gather*}
        h(p) = P_p[X=t] = \binom{n}{x} p^t {(1-p)}^{n-t} \quad\Longrightarrow\quad \ln h(p) =\ln\binom{n}{x}+ t\ln p + (n-t)\ln (1-p) \\
        \dfrac{\partial \ln h(p)}{\partial p} = \frac{t-np}{p(1-p)}
    \end{gather*}
    Observamos ahora que:
    \begin{equation*}
        t \geq \lfloor k\rfloor +1 > k \geq np_0 \geq np 
    \end{equation*}
    De donde deducimos que $h(p)$ es creciente, por lo que el supremo anterior se alcanza en $p_0$:
    \begin{equation*}
        \alpha \geq \sup_{p\in \Theta_0}\beta_\varphi(p) = \sup_{p\in \Theta_0} \left(\sum_{t=\lfloor k \rfloor+1}^{n} P_p[X=t] \right) = \sum_{t=\lfloor k \rfloor +1}^{n} P_{p_0}[X=t]
    \end{equation*} 
    En un caso concreto de esta fórmula se despeja el valor de $k$.
\end{ejercicio}

\begin{ejercicio}
    Sea $X$ una variable con función de densidad
    \begin{equation*}
        f_\theta(x) = \theta x^{\theta-1}, \qquad 0<x<1
    \end{equation*}
    Basándose en una observación de $X$, deducir el test de razón de verosimilitudes de tamaño arbitrario para contrastar
    \begin{align*}
        H_0 &: \theta \leq \theta_0 \\
        H_1 &: \theta > \theta_1
    \end{align*}~\\

    \noindent
    Tendremos:
    \begin{equation*}
        \Theta_0 = \left]0,\theta_0\right], \qquad \Theta_1 = \left]\theta_0,+\infty\right[, \qquad \Theta = \mathbb{R}^+
    \end{equation*}
    Calculamos el EMV de $\theta$ para $x\in \cc{X}$:
    \begin{gather*}
        L_x(\theta) = f_\theta(x) = \theta x^{\theta-1} \quad\Longrightarrow\quad \ln L_x(\theta) = \ln\theta + (\theta-1)\ln x \\
        \frac{\partial  \ln L_x(\theta)}{\partial \theta} = \frac{1}{\theta} + \ln x = 0 \Longleftrightarrow \theta = \frac{-1}{\ln x}
    \end{gather*}
    Y observemos que como $x\in \left]0,1\right[$ tenemos que $\ln x < 0$. El EMV es $\hat{\theta} = \nicefrac{-1}{\ln X}$. Por tanto:
    \begin{equation*}
        \sup_{\theta\in \Theta_0} L_x(\theta) = \left\{\begin{array}{ll}
            L_x(\nicefrac{-1}{\ln x}) & \text{si\ } \theta_0 \geq \nicefrac{-1}{\ln x} \\
             L_x(\theta_0)& \text{si\ } \theta_0 < \nicefrac{-1}{\ln x}
        \end{array}\right. 
    \end{equation*}

    de donde:
    \begin{equation*}
        \lm(x) = \dfrac{\sup\limits_{\theta\in \Theta_0}L_x(\theta)}{\sup\limits_{\theta\in \Theta}L_x(\theta)} = \left\{\begin{array}{ll}
            1 & \text{si\ } \theta_0 \geq \nicefrac{-1}{\ln x} \\
              \dfrac{L_x(\theta_0)}{L_x(\nicefrac{-1}{\ln x})}& \text{si\ } \theta_0 < \nicefrac{-1}{\ln x}
        \end{array}\right. 
    \end{equation*}
    En el caso $\theta_0 < \frac{-1}{\ln x}$:
    \begin{equation*}
        \dfrac{L_x(\theta_0)}{L_x(\nicefrac{-1}{\ln x})} = \dfrac{\theta_0 x^{\theta_0-1}}{\frac{-1}{\ln x}x^{\left(\frac{-1}{\ln x}-1\right)}} = -\theta_0 \ln x \cdot x^{\theta_0 + \frac{1}{\ln x}}
    \end{equation*}
    Y veamos si en este caso $\lm$ es creciente o decreciente en función de $x$. Para ello, consideramos:
    \begin{gather*}
        g(x) =  -\theta_0 \ln x \cdot x^{\theta_0 + \frac{1}{\ln x}}  \\
        \ln g(x) = \ln(-\theta_0 \ln x) + \left(\theta_0 + \frac{1}{\ln x}\right)\ln x = \ln(\theta_0) + \ln(-\ln x) + \theta_0 \ln x + 1 \\
        \dfrac{\partial \ln g(x)}{\partial  x} = \frac{-1}{x\ln x} + \frac{\theta_0}{x} = \frac{-1 + \theta_0 \ln x}{x\ln x} < 0 \qquad (x\in \left]0,1\right[ )
    \end{gather*}
    Por lo que $g$ es estrictamente decreciente, de donde si vemos $\lm$ como función de $x$ ($\theta_0 \geq \nicefrac{-1}{\ln x} \Longleftrightarrow \ln x \leq \nicefrac{-1}{\theta_0} \Longleftrightarrow x \leq e^{\nicefrac{-1}{\theta_0}}$):
    \begin{figure}[H]
        \centering
        \begin{tikzpicture}
            % Ejes
            \draw[-Stealth] (-0.2,0) -- (6,0) node[right] {$x$};
            \draw[-Stealth] (0,-0.2) -- (0,3.5) node[left] {$\lm$};

            % Recta con pendiente negativa
            \draw[thick] (3,2.65) -- (6,0.65);
            \draw[thick] (0,2.65) -- (3,2.65);

            % Punto c en el eje x
            \coordinate (C) at (4.5,0);
            \draw (C) node[below] {$k$};

            % Proyección sobre la recta
            \coordinate (K) at (4.5,1.65); % Ajustado para caer sobre la recta dibujada
            \draw[dashed] (C) -- (K);
            \draw[dashed] (K) -- (0,1.65);

            \draw[dashed] (3,0) -- (3,2.65);
            \draw (3,0) node[below] {$e^{\nicefrac{-1}{\theta_0}}$};

            % Etiqueta k en el eje y
            \draw (0,1.65) node[left] {$c$};
            \draw (0,2.65) node[left] {$1$};

            % Punto en la recta
            \fill (K) circle (2pt);
        \end{tikzpicture}
    \end{figure}
    \noindent
    Tenemos entonces que:
    \begin{equation*}
        \left.\begin{array}{l}
            \lm < c \\
            \lm \geq c
        \end{array}\right\} \Longleftrightarrow \left\{\begin{array}{l}
            x > k \\
            x \leq k
        \end{array}\right. \qquad \text{con}\quad  k\geq e^{\nicefrac{-1}{\theta_0}}
    \end{equation*}
    Y el test a considerar será:
    \begin{equation*}
        \varphi(X) = \left\{\begin{array}{ll}
            1 & \text{si\ } X>k \\
            0 & \text{si\ } X\leq k
        \end{array}\right. 
    \end{equation*}
    y determinamos el valor de $k$ imponiendo tamaño $\alpha$:
    \begin{align*}
        \alpha &= \sup_{\theta\in \Theta_0}\beta_\varphi(\theta) = \sup_{\theta\in \Theta_0} E_\theta[\varphi] = \sup_{\theta\in \Theta_0} P_\theta[X>k] = \sup_{\theta\in \Theta_0}\left( \int_{k}^{1} \theta t^{\theta-1}~dt  \right) \\
               &= \sup_{\theta\in \Theta_0}\left([t^{\theta}]_k^1\right) = \sup_{\theta\in \Theta_0}(1-k^\theta) \stackrel{(k\in [0,1])}{=} 1-k^{\theta_0}
    \end{align*}
    Luego:
    \begin{equation*}
        k^{\theta_0} = 1-\alpha \quad\Longrightarrow\quad \ln k = \frac{\ln (1-\alpha)}{\theta_0} \quad\Longrightarrow\quad k = e^{\ln\left({(1-\alpha)}^{\frac{1}{\theta_0}}\right)} = {(1-\alpha)}^{\frac{1}{\theta_0}}
    \end{equation*}
\end{ejercicio}

\begin{ejercicio}% // TODO: HACER
    Un fabricante de coches asegura que la distancia media recorrida con un galón de gasolina es al menos 30 millas. Probados 9 coches de esta fábrica, la distancia media recorrida con un galón de gasolina ha sido 26 millas, y la suma de los cuadrados 6106 millas al cuadrado.
    \begin{enumerate}[label=\alph*)]
        \item Suponiendo que la distancia recorrida por estos coches con un galón de gasolina tiene distribución normal, contrastar la hipótesis del fabricante a partir de estos datos, a nivel de significación $0.01$.
        \item ¿Qué conclusión se obtendría de estos mismos datos, al mismo nivel de significación, si se sabe que la desviación típica de la variable considerada es $5.5$?
    \end{enumerate}~\\

    \noindent
    Sale 1 si $36  > \chi^2_{25,;0.02}$, $\chi^2_{25;0.02}$. Los datos no dan evidencia para rechazar.
\end{ejercicio}

\begin{ejercicio}% // TODO: HACER
    Un fabricante de baterías asegura que la desviación típica del tiempo de vida de las mismas es, a lo sumo, 70 horas. Una muestra de 26 baterías tomadas al azar ha dado una cuasidesviación típica de 84 horas. Haciendo las hipótesis adecuadas de normalidad, ¿proporcionan los datos evidencia para rechazar la hipótesis del fabricante al nivel $0.02$?
\end{ejercicio}

\begin{ejercicio}
    Un profesor asegura que tiene un nuevo método de enseñanza mejor que el usado tradicionalmente. Para comprobar si tiene razón se selecciona de forma aleatoria e independiente dos grupos de alumnos, $A$ y $B$, utilizándose el nuevo método con el grupo $A$ y el tradicional con el $B$. $A$ final de curso se hace un examen a los alumnos, obteniéndose las siguientes puntuaciones:\newline
    Grupo $A$: 6, 5, 4, 7, 3, $5.5$, 6, 7, 6.\newline
    Grupo $B$: 5, 4, 5, 6, 4, 6, 5, 3, 7.\newline
    Supuesto que las puntuaciones de cada grupo tienen distribución normal, ¿proporcionan estos datos evidencia para rechazar el nuevo método a nivel de significación $0.05$?\\

    \noindent
    Si tomamos:
    \begin{gather*}
        X\equiv \text{``Puntuaciones del grupo A (nuevo método)''} \rightsquigarrow\cc{N}(\mu_1,\sigma_1^2) \\
        Y\equiv \text{``Puntuaciones del grupo B (método tradicional)''} \rightsquigarrow\cc{N}(\mu_2,\sigma_2^2) 
    \end{gather*}
    Consideramos la hipótesis $\mu_1>\mu_2 \Longleftrightarrow \mu_1-\mu_2>0$. Como no sabemos resolver este tipo de contrastes (donde en $H_0$ no aparece una igualdad), consideramos que esta hipótesis es $H_1$, por lo que el contraste sería el siguiente:
    \begin{equation*}
        \left\{\begin{array}{l}
            H_0 : \mu_1-\mu_2\leq 0 \\
            H_1 : \mu_1-\mu_2> 0
        \end{array}\right.
    \end{equation*} % // TODO: CAMBIAR (Hacer lo de sigma = sigma)
    En teoría vimos que se rechaza la hipótesis nula si:
    \begin{equation*}
        t_{exp} >t_{n_1+n_2-2;\alpha}
    \end{equation*}

    donde:
    \begin{equation*}
        t_{exp} = \frac{\overline{x}-\overline{y}}{s_p\sqrt{\frac{1}{n_1}+\frac{1}{n_2}}}, \qquad s_p = \sqrt{\frac{(n_1-1)s_1^2+(n_2-1)s_2^2}{n_1+n_2-2}}
    \end{equation*}
    Calculamos para ello:
    \begin{itemize}
        \item $n_1 = n_2 = 9$.
        \item $\overline{x} = 5.5$.
        \item $s_1^2 = 1.75$.
        \item $\overline{y} = 5$.
        \item $s_2^2 = 1.5$.
    \end{itemize}
    Por lo que:
    \begin{equation*}
        s_p = \sqrt{\frac{8\cdot 1.75 + 8\cdot 1.5}{9+9-2}} \approx 1.274755, \qquad 
        t_{exp} = \frac{5.5 - 5}{1.274755\cdot \sqrt{\frac{2}{9}}} \approx 0.83205
    \end{equation*}
    Y si calculamos ahora $t_{n_1+n_2-2;\alpha}$:
    \begin{equation*}
        t_{n_1+n_2-2;\alpha} = t_{9+9-2;0.05} = t_{16;0.05} = 1.7459
    \end{equation*}
    Como $t_{exp}$ no es mayor que $t_{n_1+n_2-2;\alpha}$, los datos no proporcionan evidencia suficiente para rechazar el nuevo método a nivel de significación $0.05$. %// TODO: ?
\end{ejercicio}

\begin{ejercicio}
    A partir de las siguientes observaciones de muestras independientes de dos poblaciones normales, contrastar, al nivel de significación $0.01$, si la media de la primera población supera en al menos una unidad la media de la segunda.\newline
    Muestra 1: 132, 139, 126, 114, 122, 132, 141, 126.\newline
    Muestra 2: 124, 141, 118, 116, 114, 132, 145, 123, 121.\\

    \noindent
    Planteamos el constrate de hipótesis:
    \begin{equation*}
        \left\{\begin{array}{l}
            H_0 : \mu_1-\mu_2 \geq \mu_0 = 1 \\
            H_1 : \mu_1-\mu_2 < \mu_0
        \end{array}\right.
    \end{equation*}
    Sean $X\rightsquigarrow \cc{N}(\mu_1, \sigma^2)$, $Y\rightsquigarrow\cc{N}(\mu_2,\sigma^2)$, tenemos dos observaciones: $(x_1, \ldots, x_8)$ de $X$ y $(y_1, \ldots, y_9)$ de $Y$ (tomamos $n_1=8$ y $n_2 = 9$). Usaremos el estadístico:
    \begin{equation*}
        T = \frac{\overline{X}-\overline{Y}-(\mu_1-\mu_2)}{S_p\sqrt{\frac{1}{n_1}+\frac{1}{n_2}}}\rightsquigarrow t_{n_1+n_2-2} , \qquad S_p = \sqrt{\frac{(n_1-1)S_1^2 + (n_2-1)S_2^2}{n_1+n_2-2}}
    \end{equation*}
    Por lo que calculamos $t_{exp}$: 
    \begin{equation*}
        t_{exp} = \frac{\overline{x}-\overline{y}-\mu_0}{s_p\sqrt{\frac{1}{n_1}+\frac{1}{n_2}}}
    \end{equation*}

    para lo que necesitaremos:
    \begin{equation*}
        \overline{x} = 129, \qquad \overline{y} = 126, \qquad s_1^2 = 79.142857, \qquad s_2^2 = 121
    \end{equation*}
    Por lo que:
    \begin{gather*}
        s_p = \sqrt{\frac{(n_1-1)s_1^2 + (n_2-1)s_2^2}{n_1+n_2-2}} = \sqrt{\frac{7\cdot 79.142857 + 8\cdot 121}{8+9-2}}  \approx 10.073066 \\
        t_{exp} = \frac{\overline{x}-\overline{y}-\mu_0}{s_p\sqrt{\frac{1}{n_1}+\frac{1}{n_2}}} = \frac{129 - 126 -1}{10.073066 \sqrt{\frac{1}{8} + \frac{1}{9}}} \approx 0.408611
    \end{gather*}
    Sea $\alpha = 0.01$, vimos en teoría que la región crítica era:
    \begin{equation*}
        t_{exp} < -t_{n_1+n_2-2;\alpha} = -t_{15;0.01} = 2.6025
    \end{equation*}
    Por tanto, como $0.408611 < 2.6025$ estamos en la región crítica, por lo que podemos rechazar a nivel de significación $0.01$ que la media de la primera población supera en al menos una unidad la media de la segunda.
\end{ejercicio}

\begin{ejercicio}% // TODO: HACER
    Una central lechera recibe diariamente leche de dos granjas $A$ y $B$. Para comparar la calidad de los productos recibidos se ha medido el contenido en grasa en muestras de leche tomadas al azar de cada una de las granjas, con los siguientes resultados:
    \begin{table}[H]
    \centering
    \begin{tabular}{c|cccccc}
     & \multicolumn{6}{c}{Contenido en grasa $(\%)$} \\
     \hline
        Granja A & 14 & 12 & 15 & 15 & 11 & 16 \\
        Granja B & 20 & 18 & 18 & 19 & 15 &
    \end{tabular}
    \end{table}
    \begin{enumerate}[label=\alph*)]
        \item ¿Puede suponerse, a nivel de significación $0.05$, que el contenido medio en grasa de la leche de las dos granjas es el mismo? Especificar las hipótesis bajo las que se resuelve este problema.
        \item Calcular un intervalo de confianza, a nivel de confianza $0.9$, para la varianza del contenido en grasa de la leche de la granja B. A partir de dicho intervalo, deducir si puede aceptarse que la varianza de esta población es igual a 3. Especificar el problema de contraste y el test utilizado; calcular su tamaño.
    \end{enumerate}
\end{ejercicio}

    % // TODO: Están comentadas porque no están (bien) hechas
    % \section{Regresión lineal y análisis de la varianza}

% TODO: Los ejercicios no están bien hechos
% TEngo copiadas estas soluciones de clase:
%
% 1.
% b) 
% VT = 1876.8888
% VE = 1454.3558
% 
% S_R^2 = 60.3619
% 
% c)
% yp = 62.4291
% ECM = 79.9527
% 
% d)
% Fexp = 24.0939
% F_ = 5.59
% se rechaza a ns 0.05
% En base a los datos no podemos descartarlo
% 
% 
% 2.
% a)
% xmed = 21.5
% ymed = 62
% sigmax^2 = 55.25
% sigmay^2 = 201
% cov = 84.5
% r^2 = 0.6429
% 
% y = 29.1177 + 1.5294 x
% xp = 28 => yp = 71.9404
% 
% b)
% F_... = 5.32
% Fexp = 14.41
% 
% No tiene sentido: 0e no tiene sentido
% 
% 3.
% X = calificaciones medias graduados
% Y = salarios iniciales en miles
% 
% xmed = 3.1142
% ymed = 19.3
% sigmax^2 = 0.07979
% sigmay^2 = 4.7685
% suma xy = 424.56
% cov = 0.547368
% 
% a) y = -1.9982 + 6.8388 x 
% b) S_R^2 = 7.2555/5 = 1.4509
% c) r^2 = 0.7825
% d) 20.2279
% e) 
% F_ = 16.3
% Fexp = 17.9885
% 
% 5.
% Y = tiempo de reaccion
% 5 compuestos
% n = 13
% k = 5
% Y = "tiempo de reaccion"
% 
% y1med = 
% y2med = 
% y3med = 
% y4med = 
% y5med = 
% 
% Fexp = 1.98
% VT = 9.24769
% VNE = 4.645
%

\begin{ejercicio}
    Una compañía de energía eléctrica pretende desarrollar un modelo lineal para el consumo de energía en función de la temperatura media diaria durante los meses de invierno. En 9 días elegidos al azar se obtuvo la siguiente información:
    \begin{table}[H]
    \centering
    \begin{tabular}{c|ccccccccc}
        Temperatura & 0 & 2 & 4 & 8 & 13 & -4 & -6 & -8 & -11 \\
        \hline
        Consumo & 70 & 79 & 67 & 66 & 63 & 97 & 82 & 90 & 107
    \end{tabular}
    \end{table}
    \begin{enumerate}[label=\alph*)]
        \item Obtener la recta de regresión estimada e interpretar sus coeficientes.
        \item Descomponer la variabilidad de los datos de consumo. Obtener la varianza residual, el coeficiente de determinación y dar su interpretación.
        \item Obtener la predicción del consumo de energía para un día con temperatura media 10 grados y estimar el error cuadrático medio de esta predicción.
        \item Si se suponen las hipótesis adecuadas de normalidad, ¿se puede considerar, al nivel de significación $\alpha=0.05$, que el consumo no depende linealmente de la temperatura?
    \end{enumerate}
\end{ejercicio}

\begin{ejercicio}
    Una compañía de seguros desea establecer una relación lineal, y el grado de dicha relación, para determinar el montante anual del seguro de vida del cabeza de familia en función del ingreso mensual familiar. Observada una muestra aleatoria de 10 familias elegidas de forma independiente, se obtuvo la siguiente información:
    \begin{table}[H]
    \centering
    \begin{tabular}{c|cccccccccc}
        Ingreso (cientos de euros)  & 10 & 10 & 15 & 20 & 20 & 25 & 25 & 30 & 30 & 30 \\
        \hline
        Seguro (decenas de euros) & 50 & 35 & 55 & 55 & 70 & 65 & 65 & 60 & 75 & 90
    \end{tabular}
    \end{table}
    \begin{enumerate}[label=\alph*)]
        \item ¿Cabe pensar en la citada relación lineal? En caso afirmativo, estimar la recta de regresión, interpretar los coeficientes y obtener la predicción del montante del seguro de vida para un ingreso mensual de 2800 euros.
        \item Suponiendo las hipótesis adecuadas de normalidad, realizar el contraste de regresión al nivel de significación $\alpha = 0.05$ e interpretar el resultado.
    \end{enumerate}
\end{ejercicio}

\begin{ejercicio}
    Los datos de la siguiente tabla representan las calificaciones medias ($X$) de 7 recién graduados y sus respectivos salarios iniciales ($Y$) en miles de euros:
    \begin{table}[H]
    \centering
    \begin{tabular}{c|ccccccc}
        $X$ & $2.75$ & $2.85$ & $2.95$ & $3.05$ & $3.2$ & $3.4$ & $3.6$ \\
        \hline
        $Y$ & $16.8$ & $18.8$ & $17.2$ & $17.2$ & $21.2$ & $21.5$ & $22.4$
    \end{tabular}
    \end{table}
    \begin{enumerate}[label=\alph*)]
        \item Estimar la recta de regresión de $Y$ sobre $X$ e interpretar sus coeficientes.
        \item Determinar la varianza residual.
        \item Calcular e interpretar el coeficiente de determinación.
        \item Predecir el salario inicial de un estudiante con una calificación media de $3.25$.
        \item Suponiendo las hipótesis adecuadas de normalidad, contrastar la hipótesis de que no existe relación lineal entre ambas variables al nivel de significación $\alpha = 0.01$.
    \end{enumerate}
\end{ejercicio}

\begin{ejercicio}
    En cierto estudio sobre la relación entre el diámetro de los guisantes ($X$) y el diámetro medio de sus descendientes ($Y$), Galton obtuvo los siguientes resultados:
    \begin{table}[H]
    \centering
    \begin{tabular}{c|ccccccc}
        D. Padres & 21 & 20 & 19 & 18 & 17 & 16 & 15 \\
        \hline
        D. Descendientes & $17.26$ & $17.07$ & $16.37$ & $16.4$ & $16.13$ & $16.17$ & $15.98$
    \end{tabular}
    \end{table}
    \begin{enumerate}[label=\alph*)]
        \item Determinar el modelo de regresión lineal estimado de $Y$ sobre $X$ e interpretar el valor estimado de la pendiente. Dar la predicción del diámetro de los guisantes cuyos progenitores tienen un diámetro de $18.5$. Dar una medida de la bondad del ajuste de los datos a la recta estimada.
        \item Suponiendo las hipótesis adecuadas de normalidad, ¿puede deducirse, al nivel $0.05$, que no hay relación lineal entre las variables consideradas? Relacionar este resultado con las conclusiones anteriores.
    \end{enumerate}
\end{ejercicio}

\begin{ejercicio}
    Una compañía farmacéutica investiga los efectos de 5 compuestos. El experimento consiste en inyectar los compuestos a 13 ratas de características similares y anotar los tiempos de reacción. Los animales se clasifican en 5 grupos de 4, 2, 2, 3 y 2 ratas, respectivamente, y a cada grupo se le administra un compuesto diferente, obteniéndose los resultados de la siguiente tabla:
    \begin{table}[H]
    \centering
    \begin{tabular}{c|cccc}
        Grupo & \multicolumn{4}{c}{Tiempo de reacción} \\
        \hline
        1 & $8.3$ & $7.6$ & $8.4$ & $8.3$ \\
        2 & $7.4$ & $7.1$ & & \\
        3 & $8.1$ & $6.4$ & & \\
        4 & $7.9$ & $8.5$ & $10.0$ & \\
        5 & $7.1$ & 8 & &
    \end{tabular}
    \end{table}
    \noindent
    Suponiendo que se verifican las hipótesis de normalidad, aleatoriedad, independencia e igualdad de varianzas, contrastar la hipótesis de que los tiempos medios de reacción coinciden en los cinco grupos y, por tanto, la eficacia de los cinco compuestos es la misma.
\end{ejercicio}

\begin{ejercicio}
    Se quiere estudiar la eficacia de tres fertilizantes, $A$, $B$ y $C$, en la producción de cierto fruto. Para ello se aplica el $A$ en 8 parcelas, el $B$ en 6, y el $C$ en 12 parcelas. Las parcelas son de características similares en cuanto a fertilidad, por lo que se considera que las diferencias en la producción, si las hay, serán debidas al tipo de fertilizante.  Las toneladas producidas en cada parcela en una determinada temporada son:
    \begin{table}[H]
    \centering
    \begin{tabular}{|c|cccccccccccc|}
        \hline
        $A$ & 6 & 7 & 5 & 6 & 5 & 8 & 4 & 7 & & & & \\
        $B$ & 10 & 9 & 9 & 10 & 10 & 6 & & & & & & \\
        $C$ & 3 & 4 & 8 & 3 & 7 & 6 & 3 & 6 & 4 & 7 & 6 & 3 \\
        \hline
    \end{tabular}
    \end{table}
    \noindent
    Suponiendo que las tres muestras proceden de poblaciones normales con varianzas iguales, contrastar la hipótesis de que los abonos son igualmente eficaces.
\end{ejercicio}

\begin{ejercicio}
    Los siguientes datos corresponden a observaciones del consumo medio (en Kw/h) realizado por 5 tipos de calefactores para mantener una habitación a una temperatura determinada durante todo un día:
    \begin{table}[H]
    \centering
    \begin{tabular}{c|cccccc}
        Tipo & \multicolumn{6}{c}{Consumo (en kw/h)} \\
        \hline
        1 & $14.5$ & $14.1$ & $14.6$ & $14.2$ & & \\
        2 & $13.2$ & $13.4$ & $13.0$ & & & \\
        3 & $13.7$ & $13.6$ & $14.1$ & $13.8$ & $14.0$ & \\
        4 & $12.7$ & $13.1$ & $12.8$ & $12.9$ & $13.3$ & $13.2$ \\
        5 & $14.6$ & $15.2$ & $14.4$ & $14.8$ & $14.3$ & 
    \end{tabular}
    \end{table}
    \noindent
    Contrastar la hipótesis de igualdad de los consumos medios de los diferentes tipos de calefactores. ¿Bajo qué hipótesis se puede realizar este contraste?
\end{ejercicio}

\begin{ejercicio}
    En un tratamiento contra la hipertensión se seleccionaron 35 enfermos de características similares. Los enfermos se distribuyeron en cuatro grupos de 10 ($P, A, B$ y $AB$). El grupo $P$ tomó ``placebo'' (fármaco inocuo), el grupo $A$ tomó un fármaco ``$A$'', el grupo $B$ un fármaco ``$B$'' y el grupo $AB$ una asociación entre ``$A$'' y ``$B$''. Para valorar la eficacia de los tratamientos, se registró el descenso de la presión diastólica desde el inicio del tratamiento hasta después de una semana de tratamiento. Los resultados, después de registrarse algunos abandonos, fueron:
    \begin{table}[H]
    \centering
    \begin{tabular}{|c|cccccccccc|}
        \hline
        $P$ & 10 & 0 & 15 & -20 & 0 & 15 & -5 & & & \\
        $A$ & 20 & 25 & 33 & 25 & 30 & 18 & 27 & 0 & 35 & 20 \\
        $B$ & 15 & 10 & 25 & 30 & 15 & 35 & 25 & 22 & 11 & 25 \\
        $AB$ & 10 & 5 & -5 & 15 & 20 & 20 & 0 & 10 & & \\
        \hline
    \end{tabular}
    \end{table}
    \noindent
    A la vista de estos datos, ¿Puede afirmarse que el descenso de la presión diastólica coincide en los cuatro grupos? ¿Bajo qué hipótesis?
\end{ejercicio}

    % \section{Constrastes de hipótesis no paramétricos}

\begin{ejercicio}
    A partir de los siguientes datos, que muestran el número de accidentes en un determinado regimiento del ejército durante 200 días elegidos al azar, contrastar si el número de accidentes diarios sigue una distribución de Poisson de parámetro 2.    
    \begin{table}[H]
    \centering
    \begin{tabular}{c|cccccccc}
        \hline
        Nº de accidentes & 0 & 1 & 2 & 3 & 4 & 5 & 6 & 7 \\
        \hline
        Nº de días & 22 & 53 & 58 & 39 & 20 & 5 & 2 & 1 \\
        \hline
    \end{tabular}
    \end{table}~\\

    \noindent
    Sea la variable aleatoria:
    \begin{equation*}
        X \equiv \text{``Número de accidentes en un día.''}
    \end{equation*}
    que se distribuye según una cierta función de distribución $F$, se plantea el contraste de hipótesis:
    \begin{equation*}
        \left\{\begin{array}{l}
                H_0 : F = F_{\cc{P}(2)}\\
                H_1 : F\neq F_{\cc{P}(2)}
        \end{array}\right.
    \end{equation*}
    En teoría hemos visto cómo resolver este contraste usando dos tests distintos:
    \begin{itemize}
        \item Test $\chi^2$ de Pearson.
        \item Test de Kolmogorov-Smirnov.
    \end{itemize}
    El segundo solo se puede aplicar para distribuciones continuas, por lo que no podemos usarlo en este caso. Usaremos por tanto el test $\chi^2$ de Pearson, que usa el estadístico:
        \begin{equation*}
            \chi^2(N_0,\ldots,N_k) = \sum_{i=0}^{k} \frac{{(N_i - np_i^0)}^{2}}{np_i^0} = -n + \sum_{i=0}^{k} \frac{N_i^2}{np_i^0} 
        \end{equation*}
        donde tenemos $n=200$, $k=8$ y los $N_i$ dados por la tabla anterior. Calculamos cada una de las $p_i^0 = P_{H_0}[X=i]$:
        \begin{table}[H]
        \centering
        \begin{tabular}{cccccccc}
            $p_0^0$ & $p_1^0$ & $p_2^0$ & $p_3^0$ & $p_4^0$ & $p_5^0$ & $p_6^0$ & $p_7^0$ \\
            \hline
            $0.1353$ & $0.2707$ & $0.2707$ & $0.1804$ & $0.0902$ & $0.0361$ & $0.0120$ & $0.0034$ \\
        \end{tabular}
        \end{table}
        \noindent
        de donde:
        \begin{equation*}
            \chi^2_{\text{exp}} = -n + \sum_{i=0}^{k} \frac{N_i^2}{np_i^0} = -n + \frac{1}{n}\sum_{i=0}^{k} \frac{N_i^2}{p_i^0} \approx 2.834515
        \end{equation*}
        Si calculamos ahora el $p-$valor obtenemos que:
        \begin{equation*}
            P[\chi^2(k-1)\geq \chi^2_{\text{exp}}] = P[\chi^2(7)\geq 2.834515] \approx P[\chi^2(7) \geq 2.8331] = 0.9
        \end{equation*}
        Por lo que no podemos rechazar $H_0$.
\end{ejercicio}

\begin{ejercicio}
   Una tela cuadrada tiene 60 defectos de fabricación. Con objeto de analizar la distribución de los defectos en la superficie de la tela, se ha dividido ésta en 9 zonas cuadradas exactamente iguales, observándose los siguientes defectos en cada zona: 
   \begin{table}[H]
   \centering
   \begin{tabular}{|c|c|c|}
       \hline
       8 & 7 & 3 \\
       \hline
       5 & 9 & 11 \\
       \hline
       6 & 4 & 7 \\
       \hline
   \end{tabular}
   \end{table}
   \noindent
   Contrastar, a partir de estos datos, si los defectos se distribuyen uniformemente en toda la superficie o, por el contrario, siguen algún patrón de ocurrencia.\\

   \noindent
   Sea la variable aleatoria:
   \begin{equation*}
       X\equiv \text{``Región en la que se localiza un defecto.''} 
   \end{equation*}
   que se distribuye según una función de distribución $F$, planteamos el contraste:
   \begin{equation*}
       \left\{\begin{array}{l}
           H_0: F=F_{U(1,2,\ldots,9)} \\
           H_1: F\neq F_{U(1,2,\ldots,9)}
       \end{array}\right.
   \end{equation*}
   Usaremos el test $\chi^2$ de Pearson, pues la distribución es discreta. Para ello, calculamos primero $p_i^0 = P_{H_0}[X = i]$, obteniendo en este caso que:
   \begin{equation*}
       p_i^0 = \frac{1}{9} \qquad \forall i \in \{1,\ldots,9\}
   \end{equation*}
   Por lo que ($k = 9$, $n=60$):
   \begin{equation*}
       \chi^2_{\text{exp}} = -n + \sum_{i=1}^{k} \frac{N_i^2}{np_i^0} = -n + \frac{1}{np_1^0}\sum_{i=1}^{k} N_i^2 = 7.5
   \end{equation*}
   Calculamos ahora el valor del $p-$valor:
   \begin{equation*}
       P_{H_0}[\chi^2(k-1)\geq \chi^2_{\text{exp}}] = P[\chi^2(8)\geq 7.5] \approx 0.475
   \end{equation*}
   Por lo que los datos no aportan evidencia suficiente para rechazar la hipótesis nula.
\end{ejercicio}

\begin{ejercicio}
    Un modelo genético indica que la distribución de una población de hombres y mujeres, daltónicos o no, se ajusta a probabilidades de la forma:
    \begin{table}[H]
    \centering
    \begin{tabular}{|c|c|c|}
        \hline
        & Hombres & Mujeres \\
        \hline
        No daltónicos & $\nicefrac{(1-p)}{2}$ & $\nicefrac{(1-p^2)}{2}$ \\
        \hline
        Daltónicos & $\nicefrac{p}{2}$ & $\nicefrac{p^2}{2}$ \\
        \hline
    \end{tabular}
    \end{table}
   \noindent
    Para comprobar esta teoría se examinaron 2000 individuos de la población, elegidos al azar, obteniéndose los siguientes resultados:
    \begin{table}[H]
    \centering
    \begin{tabular}{|c|c|c|}
        \hline
        & Hombres & Mujeres \\
        \hline
        No daltónicos & 894 & 1015 \\
        \hline
        Daltónicos & 81 & 10 \\
        \hline
    \end{tabular}
    \end{table}
   \noindent
    Contrastar, mediante el test de la $\chi^2$, si esta muestra concuerda con el modelo teórico.
\end{ejercicio}

\begin{ejercicio}
    Un laboratorio farmacéutico afirma que uno de sus productos confiere inmunidad a la picadura de insectos durante un tiempo exponencial de media $2.5$ horas. Probado el producto en 20 sujetos, en un ambiente con gran número de mosquitos, los instantes (en horas) en que recibieron la primera picadura fueron:
    \begin{table}[H]
    \centering
    \begin{tabular}{cccccccccc}
        $0.01$ & $0.02$ & $0.03$ & $0.23$ & $0.51$ & $0.74$ & $0.96$ & $1.17$ & $1.46$ & $1.62$ \\
        $2.18$ & $2.25$ & $2.79$ & $3.45$ & $3.82$ & $3.92$ & $4.27$ & $5.43$ & $5.79$ & $6.34$
    \end{tabular}
    \end{table}
    \noindent
    Usando el test de Kolmogorov-Smirnov, contrastar, a partir de estos datos, si puede aceptarse la afirmación del laboratorio.\\

    \noindent
    Sea la variable aleatoria:
    \begin{equation*}
        X\equiv \text{``Tiempo (en horas) en recibir la primera picadura''}
    \end{equation*}
    que se distribuye según una función de distribución $F$, planteamos el contraste:
    \begin{equation*}
        \left\{\begin{array}{l}
            H_0:F=F_0 \\
            H_1:F\neq F_0
        \end{array}\right.
    \end{equation*}
    donde $F_0 = F_{\text{exp}(2.5)}$. Como la distribución exponencial es continua podemos aplicar el Test de Kolmogorov-Smirnov, por lo que nos disponemos a calcular primero $D_{\text{exp}}$, mediante la fórmula (como las $n=20$ observaciones son distintas):
    \begin{equation*}
        D_{\text{exp}} = \max\{\max_{x_i}\{F_{X_1,\ldots,X_n}^\ast(x_i)-F_0(x_i)\}, \max_{x_i}\{F_0(x_i)-F_{X_1,\ldots,X_n}^\ast(x_i^-)\}\}
    \end{equation*}
    Para ello, usamos la tabla (abreviando con $F^\ast = F_{X_1,\ldots,X_n}^\ast$):
    \begin{table}[H]
    \centering
    \begin{tabular}{|c|c|c|c|c|c|c|}
        \hline
        $x_i$ & $n F^\ast(x_i)$ & $F^\ast(x_i)$ & $F^\ast(x_i^-)$ & $F_0(x_i)$ & $F^\ast(x_i)-F_0(x_i)$ & $F_0(x_i) - F^\ast(x_i^-)$\\
        \hline
        $0.01$ & 1 & 0.05& 0    & 0.02469 & 0.02530 & 0.02469\\
        $0.02$ & 2 & 0.1 & 0.05 & 0.04877 & 0.05122 & -0.0012\\
        $0.03$ & 3 & 0.15& 0.1  & 0.07225 & 0.07774 & -0.0277\\
        $0.23$ & 4 & 0.2 & 0.15 & 0.43729 & -0.2372 & 0.28729\\
        $0.51$ & 5 & 0.25& 0.2  & 0.72056 & -0.4705 & 0.52056\\
        $0.74$ & 6 & 0.3 & 0.25 & 0.84276 & -0.5427 & 0.59276\\
        $0.96$ & 7 & 0.35& 0.3  & 0.90928 & -0.5592 & 0.60928\\
        $1.17$ & 8 & 0.4 & 0.35 & 0.94633 & -0.5463 & 0.59633\\
        $1.46$ & 9 & 0.45& 0.4  & 0.98772 & -0.5377 & 0.58772\\
        $1.62$ & 10& 0.5 & 0.45 & 0.98257 & -0.4825 & 0.53257\\
        $2.18$ & 11& 0.55& 0.5  & 0.99570 & -0.4457 & 0.49570\\
        $2.25$ & 12& 0.6 & 0.55 & 0.99639 & -0.3963 & 0.44639\\
        $2.79$ & 13& 0.65& 0.6  & 0.99906 & -0.3490 & 0.39906\\
        $3.45$ & 14& 0.7 & 0.65 & 0.99982 & -0.2998 & 0.34982\\
        $3.82$ & 15& 0.75& 0.7  & 0.99993 & -0.2499 & 0.29993\\
        $3.92$ & 16& 0.8 & 0.75 & 0.99994 & -0.1999 & 0.24994\\
        $4.27$ & 17& 0.85& 0.8  & 0.99997 & -0.1499 & 0.19997\\
        $5.43$ & 18& 0.9 & 0.85 & 0.99999 & -0.0999 & 0.14999\\
        $5.79$ & 19& 0.95& 0.9  & 0.99999 & -0.0499 & 0.09999\\
        $6.34$ & 20& 1   & 0.95 & 0.99999 & 0       & 0.04999\\
        \hline
    \end{tabular}
    \end{table}
    \noindent
    de donde obtenemos:
    \begin{equation*}
        \left.\begin{array}{l}
            \max\limits_{x_i}\{F_{X_1,\ldots,X_n}^\ast(x_i)-F_0(x_i)\} = 0.07774 \\
            \max\limits_{x_i}\{F_0(x_i)-F_{X_1,\ldots,X_n}^\ast(x_i^-)\} = 0.60928 
        \end{array}\right\} \quad\Longrightarrow\quad D_{\text{exp}} = 0.60928
    \end{equation*}
    Calculamos el $p-$valor, usando para ello la distribución $Z$ de Kolmogorov, donde $D(X_1,\ldots,X_{20})\rightsquigarrow Z(20)$:
    \begin{equation*}
        P_{H_0}[D(X_1,\ldots,X_n)\geq D_{\text{exp}}] = P_{H_0}[D(X_1,\ldots,X_n)\geq 0.60928] < 0.001
    \end{equation*}
    obtenemos un dato mucho menor que $0.001$, por lo que podemos rechazar la afirmación del laboratorio.
\end{ejercicio}

\begin{ejercicio}
    Cierta comunidad ha modificado la procedencia del agua destinada al consumo doméstico. Se sabe que, con el antiguo suministro, la distribución de la cantidad de sodio por unidad de volumen de sangre de sus habitantes es simétrica alrededor de $3.24$ gr. Tras cierto tiempo, se quiere comprobar si la modificación ha afectado a la concentración de sodio, en el sentido de que su distribución se haya trasladado o no. Para ello, se han realizado 15 análisis con los siguientes resultados (en gr. por unidad):
    \begin{multline*}
        2.37 \quad 2.95 \quad 3.4 \quad 2.64 \quad 3.66 \quad 3.18 \quad 2.72 \quad 3.61 \quad 3.87 \quad 1.97 \quad 1.66 \quad 3.72 \quad 2.10 \\ 1.83 \quad 3.03
    \end{multline*}
    ¿Se puede afirmar, al nivel de significación $0.1$, que la distribución de la cantidad de sodio no ha variado?
\end{ejercicio}

\begin{ejercicio}
    La siguiente tabla presenta las presiones sanguíneas sistólicas de 10 individuos antes y después de haber dejado la bebida.
    \begin{table}[H]
    \centering
    \begin{tabular}{c|cccccccccc}
        \hline
        $A$ & 140 & 165 & 160 & 160 & 175 & 190 & 170 & 175 & 155 & 160 \\
        \hline
        $D$ & 145 & 150 & 150 & 160 & 170 & 175 & 160 & 165 & 145 & 170 \\
        \hline
    \end{tabular}
    \end{table}
    \noindent
    ¿Se puede afirmar a partir de los datos que el abandono de la bebida no disminuye la presión sanguínea? ¿Bajo qué hipótesis?
\end{ejercicio}

\begin{ejercicio}
    En cierta comunidad de E.E.U.U. se realizó un estudio para investigar si el sueldo anual de las familias influía en los hijos para la elección de los diferentes cursos de enseñanza secundaria (Preparatorio, General y Comercial). Para ello, se hizo una clasificación de los sueldos en cuatro niveles (I, II, III y IV), y se tomó una muestra aleatoria simple de 390 estudiantes, obteniéndose la siguiente tabla de frecuencias:
    \begin{table}[H]
    \centering
    \begin{tabular}{|c|c|c|c|c|}
        \hline
        Sueldo & I & II & III & IV \\
        \hline
        Preparatorio & 23 & 40 & 16 & 2 \\
        \hline
        General & 11 & 75 & 107 & 14 \\
        \hline
        Comercial & 1 & 31 & 60 & 10 \\
        \hline
    \end{tabular}
    \end{table}
    \noindent
    A la vista de los datos, decidir, al nivel de significación $0.01$, si se acepta que el nivel económico familiar no influye en la decisión de los estudiantes a la hora de elegir curso.
\end{ejercicio}

\begin{ejercicio}
    En un estudio sociológico sobre la polución atmosférica se entrevistó a 40 residentes de cada una de tres zonas residenciales en Gran Bretaña. La siguiente tabla muestra las respuestas a la pregunta: ¿Hay problema de polución en su barrio?
    \begin{table}[H]
    \centering
    \begin{tabular}{|c|c|c|c|c|}
        \hline
        Zona residencial & No & Sí & No sabe & No contesta \\
        \hline
        1 & 5 & 31 & 2 & 2 \\
        \hline
        2 & 10 & 21 & 4 & 5 \\
        \hline
        3 & 11 & 20 & 7 & 2 \\
        \hline
    \end{tabular}
    \end{table}
    \noindent
    Contrastar si las tres poblaciones de residentes pueden considerarse homogéneas con respecto a su opinión sobre la polución.
\end{ejercicio}

\begin{ejercicio}
    Para determinar si diferentes tipos de profesiones de los individuos activos de cierto colectivo afectan a la tensión arterial, se clasificó a los individuos en cuatro grupos, atendiendo a su profesión, y se midió la tensión a una muestra de individuos elegidos de forma aleatoria. Clasificando la tensión en los niveles ``Bajo'', ``Normal'' y ``Alto'', se obtuvo los siguientes resultados, que muestran el número de individuos de cada tipo de profesión con los distintos niveles de tensión:
    \begin{table}[H]
    \centering
    \begin{tabular}{|c|c|c|c|}
        \hline
        Profesión & Bajo & Normal & Alto \\
        \hline
        I & 8 & 4 & 3 \\ 
        \hline
        II & 5 & 7 & 7 \\
        \hline
        III & 4 & 8 & 8 \\
        \hline
        IV & 5 & 7 & 8 \\
        \hline
    \end{tabular}
    \end{table}
    \noindent
    ¿Qué conclusión acerca del problema planteado se obtiene a la vista de estos datos?  Especificar las hipótesis nula y alternativa que se contrastan.
\end{ejercicio}

\begin{ejercicio}
    Para determinar si las calificaciones de los alumnos en selectividad son independientes de las calificaciones en bachiller, se eligió de forma aleatoria una muestra de alumnos, a los que se preguntó ambas calificaciones, obteniendo los siguientes resultados:
    \begin{table}[H]
    \centering
    \begin{tabular}{|c|c|c|c|c|}
        \hline
        Bachiller & Suspenso & Aprobado & Notable & Sobresaliente \\
        \hline
        Aprobado & 10 & 6 & 4 & 5 \\
        \hline
        Notable & 7 & 9 & 9 & 4 \\
        \hline
        Sobresaliente & 6 & 10 & 10 & 6 \\
        \hline
    \end{tabular}
    \end{table}
    \noindent
    ¿Qué conclusión acerca del problema planteado se obtiene a la vista de estos datos?  Especificar las hipótesis nula y alternativa que se contrastan.
\end{ejercicio}

\end{document}
